\section{Involutions of the three-cusp group}

Throughout this part \(F=\widehat F(x,y)\), \(z=(xy)^{-1}\), and
\(\mathcal P\) denotes the subgroup of \(\operatorname{Out}(F)\)
preserving each named peripheral conjugacy class up to a common
unit power. All group homomorphisms are continuous. We use the
normalized lift
\[
 (\lambda,f):\quad x\longmapsto x^\lambda,\qquad
 y\longmapsto f^{-1}y^\lambda f,\qquad f\in F',
\]
where \(F'\) is the closed derived subgroup. This normalization is
a homomorphic section, as in the refereed LM4 D.32. Write
\(c=(-1,1)\). We use R1--R3 only with their stated scopes.

\statement{A.1}{Known}{Torsion in an acting group and its action on torsion are different assertions}
Lochak and Schneps (1997, Proposition 4(i)) prove
\(C_F(c)=1\); Proposition 4(ii) proves that \(c\) is
self-centralizing in \(\widehat{GT}\).
Schneps (1997, \S1.4(4)) and Lochak and Schneps (2006,
Question 1.4) ask whether all nonidentity torsion in
\(\widehat{GT}\) consists of conjugates of \(c\).
Harbater and Schneps (2000, Main Theorem(a), proved by
Propositions 3--4 and Theorem 2) identify
\(\widehat{GT}_0\) with the centralizer in \(\mathcal P\) of
the geometric \(S_3\).
Collas's genus-zero Theorems 3.3 and 3.10 identify prime-order
torsion in \(\widehat\Gamma_{0,[n]}\) with discrete torsion
and prove the \(\lambda\)-conjugacy action of \(\widehat{GT}\).
His genus-one Theorems 6.3 and 6.5 give the corresponding
statements with acting group \(\widehat{GR}\), isomorphic
to the subgroup denoted \(I\Gamma\) in that article;
replacing it there by all of \(\widehat{GT}\) is not
justified by those theorems.\cite{A_LS1997,A_Schneps1997,A_LS2006,A_HS2000,A_Collas2012a,A_Collas2012b}

\begin{proof}
For the centralizer comparison identify \(K(0,4)\) with
\(\widehat F(x,y)\). The natural \(S_4\)-action on its outer
automorphism group factors through the anharmonic \(S_3\).
Harbater--Schneps's quasi-special condition is precisely the
three peripheral power condition here; their symmetric
condition is commutation with that \(S_3\). Their proof
translates commutation with the transposition into the
two-cycle relation and commutation with the three-cycle into
the hexagon. The identification has the required source
and target, not merely an abstract isomorphism.

In Collas's formulas the variable of prime order lies in a
completed mapping class group; \((\lambda,f)\) is an arbitrary
element of the acting group. The conclusion
\(F(\gamma)=u\gamma^\lambda u^{-1}\) does not classify
finite-order elements \(F\) of the acting group. The subgroup
\(\widehat{GR}\) in the genus-one article is retained in its
Theorem 6.5. Neither action theorem therefore proves task A.8.
\end{proof}

The opened sources establish no literature antecedent for the
full order bound in R1. Searches included the authors' source
lists, the cited torsion papers, and recent work on profinite
curve-complex automorphisms and Galois--Teichm\"uller theory.
This records the search, not a priority claim or a proof that
no antecedent exists.

\statement{A.2}{Proved}{Relative pronilpotent levels recover the full peripheral group}
For an open characteristic subgroup \(U\) of \(F\), put
\[
 R_U=R_{\mathrm{nil}}(U),\quad F_U=F/R_U,\quad
 P_U=\operatorname{Out}^{\sharp}(F_U).
\]
For \(V\subset U\), the quotient \(F_V\to F_U\) induces
a continuous epimorphism \(P_V\to P_U\). There is an
isomorphism of profinite groups
\[
                 \mathcal P \xrightarrow{\sim}\varprojlim_U P_U.
\]
Two elements of \(\mathcal P\) are conjugate in \(\mathcal P\)
if and only if their images are conjugate in every \(P_U\).
This proves Lead 1.

\begin{proof}
Restricting finite nilpotent quotients of \(U\) to \(V\)
shows \(R_V\subset R_U\). Each \(R_U\) is characteristic in
\(F\). Since \(R_U\subset U\) and open characteristic
subgroups are cofinal in the finitely generated profinite
group \(F\), we have \(\bigcap_U R_U=1\).

We first prove surjectivity onto \(P_U\). Choose a representative
of a peripheral outer class with common multiplier \(\lambda\),
and a Galois element with cyclotomic multiplier \(\lambda\);
the cyclotomic character is onto
\(\widehat{\mathbb Z}^{\times}\). The latter acts on \(F_U\).
Composing with its inverse gives a multiplier-one peripheral
automorphism. R2 lifts this automorphism, with chosen
conjugators, to \(F\). Composing back proves surjectivity.
In particular every peripheral automorphism of \(F_V\)
lifts to \(F\), so preserves the image of \(R_U\).
This proves that the transition maps are defined on the
entire displayed outer groups. Their surjectivity follows
by lifting first to \(\mathcal P\).

Here the automorphism groups are compact: intersecting
kernels of all maps to finite groups of bounded order
gives a cofinal system of open fully invariant subgroups.
The conditions that a continuous endomorphism induces
bijections on these finite quotients are closed conditions
on its generator images. Consequently \(\operatorname{Aut}(F)\)
is profinite. Inner automorphisms form a compact subgroup.
The peripheral conditions are closed, since conjugators
and the common multiplier range over compact groups.
The same argument applies to the finitely generated
quotients \(F_U\).

If \([a]\in\mathcal P\) dies in every \(P_U\), the sets
\[
 \{g\in F:a(x)\equiv gxg^{-1},\
              a(y)\equiv gyg^{-1}\pmod {R_U}\}
\]
are nonempty compact sets with the finite intersection
property. Their intersection gives an actual inner
representative for \(a\), proving injectivity.
For a compatible family of classes in \(P_U\), their
nonempty compact inverse images in \(\mathcal P\) have
the same finite intersection property, proving surjectivity.
Finally apply this argument to the transporter sets
\(\{g:g_Ua_Ug_U^{-1}=b_U\}\) in \(\mathcal P\).
Surjectivity onto \(P_U\) makes these nonempty whenever
the two images are conjugate. Their intersection is the
required full conjugator.
\end{proof}

\statement{A.3}{Proved}{A central quotient is the exact obstruction to a peripheral lift}
Let \(N\) be a closed characteristic subgroup of \(F\), and
let \([N,F]\) denote the closed subgroup generated by
commutators between \(N\) and \(F\).
An automorphism \(a\) of \(F/N\), with specified conjugators
\(b_i\) satisfying
\[
 a(\bar x_i)=b_i\bar x_i b_i^{-1},\qquad
 (x_1,x_2,x_3)=(x,y,z),
\]
lifts to a peripheral automorphism of \(F\) with these
conjugators if and only if it lifts to a peripheral
\emph{automorphism} of \(F/[N,F]\) with conjugators above
the \(b_i\). This proves Lead 2.

\begin{proof}
Necessity follows from characteristicity. For sufficiency
let \(a^\#\) and \(b_i^\#\) be the data on \(F/[N,F]\).
We use the finite central-refinement count in the refereed
LM7\_PUSH D.50, equations (D.50.6)--(D.50.16).
For any actual quotient \(q_E:F\twoheadrightarrow E\)
and normal subgroup \(A\triangleleft E\), that count
constructs an actual quotient
\(q^\dagger:F\twoheadrightarrow D^\dagger\), with central
kernel over \(D=E/A\). A second epimorphism to
\(D^\dagger\) with specified peripheral conjugators has
the same number of epimorphism lifts to \(E\), with the
transported \(A\)-conjugacy classes on the boundary.
This counting assertion uses only normality of \(A\).
The hypothesis \([E,A]=A\) occurs in the separate
joint-image assertion (D.50.17), not used here.

Take \(A=q_E(N)\). The image \(q^\dagger(N)\) is central,
so \(q^\dagger\) kills \([N,F]\).
Thus \(q^\dagger a^\#\) is an epimorphism from the
same source to \(D^\dagger\), with conjugators the
images of \(b_i^\#\). The finite count gives an epimorphism
to \(E\) whose conjugators agree with the \(b_i\) in \(D\);
the reference count is positive since it includes \(q_E\).

In the compact space \(F^3\) of conjugators impose the
boundary product relation, generation in every finite
quotient, and agreement with each \(b_i\) in every finite
quotient of \(F/N\). A finite list of conditions lies in
one joint finite quotient \(E\), and the preceding
construction satisfies that list. Compactness gives
full conjugators satisfying all conditions. Their
endomorphism is surjective since its compact image
maps onto every finite quotient; Hopficity makes it
an automorphism. Its generator images induce exactly
\(a\) modulo \(N\).
\end{proof}

\statement{A.4}{Proved}{A degree-three coefficient imposed by the third cusp}
Over a field \(k\) of characteristic zero or greater
than three, use the free class-three Lie algebra on
\(X,Y\), and set
\[
 Z=[X,Y],\quad U_3=[X,Z],\quad V_3=[Y,Z],\quad
 x=\exp X,\quad y=\exp Y .
\]
If \(x\mapsto x^\mu,\ y\mapsto g^{-1}y^\mu g\), with
\(\mu\in k^\times\) and
\(\log g=tZ+\text{central terms}\), sends \(xy\)
to a conjugate of \((xy)^\mu\), then
\(t=(\mu^2-1)/24\).

\begin{proof}
The Baker--Campbell--Hausdorff formula gives
\[
 W=\log(xy)=X+Y+\tfrac12Z+\tfrac1{12}(U_3-V_3).
\]
The logarithm of the image is
\[
 \mu(X+Y)+\tfrac{\mu^2}{2}Z+
       \tfrac{\mu^3}{12}(U_3-V_3)+\mu tV_3 .
\]
Write the logarithm of a conjugator as
\(aX+bY+dZ+\text{central terms}\).
In \(\exp(\operatorname{ad}(aX+bY+dZ))(\mu W)\),
comparison of degree two gives \(m=a-b=(\mu-1)/2\).
The difference of the \(U_3\)- and \(V_3\)-coefficients,
divided by \(\mu\), then gives
\[
 \frac{\mu^2}{6}-t=\frac16+\frac m2+\frac{m^2}{2}
                  =\frac{3\mu^2+1}{24}.
\]
The result follows. The coefficient \(d\) and all
central coordinates cancel.
\end{proof}

\statement{A.5}{Refuted}{The specified conjugator in task A.3 need not preserve the third cusp}
There is an involution \(\alpha\in\widehat{GT}\), conjugate
to \(c\) by a Galois element, for which the unique solution
\(h\in F'\) of \(f_\alpha=h^{-1}\alpha(h)\) gives
\[
                 \psi(x)=x,\qquad\psi(y)=hyh^{-1}
\]
that does not preserve the conjugacy class of \(z\)
on \(F/R_{\mathrm{nil}}(F)\). The failure occurs in
the class-three exponent-seven quotient of order \(7^5\).

\begin{proof}
Choose \(\theta=(\mu,g)\) in the Galois image with
\(\mu\equiv2\pmod7\), possible by surjectivity of the
cyclotomic character modulo seven, and put
\(\alpha=\theta c\theta^{-1}\).
Normalization is a homomorphic section, so these
are actual automorphisms with that equality.
From \(\alpha\theta(y)=\theta(y)^{-1}\) we obtain
\[
 \alpha(y)=\alpha(g)g^{-1}y^{-1}g\alpha(g)^{-1},
 \qquad f_\alpha=g\alpha(g)^{-1}.
\]
Thus \(h=g^{-1}\) is a solution. Also
\(C_F(\alpha)=\theta(C_F(c))=1\) by
Lochak--Schneps Proposition 4(i). The quotient of two
solutions is fixed by \(\alpha\), so \(h\) is unique.

Let \(Q=\exp L\) for the free class-three Lie algebra
\(L\) on two generators over \(\mathbb F_7\).
It has order \(7^5\), exponent seven, and is generated
by \(x=\exp X,y=\exp Y\); hence it is a quotient of
\(F^{(7)}\) and of \(F/R_{\mathrm{nil}}(F)\).
The third-cusp condition on \(\theta\) and A.4 give
\[
       \operatorname{coeff}_Z(\log\bar g)
                    =\frac{2^2-1}{24}=1 .
\]
But \(\psi\) has multiplier one and the same conjugator
\(g\). Peripheral preservation for \(\psi\) would force
this coefficient to be zero by A.4, a contradiction.

The exact certificate \texttt{check\_A3\_class3.py}
realizes \(Q\) in the units of
\[
 \mathbb F_7\langle X,Y\rangle/(\text{words of length }4).
\]
It enumerates all \(16807\) elements; the conjugacy
orbit of \(xy\) has \(49\) elements. For
\(g=\exp[X,Y]\) it verifies
\[
 xg^{-1}yg\not\sim xy,\qquad
 x^2g^{-1}y^2g\sim(xy)^2 .
\]
The actual Galois conjugator differs only by central
degree-three coordinates relevant to this calculation,
which do not affect it. This certifies a failure in
an actual full profinite example; no order-preserving
lift of a finite involution is being assumed.
\end{proof}

The correct universal statement from R1 is that this
\(\psi\) intertwines \(\alpha,c\) on the maximal
pro-\(p\) quotients as automorphisms of free groups.
Peripheral preservation of this specific \(\psi\)
is an additional condition and is false.

\statement{A.6}{Proved}{The pure centralizer of standard conjugation has two elements}
\[
 C_{\mathcal P}(c)=N_{\mathcal P}(\langle c\rangle)=\{1,c\}.
\]
If \(\alpha\in\widehat{GT}_0\), every conjugator in
\(\mathcal P\) from \(c\) to \(\alpha\) belongs
to \(\widehat{GT}_0\).

\begin{proof}
A normalized \((\mu,g)\) commuting with \(c\) satisfies
\(c(g)=g\), by the normalized group law. A.1 gives
\(g=1\). In particular the multiplier-one centralizer
is trivial. Let \(\Sigma=S_3\) be the geometric
permutation group. If
\(\beta c\beta^{-1}=\alpha\) and \(\alpha,c\) commute
with \(\Sigma\), then
\(r_\sigma=\beta^{-1}\sigma\beta\sigma^{-1}\)
is pure, has multiplier one, and centralizes \(c\).
Thus \(r_\sigma=1\), and the Harbater--Schneps
identification in A.1 puts \(\beta\) in
\(\widehat{GT}_0\).

Apply this to a centralizing \((\mu,1)\).
The hexagon, with \(m=(\mu-1)/2\), reads
\((xy)^m=y^mx^m\). In the finite Heisenberg groups
with multiplication
\((a,b,d)(a',b',d')=(a+a',b+b',d+d'+ab')\),
the final coordinates of these two words are
\(m(m+1)/2\) and zero. Hence \(\mu^2=1\) in
\(\widehat{\mathbb Z}\). The actual automorphism
\((\mu,1)\) then has square one. R1 forces its
multiplier to be globally \(1\) or \(-1\), proving
the centralizer assertion. A normalizer of
\(\langle c\rangle\) centralizes its unique
nonidentity element.
\end{proof}

\statement{A.7}{Proved}{The multiplier-one cohomology set already has continuum many classes}
For \(G=\mathcal P,\widehat{GT}_0\), or \(\widehat{GT}\),
put \(H_G=\ker\lambda\), with \(C_2\) acting by \(c\).
There is a canonical closed injection of profinite spaces
\[
 \widehat{\mathbb Z}^{\times}/\{1,-1\}
 \longrightarrow H^1(C_2,H_G),\qquad
 \mu\longmapsto[\theta_\mu c\theta_\mu^{-1}c],
\]
where \(\theta_\mu\) is any element of multiplier \(\mu\).
The set \(H^1(C_2,H_G)\) has cardinality \(2^{\aleph_0}\).
For \(G=\widehat{GT}_0\), \(\Delta\) vanishes on this
closed family. All these \(H_G\)-conjugacy classes form
one \(G\)-conjugacy class.

\begin{proof}
A cocycle is \(a\in H_G\) with \(a\,c(a)=1\); it
corresponds to the involution \(ac\). Cohomology
identifies \(a\) and \(h^{-1}a\,c(h)\), exactly
conjugacy by \(H_G\). The Galois image makes
\(\lambda:G\to\widehat{\mathbb Z}^{\times}\) onto.
The element \(a_\mu=\theta_\mu c\theta_\mu^{-1}c\)
is a cocycle, and replacing \(\theta_\mu\) by a
different lift conjugates its involution by \(H_G\).

Two displayed involutions are \(H_G\)-conjugate
if and only if
\(\theta_\nu^{-1}h\theta_\mu\in C_G(c)=\{1,c\}\)
for some \(h\in H_G\), by A.6. Taking multipliers
gives \(\mu=\pm\nu\); these conditions are sufficient
as well. This proves the injection.
The cocycle space is compact, and its quotient by
the compact conjugating group is Hausdorff.
The orbit map from \(G\) factors continuously
through the multiplier quotient and then its
two signs, so the injection is continuous and closed.
Its domain has cardinality continuum. Conversely
\(G\) is contained in the metrizable profinite group
\(\operatorname{Out}(F)\), of cardinality at most
continuum. This proves equality.
Finally every represented involution is conjugate
to \(c\); the character \(\Delta\) is conjugacy
invariant and \(\Delta(c)=0\).
\end{proof}

\statement{A.8}{Proved}{Full conjugacy is an additional quotient of nonabelian cohomology}
With the notation of A.7, the full \(G\)-conjugacy
classes of nonidentity involutions are
\[
 \widehat{\mathbb Z}^{\times}\backslash H^1(C_2,H_G).
\]
The action of \(\mu\), using any lift \(\theta_\mu\), is
\[
 [a]\longmapsto[\,\theta_\mu a\,c(\theta_\mu^{-1})\,].
\]
For \(\mathcal N=\ker(\mathcal P\to P_F)\), the
\(\mathcal N\)-conjugacy classes of involutions with
image exactly \(c_F\) are \(H^1(C_2,\mathcal N)\).
Full \(\mathcal P\)-conjugacy on that fibre further
quotients by \(C_{P_F}(c_F)\), acting by lifted
conjugation, which need not preserve the distinguished
cohomology class.

\begin{proof}
By R1 every nonidentity involution has multiplier
\(-1\), so is uniquely \(ac\), \(a\in H_G\).
Its square-one condition is the cocycle equation.
Conjugation by \(\theta_\mu\) gives the formula.
Different lifts differ by \(H_G\), already divided
out in cohomology, proving the first assertion.

The same calculation for \(\mathcal P/\mathcal N\)
shows that involutions over the literal \(c_F\)
are \(ac\), \(a\in\mathcal N\). A conjugation
preserves this fibre exactly when its quotient
centralizes \(c_F\). Changing a lift changes
the action only by \(\mathcal N\)-conjugacy.
Such a lift need not centralize \(c\), explaining
the last qualification.
\end{proof}

\statement{A.9}{Proved}{Odd-prime conjugacy over a relative Frattini level}
Let \(U\) be open characteristic in \(F\), let \(p\)
be odd, put \(E_p=F/O^p(U)\), \(V=U/O^p(U)\),
and consider the peripheral outer subgroup preserving
\(V\). Its kernel on
\[
           E_p/\Phi(V)=F/\overline{U'U^p}
\]
is pro-\(p\). Two involutions with the same image
in this Frattini quotient are conjugate by the kernel.
This applies to the fibre over \(c\), where that
fibre is defined by descent to \(F/U\).

\begin{proof}
\(V\) is a nonabelian finitely generated free
pro-\(p\) group. Let \(J\) be the actual
automorphisms of \(E_p\) inducing identity on
\(E_p/\Phi(V)\). Restriction maps \(J\) into
the Frattini kernel of \(\operatorname{Aut}(V)\).
That kernel is pro-\(p\): on finite lower
\(p\)-central quotients its successive kernels
embed into homomorphism groups of elementary
abelian \(p\)-groups, and those quotients are cofinal.

Restriction is injective. If an element of \(J\)
fixes \(V\) pointwise, write \(a(g)=gu_g\) with
\(u_g\in\Phi(V)\). Comparing the conjugation
actions on \(V\) gives \(u_g\in C_V(V)=1\).
Thus \(J\) is pro-\(p\). If an outer automorphism
acts inner on \(E_p/\Phi(V)\), lift its inner
conjugator to \(E_p\) and correct it to an
element of \(J\). Consequently the outer
kernel is a continuous quotient of a closed
subgroup of \(J\), hence pro-\(p\).

The common image of the two involutions has
order two, since a pro-\(p\) group for odd \(p\)
has no involution. In its inverse image the
two cyclic subgroups are complements to the
normal pro-\(p\) kernel. Finite
Schur--Zassenhaus gives conjugators in every
finite quotient; compactness gives a conjugator
in the kernel itself.
\end{proof}

For arbitrary quotient automorphisms, preservation
of the image of a merely characteristic \(U\)
must be checked. It is included in the phrase
over the class of \(c\). For fully characteristic
\(U\) it is automatic: any quotient endomorphism
lifts to an endomorphism of the free profinite
group and preserves \(U\). Such \(U\) are cofinal.

\statement{A.10}{Proved}{Every negative pure pro-\(p\) involution is locally conjugate for odd \(p\)}
Let \(P^{(p)}\) be the pure peripheral outer group
of \(F^{(p)}\), with multiplier in
\(\mathbb Z_p^\times\), for odd \(p\).
Every involution of multiplier \(-1\) is
conjugate to \(c\) by a multiplier-one element.

\begin{proof}
The multiplier kernel \(H^{(p)}\) acts trivially
on abelianization, so lies in the pro-\(p\)
Frattini kernel used in A.9. The inverse image
of \(\{1,-1\}\) under the multiplier is
\(H^{(p)}\rtimes\langle c\rangle\).
Every negative involution defines a complement;
Schur--Zassenhaus gives the required conjugacy.
\end{proof}

This does not rescue the particular map of A.5.
Its local coboundary can be changed by a local
fixed element, whereas its full coboundary is
unique. The comparison of fixed subgroups
need not be surjective.

\statement{A.11}{Proved}{A peripheral reflection triangle is equivalent to full conjugacy}
For a nontrivial involution \(\alpha\in\mathcal P\)
let \(E=F\rtimes_\alpha C_2\). Then \(\alpha\)
is conjugate to \(c\) in \(\mathcal P\) exactly
when there are involutions \(t_0,t_1,t_2\)
representing its three section classes such that
\[
 a=t_1t_0,\qquad b=t_0t_2,\qquad d=t_2t_1
\]
generate conjugates of the respective inertia
groups of \(x,y,z\), and the \(t_i\) generate
\(E\) topologically. The last condition is
equivalent to generation in every \(E/R_U\).

\begin{proof}
For standard \(c\) take
\(t_0=c,t_1=xc,t_2=y^{-1}c\); their products
are \(x,y,z\) and they generate the extension.
A peripheral conjugator transports this triple.

Conversely the universal property gives a
surjection from the profinite coproduct
\(C_2^{(0)}\amalg C_2^{(1)}\amalg C_2^{(2)}\)
to \(E\), sending the factors to the named
involutions. Its parity kernel is free
profinite on \(a,b\). The restriction maps
onto \(F\): every even word in three
reflections is a word in the two edge products.
Both kernels are free profinite of rank two,
so Hopficity makes this restriction an
isomorphism. The equality \(abd=1\) is literal.
Abelianization makes the unit exponents in
the three peripheral conditions equal,
so the induced outer map belongs to \(\mathcal P\).
Conjugation by \(t_0\) inverts both \(a\) and \(b\).
Its outer action on \(F\) is \(\alpha\), since
\(t_0\) and the named splitting element have the
same image in \(E/F\). Hence the resulting outer
map conjugates \(c\) to \(\alpha\); no choice of
an actual splitting representative is needed.

Finally put \(H=\overline{\langle t_0,t_1,t_2\rangle}\).
If \(HR_U=E\) for all \(U\), then for every \(e\in E\)
the compact sets \(H\cap eR_U\) are nonempty
and have the finite intersection property.
Their intersection and \(\bigcap R_U=1\)
imply \(e\in H\). Thus \(H=E\).
\end{proof}

The transcription of cote 145, scan 23, lists
\(l_0=t_\infty t_1,\ l_1=t_0t_\infty,\
l_\infty=t_1t_0\). With ordinary multiplication
these satisfy \(l_0l_\infty l_1=1\).
The cyclic order or path-composition convention
must therefore be made explicit before
identifying them with \(x,y,z\).
We use \(abd=1\), and do not emend an unseen scan.

\statement{A.12}{Open}{Dyadic and mixed-prime involution obstructions}
A.9 does not determine all two-primary classes or all
gluing classes over \(U=F'F^3\) and \(U_{12}\).
A.25 below constructs an actual involution with
parity one in the image of \(\widehat{GT}_0\) at
\(U_{12}\), and A.24 proves that this parity
descends to that outer image. Thus the existence
assertion at this specified relative level is proved.
Lead 3 still needs a compatible order-two full lift.
Neither R2 nor A.3 supplies that further equivariance.

For a surjection \(G\to Q\) with kernel \(L\),
a chosen lift \(t\) of an involution can be
corrected to an involution \(ut\), \(u\in L\),
exactly when
\[
                    u(tut^{-1})=t^{-2}.
\]
Surjectivity asserts the existence of \(t\),
not solvability of this nonabelian norm equation.
The map \(C_4\to C_2\) disproves that implication
in general, but is not a counterexample in
the peripheral groups. The residual task is
to solve the norm equations compatibly in
the actual relative-level kernels with the
specified frames and \(\Delta=1\), or to
prove they are insoluble.

\statement{A.13}{Open}{Parity on arbitrary involutions}
For an involution \(\tau=(-1,f)\) in
\(\widehat{GT}_0\), R3 gives
\(\Delta(\tau)=\delta((1,f))\).
It is zero for A.7's continuum family and
for every \(\widehat{GT}_0\)-conjugate of \(c\).
The cocycle relation alone gives only
\(2\Delta(a)=0\) in \(\mathbb F_2\).
A.25 gives such an involution in an actual relative
outer image of \(\widehat{GT}_0\), but no full
involution of parity one or universal full vanishing
proof is obtained, in \(\widehat{GT}_0\) or
\(\widehat{GT}\).

\statement{A.14}{Open}{The three full conjugacy assertions}
We have not decided whether every involution
of \(\mathcal P,\widehat{GT}_0,\widehat{GT}\)
is conjugate to \(c\) in its respective group.
A.2 detects conjugacy of a fixed pair on
all relative levels; A.11 identifies the
equivalent peripheral generation condition.
Neither supplies the compatible local
conjugators or reflection triangles whose
existence is required.

For \S31, folios 198--199 verso, scans
391--394, the full finite-subgroup conclusion
proved by R1 is that the only abstract
possibilities in \(\mathcal P\) are \(1,C_2\),
and \(C_2=\langle c\rangle\) occurs.
Exhaustion of its embedding conjugacy
classes is exactly the remaining assertion.

\statement{A.15}{Proved}{The pronilpotent lifting kernel also has continuum many cohomology classes}
Put \(Q=\mathcal P_F\), \(\mathcal N=\ker(\mathcal P\to Q)\),
and write \(c_Q\) for the image of \(c\). There is a canonical
closed injection
\[
 C_Q(c_Q)/\langle c_Q\rangle
 \longrightarrow H^1(C_2,\mathcal N),\qquad
 \bar\beta\longmapsto[\beta c\beta^{-1}c].
\]
Its image is exactly the \(\mathcal N\)-conjugacy classes in the
fibre over \(c_Q\) which are conjugate to \(c\) in \(\mathcal P\).
Both its domain and \(H^1(C_2,\mathcal N)\) have cardinality
\(2^{\aleph_0}\). Thus the large cohomology set does not itself
provide an additional full involution-conjugacy class.

\begin{proof}
The pronilpotent quotient is the product of its maximal
pro-\(p\) quotients. Every factor is its characteristic
Sylow subgroup; conversely, a tuple of continuous
automorphisms of the factors gives a continuous automorphism
of their product. Inner automorphisms and the three peripheral
conditions are coordinatewise. Therefore
\[
 Q=\prod_p\mathcal P^{(p)},\qquad
 C_Q(c_Q)=\prod_p C_{\mathcal P^{(p)}}(c).
\]
For odd \(p\), the multiplier map on the latter centralizer
is onto \(\mathbb Z_p^\times\). Indeed choose an element
\(\theta\) of the desired multiplier using the Galois image.
By A.10, some multiplier-one \(v\) conjugates
\(\theta c\theta^{-1}\) to \(c\); then \(v\theta\)
centralizes \(c\) and retains that multiplier.
Consequently the displayed centralizer, and its quotient
by the finite group \(\langle c_Q\rangle\), have at least
continuum many elements.

By A.2 every \(\bar\beta\in C_Q(c_Q)\) lifts to
\(\beta\in\mathcal P\). Its displayed commutator with
\(c\) lies in \(\mathcal N\) and is a cocycle.
Changing the lift changes its class only by
\(\mathcal N\)-conjugacy. Two lifts \(\beta_1,\beta_2\)
give the same class exactly when
\(\beta_2^{-1}n\beta_1\in C_{\mathcal P}(c)=\{1,c\}\)
for some \(n\in\mathcal N\). On taking the quotient,
this is exactly equality of the two cosets modulo
\(\langle c_Q\rangle\); the converse follows by lifting
that equality. This proves injectivity and the asserted
description of the image.
Compactness makes the injection continuous and closed,
as in A.7. All groups and orbit spaces here are compact
metrizable, so their cardinalities are at most continuum.
The preceding lower bound proves the asserted equalities.
\end{proof}

If the same construction is made in \(\widehat{GT}_0\),
using its actual image in any \(c\)-stable quotient,
every class in the connecting-map image has \(\Delta=0\).
This follows from conjugacy invariance of the character;
no surjectivity onto the entire peripheral quotient is
required for this observation.

\statement{A.16}{Proved}{A prescribed reflection splits every pro-two extension of the symmetric group}
Let \(1\to K\to E\xrightarrow{q}S_3\to1\) be an extension
of profinite groups with \(K\) pro-two. If a transposition
\(t\in S_3\) has an involution lift \(\widetilde t\in E\),
then there is a continuous section \(j:S_3\to E\) satisfying
\(j(t)=\widetilde t\).

\begin{proof}
For every finite \(\mathbb F_2[S_3]\)-module \(M\), restriction
\[
 H^n(S_3,M)\longrightarrow H^n(\langle t\rangle,M)
\]
is an isomorphism for \(n\geq1\). Indeed the double-coset
formula for restriction after corestriction has the identity
term and one term through the trivial intersection
\(\langle t\rangle\cap r\langle t\rangle r^{-1}=1\),
where \(r\) is a three-cycle. The second term is zero
in positive degrees. Corestriction after restriction is
multiplication by the odd index three, hence the identity
on these \(\mathbb F_2\)-vector spaces.

First suppose \(K\) is finite. Induct on its order.
When \(K\ne1\), the subgroup
\(A=\{z\in Z(K):z^2=1\}\) is a nonzero elementary
abelian characteristic subgroup of \(K\), hence normal
in \(E\). Induction gives a section of \(E/A\to S_3\)
with the prescribed image of \(\widetilde t\).
Pull back that section to an extension with kernel \(A\).
Its class in \(H^2(S_3,A)\) restricts to zero, since
\(\widetilde t\) is an involution. The proved injectivity
of restriction makes the extension split.
The difference between any splitting on \(\langle t\rangle\)
and the specified lift is an \(A\)-valued one-cocycle.
Surjectivity in degree one extends its cohomology class to
\(S_3\); changing the section by that cocycle and then by
an \(A\)-conjugation makes its value at \(t\) literally
\(\widetilde t\). This proves the finite assertion.

For general pro-two \(K\), take all open normal subgroups
\(L\) of \(E\) contained in \(K\). Each finite quotient
has such a section with the prescribed transposition.
The corresponding closed subsets of the compact space
of tuples \((e_g)_{g\in S_3}\in E^{S_3}\), imposing the
section identities modulo \(L\), are nonempty and have
the finite intersection property. Their intersection
gives a section with all identities literal.
\end{proof}

\statement{A.17}{Proved}{The parity-one relative frame can be chosen anti-invariant under conjugation}
In the notation of R2, put \(H_0=F'F^3\),
\(\Gamma_2=\widehat B/O^2(H_0)\), and let
\(d=\operatorname{Inn}(s)\circ c\), so that
\(d(s)=s\), \(d(r)=r^{-1}\), \(d(x)=y^{-1}\),
and \(d(y)=x^{-1}\). There is an actual element
\[
 b\in C_{H_0^{(2)}}(r),\qquad q(b)=[5I],\qquad d(b)=b^{-1}.
\]
Consequently the identity automorphism of
\(\widehat B/R_{\mathrm{nil}}(H_0)\), with \(s\)- and
\(\tau\)-frames one and \(r\)-frame equal to \(b\) at two
and one at all other primes, is a framed cocycle for
conjugation by \(c\). Its modulo-twelve frame parity is one.
This establishes local equivariance of these data; it does
not assert an equivariant full lift.

\begin{proof}
Use R2's actual joint quotient
\[
 J=\operatorname{im}\bigl(\widehat B\to Q\times D_0\bigr),
 \qquad Q=\mathrm{SL}_2(\mathbb Z/12)/\{\pm I\},\quad
 D_0=\widehat B/H_0.
\]
The map \(\Gamma_2\to J\) has pro-two kernel; both factors
and the joint quotient are preserved by \(d\).
The element \(z_0=([5I],1)\in J\) is central of order two
and is fixed by \(d\). Thus
\(S=\langle\bar r,z_0d\rangle\) in \(J\rtimes\langle d\rangle\)
is an actual copy of \(S_3\), with its named three-cycle
and transposition.

There is a particularly small literal lift of that
transposition as an involution. Set
\[
 C=[x,y],\qquad h=CxCx^{-1}\in F'\subseteq H_0.
\]
In R2's signed matrix coordinates modulo twelve,
\[
 q(h)=\begin{pmatrix}9&4\\8&9\end{pmatrix},\qquad
 D=\begin{pmatrix}0&1\\1&0\end{pmatrix},\qquad
 q(h)\bigl(Dq(h)D^{-1}\bigr)^{-1}=5I.
\]
Conjugation by \(D\) represents \(d\) on \(Q\).
The \(D_0\)-coordinate of \(h\) is one, so the actual
involution \(hdh^{-1}\) maps to \(z_0d\).
Apply A.16 to the preimage of \(S\) in
\(\Gamma_2\rtimes\langle d\rangle\), retaining this
involution. Its lifted three-cycle is conjugate to
the named \(r\) by the pro-two kernel, by the
coprime complement-conjugacy argument of A.9.
After that conjugation write the lifted transposition
as \(bd\). Its image in \(D_0\rtimes\langle d\rangle\)
is \(d\), so \(b\in H_0^{(2)}\), and its \(Q\)-coordinate
gives \(q(b)=[5I]\). Its order two gives \(b\,d(b)=1\).
Since it inverts the named \(r\), while \(d\) also
inverts \(r\), the element \(b\) centralizes \(r\).

All components of the prescribed relative frame have
image one in \(D_0\), so they give an actual element
of the relative fibre product. For a multiplier-one
normalized automorphism \(\Phi\) of \(\widehat B\),
write \(\Phi(s)=\beta s\beta^{-1}\) and
\(\Phi(r)=A rA^{-1}\). Direct substitution from
\(c(s)=s\), \(c(r)=sr^{-1}s\) gives the frames of
\(c\Phi c\):
\[
 \beta^c=c(\beta),\qquad
 A^c=c(\beta)s\,c(\beta)^{-1}c(A)s.
\]
For the specified local data \(\beta=1\), the second
formula is \(A^c=d(A)=A^{-1}\). These are exactly
the frames of the inverse of the identity framed
automorphism. The parity is the same as in R2's E.58:
the signed \(r\)-frame is \(5I\), whose modulo-three
value is \(-I\), and the \(s\)-frame is one.
The independently executed certificate
\texttt{check\_A17\_frame.py} verifies the literal
matrix word and the covariance identities just used.
\end{proof}

The remaining full obstruction can now be placed in a
smaller, specified kernel. Let \(L\) consist of the
multiplier-one elements whose normalized action is
identity on \(\widehat B/R_{\mathrm{nil}}(H_0)\) and whose
two pure frames both have image one there. It is the
kernel of the framed action, is stable under \(c\), and
\(\Delta(L)=0\). R2 supplies a full \(\sigma\) lifting
A.17's data, with \(\Delta(\sigma)=1\). Put \(a=\sigma c\).
Then \(a^2\in L\). A parity-one involution in this
specified fibre exists exactly when some \(v\in L\)
solves
\[
                   v\,(ava^{-1})=a^{-2}.
\]
The local construction proves that the right-hand side
lies in this framed kernel. Solvability of this last full
norm equation remains open.

\statement{A.18}{Refuted}{R2 does not admit an unrestricted order-preserving refinement}
There is an actual involution of
\(\operatorname{Out}^{\sharp}(F/R_{\mathrm{nil}}(F))\)
which lifts to a full peripheral automorphism by R2 but
has no involution lift in \(\mathcal P\).

\begin{proof}
Use the actual product
\(F/R_{\mathrm{nil}}(F)=\prod_pF^{(p)}\).
Define \(\bar a\) to be standard conjugation \(c\)
on \(F^{(3)}\) and identity on every other factor.
This is a continuous automorphism of order two,
preserving each of the three peripheral conjugacy
classes up to the common unit
\[
 \lambda_3=-1,\qquad \lambda_p=1\quad(p\ne3).
\]
Its outer class is nonidentity, since its action on
the three-primary abelianization is \(-I\).
By A.2, which applies R2 after Galois translation
of the multiplier, this class has a lift in
\(\mathcal P\). Every lift has the displayed
multiplier, because the product quotient retains
the entire abelianization \(\widehat{\mathbb Z}^2\).
By R1, a nonidentity full involution must instead
have multiplier \(-1\) at every prime. Hence no
lift of \(\bar a\) can be an involution.
\end{proof}

This counterexample is in the actual peripheral level
groups, rather than in an unrelated finite extension.
It excludes a blanket equivariant strengthening of
the central-refinement count. A.17's remaining fibre
has multiplier globally \(-1\), so this orientation
obstruction does not decide that fibre.

\statement{A.19}{Proved}{The square obstruction at an abelian refinement and the character-lifting test}
Suppose \(a\) is a lift of a quotient involution,
\(a^2\in L\), and \(L_0\triangleleft\langle L,a\rangle\)
is contained in \(L\), with \(M=L/L_0\) finite abelian.
Conjugation by \(a\) makes \(M\) a \(C_2\)-module.
An involution exists above that quotient element in
\(\langle L,a\rangle/L_0\) exactly when
\[
 [a^2]=0\quad\hbox{in}\quad
 H^2(C_2,M)=M^{a}/(1+a)M.
\]
This obstruction vanishes when \(M\) has odd order.
For elementary abelian two-primary \(M\), its target is
\(\ker(a+1)/\operatorname{im}(a+1)\).

If \(\Delta\) has a continuous lift to a character
\(\widehat{GT}_0\to\mathbb Z/4\), then it vanishes on
every involution. The same conclusion follows from
a continuous lift of \(\delta\) on \(\mathcal H\)
which is invariant under conjugation by \(c\).
Existence of either lift is not established here.

\begin{proof}
The square of the conjugation action is conjugation
by \(a^2\in L\), which is trivial on \(M\).
Also \(a^2\) is invariant under conjugation by \(a\).
Replacing \(a\) by \(va\), with \(v\in L\), changes
its square in additive notation by \((1+a)\bar v\).
This proves the quotient criterion, which is the
usual cyclic two-cocycle description of \(H^2\).
For an invariant \(u\) of odd order, \(u/2\) is
defined and \((1+a)(u/2)=u\), proving vanishing.
In characteristic two the invariants are
\(\ker(a+1)\), giving the last formula.

A character to \(\mathbb Z/4\) sends an involution
to zero or two; both reduce to zero modulo two.
Alternatively, if \(b\in\mathcal H\) is a cocycle
for \(c\) and \(\theta\) is a \(c\)-invariant
lift of \(\delta\), then
\(0=\theta(bc(b))=2\theta(b)\), so
\(\delta(b)=0\).
\end{proof}

For the actual \(a=\sigma c\) of A.17, any proposed
global character lift \(\Theta\) would necessarily satisfy
\(\Theta(a^2)=2\), although \(a^2\in L\) and
\(\Delta(L)=0\). Thus such a character must detect the
residual square inside the framed kernel; it cannot be
inferred from the local frames alone. Conversely A.17's
local parity-one framed involution rules out a character
lift to \(\mathbb Z/4\) on that local framed group itself.

\statement{A.20}{Proved}{The locally glued symmetric subgroup does not lift in the standard extension}
Let \(H_0\), \(d\), and \(b\) be as in A.17. In the standard
relative extension
\[
 \bigl(\widehat B/R_{\mathrm{nil}}(H_0)\bigr)
                  \rtimes\langle d\rangle
\]
there is a subgroup \(S_{\mathrm{loc}}\simeq S_3\), generated
by the named \(r\) and a transposition with component \(bd\)
at two and \(d\) at every odd prime. No finite subgroup
of the standard full extension
\(\widehat B\rtimes\langle d\rangle\) maps onto
\(S_{\mathrm{loc}}\). The obstruction is visible already in
\(Q\rtimes\langle d\rangle\), with \(|Q|=576\).

\begin{proof}
All local components of the transposition have the same
image \(d\) in \(D_0\rtimes\langle d\rangle\).
They therefore form an actual element of the relative fibre
product. A.17 gives its order two and its inversion of \(r\).
Their images in \(D_0\rtimes\langle d\rangle\) generate
\(S_3\), so the resulting subgroup has order exactly six.

In the full free product \(\widehat B=C_2\amalg C_3\),
every element of order three is conjugate into the named
\(C_3\), and that factor is malnormal. These are
Herfort--Ribes Theorems A' and B'.\cite{A_HR1985}
In particular \(N_{\widehat B}(\langle r\rangle)=\langle r\rangle\).
Since \(d\) inverts \(r\),
\[
 N_{\widehat B\rtimes\langle d\rangle}(\langle r\rangle)
                         =\langle r,d\rangle\simeq S_3.
\]
Every full subgroup isomorphic to \(S_3\) intersects
\(\widehat B\) in its unique subgroup of order three:
it cannot be contained in \(\widehat B\), whose finite
subgroups conjugate into its factors by the same paper's
Theorem 1. Conjugating that order-three subgroup into
\(\langle r\rangle\) proves that every full \(S_3\)
is conjugate to \(\langle r,d\rangle\).

The image of \(S_{\mathrm{loc}}\) in the actual finite
quotient \(Q\rtimes\langle d\rangle\) is
\[
                   \langle R,[5I]d\rangle .
\]
It is not conjugate there to \(\langle R,d\rangle\).
Indeed a conjugator may be taken in \(Q\), since \(d\)
normalizes the latter subgroup, and it must normalize
its unique order-three subgroup. Directly solving the
two-by-two equations in \(\mathrm{SL}_2(\mathbb F_3)\)
shows that no determinant-one matrix conjugates \(R\)
to \(R^{-1}\), and its centralizer is \(\{\pm R^j\}\).
The centralizer modulo four is also \(\{\pm R^j\}\).
Traces exclude a negative sign in the projective
commutation equation. Thus a representative of a
normalizing element has coordinates
\(\varepsilon_4R^i\) modulo four and
\(\varepsilon_3R^j\) modulo three.
Conjugation by \(D\) fixes both signs and inverts \(R\).
Its \(d\)-norm consequently has positive signs in both
coordinates. It cannot lie in \([5I]\langle R\rangle\):
one global choice of representative has negative sign
modulo three, and the other has negative sign modulo four.
This proves the nonconjugacy and hence the asserted
nonlifting of an \(S_3\) subgroup.

For completeness, a finite subgroup mapping onto
\(S_{\mathrm{loc}}\) has trivial kernel: the kernel lies
in the torsion-free free profinite group \(H_0\).
It would therefore be an \(S_3\), already excluded.
The exact exhaustive certificate
\texttt{check\_A20\_standard\_extension.py} generates all
\(576\) elements of \(Q\), obtains normalizer and
centralizer orders \(18\), lists all nine \(d\)-norm
values and all three elements of the target coset, and
checks that the lists are disjoint. It independently
enumerates the two signed centralizers used above.
\end{proof}

A.20 blocks lifting this subgroup inside the \emph{standard}
extension. A hypothetical new involution of
\(\widehat{GT}_0\) would change the action defining the
full extension. No obstruction in that unknown twisted
extension is inferred from A.20.

\statement{A.21}{Refuted}{Real projectivity and a rank-two open free subgroup do not imply finite graph recognition}
There exists a finitely generated real-projective profinite
group \(E_*\) with an open normal subgroup isomorphic to
\(\widehat F_2\), quotient \(C_2\), and self-centralizing
involutions, which is not the profinite fundamental group
of a finite graph of finite groups. It has one conjugacy
class of involutions, so it does not satisfy R1's
three-section conclusion.

\begin{proof}
Put \(A=\prod_{p\text{ odd}}\mathbb Z_p\) and define
\[
 \Gamma_*=A\rtimes\mathbb Z_2,
 \qquad u\cdot a=(-1)^{u\bmod2}a.
\]
All groups and actions here are profinite and continuous.
The kernel of the parity map \(\Gamma_*\to C_2\) is
\(A\times2\mathbb Z_2\simeq\widehat{\mathbb Z}\).
The group \(\Gamma_*\) is torsion-free: a finite-order
element projects trivially to \(\mathbb Z_2\), and then
belongs to the torsion-free group \(A\).
Set
\[
 E_*=C_2\amalg\Gamma_*,\qquad
 \varepsilon:E_*\longrightarrow C_2,
\]
where \(\varepsilon\) is identity on the first factor
and parity on the second. In the open-subgroup Kurosh
decomposition of \(\ker\varepsilon\), the two vertex
groups are \(1\) and \(\widehat{\mathbb Z}\), and the
quotient graph has two vertices joined by two edges.
Consequently
\[
 \ker\varepsilon\simeq
 \widehat{\mathbb Z}\amalg\widehat{\mathbb Z}
 \simeq\widehat F_2.
\]
The open Kurosh theorem in precisely this full-profinite
form is recalled in Herfort--Ribes, Section~2.
Their Theorems A', B', and 1 show that all nonidentity
torsion lies in conjugates of the single \(C_2\)
factor, and that each such involution has centralizer
exactly that factor.\cite{A_HR1985}
Haran's characterization of real projectivity therefore
applies: an index-at-most-two subgroup of cohomological
dimension at most one, together with self-centralizing
involutions, suffices. We use its exact formulation in
P\'al--Quick, Theorem~2.1.\cite{A_PalQuickFormal2025}

Suppose \(E_*\) were the profinite fundamental group of
a finite graph of finite groups. Its vertex groups have
order at most two. Collapse each nonloop edge whose
nontrivial edge group is equal to an endpoint group.
In the resulting reduced graph, a loop with edge group
\(C_2\) would make its stable letter centralize that
\(C_2\); the stable letter has infinite order, as is
seen by killing the vertex groups and mapping that
loop to \(\widehat{\mathbb Z}\). This contradicts the
centralizer just proved. All remaining edges therefore
have trivial group. We obtain
\[
 E_*\simeq\bigl(\mathop{\amalg}_{j=1}^{k}C_2\bigr)
                   \amalg\widehat F_r .
\]
The single torsion class forces \(k=1\). The
torsion-free index-two subgroup of this free product
has rank \(2r+k-1\), by its two-sheeted graph or by
open Kurosh. The rank two already computed forces
\(r=1\).

On the other hand \([\Gamma_*,\Gamma_*]=A\), since
multiplication by two is invertible on \(A\).
Thus
\[
 E_*^{\mathrm{ab}}=C_2\times\mathbb Z_2,
\]
whereas \((C_2\amalg\widehat{\mathbb Z})^{\mathrm{ab}}
=C_2\times\widehat{\mathbb Z}\). The latter has
nontrivial odd-primary quotients and the former does not.
This contradiction proves the claim.
\end{proof}

The construction uses the torsion-free virtually
procyclic mechanism already present in Haran's
counterexamples to full-profinite virtual freeness
recognition.\cite{A_Haran1985}
It identifies a genuine limitation of that method;
it is not a new peripheral involution. For the groups
in A.14 the three-section hypothesis, and ultimately
the marked cusp products in A.11, still have to be used.

\statement{A.22}{Proved}{The three distinguished reflections normally generate, and a proper generated subgroup has infinite index}
Let \(\alpha\) and \(h\) be as in R1, and put
\(E=F\rtimes_\alpha\langle t\rangle\), \(b=hyh^{-1}\).
The three involutions
\[
       t_0=t,\qquad t_1=xt,\qquad t_2=b^{-1}t
\]
represent the three section classes and normally
generate \(E\). If their closed generated subgroup
\(J\) is proper, then \([E:J]\) is infinite.
Normal generation alone does not supply the full
generation required in A.11.

\begin{proof}
The equality \(f=h^{-1}\alpha(h)\) gives
\(\alpha(b)=b^{-1}\), so all three displayed
elements have order two. In
\(E^{\mathrm{ab}}\simeq C_2^3\), their images are
\(t,x+t,y+t\); these are distinct and hence the
elements are pairwise nonconjugate. R1 states that
there are exactly three classes. Quotienting \(E\)
by their closed normal closure kills \(t\), then
\(x=t_1t\), then \(b^{-1}=t_2t\). Since \(b\) is a
conjugate of \(y\), it kills \(y\), and the quotient
is trivial.

There is a continuous epimorphism
\(C_2\amalg C_2\amalg C_2\to J\) sending the three
factors to the displayed involutions. Its parity
kernel is free profinite of rank two, so \(J\cap F\)
is generated by at most two elements. If \(J\)
had finite index in \(E\), then \(J\cap F\) would
be open in \(F\) and the profinite Schreier formula
would give
\[
 d(J\cap F)=1+[F:J\cap F].
\]
It follows that \(J\cap F=F\), and \(t\in J\)
then gives \(J=E\). This proves the assertion.
\end{proof}

\statement{A.23}{Proved}{Torsion and cusp centralizers in the relative primary modular groups}
Let \(q:\widehat B=C_2\amalg C_3\to D\) be a finite
quotient injective on the two factors, with torsion-free
kernel \(U\), and write
\[
 \Gamma_p=\widehat B/O^p(U),\qquad V_p=U^{(p)}.
\]
For \(a=s\) or \(r\), of prime order \(\ell\), the
following statements hold:
\[
 \begin{array}{ll}
 p\ne\ell:&C_{\Gamma_p}(a)\longrightarrow C_D(q(a))
                    \text{ is onto},\\[2pt]
 p=\ell:&C_{\Gamma_p}(a)=\langle a\rangle.
 \end{array}
\]
For \(U=U_{12}=\ker q_{12}\) in the projective modular
quotient, put \(\Gamma=\widehat B/R_{\mathrm{nil}}(U)\).
Then
\[
 C_{\Gamma_p}(\tau)=\overline{\langle\tau\rangle},
 \qquad C_\Gamma(\tau)=\overline{\langle\tau\rangle}
                    \simeq\widehat{\mathbb Z}.
\]
In the second formula the closure is taken in \(\Gamma\);
in the first it is taken in \(\Gamma_p\).

\begin{proof}
If \(p\ne\ell\), choose a lift \(g\in\Gamma_p\)
of an element of \(C_D(q(a))\). The subgroups
\(\langle a\rangle\) and
\(\langle gag^{-1}\rangle\) are complements to
\(V_p\) in the inverse image of \(\langle q(a)\rangle\).
Coprime Schur--Zassenhaus, first in finite quotients
and then by compactness, conjugates the second to
the first by an element of \(V_p\). The generators
with image \(q(a)\) then agree. Adjusting \(g\)
by that element gives the asserted centralizer lift.

Now let \(p=\ell\), and put
\(J=q^{-1}(\langle q(a)\rangle)\subset\widehat B\).
The open Kurosh decomposition expresses \(J\) as
a finite free profinite product of copies of
\(C_p\), together with a finitely generated free
profinite factor. It has no factor of the other
prime order. Since \(J/U=C_p\),
\[
          J/O^p(U)=J^{(p)}.
\]
Indeed the left side is pro-\(p\), and every
homomorphism from \(J\) to a finite \(p\)-group
kills \(O^p(U)\). Its maximal pro-\(p\) quotient
is the corresponding free pro-\(p\) product.

If \(g\in C_{\Gamma_p}(a)\), choose
\(g_0\in\widehat B\) with the same \(D\)-image.
The elements \(a\) and \(g_0ag_0^{-1}\) belong
to \(J\) and become conjugate by \(V_p\) in
\(J^{(p)}\). Distinct Kurosh factors remain
nonconjugate in that free pro-\(p\) product,
by Herfort--Ribes Theorem B'. Their double-coset
labels in \(J\backslash\widehat B/\langle a\rangle\)
must consequently agree. Thus \(g_0\in J\),
so \(g\in J^{(p)}\). The same theorem gives
\(C_{J^{(p)}}(a)=\langle a\rangle\).
This proves the prime-primary assertion without
assuming a full-profinite Kurosh theorem for
arbitrary closed subgroups.\cite{A_HR1985}

For the cusp assertion, \(|D|=576\) and
\(q(\tau)\) has order twelve. The discrete kernel
has 48 cusps, genus 25, and free rank 97:
the cusp count is \(576/12\), and its Euler
characteristic is \(576(-1/6)=-96\).
In \(H_1(U,\mathbb F_p)\), the 48 oriented cusp
vectors have precisely the relation that their
sum is zero. In particular they are pairwise
distinct for every prime, including two and three.
The \(D\)-stabilizer of the vector represented
by \(\tau^{12}\) is exactly
\(\langle q(\tau)\rangle\): deck transformations
permute these vectors by their cusp permutation.
These assertions follow directly from the usual
oriented surface presentation of this torsion-free
finite cover of the modular orbifold.

If \(g\in\Gamma_p\) centralizes \(\tau\), it
centralizes \(\tau^{12}\in V_p\). Its image in
\(D\) therefore belongs to the indicated cusp
stabilizer. After multiplying \(g\) by an integral
power of \(\tau\), it belongs to \(V_p\).
The boundary element \(\tau^{12}\) is primitive
in \(U\), since another cusp can be omitted in
the surface free basis, and remains primitive
in \(V_p\). Herfort--Ribes Theorem B' in that
free pro-\(p\) group gives
\(C_{V_p}(\tau^{12})=\overline{\langle\tau^{12}\rangle}\).
This proves the first cusp equality. The fibre
product of these procyclic groups over
\(\langle q(\tau)\rangle=C_{12}\) is exactly
\(\widehat{\mathbb Z}\): its \(p\)-coordinate
is the infinite \(p\)-primary exponent, and the
remaining finite coordinates are the common
residue modulo twelve. This proves the second.
The same proof applies to \(x=\tau^{-2}\): its
negative sixth power is \(\tau^{12}\), so a
centralizer of \(x\) fixes the same cusp vector.
Conjugating by \(s\) treats \(y=sxs\). Thus
\(C_\Gamma(x)=\overline{\langle\tau\rangle}\)
and \(C_\Gamma(y)=s\overline{\langle\tau\rangle}s\),
as used below.
\end{proof}

\statement{A.24}{Proved}{The parity character descends to the actual outer image at level twelve}
Let
\[
 F_{12}=F/R_{\mathrm{nil}}(U_{12}),\qquad
 G_{12}=\operatorname{im}\!\left(
 \widehat{GT}_0\longrightarrow
               \operatorname{Out}^{\sharp}(F_{12})\right).
\]
The map is well-defined by the finite covariance in
refereed LM7\_PUSH E.65, which preserves \(U_{12}\).
There is a unique continuous character
\(\Delta_{12}:G_{12}\to\mathbb F_2\) such that
\(\Delta=\Delta_{12}\circ\pi_{12}\).

\begin{proof}
First consider a normalized multiplier-one action
\(\Phi\) which is identity on
\(\Gamma=\widehat B/R_{\mathrm{nil}}(U_{12})\).
Its pure \(s\)-frame, when projected to \(\Gamma_2\),
centralizes \(s\), so A.23 puts it in
\(\langle s\rangle\). Purity makes this projection
identity, since \(F\cap\langle s\rangle=1\)
already in the quotient \(S_3\). Hence its
\(Q_{12}\)-image is identity. At \(p=3\) the
same argument applies to the pure \(r\)-frame.
Moreover \(\Phi(s)=sf\) gives \(q_{12}(f)=1\).
The explicit covariance formula of LM7\_PUSH
E.65.9 then gives \(\rho_\Phi\equiv0\pmod3\),
and E.65.11 gives \(\rho_\Phi\in2\widehat{\mathbb Z}\).
Thus \(V_\Phi=I\) modulo twelve. Both parity
bits are zero. This proves that \(\Delta\)
vanishes on the kernel of the actual \(B\)-action.

We check that passing to an outer action on
\(F_{12}\) adds no kernel. Its abelianization is
\(\widehat{\mathbb Z}^2\), because
\(R_{\mathrm{nil}}(U_{12})\subset U_{12}'\subset F'\).
An outer-trivial action therefore has multiplier
one. Suppose its restriction is conjugation by
\(u\in F_{12}\). The centralizer of \(F_{12}\)
in \(\Gamma\) is trivial. To see this, the cusp
argument in A.23 also gives
\(C_\Gamma(x)=\overline{\langle\tau\rangle}\)
and
\(C_\Gamma(y)=s\overline{\langle\tau\rangle}s\).
Their images in \(S_3\) are distinct order-two
subgroups; an element of their intersection lies
in \(F_{12}\), and then in both
\(\overline{\langle x\rangle}\) and
\(\overline{\langle y\rangle}\). Abelianization
shows that their intersection is trivial.

An automorphism of \(\Gamma\) fixing the open
normal subgroup \(F_{12}\) pointwise is consequently
identity: for each \(g\in\Gamma\), its discrepancy
\(\Phi(g)g^{-1}\) centralizes that subgroup.
Thus the given \(\Phi\) is conjugation by \(u\)
on all of \(\Gamma\). It fixes \(\tau\), so A.23
gives \(u\in\overline{\langle\tau\rangle}\cap F_{12}
=\overline{\langle x\rangle}\). Write \(u=x^a\).
The normalized formula \(\Phi(s)=sf\) then gives
\(f=y^a x^{-a}\). Since \(f\in F_{12}'\), its
abelian coordinates \((-a,a)\) vanish; hence
\(a=0\). The actual action was identity after all.

We have proved kernel containment for \(\Delta\).
It therefore factors through \(G_{12}\); continuity
follows because a continuous surjection of compact
groups to a Hausdorff group is a quotient map.
\end{proof}

\statement{A.25}{Proved}{An actual parity-one involution exists at the relative level twelve}
The group \(G_{12}\) of A.24 contains an involution
\(a_{12}\) with
\[
             \lambda(a_{12})=-1,
             \qquad\Delta_{12}(a_{12})=1.
\]
It is not conjugate to the image of \(c\) in
\(G_{12}\). Thus the first, specified-relative-level
assertion of Lead~3 holds at \(U_{12}\).
The existence of an involution lifting it in the
full group is not asserted.

\begin{proof}
Let \(r_0=(3/4)^{1/4}>0\). The splitting field
\(L=\mathbb Q(i,r_0)\) has degree eight: the
polynomial \(4X^4-3\) is Eisenstein at three,
and \(\mathbb Q(r_0)\) is real. Choose a Galois
element \(\gamma\) whose restriction to \(L\)
fixes \(i\) and sends \(r_0\) to \(ir_0\).
Such a restriction is an automorphism of \(L\)
and extends to \(\overline{\mathbb Q}\).
Use its actual Galois image in \(\widehat{GT}\),
and put
\[
 \sigma_0=\gamma c\gamma^{-1}c,
 \qquad \sigma=\sigma_0^3.
\]
Then \(\lambda(\sigma_0)=1\) and
\(c\sigma_0c=\sigma_0^{-1}\), so
\(\sigma c\) is an actual full involution.
Also \(\sigma_0(r_0)=-r_0\). The Galois comparison
in refereed LM7\_PUSH E.65.13, including the model
and sign calculation supplied by its referee,
makes both normalized CM bits of \(\sigma_0\) equal
to one. More precisely these are
\(\varepsilon_3(T_{\sigma_0})\) and
\(\varepsilon_2(W_{\sigma_0})\), with the shear
\(V\) included in the definitions of \(T,W\).
LM7\_PUSH E.65.12 proves that \(T,W\) are
homomorphisms on \(\mathcal H\); composing with
the two indicated characters proves that both
bits are characters on all of \(\mathcal H\).
The same bits therefore hold for \(\sigma\).

The cubing calculation in the last paragraph of
that refereed proof gives
\[
 q\Phi_\sigma=q,
 \qquad q_{12}(f_\sigma)=I,
 \qquad V_\sigma=I.
\]
Let \(A,\beta\in F\) be its actual pure frames,
and write \(a=q(A),\ b=q(\beta)\). The r-bit
and s-bit are both one, while \(a\) and \(b\)
centralize \(q(r)\) and \(q(s)\), respectively.
More precisely, cubing the two finite cyclic
characters gives
\[
 a=[5I],\qquad
 b=[B]\ \text{or}\ [B^{-1}],\qquad
 B=\begin{pmatrix}9&4\\8&9\end{pmatrix}\pmod{12}.
\]
Here \(B\equiv I\pmod4\), \(B\equiv S\pmod3\),
and \(B^2=5I\). The choice between the two
possibilities will not matter.

On \(\Gamma_p=\widehat B/O^p(U_{12})\), define
\[
 \phi_2=\Phi_\sigma,\qquad
 \phi_p=\operatorname{id}\quad(p\ne2).
\]
All induce identity on the shared finite quotient
\(Q_{12}\). They therefore define an actual
continuous automorphism \(\phi\) of their fibre
product \(\Gamma\). Each fixes \(\tau\).
The automorphism \(\phi c\) has square one:
at two it is \(\Phi_\sigma c\), and at every
other prime it is \(c\).

We now verify full peripheral gluing with literal
conjugators, the essential point of the construction.
At two, A.23 supplies \(k_2\in C_{\Gamma_2}(r)\)
with \(q(k_2)=a^{-1}\). Replace the r-frame
\(A_2\) by \(A_2k_2\). Its \(Q_{12}\)-image
is one. At every other prime take the identity
r-frame. These give a single pure r-frame of
\(\phi\), with \(Q_{12}\)-image one.

At two retain the s-frame \(\beta_2\), with
\(Q_{12}\)-image \(b\). For every odd prime,
A.23 supplies \(k_p\in C_{\Gamma_p}(s)\)
with \(q(k_p)=b\). These are frames for the
identity automorphism at those primes. Their
common image \(b\in q(F)\) makes the resulting
global s-frame pure. The \(\tau\)-frame is one
at every prime. Thus \(\phi\) preserves all
three conjugacy classes \(s,r,\tau\), with
the stated common quotient frames.

Apply R2 with these literal conjugators. It gives
an automorphism of \(\widehat B\) fixing \(\tau\)
and lifting \(\phi\), with pure frames whose
\(Q_{12}\)-images are \((1,b)\) in the order
\((r,s)\). This lift belongs to \(\mathcal H\).
Indeed write it as \(s\mapsto sf\),
\(r\mapsto f^{-1}r\). Its reduction has
\(f_2=f_\sigma\) and \(f_p=1\) at odd primes,
so its image in every primary abelianization of
\(F\) is zero. Since
\(R_{\mathrm{nil}}(U_{12})\subset F'\), the
full \(f\) lies in \(F'\). The two torsion
relations are exactly the gentle relations, and
R2 supplies invertibility.

Denote this full multiplier-one lift by
\(\widetilde\sigma\). Its action on \(Q_{12}\)
is identity and its \(f\)-image there is one;
the covariance formula makes \(V=I\). Its
r-bit is zero and its s-bit is one. Hence
\(\Delta(\widetilde\sigma)=1\). The outer
image
\[
                  a_{12}=\pi_{12}(\widetilde\sigma c)
\]
is precisely the involution already constructed.
Its action on \(F_{12}^{\mathrm{ab}}\) is \(-I\),
so it has order exactly two. A.24 gives
\(\Delta_{12}(a_{12})=1\), whereas
\(\Delta_{12}(\pi_{12}(c))=0\). The character
separates their conjugacy classes.
\end{proof}

The exact finite computations underlying the construction
are retained in \path{check_A25_level.py} and its output:
the quotient has order 576, the pure subgroup has order 96,
the cusp orbit has size 48, and the frame matrix has
square \([5I]\). The script also verifies the dihedral
action on the four roots used above. The centralizer
and lifting assertions are proved in the text, not
inferred from this finite computation.

Unlike A.17, A.25 is an actual involution in an
outer image of the relative level. It uses the
different primes where torsion centralizers acquire
extra lifts, and retains their common finite quotient
frames throughout. Let
\(L_{12}=\ker(\widehat{GT}_0\to G_{12})\) and
\(a=\widetilde\sigma c\). Then
\(a^2\in L_{12}\) and \(\Delta(L_{12})=0\).
A full involution above this specified class exists
if and only if
\[
          v(ava^{-1})=a^{-2}
          \quad\text{has a solution }v\in L_{12}.
\]
A.19 describes every abelian refinement of this
remaining square obstruction. No solution or
incompatible finite refinement has yet been proved.

\statement{A.26}{Proved}{Finite frame visibility prevents arithmetic-only continuation}
Put
\[
 U=U_{12},\qquad V=\overline{U'U^6},\qquad D_6=\widehat B/V.
\]
This is a finite quotient of order \(576\cdot6^{97}\),
and \(V\) is characteristic in \(U\), normal in
\(\widehat B\), and preserved by \(\widehat{GT}_0\).
The two normalized CM bits of an element of
\(\mathcal H\) are determined by its actual action
on \(D_6\). In particular both bits are zero for
an element acting identically on \(D_6\).

Let \(H_6\) be the actual image of \(\mathcal H\)
in \(\operatorname{Aut}(D_6)\), and
\(J_6=\langle H_6,c_6\rangle\). There is a
character \(\Delta_6:J_6\to\mathbb F_2\) induced
by \(\Delta\) on
\(\mathcal H\rtimes\langle c\rangle\).
The involution of A.25 induces an actual involution
\(a_6\in J_6\) with \(\Delta_6(a_6)=1\).
No Galois element of multiplier \(-1\) has the
same actual action on \(D_6\).

\begin{proof}
For \(p=2,3\), write
\(D_p=\widehat B/\overline{U'U^p}\).
If \(a=s\) when \(p=2\), or \(a=r\) when
\(p=3\), then
\[
  \operatorname{im}\bigl(C_{D_p}(a)\to Q_{12}\bigr)
                      =\langle q(a)\rangle.
\]
To prove this, let \(g\) be a centralizing element
and choose a lift \(g_0\in\widehat B\) of its
\(Q_{12}\)-image. In
\(J=q^{-1}(\langle q(a)\rangle)\), the elements
\(a\) and \(g_0ag_0^{-1}\) become conjugate
by an element of \(U/\overline{U'U^p}\).
If their Kurosh factor labels were different,
the free-product universal property would give
a continuous character \(J\to C_p\) sending
one to one and the other to zero. Every such
character kills \(\overline{U'U^p}\), and
characters are constant on conjugacy classes.
This is a contradiction. Therefore the labels
coincide and \(q(g_0)\in\langle q(a)\rangle\),
as claimed. The reverse inclusion is immediate.

Any two pure s-frames for the same actual action
on \(D_6\) differ by a centralizer element.
The displayed assertion at two, followed by
intersection with the pure subgroup, shows that
their \(Q_{12}\)-images agree. The assertion at
three does the same for r-frames. The shear
\(V_\Phi\) modulo twelve is determined by the
\(Q_{12}\)-image of \(f\), by the explicit
covariance formula and the evenness of \(\rho\).
Thus both normalized bits are determined.

The kernel of \(\mathcal H\to H_6\) therefore
lies in \(\ker\delta\), giving a character
\(\delta_6\) on \(H_6\). It is invariant under
\(c_6\) because \(\delta\) is c-invariant.
The element \(c_6\) is not in \(H_6\):
on the retained abelian quotient
\(\widehat B^{\mathrm{ab}}=C_6\), its action
is inversion, whereas \(H_6\) acts identically.
Hence \(J_6=H_6\rtimes\langle c_6\rangle\),
and
\(\Delta_6(hc_6^e)=\delta_6(h)\) is well-defined.

The group \(U/R_{\mathrm{nil}}(U)\) has
characteristic subgroup whose quotient is
\(U/\overline{U'U^6}\). The actual involution
\(\phi c\) in A.25 therefore induces an actual
involution on \(D_6\). It belongs to \(J_6\),
has \(\Delta_6=1\), and is nonidentity since
it inverts the abelian quotient \(C_6\).
If a Galois element of multiplier \(-1\)
induced this action, it would belong to
\(\mathcal H\rtimes\langle c\rangle\) and
its \(\Delta\)-value would be one. This
contradicts R3's Galois vanishing.
\end{proof}

This is an obstruction to continuing the
\emph{arithmetic-only construction}, not an
obstruction to lifting \(a_{12}\) at all.
At any relative level below \(V\), the common
finite action must be \(a_6\). Thus no primary
component of such a continuation can be obtained
from a full Galois element of multiplier \(-1\).
In particular one cannot choose identity-on-
\(D_6\) Galois c-cocycles with the two CM bits
still equal to one. Nongalois relative lifts,
or a solution of the norm equation in
\(L_{12}\), remain possible. This conclusion
does not identify the entire character
\(\Delta\) on the image of arbitrary multipliers
in \(\operatorname{Aut}(D_6)\); the precise
finite domain here is \(J_6\).

\statement{A.27}{Proved}{The exact finite real-reflection count and its properness condition}
Let \(\pi:E\twoheadrightarrow D\) be an epimorphism of
finite groups. Fix involutions \(d_0,d_1,d_2\in D\)
which generate \(D\), and subsets
\(T_{ij}\subseteq\pi^{-1}(d_id_j)\), for
\((i,j)=(0,1),(1,2),(2,0)\). For every subgroup
\(J\leq E\), put
\[
 I_i(J)=\{u\in J:\pi(u)=d_i,\ u^2=1\},
 \qquad
 M_{ij}(J)_{u,v}=\mathbf1_{T_{ij}}(uv).
\]
Thus \(M_{ij}(J)\) is a rectangular matrix with rows
indexed by \(I_i(J)\) and columns indexed by
\(I_j(J)\). The number of triples of involutions
in \(E\) with the specified images and pair products,
and which generate \(E\), is
\[
 \sum_{J\leq E}\mu(J,E)
       \operatorname{tr}\bigl(M_{01}(J)M_{12}(J)M_{20}(J)\bigr),
 \tag{A.27.1}
\]
where \(\mu\) is the ordinary subgroup-lattice
M\"obius function. Empty index sets contribute zero.
This count concerns actual homomorphisms from
\(C_2*C_2*C_2\), with actual quotient data.

\begin{proof}
Expanding the trace gives exactly one summand for
\((u_0,u_1,u_2)\in I_0(J)\times I_1(J)\times I_2(J)\),
and that summand is one exactly when all three
pair-product conditions hold. Their ordered product
is automatically one:
\((u_0u_1)(u_1u_2)(u_2u_0)=1\).
Every counted triple has a unique generated subgroup.
M\"obius inversion on the finite subgroup lattice
therefore gives (A.27.1). The assumption that the
\(d_i\) generate makes every subgroup not mapping
onto \(D\) contribute zero.
\end{proof}

The matrices in (A.27.1) retain the common reflection
at the ends of adjacent edges. The ordinary
boundary-class trace sums of R2 do not retain this
incidence datum. Equality of those ordinary sums
therefore does not by itself establish equality of
(A.27.1). A concrete perfect-kernel control is
\(S_5\twoheadrightarrow C_2\), with kernel \(A_5\),
and all three \(d_i\) nonidentity. Every \(d_i\)
has an involution lift, and \(S_5\) has three
ordinary generators, but three such lifts cannot
generate \(S_5\): they are transpositions, whose
three edges cannot connect a graph on five vertices.
The certificate \path{check_A27_real_count.py}
checks perfectness of the kernel and enumerates
all 1000 triples. The generated-subgroup orders
and multiplicities are
\[
 (2,10),\quad(4,90),\quad(6,240),\quad
 (12,180),\quad(24,480).
\]
This is a control on a proposed general real
embedding theorem, not a counterexample to the
negative pants problem. Indeed a product of two
odd involutions in \(S_5\) has order at most three;
no generating pants triple of \(A_5\) has all
three orders at most three. The same certificate
checks all 816 candidate ordered pairs and finds
no pair generating \(A_5\).

\statement{A.28}{Proved}{A negative finite peripheral action is stopped by a central spin obstruction}
There is an epimorphism \(q:F\twoheadrightarrow A_6\)
and an involutory automorphism \(\vartheta\) of
\(A_6\) which sends each of the three labelled
cusp images to an \(A_6\)-conjugate of its inverse,
but no multiplier-minus-one peripheral automorphism
of \(F\) satisfies
\[
                       q\Phi=\vartheta q.
\]
In particular no full negative peripheral involution
induces this finite action. The obstruction is
already visible in a central double cover; it is
not a counterexample to R2 or A.3.

\begin{proof}
We use rightmost-first permutation composition and put
\[
\begin{aligned}
 X_0&=(1\,2\,3\,4)(5\,6),&
 Y_0&=(1\,2\,4\,6)(3\,5),\\
 Z_0&=(1\,2\,6\,3)(4\,5),&
 t&=(2\,4).
\end{aligned}
\]
The certificate \path{check_A28_s6_passport.py}
verifies \(X_0Y_0Z_0=1\) and
\(\langle X_0,Y_0\rangle=A_6\), giving
\(q(x)=X_0,q(y)=Y_0,q(z)=Z_0\).
Let \(\vartheta=\operatorname{Ad}(t)\).
The required inverse-class conjugators are,
in the displayed cusp order,
\[
                    1,\qquad(2\,6\,4),\qquad(2\,3\,4).
\]
In particular these conjugators are even.
Each cusp also has an odd involution which
inverts it: respectively \((2\,4),(2\,6),(2\,3)\).
The permutation \((3\,4\,6)\), of order three
in \(A_6\), cyclically permutes the three
cusp images by conjugation.

For completeness there is a short section-incidence
explanation of why an order-two full lift is impossible.
Odd involutions of \(S_6\) have just two cycle types,
\((2)\) and \((2,2,2)\). The product of two of
one type cannot have cycle type \((4,2)\).
For transpositions its order is at most three.
For two fixed-point-free involutions, their two
perfect matchings decompose into alternating cycles;
on six letters the possible product orders are
again at most three. Thus every \((4,2)\) cusp
requires opposite types at its two reflection ends.

In a hypothetical full negative pants extension,
R1 gives exactly three section classes. Their
images in
\(E^{\mathrm{ab}}=C_2^3\) are
\(\bar t,\bar x+\bar t,\bar y+\bar t\).
The x-cusp joins the first two, the y-cusp the
first and third, and the z-cusp the last two:
its two ends differ by \(\bar x+\bar y\).
Here each cusp has genuine involution ends.
Indeed an element of the odd coset which inverts
a full cusp has square in its procyclic centralizer;
inversion fixes that square, so the square is one.
Their three finite class labels would therefore
colour a triangle with two colours, with opposite
colours on each edge. This is impossible. No
assumption that the three reflections generate
is made in this argument.

A stronger obstruction excludes even infinite-order
full peripheral lifts of multiplier \(-1\).
We give a finite certificate rather than assume
that an arbitrary finite peripheral automorphism
lifts. The two degree-80 permutations in the
ATLAS representation \texttt{2S6G1-p80B0}
are opened source data. The code
\path{check_A28_double_cover.py} enumerates their
joint group with the natural standard generators
\((1\,2),(2\,3\,4\,5\,6)\) of \(S_6\).
It verifies an epimorphism
\[
 \widetilde S\twoheadrightarrow S_6,
 \qquad |\widetilde S|=1440,
 \qquad \ker=\langle\epsilon\rangle\cong C_2,
\]
with central kernel. Its even preimage
\(\widetilde A\) has order 720. Choose arbitrary
lifts \(X,Y\in\widetilde A\) of \(X_0,Y_0\),
and put \(Z=(XY)^{-1}\). Choose an involution
\(T\in\widetilde S\) above \(t\).
The enumeration of all conjugators in
\(\widetilde A\) proves, for each
\(Q\in\{X,Y,Z\}\),
\[
 TQT^{-1}\not\sim_{\widetilde A}Q^{-1},
 \qquad
 TQT^{-1}\sim_{\widetilde A}\epsilon Q^{-1}.
 \tag{A.28.1}
\]
Each lift has order eight. More precisely the
unit exponents \(k\pmod8\) for which
\(TQT^{-1}\sim Q^k\) are exactly \(3,5\).
The same statements hold for either central
choice of each lift. All 1440 elements, the
central kernel, the 720 even conjugators and
the stated class tests are checked by the code;
the group order is not inferred solely from
the ATLAS label.\cite{AAtlasA6}

Let \(q_2:F\to\widetilde A\) send \(x,y\)
to \(X,Y\). If \(\Phi\) as in the statement
existed, the two homomorphisms
\(q_2\Phi\) and \(\operatorname{Ad}(T)q_2\)
would have the same projection to \(A_6\).
Their ratio is consequently a continuous central
character \(\eta:F\to\langle\epsilon\rangle\).
The multiplier-minus-one cusp conditions and
(A.28.1) require
\[
             \eta(x)=\eta(y)=\eta(z)=\epsilon.
\]
But \(xyz=1\), contradicting
\(\eta(x)\eta(y)\eta(z)=1\).
This proves the stronger assertion.
\end{proof}

The kernel of \(q\) is open normal; it is not
asserted to be characteristic. Any characteristic
refinement retaining this quotient and this action
inherits the displayed obstruction. Conversely,
we have not constructed an automorphism of such a
characteristic refinement or a relative pronilpotent
level inducing \(\vartheta\). The central calculation
explains why the finite inverse-class data alone
cannot be inserted into R2. The dyadic class
obstruction, rather than perfectness alone, is
what is retained here.

\statement{A.29}{Proved}{The new relative involution has exactly three section classes}
For the involution \(a_{12}\) constructed in A.25,
choose its actual representative \(\phi c\) on
\(F_{12}=F/R_{\mathrm{nil}}(U_{12})\).
Then
\[
                \bigl|H^1(C_2,F_{12})\bigr|=3,
\]
where the action is this specified representative.
Its relative dyadic component likewise has exactly
three section classes and has trivial fixed subgroup
in its orientation kernel. These are section
classes of the group extension, not conjugacy
classes of peripheral automorphisms.

\begin{proof}
First let
\(W=C_2\amalg C_2\amalg C_2=F\rtimes\langle c\rangle\),
and let \(U\) be any open normal subgroup of
\(W\) contained in \(F\). In
\(W/O^2(U)\), every nonidentity torsion element
of order two has nonidentity image in \(W/U\),
because \(U^{(2)}\) is torsion-free. If that
image generates \(C\cong C_2\), let
\(J\leq W\) be its inverse image. The open
Kurosh theorem expresses \(J\) as a free
profinite product of order-two factors and a
free profinite group. Since \(J/U=C_2\),
\[
                     J/O^2(U)=J^{(2)}.
\]
Indeed the left side is pro-2, and every finite
2-group quotient of \(J\) kills \(O^2(U)\).
The torsion theorem for the resulting free pro-2
product shows that the given involution is conjugate
to an image of a full factor of \(W\).
The three full factors remain nonconjugate,
because \(W^{\mathrm{ab}}=C_2^3\) survives this
quotient. Conjugacy under the even subgroup
is the same as full conjugacy for odd involutions:
one may multiply an odd conjugator by its
source involution to make it even.

The same open-product argument proves that every
such torsion centralizer is exactly the cyclic
subgroup it generates. Explicitly, if a centralizing
element has finite image \(g\), conjugating its
factor by a lift of \(g\) gives two factors of
\(J\) that become conjugate in \(J^{(2)}\).
The free pro-2 torsion theorem forces equality
of their Kurosh labels, so \(g\in C\).
The centralizer computation now takes place inside
\(J^{(2)}\), where the free-factor centralizer
is that factor. Its intersection with the even
subgroup is trivial. This uses the free-factor
case of Herfort--Ribes, not a closed-subgroup
Kurosh assertion.\cite{A_HR1985}

Apply this to \(U=U_{12}\), which is preserved
by the actual Galois action used in A.25.
The dyadic involution \(\Phi_\sigma c\) is
Galois-conjugate to \(c\), so its section set
has the same three classes and trivial even
centralizers. The choice of conjugating Galois
element preserves \(U_{12}\), by R3's quotient
covariance.

At each odd prime the kernel above the common
finite quotient is pro-\(p\). Every involution
in the specified quotient coset has an involution
lift, and any two such lifts with the same finite
image are conjugate by that kernel, by
Schur--Zassenhaus. Moreover the centralizer of
an involution surjects onto the centralizer of
its finite image: conjugate a lift by any chosen
finite centralizer lift, then apply the same
complement-conjugacy assertion in the kernel.

It follows that projection to the dyadic component
induces a bijection on section-conjugacy classes
for the complete fibre product. Surjectivity
follows by choosing odd-prime lifts with the
same finite image. For injectivity, choose a
dyadic conjugator and its finite image; at every
odd prime use centralizer surjectivity to adjust
a conjugator to this very same finite image.
The resulting tuple is an element of the fibre
product and conjugates the two sections. This
proves the asserted cardinality three.
\end{proof}

The unmarked pro-2 recognition theorem is fully
consistent with A.29. Herfort--Zalesskii's
Theorem~6.2 conjugates a finite 2-subgroup of
\(\operatorname{Aut}(F^{(2)})\) into the
automorphisms of a dense abstract free group.
It does not retain three prescribed peripheral
classes.\cite{AHerfortZalesskii2013}
Thus neither that theorem nor the section count
proves the peripheral conjugacy statement left
open in A.12. The actual involution of A.25
already passes the three-section test.

\statement{A.30}{Proved}{The level-twelve frames can also satisfy the involution relations}
The construction in A.25 can be made with its literal
relative frames \(A,\beta\in\Gamma\) so that the
associated r-normalizer and s-normalizer in
\(\Gamma\rtimes\langle\phi c\rangle\) are actual
subgroups \(S_3\) and \(C_2\times C_2\), with the
finite frame values still
\[
                         q(A)=1,\qquad q(\beta)=b.
\]
Consequently an R2 lift \(a=\widetilde\sigma c\)
can be chosen with
\[
 a^2\in L_{12}^{\mathrm{fr}},\qquad
 L_{12}^{\mathrm{fr}}
 =\{k\in\mathcal H:\Phi_k|_\Gamma=1,
                  \ A_k|_\Gamma=\beta_k|_\Gamma=1\}.
 \tag{A.30.1}
\]
This is a smaller kernel than the unframed
\(L_{12}\). It does not assert a lift at a
smaller characteristic level or solve the remaining
norm equation.

\begin{proof}
Write \(a_2=\Phi_\sigma c\) for the genuine
Galois involution used at the dyadic component,
and write \(A_2,\beta_2\) for its underlying
multiplier-one element's frames. Put
\[
 J_2=\Phi_\sigma(s)A_2,
 \qquad D_2=J_2^{-1}a_2,
 \qquad T_2=\beta_2^{-1}a_2.
\]
The element \(D_2\) inverts \(r\), because
\(c(r)=sr^{-1}s\). In the full Galois extension,
its square lies in \(C_{\widehat B}(r)=\langle r\rangle\)
and is fixed by inversion; hence \(D_2^2=1\).
Similarly \(T_2\) centralizes \(s\) and its
square lies in \(\langle s\rangle\). That square
cannot be \(s\): the abelianization of
\(\widehat B\rtimes\langle\Phi_\sigma c\rangle\)
is \(C_2^2\), and \(s\) has nonzero image.
Thus \(T_2^2=1\). These equations descend to
\(\Gamma_2\rtimes\langle a_2\rangle\).

In the common finite extended quotient, put
\(d=sc\) and \(z_0=[5I]\). Since
\(q\Phi_\sigma=q\), the frame values give
\[
                       q(D_2)=z_0d,
                       \qquad q(T_2)=b^{-1}c.
\]
Choose a pure \(h\in\Gamma_2\) with finite
image \([B]\), where \(B\) is the matrix of
A.25. Direct multiplication gives
\[
             d(B)=B^{-1},\qquad B^2=5I.
\]
It follows that the involution \(hD_2h^{-1}\)
has finite image \(d\). The inverse image of
\(\langle q(r),d\rangle\cong S_3\) in this
relative extension has pro-2 kernel. A.16 supplies
a section retaining this exact involution. Its
order-three subgroup is conjugate to \(\langle r\rangle\)
by an element of \(U_{12}^{(2)}\), by coprime
complement conjugacy. After this conjugation the
section has an involution \(D'_2\) satisfying
\[
        D'_2r(D'_2)^{-1}=r^{-1},\qquad q(D'_2)=d.
\]
Then \(k_2=D_2D'_2\) centralizes \(r\), belongs
to \(\Gamma_2\), and has finite image \(z_0\).
Replace \(A_2\) by \(A'_2=A_2k_2\).
The r-frame now has finite image one, and its
associated r-normalizing reflection is \(D'_2\).
At every odd prime retain the identity r-frame
and the standard r-normalizing reflection \(sc\).
Thus the r-frames glue, and their normalizers
glue to an actual finite \(S_3\).

At two retain \(\beta_2\) and \(T_2\).
For an odd prime \(p\), A.23 gives a surjection
\[
             C_{\Gamma_p}(s)\longrightarrow C_{Q_{12}}(q(s))
\]
with pro-\(p\) kernel, equivariant for \(c\).
The element \(b^{-1}c\) has order two in the
finite semidirect product, since \(c(b)=b^{-1}\).
Schur--Zassenhaus supplies an involution above it
in \(C_{\Gamma_p}(s)\rtimes\langle c\rangle\).
Write this involution as \(\beta_p^{-1}c\).
Then \(q(\beta_p)=b\), and \(\beta_p\) is a
frame for the identity action at that prime.
These frames glue with \(\beta_2\); their
involutions glue to \(T=\beta^{-1}(\phi c)\),
which centralizes \(s\). Since \(T\) is outside
\(\Gamma\), it and \(s\) generate \(C_2^2\).
The unchanged \(\tau\)-frame is one.

Apply R2 to these actual frames, and call the full
lift \(a=\widetilde\sigma c\). It still has
\(\Delta(a)=1\) and the same relative action as
A.25. For its s-frame, applying \(a\) twice
gives the pure conjugator \(a(\beta)\beta\).
Its relative image is one, exactly because
\((\beta^{-1}(\phi c))^2=1\).
For the r-frame put \(J=\Phi_{\widetilde\sigma}(s)A\).
Applying \(a\) twice gives the pure conjugator
\(a(J)J\), whose relative image is one because
\((J^{-1}(\phi c))^2=1\). Uniqueness of the full
pure frames identifies these with the frames of
\(a^2\). Together with its identity relative
action, this proves (A.30.1).
\end{proof}

The kernel \(L_{12}^{\mathrm{fr}}\) is invariant
under conjugation by this chosen \(a\). One may
verify this by composing actual peripheral actions
and their literal conjugators: an identity action
with both conjugators equal to one remains such
under framed conjugation. In the specified
\emph{framed} fibre, an order-two full lift exists
exactly when
\[
             v(ava^{-1})=a^{-2}
             \quad\text{for some }v\in L_{12}^{\mathrm{fr}}.
 \tag{A.30.2}
\]
Solving (A.30.2) would solve Lead~3. Failure to
solve it would exclude only this framed fibre;
a full involution over the unframed \(a_{12}\)
could still use different relative frames, and
is governed by A.25's equation in \(L_{12}\).
Neither equation is declared solvable or insoluble.

\statement{A.31}{Proved}{Parity-one involutions are compatible through the entire relative principal-congruence tower}
For every positive integer \(N\) divisible by twelve, put
\[
 \begin{gathered}
 Q_N=\operatorname{SL}_2(\mathbb Z/N)/\{\pm I\},\qquad
 U_N=\ker(\widehat B\longrightarrow Q_N),\\
 \Gamma_N=\widehat B/R_{\mathrm{nil}}(U_N),\qquad
 F_N=F/R_{\mathrm{nil}}(U_N).
 \end{gathered}
\]
These kernels are preserved by the actual normalized
\(\widehat{GT}_0\)-action. Let \(G_N\) be its actual outer
image on \(F_N\). There is a single full element
\(a\in\widehat{GT}_0\) such that
\[
 \lambda(a)=-1,\qquad \Delta(a)=1,\qquad
 a_N^2=1\quad\hbox{for every }12\mid N,
 \qquad a_N=\pi_N(a)\in G_N.
 \tag{A.31.1}
\]
The \(a_N\) are compatible and each has order exactly two.
The character \(\Delta\) descends to every \(G_N\), and
\(\Delta_N(a_N)=1\). In particular their inverse limit
is an involution of parity one in the image of the full
group on this entire tower.

The assertion is restricted to this tower. The full
element \(a\) is not proved to have order two, and the
relative principal-congruence tower is not cofinal in
the tower in A.2.

\begin{proof}
We first choose the arithmetic cocycle more carefully
than in A.25. Fix an embedding
\(\overline{\mathbb Q}\hookrightarrow\overline{\mathbb Q}_3\)
and write \(r_0=(3/4)^{1/4}>0\). The field
\(\mathbb Q_3(i)\) is unramified quadratic over
\(\mathbb Q_3\); over it the polynomial
\(4X^4-3\) is Eisenstein. Since it contains the fourth
roots of unity, adjoining \(r_0\) gives a cyclic,
totally and tamely ramified extension of degree four.
Its inertia group therefore contains the automorphism
\(r_0\mapsto ir_0\), fixing \(i\). Choose a lift
\(\gamma\) in the absolute inertia subgroup at three.

Every splitting field obtained by adjoining roots of
unity of 2-power order and a 2-power root of two is
unramified at three: the discriminant of
\(X^{2^m}-2\) is a power of two, and the same holds
for the relevant cyclotomic extensions. Thus
\(\gamma\) fixes all these roots. Let
\[
 \sigma_0=\gamma c\gamma^{-1}c.
\]
It has multiplier one, satisfies
\(c\sigma_0c=\sigma_0^{-1}\), sends \(r_0\) to
\(-r_0\), and has zero dyadic Kummer character of two.
The last statement also follows directly because both
\(\gamma\) and \(c\) fix each positive real
\(2^{1/2^m}\).

Let \(e_2\in\widehat{\mathbb Z}\) be the idempotent
whose 2-adic component is one and whose odd-primary
components are zero. Define, in the procyclic closure
of \(\sigma_0\),
\[
                         \sigma=\sigma_0^{e_2}.
\]
The complete Kummer character \(\kappa_2\) is additive
on the multiplier-one Galois subgroup. Hence
\[
 \kappa_2(\sigma)=e_2\kappa_2(\sigma_0)=0,
 \qquad c\sigma c=\sigma^{-1},
 \qquad \sigma(r_0)=-r_0.
 \tag{A.31.2}
\]
The Galois comparison in R3 again makes both normalized
CM bits equal to one. Moreover \(\sigma c\) is a
Galois conjugate of \(c\): if
\(b=(e_2-1)/2\in\widehat{\mathbb Z}\), then
\[
 \sigma c=(\sigma_0^b\gamma)c(\sigma_0^b\gamma)^{-1}.
\]

Marin's modular representation calculation identifies
its congruence shear on a Galois element with
\(-8\kappa_2\), prime by prime
\cite[\S4, p.~776]{AMarin2008}.
Consequently
\[
                         q_N\Phi_\sigma=q_N
                         \quad\text{for every }N.
 \tag{A.31.3}
\]
Here the sign or left-versus-right normalization of
the shear is irrelevant to its vanishing. The full
covariance formula in refereed E.65.15 proves that
every \(U_N\) is preserved by
\(\widehat{GT}_0\). We do not assert that \(U_N\)
is characteristic under all automorphisms of \(F\).

Write \(A_\sigma,\beta_\sigma\) for the full pure
r- and s-frames of \(\sigma\), and put
\(a_N'=q_N(A_\sigma)\), \(b_N=q_N(\beta_\sigma)\).
At twelve, the signed r-character takes values in
\(C_6\), and the signed s-character in \(C_4\).
Since \(e_2\equiv3\pmod6\) and
\(e_2\equiv1\pmod4\), (A.31.2)--(A.31.3) give
\[
 a_{12}'=[5I],\qquad b_{12}=[B]^{\pm1},\qquad
 B=\begin{pmatrix}9&4\\8&9\end{pmatrix}\pmod{12}.
 \tag{A.31.4}
\]

For fixed \(N\), use the fibre-product decomposition
\[
 \Gamma_N=\mathop{\prod}_{Q_N,p}
          \Gamma_{N,p},\qquad
 \Gamma_{N,p}=\widehat B/O^p(U_N).
\]
Define the actual automorphism \(\phi_N\) by
\(\Phi_\sigma\) on its dyadic component and identity
on every odd-primary component. Equation (A.31.3)
gives the common identity action on \(Q_N\), so
these actions glue; \((\phi_Nc)^2=1\).
They are compatible as \(N\) varies, since the same
full \(\sigma\) is used at each level.

The literal frames glue as well. At two, the
coprime-centralizer assertion of A.23 gives
\(k_2\in C_{\Gamma_{N,2}}(r)\) of finite image
\((a_N')^{-1}\). Replace the dyadic r-frame by
\(A_{\sigma,2}k_2\); its finite image is one.
Use the identity r-frame at every odd prime.
For s, retain \(\beta_{\sigma,2}\) at two.
At every odd prime A.23 supplies an s-centralizing
frame with finite image \(b_N\). All these frames
are pure because their common finite images are
in \(q_N(F)\) and \(U_N\subset F\).
The \(\tau\)-frame is one. R2 therefore supplies
a full lift in \(\mathcal H\), just as in A.25.
The commutator normalization follows because its
\(f\)-image has zero abelianization at every prime
and \(R_{\mathrm{nil}}(U_N)\subset F'\).

For compatibility of the full lifts, define the closed
subsets of the compact group \(\mathcal H\)
\[
 \mathcal C_N=\{h:\ \Phi_h|_{\Gamma_N}=\phi_N,
       \ q_N(A_h)=1,\ q_N(\beta_h)=b_N\}.
\]
The preceding construction proves each nonempty.
For any finite set of moduli, a common multiple
\(M\) satisfies \(\mathcal C_M\subset\bigcap_N
\mathcal C_N\). Compactness gives
\(h\in\bigcap_{12\mid N}\mathcal C_N\).
This argument requires no surjectivity assertion
for transition maps between centralizers.
Set \(a=hc\). Its action on every \(\Gamma_N\)
is \(\phi_Nc\), of square one. The frame values
at twelve in (A.31.4), now changed to \((1,b_{12})\),
give \(\Delta(h)=1\); also \(\Delta(c)=0\).
Every \(a_N\) acts by \(-I\) on
\(F_N^{\mathrm{ab}}=\widehat{\mathbb Z}^{,2}\),
so it has order exactly two. Since \(U_N\subset U_{12}\),
the outer image \(G_N\) maps to \(G_{12}\).
A.24 proves that \(\Delta_N\) is well defined.
This proves (A.31.1).

We finally verify the precise limitation on the tower.
Every finite quotient of \(\Gamma_N\) is an extension
of a quotient of \(Q_N\) by a finite nilpotent group.
Its nonabelian composition factors are therefore of
the form \(\operatorname{PSL}_2(\mathbb F_p)\),
with \(p\) a prime. The same restriction holds for
finite quotients of the open normal subgroup \(F_N\):
given such a quotient, take the finite intersection
of the conjugates of its kernel in \(\Gamma_N\),
then use the resulting finite quotient of
\(\Gamma_N\). A normal subgroup and its quotient
cannot acquire new composition factors.

There is an epimorphism \(F\twoheadrightarrow A_7\),
sending \(x\) to \((1\,2\,3\,4\,5\,6\,7)\)
and \(y\) to \((1\,2\,3)\). These generate
\(A_7\), which is simple of order 2520.
It is not \(\operatorname{PSL}_2(\mathbb F_p)\):
for \(p\ge5\), equality of orders would require
\(p\mid2520\), leaving \(p=5,7\), whose respective
orders are 60 and 168. The cases \(p=2,3\)
are not nonabelian simple. The finite certificate
\path{check_A31_congruence_scope.py} verifies
generation, simplicity by all conjugacy-class normal
closures, and these order calculations.

Put \(T=\bigcap_{12\mid N}R_{\mathrm{nil}}(U_N)\).
Compactness gives
\(F/T\cong\varprojlim_NF_N\), and every finite
quotient of this inverse limit factors through one
\(F_N\). Thus \(A_7\) is not its quotient.
The image of \(T\) under the displayed epimorphism
to \(A_7\) is consequently nontrivial, and hence
is all of \(A_7\). In particular \(T\ne1\),
proving noncofinality. This group-theoretic kernel
statement alone makes no assertion about the
faithfulness of the \(\widehat{GT}_0\)-action
on the quotient.
\end{proof}

The sharper arithmetic-only obstruction of A.26 is
consistent with this construction. Its finite
quotient \(D_6\) is not a quotient of any of the
common bases \(Q_N\) through the named quotient
maps: the \(\sigma\) above acts identically on all
\(Q_N\), but its two nonzero normalized CM bits
prevent identity on \(D_6\) by A.26.
Thus no \(U_N\) is contained in
\(\overline{U_{12}'U_{12}^6}\).
This concerns the finite common bases; \(D_6\)
is already a quotient of the larger relative group
\(\Gamma_{12}\). A continuation below that kernel
would require additional information beyond A.31.

\statement{A.32}{Proved}{Congruence CM refinements cannot supply the invariant fourth-root character obstruction}
Let \(\mathcal H_{\mathrm{cong}}\) be the image of
\(\mathcal H\) on the family of actual groups
\(\Gamma_N\), for \(12\mid N\). The character
\(\delta=\Delta|_{\mathcal H}\) descends to this
image. There is no continuous homomorphism
\[
 \eta:\mathcal H_{\mathrm{cong}}\longrightarrow
             \mathbb Z/4
\]
whose reduction modulo two is \(\delta\) and
which is invariant under conjugation by \(c\).
The same exclusion applies to a proposed character
on \(\mathcal H\) depending only on its congruence
CM-frame data, including any refinement at moduli
24 or 48, or continuously on their entire tower.

\begin{proof}
The descent of \(\delta\) follows from A.24.
Let \(h\) be the element constructed in A.31,
and let \(\bar h\) be its actual image on the
entire relative congruence tower. There
\(\bar h(c\bar h c)=1\), while
\(\delta(\bar h)=1\). A putative \(\eta\)
would give
\[
 0=\eta\bigl(\bar h(c\bar h c)\bigr)
   =2\eta(\bar h)=2\quad\text{in }\mathbb Z/4,
\]
a contradiction.

For the frame formulation, two normalized full
elements having the same actual action on
\(\Gamma_N\) have r-frames differing by an
r-centralizer there. Projecting to
\(\Gamma_{N,3}\), A.23 identifies that centralizer
with \(\langle r\rangle\); purity makes the
difference trivial. Their finite r-frame images
therefore agree. The same argument at two works
for the s-frames. The finite signed frame is
determined by its projective class, since its
reduction modulo four lies in the pure four-group,
which does not contain \(-I\).
The multiplier-one congruence shear is likewise
determined by the actual \(Q_N\)-action: a shear
\(U^t\) centralizing the projective image of
\(S\) has \(t=0\pmod N\), by its lower-left
and diagonal entries. Hence the normalized CM
values formed from these frames and the shear
are functions of the actual action.
Any character depending only on those values
factors through \(\mathcal H_{\mathrm{cong}}\).
The compact quotient topology makes this factor
continuous, so the preceding contradiction applies.
\end{proof}

A global \(c\)-invariant \(\mathbb Z/4\)-lift of
\(\delta\), if it exists, must therefore use data
not determined by this relative congruence tower.
Neither its existence nor a solution of the full
norm equation is proved here.

\statement{A.33}{Proved}{The compatible congruence family also has equivariant literal frames}
The full element \(a=hc\) in A.31 can be chosen so that
its square belongs to
\[
 L_{\mathrm{cong}}^{\mathrm{fr}}=
 \bigcap_{12\mid N}\{k\in\mathcal H:
       \Phi_k|_{\Gamma_N}=1,\quad
       A_k|_{\Gamma_N}=\beta_k|_{\Gamma_N}=1\}.
 \tag{A.33.1}
\]
At every level, the literal r- and s-frames can therefore
be retained with the actual finite normalizers \(S_3\)
and \(C_2^2\) from A.30. This strengthens the compatible
actual-action assertion without asserting a full
order-two lift.

\begin{proof}
Use the arithmetic \(\sigma\) of A.31 and set
\(u_N=q_N(A_\sigma)\), \(b_N=q_N(\beta_\sigma)\).
We first prove the matrix fact replacing the special
modulo-twelve adjustment in A.30. In signed matrices put
\[
 R=\begin{pmatrix}0&-1\\1&-1\end{pmatrix},\qquad
 D=\begin{pmatrix}0&1\\1&0\end{pmatrix}.
\]
The signed representative \(u_N\) centralizes \(R\)
and is identity modulo four. Indeed projective
centralization cannot replace \(R\) by \(-R\):
their traces are \(-1\) and 1, different modulo
every \(N\) under consideration. The signed
modulo-four assertion is part of R3's pure-frame
normalization. Since \(R\) is cyclic, its matrix
centralizer consists of \(vI+wR\); for a
determinant-one element \(u\) of this centralizer,
\[
              DuD=u^{-1},\qquad
              (u^{-1}D)^2=I,
              \operatorname{tr}(u^{-1}D)=0.
 \tag{A.33.2}
\]

The reflection \(u^{-1}D\) is conjugate to \(D\)
by an element of \(\operatorname{SL}_2(\mathbb Z_p)\)
at every prime, with conjugator identity modulo two
at the dyadic prime. For odd \(p\), its two eigenvalues
are 1 and \(-1\), and its eigensummands are free
rank-one direct summands. Thus a general linear
conjugator exists. The determinant map on
\(C_{\operatorname{GL}_2(\mathbb Z_p)}(D)\) is onto,
so its determinant can be corrected to one.
For \(p=2\), put \(T=u^{-1}D\equiv D\pmod4\)
and take the cyclic-basis matrix
\[
                   H_0=(e_1,Te_1).
\]
It satisfies \(H_0DH_0^{-1}=T\),
\(H_0\equiv I\pmod2\), and
\(d_0=\det H_0\equiv1\pmod4\).
The matrix
\[
 K=\begin{pmatrix}v&w\\w&v\end{pmatrix},\qquad
 v=(d_0^{-1}+1)/2,\quad w=(d_0^{-1}-1)/2,
\]
centralizes \(D\), is identity modulo two,
and has determinant \(d_0^{-1}\).
Consequently \(H_0K\) is the required special
linear conjugator. Reduction and the Chinese
remainder theorem give the finite pure conjugator
at every \(N\). It lifts to a pure element of
\(\widehat B\), because reduction modulo two
is the quotient \(\widehat B/F\cong S_3\).

At the dyadic relative level, retain the full
Galois involution \(a_2=\Phi_\sigma c\).
Its r-normalizing involution
\(D_2=J_2^{-1}a_2\), with
\(J_2=\Phi_\sigma(s)A_{\sigma,2}\), has finite
image \(u_N^{-1}d\), where \(d=sc\).
The matrix calculation provides a pure conjugate
of \(D_2\) with finite image \(d\).
Apply A.16 to the inverse image of
\(\langle q_N(r),d\rangle\cong S_3\), retaining
that involution. Coprime conjugacy of its
order-three complement with \(\langle r\rangle\)
then gives \(D'_2\) inverting \(r\) and having
finite image \(d\). The correction
\(D_2D'_2\) centralizes \(r\) and has finite
image \(u_N^{-1}\). It changes the r-frame
to finite image one while keeping its associated
r-normalizer involutory. At odd primes use the
identity frame and \(sc\), as in A.30.

For the s-frame, \(b_N\) centralizes the signed
matrix \(S\). Conjugation by
\(C=\operatorname{diag}(-1,1)\) sends it to its
inverse, because \(CSC=S^{-1}\) and the
determinant-one \(S\)-centralizer is the
corresponding quadratic norm-one group.
Thus \(b_N^{-1}c\) is an involution in the
common finite extension. At each odd prime the
centralizer epimorphism in A.23, followed by
Schur--Zassenhaus, supplies an involution above
it. At two keep the full Galois s-normalizer.
This gives the same equivariant s-frame gluing
as in A.30.

R2 now supplies a full lift at each \(N\) with
these literal frames. The proof of A.30 shows
that its square has identity actual action and
both relative frames equal to one. Restrict
the nonempty compact sets \(\mathcal C_N\)
from A.31 by these additional closed conditions
on the frames of \((hc)^2\). A common multiple
again provides the finite-intersection property,
so one full \(h\) satisfies the conditions for
all \(N\). This proves (A.33.1).
\end{proof}

The matrix certificate
\path{check_A33_equivariant_tower.py} constructs
the pure special-linear conjugators for every
allowed r-Cartan value at moduli
12, 24, 48, 60, 84, 120, and 132. The proof above
establishes all moduli. In the specified full
framed fibre the remaining equation is
\[
            v(ava^{-1})=a^{-2},
            \qquad v\in L_{\mathrm{cong}}^{\mathrm{fr}}.
\]
This is again a sufficient route to a full
involution, not a claimed solution or a necessary
restriction on differently framed lifts.

\statement{A.34}{Proved}{The full peripheral lifting theorem holds over any closed normal relative kernel}
Let \(\Pi\) be either the full free three-cusp group
\(F\), or the full modular group \(\widehat B\), with
the respective three peripheral elements
\((x,y,z)\) or \((s,r,\tau^{-1})\).
Let \(V\triangleleft\Pi\) be any closed normal subgroup,
without an openness or characteristicity assumption,
and put \(R=R_{\mathrm{nil}}(V)\).
Every continuous automorphism of \(\Pi/R\) preserving
the three positive peripheral conjugacy classes, with
specified literal conjugators there, lifts to a full
peripheral automorphism of \(\Pi\) with those conjugators.
The lift stabilizes \(R\). If a specified peripheral
conjugator is one, that peripheral can be fixed literally.

For a descending directed family of closed subgroups
\((V_i)\) of any profinite group, one also has
\[
 R_{\mathrm{nil}}\Bigl(\bigcap_i V_i\Bigr)
                  =\bigcap_i R_{\mathrm{nil}}(V_i).
 \tag{A.34.1}
\]
In particular, if \(K_{\mathrm{cong}}=\bigcap_NU_N\),
the group kernel \(T\) in A.31 is exactly
\(R_{\mathrm{nil}}(K_{\mathrm{cong}})\).

\begin{proof}
For every profinite group \(V\), with all commutator
subgroups taken closed,
\[
 R_{\mathrm{nil}}(V)=\bigcap_{n\ge1}\gamma_n(V),
 \qquad
 q(R_{\mathrm{nil}}(V))=
                   \gamma_\infty(q(V))
 \tag{A.34.2}
\]
for every finite continuous image \(q(V)\).
The first equality follows because bounded-class
nilpotent profinite quotients are separated by their
finite nilpotent quotients. For the second, a
descending family of nonempty compact fibres of
a fixed finite image value has nonempty intersection.
The finite lower central series stabilizes, so
\[
 q([R,V])=[\gamma_\infty(q(V)),q(V)]
          =\gamma_\infty(q(V))=q(R).
\]
Finite quotients separate closed subgroups. Therefore
\([R,V]=R\), and normality gives
\[
                             [R,\Pi]=R.
 \tag{A.34.3}
\]

Apply the finite scalar-cover and complete proper-fibre
count used in A.3. For a finite quotient
\(q_E:\Pi\twoheadrightarrow E\), take
\(A=q_E(R)\). The auxiliary actual source map
\(q^\dagger:\Pi\twoheadrightarrow D^\dagger\)
has central image of \(R\), because its composite
to \(E/A\) kills \(R\).
Equation (A.34.3) implies that \(q^\dagger\) itself
kills \(R\). Thus the specified automorphism of
\(\Pi/R\) composes with this actual quotient map
to give the required new epimorphism to
\(D^\dagger\), with the prescribed conjugators.
The finite count gives a positive proper fibre over
every joint finite constraint, including all the
specified conjugator coordinates.

Compactness of \(\Pi^3\) now supplies full
conjugators, the peripheral product relation, and
generation in every finite quotient, exactly as
in A.3. The reference D.50 count explicitly applies
to both sources: in the modular case conjugacy
to \(s,r\) supplies their defining orders, so no
homomorphism from the wrong free source is inserted.
The resulting continuous endomorphism of the
finitely generated group \(\Pi\) is surjective
and hence an automorphism. If it is denoted by
\(\Phi\), and \(a\) is the specified quotient
automorphism, then \(q_R\Phi=a q_R\).
Since both maps are automorphisms, this implies
\(\Phi^{-1}(R)=R\) and hence \(\Phi(R)=R\).
Thus preservation is a conclusion, not a missing
characteristicity assumption. If one conjugator
is one in the retained quotient, its full lift
lies in \(R\); the usual inner correction fixes
that peripheral literally without changing the
retained map or any prescribed quotient conjugator.

Finally take any finite image \(q\) of the ambient
group containing the \(V_i\). Compactness gives
\(q(\bigcap_iV_i)=\bigcap_iq(V_i)\).
The descending directed family of subgroups of
this finite group stabilizes. Its lower central
residual consequently commutes with that intersection.
Using (A.34.2) on each subgroup gives equal finite
images for the two sides of (A.34.1). Both sides
are closed, proving equality. Apply this to
the family \(U_N\) to obtain the final assertion.
\end{proof}

The same audit extends A.3's central-quotient criterion
from characteristic to closed normal \(N\), when
``lift'' explicitly means an automorphism inducing
the given automorphism of \(\Pi/N\). Necessity
uses its resulting preservation of \(N\) and
\([N,\Pi]\); sufficiency is the unchanged finite
count. These extensions concern ordinary peripheral
automorphisms. They do not solve an involution
norm equation or construct a desired automorphism
of an unspecified infinite quotient.

\statement{A.35}{Proved}{The full framed congruence correction kernel contains a closed dyadic procyclic subgroup}
Let \(L_{\mathrm{cong}}^{\mathrm{fr}}\leq\mathcal H\)
be the kernel in A.33, retaining the actual action and
both literal full relative frames at every principal
congruence level. Then
\[
             \mathbb Z_2\ \lhook\joinrel\longrightarrow
                    L_{\mathrm{cong}}^{\mathrm{fr}}
 \tag{A.35.1}
\]
is a continuous closed embedding. The image may also
be required to lie in all three corrected resolvent
kernels of B.53. In particular the remaining norm
kernel of A.33 is not pro-odd. This statement alone
neither obstructs nor solves its twisted norm equation.

\begin{proof}
Write
\[
 R=R_{\mathrm{nil}}(K_{\mathrm{cong}})
     =\bigcap_{12\mid N}R_{\mathrm{nil}}(U_N),
       \qquad \Pi=\widehat B.
\]
The equality is A.34. Use the specified tripod map
\(q_0:\widehat K\twoheadrightarrow A_9\) and the
fifteen normal filling kernels \(N_i\) of B.57.
We first check the extra retained kernel:
\[
                         q_0(R)=A_9.
 \tag{A.35.2}
\]
Indeed \(R\) is normal in \(\Pi\) and contained in
\(\widehat K\). Its image is normal in the simple
group \(A_9\). If that image were trivial, \(q_0\)
would factor through \(\widehat K/R\), an open normal
subgroup of \(\Pi/R\). Taking the core of the finite
quotient kernel would make \(A_9\) a composition
factor of a finite quotient of \(\Pi/R\).
Every such quotient factors through one of the relative
congruence groups in A.31. Their nonabelian composition
factors are only \(\operatorname{PSL}_2(\mathbb F_p)\)
for primes \(p\). None is \(A_9\): its order is
\(181440\), and a prime \(p\geq5\) dividing that
order is only \(5\) or \(7\), whose corresponding
simple groups have orders \(60\) and \(168\).
This contradiction proves (A.35.2).

Set
\[
                 V=R\cap\bigcap_{i=1}^{15}N_i.
\]
Each of these sixteen closed normal subgroups of
\(\widehat K\) has full \(q_0\)-image, by (A.35.2)
and B.57. Their iterated closed commutator lies in
\(V\) and still has image \(A_9\), since \(A_9\)
is perfect. Thus \(q_0(V)=A_9\). The intersection
is normal in \(\Pi\). The induced sixty-coordinate
quotient of B.57 consequently again has
\[
 q(V)=M=A_9^{60},\qquad
 \Pi/(V\cap\ker q)=M\rtimes(\Pi/V).
\]
All subsequent central-refinement arguments in B.57
apply with this smaller \(V\). The diagonal odd
transposition on \(M\) extends by the identity on
\(\Pi/V\), and extends to the actual central
refinement while preserving the classes of \(s,r,\tau\).
Moreover the three conjugators constructed there
belong to the perfect subgroup \(M^\#\). They
therefore have image one on \(\Pi/V\).

Apply A.3 and A.34, retaining these projected
conjugators. The full cusp conjugator belongs to
\(V\); conjugating the resulting automorphism by its
inverse normalizes it to fix \(\tau\) literally.
The resulting \(h\in\mathcal H\) still induces the
identity on \(\Pi/V\), and its two unique pure
frames belong to \(V\). Since \(V\subseteq R\),
it induces the identity on \(\Pi/R\) and its two
frames have image one there. In particular
\(h\in L_{\mathrm{cong}}^{\mathrm{fr}}\).
Its three resolvent images are one, because \(V\)
is contained in the original fifteen filling kernels.

The tripod outer action of \(h\) on the same
\(A_9\) quotient is the nontrivial class in
\(\operatorname{Out}(A_9)=C_2\). Cusp normalization
changes that action only by an inner automorphism:
its conjugator lies in \(V\subseteq\widehat K\).
As in B.57, this gives a continuous quotient of order
two of the closed procyclic group
\(\overline{\langle h\rangle}\). R1 says that
\(\mathcal H\) is torsion-free. A torsion-free
procyclic profinite group with a quotient of order
two has dyadic factor \(\mathbb Z_2\), proving
(A.35.1). Closedness of the action and frame
conditions keeps the entire procyclic group in
\(L_{\mathrm{cong}}^{\mathrm{fr}}\).
\end{proof}

The construction supplies no formula for the action
of the chosen negative lift \(a\) on this procyclic
subgroup, nor a value of \(a^2\) in its abelian
quotients. It therefore does not replace the
calculation of \(H^2(C_2,M)\) in A.19. Its exact
consequence for Lead~3 is that a coprime splitting
argument cannot be justified by declaring the
noncongruence correction kernel pro-odd.


\statement{A.36}{Proved}{A peripheral conjugator to a Grothendieck--Teichm\"uller involution satisfies the pentagon}
Let \(\alpha\in\widehat{GT}\) be an involution and let
\(g\in\mathcal P\) satisfy \(\alpha=g c g^{-1}\).
Then \(g\in\widehat{GT}\). Consequently
\[
 \widehat{GT}\cap\{g c g^{-1}:g\in\mathcal P\}
   =\{g c g^{-1}:g\in\widehat{GT}\}.
 \tag{A.36.1}
\]
For an involution already in \(\widehat{GT}\), conjugacy
to \(c\) in any of \(\mathcal P,\widehat{GT}_0,
\widehat{GT}\) is therefore equivalent. This statement
does not assert that every involution is so conjugate.

\begin{proof}
A.6 first puts \(g\) in \(\widehat{GT}_0\). Write
\(\alpha=(-1,f)\) and \(g=(\mu,h)\) in our normalized
convention. The equality \(\alpha g=g c\), computed
on \(y\), and uniqueness of the normalized conjugator
give
\[
                         h=f\,\alpha(h).
 \tag{A.36.2}
\]
We shall use the actual action of \(\alpha\) on the
completed parenthesized braid operad
\(\widehat{\mathrm{PaB}}\). We do not presume that
\(g\) acts on its four-strand component.

Here is the groupoid convention explicitly. Objects are
labelled binary parenthesizations. Composition is
written right to left, and the endomorphism group of
any object with \(n\) labels is \(\widehat P_n\).
Let
\[
 b=((12)3),\qquad b'=(1(23)),\qquad A:b\longrightarrow b'
\]
be the associativity arrow with underlying identity
braid. Let \(x\) be the full twist of strands one and
two at \(b\), and let
\(y=A^{-1}y' A\), where \(y'\) is the full twist of
strands two and three at \(b'\). These give the free
factor \(\widehat F(x,y)\) in
\(\widehat P_3=\widehat{\mathbb Z}\times\widehat F_2\).
Choose the customary source-based convention
\[
             \alpha(A)=A f,
 \qquad \alpha(x)=x^{-1},
 \qquad \alpha(y)=f^{-1}y^{-1}f.
 \tag{A.36.3}
\]
The last identity follows directly by applying
\(\alpha\) to \(A^{-1}y'A\), so this convention
agrees with the normalization at the start of this
part. The associator parameter has no central full-twist
coordinate and belongs to the commutator subgroup;
these are exactly the normalization conditions in the
parenthesized-braid presentation
\cite[Definitions 6.10--6.11, Lemma 7.1 and Proposition 7.3]{CEBHR2019}.

Consider the genuine arrow
\[
                     \theta=A h:b\longrightarrow b'.
\]
It exists for every \(h\in\widehat F(x,y)'\), whether
or not that parameter satisfies the pentagon. Equations
(A.36.2)--(A.36.3) imply
\[
             \alpha(\theta)=A f\alpha(h)=A h=\theta.
 \tag{A.36.4}
\]
Now form its five operadic insertions between the
following actual four-strand objects:
\[
 \begin{aligned}
 b_0&=(((12)3)4),& b_1&=((1(23))4),\\
 b_2&=(1((23)4)),& b_3&=(1(2(34))),\\
 b_4&=((12)(34)).
 \end{aligned}
\]
The three-edge path from \(b_0\) to \(b_3\) consists of
\[
 \begin{aligned}
 e_1&=\theta^{1,2,3}:b_0\to b_1,\qquad
 e_2=\theta^{1,23,4}:b_1\to b_2,\\
 e_3&=\theta^{2,3,4}:b_2\to b_3.
 \end{aligned}
\]
The two-edge path consists of
\[
 d_1=\theta^{12,3,4}:b_0\to b_4,\qquad
 d_2=\theta^{1,2,34}:b_4\to b_3.
\]
The superscripts mean literal cabling by identity arrows
of the indicated binary objects; an unused outer
strand is also inserted by an identity arrow. Since
\(\alpha\) is an object-fixing operad automorphism,
it fixes those identity arrows and commutes with these
insertions. Equation (A.36.4) therefore fixes all five
arrows, not merely their classes. Their discrepancy is
the based pure braid
\[
                D=(d_2d_1)^{-1}(e_3e_2e_1)
                   \in\operatorname{End}(b_0)=\widehat P_4,
 \qquad \alpha(D)=D.
 \tag{A.36.5}
\]

C.28 gives \(\operatorname{Fix}_{\widehat P_4}(\alpha)=1\).
For clarity its application here uses actual strand
forgetting, compatible with the unital operadic action
\cite[Proposition 11.5 and Corollary 11.6]{CEBHR2019}.
The full profinite pure-braid sequences have free
kernels of ranks two and three. Each finite
point-pushing cusp is mapped to a conjugate of its
inverse. The infinity cusp is represented, up to
inverting the peripheral generator, by
\(z_m z_{m-1}^{-1}\), where \(z_j\) is the full
central twist of the first \(j\) strands. Both twists
are sent to their inverse; choose a parenthesization
grouping the first \(m-1\) strands to make this
literal. Thus the kernel action has all the negative
peripheral classes required by C.28. Induction from
\(\widehat P_2=\widehat{\mathbb Z}\), on which the
action is inversion, proves the displayed fixed-subgroup
assertion. In particular (A.36.5) has \(D=1\).

This is precisely the pentagon for the parameter
\(h\). The two hexagons already hold because
\(g\in\widehat{GT}_0\). If the half-braiding arrow is
denoted \(R:(12)\to(21)\), its parameter \(\mu\) is
implemented by the well-typed arrow
\[
 R_\mu=R\,(R_{\mathrm{op}}R)^{(\mu-1)/2}:(12)\to(21),
 \qquad R_{\mathrm{op}}:(21)\to(12),
\]
rather than by an untyped groupoid power. Every unit
\(\mu\in\widehat{\mathbb Z}^{\times}\) is odd.
The parenthesized-braid presentation identifies these
hexagons and this pentagon with the defining
Grothendieck--Teichm\"uller equations
\cite[Lemma 7.1 and Proposition 7.3]{CEBHR2019}.
Since \(g\) already is an invertible normalized
pair in \(\widehat{GT}_0\), the extra pentagon puts
it in \(\widehat{GT}\). More explicitly, the pentagon
solutions inside the profinite group \(\widehat{GT}_0\)
form a closed submonoid by the operadic presentation.
Every closed submonoid of a profinite group is a
subgroup: the powers \(g^{n!-1}\) converge to
\(g^{-1}\). This also checks the unit condition in
the operadic monoid definition of \(\widehat{GT}\).
\end{proof}

The decisive fixed-point argument used only the honest
operadic action of the given \(\alpha\in\widehat{GT}\).
It does not construct a five-point action of an
arbitrary element of \(\widehat{GT}_0\), nor prove
that a general involution there has such an action.
Thus A.36 is a descent of an existing conjugator,
leaving the existence questions in A.13--A.14 open.

\statement{A.37}{Proved}{Every full Grothendieck--Teichm\"uller involution gives a unique pentagon solution}
Let \(\alpha=(-1,f)\in\widehat{GT}\) be an involution.
Let \(k\in F'\) be the unique solution supplied by R1
of \(f=k^{-1}\alpha(k)\), and put \(u=k^{-1}\).
Then \(u\) satisfies the full profinite pentagon.
Moreover \(\alpha\) is conjugate to \(c\) in
\(\widehat{GT}\) if and only if there exists
\(\mu\in\widehat{\mathbb Z}^{\times}\) such that
\[
                 x^\mu u^{-1}y^\mu u\sim_F(xy)^\mu.
 \tag{A.37.1}
\]

\begin{proof}
The equation for \(k\) is equivalent to
\(u=f\alpha(u)\). Therefore the arrow \(A u\)
is fixed by the genuine action of \(\alpha\),
exactly as in (A.36.4). Its five inserted arrows
have a fixed pentagon discrepancy in \(\widehat P_4\),
which C.28 makes one. This proves the pentagon
without assuming that a pair \((\mu,u)\) is an
endomorphism of the operad.

If (A.37.1) holds, the normalized substitution
\(g=(\mu,u)\) is an endomorphism preserving the
three named cusp classes with the common unit
multiplier \(\mu\). C.29 makes it a full
automorphism, so \(g\in\mathcal P\).
The identity \(u=f\alpha(u)\) gives
\(\alpha g=g c\) on both free generators, and
A.36 puts \(g\) in \(\widehat{GT}\).
Conversely any such conjugator \(g=(\mu,v)\)
has \(v=f\alpha(v)\). The uniqueness in R1
implies \(v=u\), and its peripheral property gives
(A.37.1).
\end{proof}

The remaining condition is the full third cusp
with a suitable unit multiplier. Full generation
follows once that equation holds. This automatic
generation and the general peripheral conjugator
criterion already occur in the supplied LM4 D.35;
C.29 independently proves the needed generation
statement by finite-cover boundary homology. The
additional result here is that the canonical \(u\)
already satisfies the pentagon. A.5 shows that the
third cusp need not hold with multiplier one.
The generation condition is independent of \(\mu\):
unit powering preserves each of the two closed
procyclic subgroups. The choice of multiplier
affects the third-cusp condition.
No prounipotent pentagon-to-hexagon implication
has been transferred to the full profinite setting.
Furusho's pro-$\ell$ and pronilpotent conclusions
allow a quadratic extension for the multiplier;
the full profinite implication is posed separately
as Question~14 in that paper
\cite[\S3, Corollaries 12--13 and Question 14]{Furusho2010}.

\statement{A.38}{Proved}{The prescribed coboundary gives peripheral conjugators at two and three}
Let \(\alpha\in\widehat{GT}\) be an involution, and let
\(u\in F'\) be its canonical parameter from A.37.
Use the lower central series
\(\gamma_2(F)=[F,F]\), \(\gamma_3(F)=[\gamma_2(F),F]\)
and define \(t\in\widehat{\mathbb Z}\) by
\[
              u\equiv[x,y]^t\pmod{\gamma_3(F)},
              \qquad D=1+24t.
 \tag{A.38.1}
\]
The commutator convention is the one whose logarithm
has leading term \([X,Y]\), as in A.4.
For a prime \(p\), there exists a peripheral
conjugator from \(c_p\) to \(\alpha_p\) of the
specific normalized form
\[
 g_p(x)=x^{\mu_p},\qquad
 g_p(y)=u_p^{-1}y^{\mu_p}u_p,
 \qquad \mu_p\in\mathbb Z_p^{\times},
 \tag{A.38.2}
\]
if and only if \(D_p\) is a square unit in
\(\mathbb Z_p\). This condition always holds at
\(p=2\) and \(p=3\). For \(p\geq5\), it is equivalent
to \(D_p\bmod p\) being a nonzero quadratic residue.
In particular every involution already in full
\(\widehat{GT}\) is peripherally conjugate to the
standard involution on its maximal pro-\(2\) quotient.

\begin{proof}
A.37 proves the full pentagon for \(u\), hence its
pro-\(p\) pentagon. Furusho's Corollary~12 gives the
two hexagons with parameter
\(\mu_p=\pm\sqrt{1+24t_p}\)
\cite[Theorem 2 and \S3, Corollary 12]{Furusho2010}.
The theorem generally permits a quadratic extension;
we use it only when the chosen root belongs to
\(\mathbb Z_p^{\times}\). Then the half-parameter
\(m_p=(\mu_p-1)/2\) belongs to \(\mathbb Z_p\):
for odd \(p\) this is immediate, and for \(p=2\)
it follows because \(\mu_p\) is odd. Consequently
all words in the hexagons lie in \(F^{(p)}\).
The injective comparison
\(F^{(p)}\hookrightarrow F_2(\mathbb Q_p)\)
used in Furusho's proof therefore makes the
hexagons identities in the actual pro-\(p\) group,
not merely in its prounipotent completion.
They give the three peripheral power classes.
The substitution (A.38.2) has mod-\(p\) homology
matrix \(\mu_p I\), hence is an automorphism by the
pro-\(p\) Frattini criterion. The exact identity
\(u=f\alpha(u)\) from A.37 gives
\(\alpha_p g_p=g_p c_p\) on both free generators.

Conversely, apply the necessary third-cusp coefficient
calculation A.4 in the class-three
\(\mathbb Q_p\)-Malcev quotient. This calculation is
valid for every prime after passing to \(\mathbb Q_p\),
and gives
\(\mu_p^2=1+24t_p\). A unit multiplier therefore
requires a square unit. Finally
\[
       D_2\in1+8\mathbb Z_2,
       \qquad D_3\in1+3\mathbb Z_3.
\]
Every element of \(1+8\mathbb Z_2\) is a square
unit, and Hensel's lemma gives a square root in
\(\mathbb Z_3^{\times}\) of every element of
\(1+3\mathbb Z_3\). For odd \(p\), the unit-square
criterion is exactly the nonzero-square criterion
in the residue field. This proves all assertions.
\end{proof}

The parameter in (A.38.1) uses the lower-central
quotient \(\gamma_2(F)/\gamma_3(F)\), not the
abelianization of the derived group. This is also
the integral coefficient represented by the
\(XY\)-coefficient in Furusho's exponential
notation. The displayed derived-series expression
preceding Question~14 on printed page~555 of that
paper has a typographical error: its asserted
cyclic quotient is the lower-central one used here.

A.38 does not settle the dyadic case for arbitrary
\(\mathcal P\)- or \(\widehat{GT}_0\)-involutions:
the genuine full \(\widehat{GT}\)-action supplied
A.37's pentagon. It also concerns maximal pro-\(p\)
quotients, not all relative quotients \(F/U^{(p)}\).
For other local conjugators, whose parameter need
not be \(u_p\), the necessity assertion about
\(D_p\) is not being made.
For the actual pentagon word supplied by a full
\(\widehat{GT}\)-involution, A.42 below proves
that this square-unit criterion is always satisfied.

\statement{A.39}{Proved}{The audited full conjugacy question is one unit third-cusp equation}
Let \(\alpha=(-1,f)\in\mathcal P\) be a nonidentity
involution. Let \(u\in F'\) be the unique solution
of \(u=f\alpha(u)\), supplied by R1 after inverting
its coboundary parameter. Define
\[
 \mathcal M_\alpha=
 \{\mu\in\widehat{\mathbb Z}^{\times}:
       x^\mu u^{-1}y^\mu u\sim_F(xy)^\mu\}.
 \tag{A.39.1}
\]
Then \(\alpha\) is conjugate to \(c\) in
\(\mathcal P\) if and only if
\(\mathcal M_\alpha\ne\varnothing\).
This set is either empty or exactly
\(\{\mu,-\mu\}\) for one unit \(\mu\).
If \(\alpha\) belongs to \(\widehat{GT}_0\),
respectively \(\widehat{GT}\), both resulting
conjugators belong to that same group.
The existence of a unit satisfying (A.39.1) is
detected by all characteristic finite quotients of
\(F\). For a specified unit \(\mu\), the validity
of (A.39.1) is not detected by the maximal
pro-\(p\) quotients alone, even with \(u=1\).

The basic one-equation criterion and automatic
generation were asserted in the supplied LM4 D.35;
D.49 asserted the two-element conjugator fibres.
This statement audits those conclusions using
C.29's independent covering proof, R1's now-unique
parameter, and A.6. A.36 supplies the additional
descent to full \(\widehat{GT}\).

\begin{proof}
If \(\mu\in\mathcal M_\alpha\), the free universal
property defines the endomorphism
\[
       g_\mu(x)=x^\mu,\qquad
       g_\mu(y)=u^{-1}y^\mu u.
\]
Equation (A.39.1) makes it peripheral at all three
cusps with the common unit power. C.29 makes it
automatically onto, hence an automorphism.
The identity \(u=f\alpha(u)\) gives
\(\alpha g_\mu=g_\mu c\). Thus it is the desired
conjugator in \(\mathcal P\). Conversely every
normalized conjugator has this parameter \(u\)
by its uniqueness in R1 and satisfies (A.39.1).
A.6 and A.36 give the two membership assertions.

If \(\mu,\nu\in\mathcal M_\alpha\), then
\(g_\nu^{-1}g_\mu\in C_{\mathcal P}(c)=\{1,c\}\)
by A.6. Its multiplier gives \(\mu=\nu\) or
\(\mu=-\nu\). Both signs occur because
\(g_\mu c=g_{-\mu}\), and these two units are
distinct. This proves the asserted solution set.

For finite detection, let \(Q=F/V\) range over
characteristic finite quotients, and let
\(e_Q\) be its exponent. Test whether there is
\[
 a\in(\mathbb Z/e_Q\mathbb Z)^{\times}
 \quad\hbox{with}\quad
 \bar x^a\bar u^{-1}\bar y^a\bar u
          \sim_Q(\bar x\bar y)^a.
 \tag{A.39.2}
\]
Every such residue \(a\) lifts to a profinite unit,
and every conjugator in \(Q\) lifts to \(F\).
Thus (A.39.2) is the nonemptiness of a closed subset
of the compact space
\(\widehat{\mathbb Z}^{\times}\times F\), where
we retain the unit and the third-cusp conjugator.
A common characteristic refinement gives the
finite-intersection property. Nonemptiness at
every finite quotient is therefore equivalent
to a full solution, by compactness.

For the final assertion take \(\alpha=c\), so
\(u=1\), and choose the profinite unit with
\[
          \mu_2=-1,\qquad \mu_p=1\quad(p\text{ odd}).
 \tag{A.39.3}
\]
At every odd prime the third-cusp equation is
literal. At two, \(x^{-1}y^{-1}\) is conjugate
to \((xy)^{-1}\), so it holds there as well.
But \(\mu\equiv7\pmod{12}\). In the actual
quotient \(F\twoheadrightarrow S_4\) given by
\[
                x=(234),\qquad y=(1234),
\]
with products acting right to left, direct
multiplication gives
\[
     x^7y^7=(12),\qquad (xy)^7=(1423).
 \tag{A.39.4}
\]
A transposition is not conjugate to a four-cycle.
The script \path{check_A39_mixedpower.py} verifies
that the displayed pair generates all twenty-four
elements and checks the full conjugacy class.
This is a finite mixed-prime obstruction to the
particular multiplier (A.39.3), despite every
maximal pro-\(p\) equation holding.
\end{proof}

For \(\alpha\in\widehat{GT}\), A.37 adds the
pentagon for its prescribed \(u\). The later
A.42 proves that full unit roots of
\(\mu^2=1+24c_2(u)\) exist and that every local
hexagon holds. The remaining question is whether
\emph{some} full root gives the full hexagon.
This existential assertion would settle conjugacy:
the two-cycle already follows from the pentagon,
the three-cycle gives the third cusp by C.30,
and C.29 supplies full generation. C.34 explains
why the literal prescribed-root version of
Furusho's Question~14 fails: the mixed-sign
example (A.39.3)--(A.39.4) also fails its full
hexagon. It does not refute existence of a
suitable root, since the same example has the
two standard roots \(1,-1\).

\statement{A.40}{Proved}{A parity-one involution survives an actual noncongruence perfect-socle refinement}
Put
\[
 \Pi=\widehat B,\qquad
 R=R_{\mathrm{nil}}(K_{\mathrm{cong}}),\qquad C=\Pi/R,
 \qquad \widehat K=\ker(\Pi\longrightarrow A_5).
\]
There is a finite quotient
\[
 q:\Pi\twoheadrightarrow Q=M\rtimes A_5,
 \qquad M=A_9^t\quad(t\geq60),
 \tag{A.40.1}
\]
whose kernel is invariant under the full
\(\widehat{GT}_0\)-action, and for which the
\(A_5\)-action on the simple factors is a union of
regular orbits. Let
\[
 W=R\cap\ker q,
 \qquad D=\Pi/W,
 \qquad F_D=F/W.
\]
There is a full \(a\in\widehat{GT}_0\) with
\(\lambda(a)=-1\) and \(\Delta(a)=1\) whose actual
induced automorphism on \(D\) has order two.
Its outer image on \(F_D\) has order two and
projects to the entire compatible family of A.31.
The literal \(C\)-frames from A.33 may also be retained.

This quotient genuinely extends the relative
principal-congruence tower: \(D\) maps onto \(C\)
and has an \(A_9\) composition factor, which no
finite quotient of \(C\) has. The statement does
not assert \(a^2=1\) in the full group, nor that
\(W\) is characteristic under all automorphisms
of the underlying free group.

\begin{proof}
\emph{The finite invariant family.}\par
The twelve cusps of \(\widehat K\) are the
\(A_5/C_5\) cosets. The antipodal relation is
intrinsically given by the nontrivial element of
\(N_{A_5}(C_5)/C_5\). Consider all kernels of
epimorphisms \(\widehat K\twoheadrightarrow A_9\)
which factor through a three-cusp filling, whose
retained triple consists of an antipodal pair and
one further cusp, and which send each retained
positive meridian to a seven-cycle. This is a finite
family because \(\widehat K\) is finitely generated.
It is nonempty by the exact generating triple
(B.57.5)--(B.57.6).

Congruence covariance, as used in A.31, preserves
\(\widehat K\). A peripheral automorphism permutes
its cusps and preserves the stated normalizer
description of antipodes. Its common unit powers
also preserve seven-cycle orders. Precomposition
therefore permutes the entire defined family of
kernels. Each kernel has trivial \(A_5\)-stabilizer:
its stabilizer preserves its nonzero cusp support,
and such a triple has trivial stabilizer by B.57.
Thus every orbit has sixty elements.

Choose an epimorphism for one kernel in each orbit,
and induce it from \(\widehat K\) to \(\Pi\) by
the regular wreath-product construction of B.57.
Take the product of these induced maps over their
common \(A_5\)-quotient. The coordinate kernels
on \(\widehat K\) are exactly the distinct kernels
in the family. Their image is the full product:
for a chosen coordinate, every other distinct
maximal normal kernel maps onto its simple target;
an iterated commutator of those other kernels gives
that factor alone. This proves (A.40.1), including
the asserted regular action and invariance of its
kernel.

Every coordinate sends \(R\) onto \(A_9\).
Otherwise it would factor through
\(\widehat K/R\), an open normal subgroup of
\(C\), making \(A_9\) a composition factor of a
finite quotient of \(C\). This contradicts
A.31 and A.35. Normality of \(q(R)\) in the
product \(M\), together with surjectivity onto
each simple factor, gives
\[
 q(R)=M,\qquad
 D=Q\times_{A_5}C\cong M\rtimes C.
 \tag{A.40.2}
\]
For the second equality, lift a prescribed
\(C\)-coordinate to \(\Pi\), then correct its
finite \(M\)-coordinate by an element of \(R\).
The action of \(C\) on \(M\) factors through
\(A_5\).

\emph{The actual central refinement.}\par
Set
\[
 E^\#=\Pi/[W,\Pi],\qquad Z=W/[W,\Pi],
 \qquad R^\#=R/[W,\Pi].
\]
Thus \(Z\) is central in \(E^\#\), and
\(R^\#\) is the full preimage of \(M\).
The canonical image \(M^\#\) of the universal
central cover \((2.A_9)^t\) in \(R^\#\) is
finite and perfect, maps onto \(M\), and is
normal in \(E^\#\). The universal-cover input is
the same primary Schur-multiplier datum used in
B.57.\cite{CEAtlasA9}
Uniqueness of lifts from a perfect group makes
this canonical image invariant under conjugation.
Centrality gives \(R^\#=M^\# Z\). A.34 gives
\(R=[R,\Pi]\), and consequently
\[
 R^\#=[R^\#,E^\#]=[M^\# Z,E^\#]\subseteq M^\#.
\]
Hence
\[
                   R^\#=M^\#,
                   \qquad Z\subseteq M^\#.
 \tag{A.40.3}
\]

Let \(C^\#\) be the preimage in \(E^\#\) of
the specified section \(C\leq M\rtimes C\).
In \(D\), the actual elements \(s,r\) are
\(M\)-conjugate to the pure elements
\(s_C,r_C\) of this section. Indeed, in each
regular wreath orbit their permutations consist
only of two-cycles, respectively three-cycles.
Every base cycle product is one, since the actual
elements have orders two and three. Solving the
coordinate conjugator equations therefore removes
their entire base components, as in B.57.

Lift these conjugators to \(M^\#\). The
corresponding conjugates \(S,R_3\) of the actual
\(s,r\) in \(E^\#\) lie in \(C^\#\), have
orders dividing two and three, and project to
the named \(s_C,r_C\). The universal property
of \(\Pi=C_2\amalg C_3\) gives a homomorphism
\(j:\Pi\to C^\#\) sending \(s,r\) to these
elements. Its projection to \(C\) is the
canonical quotient map. Thus \(j(R)\subseteq Z\)
is central. Since \(R=[R,\Pi]\), it follows
that \(j(R)=1\). The map factors to a continuous
section \(i:C\to C^\#\), proving
\[
                         E^\#=M^\#\rtimes C.
 \tag{A.40.4}
\]

The action of \(C\) on \(M^\#\) factors
through \(A_5\). On \(M\), it is the specified
factor permutation. An automorphism of the
perfect group \(M^\#\) inducing the identity
on its central quotient \(M\) is the identity:
its difference is a homomorphism from \(M^\#\)
to the central abelian kernel. This proves the
assertion for every element in the kernel of
\(C\to A_5\), and also proves uniqueness of the
lifted permutation action. Put
\[
 Q^\#=M^\#\rtimes A_5,
 \qquad E^\#=Q^\#\times_{A_5}C.
 \tag{A.40.5}
\]
For this chosen section the images of the actual
\(s,r\) in \(Q^\#\) are \(M^\#\)-conjugate
to their pure permutations: the section was built
from precisely those chosen conjugates.

\emph{A positive lift retaining the whole congruence quotient.}\par
Let \(\beta_C\) be A.31's actual order-two
automorphism of \(C\), with the literal frames
of A.33. It literally inverts \(\tau_C\), and
its induced action on \(A_5\) is the standard
one. Write \(c_C\) for genuine standard
conjugation and put
\[
                          h_C=\beta_Cc_C.
\]
The map \(h_C\) acts identically on \(A_5\),
fixes \(\tau_C\), and preserves the positive
\(s,r\) classes with the prescribed pure frames.
The formula
\[
 \psi^\#=\operatorname{id}_{M^\#}\rtimes h_C
        =(\operatorname{id}_{Q^\#},h_C)
       \quad\text{on }E^\#=Q^\#\times_{A_5}C
 \tag{A.40.6}
\]
is a continuous automorphism. It fixes the actual
\(\tau\) in \(E^\#\) literally: its finite
\(Q^\#\)-coordinate is unchanged and its
\(C\)-coordinate \(\tau_C\) is fixed. No
assumption about the order of a finite cusp lift
is needed.

To check \(s\), write \(A=q^\#(s)\) and let
\(d_s\in C\) be its specified positive frame,
so that \(h_C(s_C)=d_s s_Cd_s^{-1}\).
Its \(A_5\)-image centralizes \(\bar s\),
because \(h_C\) acts identically on \(A_5\).
The map
\[
          C_{Q^\#}(A)\longrightarrow C_{A_5}(\bar s)
 \tag{A.40.7}
\]
is onto. Conjugate \(A\) to its pure permutation
by the \(M^\#\)-conjugator above; every pure
permutation in \(C_{A_5}(\bar s)\) then
centralizes it, and conjugating back preserves
the \(A_5\)-coordinate. Choose a lift
\(e_s\in C_{Q^\#}(A)\) of the image of
\(d_s\). The pair \((e_s,d_s)\) lies in the
actual group \(E^\#\) and conjugates \(s\)
to \(\psi^\#(s)\), retaining exactly the
specified \(C\)-frame. The identical argument
applies to \(r\).

We have therefore constructed an automorphism of
the actual central refinement
\(\Pi/[W,\Pi]\) preserving all three positive
peripheral classes with literal conjugators.
Apply the closed-normal criterion A.3/A.34.
It gives a full positive peripheral lift inducing
\[
             \psi_D=(\operatorname{id}_Q,h_C)
                \quad\text{on }D=Q\times_{A_5}C,
 \tag{A.40.8}
\]
and retaining the specified quotient conjugators.
Its full cusp conjugator lies in \(W\). An inner
correction by its inverse makes \(\tau\) literally
fixed without changing (A.40.8) or the \(C\)-frames.
Because \(W\subseteq R\subseteq U_{12}\), the
retained pure-frame conditions also ensure purity
of these full frames. The standard modular
identification gives a full \(h\in\mathcal H\).

Finally set \(a=hc\). The kernel of \(q\) and
\(R\) are invariant under genuine \(c\), whose
action on \(Q\) is denoted by \(c_Q\). Thus
\[
       a_D=(c_Q,h_Cc_C)=(c_Q,\beta_C),
       \qquad a_D^2=1.
 \tag{A.40.9}
\]
Both coordinates induce the same action on
\(A_5\), so this is an actual fiber-product
automorphism. It has order exactly two, for
example from its nontrivial inversion on the
order-three quotient of \(\Pi^{\mathrm{ab}}\).
Its outer action on \(F_D\) also has order two,
since its multiplier is \(-1\). Its projection
at level twelve is the one in A.31, so A.24
gives \(\Delta(a)=1\). The character descends
to this actual outer image by projection to
\(G_{12}\); hence this involution is not
conjugate to the standard one in that image.
The full lift's square has only been shown to
vanish on \(D\), as asserted.
\end{proof}

\statement{A.41}{Proved}{The parity-one involution is compatible through a whole semisimple tripod tower}
Keep \(R,C,\widehat K\) from A.40. Let
\(\mathscr S\) range over finite nonabelian simple
groups not isomorphic to
\(\operatorname{PSL}_2(\mathbb F_p)\) for any
prime \(p\). For each such group take all kernels
of epimorphisms from the antipodal-pair tripod
fillings of \(\widehat K\) onto that group.
For every finite collection of targets let
\(q_J:\Pi\twoheadrightarrow Q_J\) be the joint
induced quotient over \(A_5\), and put
\[
 W_J=R\cap\ker q_J,
 \qquad W_* =\bigcap_J W_J,
 \qquad D_* =\Pi/W_*.
 \tag{A.41.1}
\]
There is one full \(a\in\widehat{GT}_0\) with
\(\lambda(a)=-1\) and \(\Delta(a)=1\) whose
actual action on \(D_*\) has order two.
It retains A.33's literal \(C\)-frames and
induces standard conjugation on every \(Q_J\).
This is an assertion about the specified joint
tower; no cofinality or full order-two lift is
asserted.

\begin{proof}
For each fixed finite target there are finitely
many kernels, and the whole family is invariant
under \(\widehat{GT}_0\). A retained cusp cannot
map trivially: killing one of the three cusps
would make the remaining two-cusp quotient
procyclic, which cannot surject onto a nonabelian
simple group. Thus the support is exactly the
retained triple, so every \(A_5\)-orbit is free.
The full-product argument in A.40 now gives
\[
 Q_J=M_J\rtimes A_5,
 \qquad M_J=\prod_{i\in I_J}S_i,
 \qquad q_J(R)=M_J.
 \tag{A.41.2}
\]
The last equality follows as before because none
of the \(S_i\) occurs as a composition factor
of a finite quotient of \(C\).

The only additional input for repeating A.40's
central-refinement argument is finiteness of the
universal central cover of \(M_J\). Here is the
elementary finiteness check. The integral bar
complex of a finite group has finite rank in
each degree, so its second homology is finitely
generated. Transfer from the trivial subgroup
shows that multiplication by the group order
annihilates this homology.\cite[Section~2 and
Corollary~3.4]{ABrown2010} It is therefore finite.
For a perfect group the kernel of its universal
central cover is precisely this second homology,
by the free-presentation construction and Hopf's
formula. The cover is consequently finite and
perfect. Thus every step of (A.40.3)--(A.40.8)
applies to (A.41.2), with no restriction on the
prime divisors of its Schur multiplier.

For a finite collection \(J\), let
\(\mathscr C_J\subseteq\mathcal H\) consist of
the full normalized positive elements whose
actual action on \(\Pi/W_J\) is
\((\operatorname{id}_{Q_J},h_C)\) and whose
two literal \(C\)-frames are the prescribed
ones. These sets are nonempty by the preceding
lifting argument. They are closed: quotient
actions and the normalized frame coordinates
are continuous, and equality of an infinite
quotient action or frame is the intersection
of its closed finite-coordinate equalities,
as in A.31--A.33. For finitely many \(J\),
their union is a single finite collection whose
solution set lies in the intersection. Compactness
of \(\mathcal H\) therefore gives
\[
                         h\in\bigcap_J\mathscr C_J.
\]
Put \(a=hc\). Its action on each joint quotient
is \((c_{Q_J},\beta_C)\), of order two.
These are the actual images of one element,
so their compatibility does not depend on a
surjectivity assertion for torsion transition
maps. Passing to the inverse limit proves the
claim. The image remains nontrivial and of
parity one already at level twelve.
\end{proof}

The finite-family argument controls direct products
of simple factors in free \(A_5\)-orbits, including
their actual scalar central refinements. It does
not control arbitrary nonsimple chief layers of
the remaining kernel. In particular it neither
annihilates every \(H^2(C_2,M)\) in A.19 nor
solves the full nonabelian norm equation.

\statement{A.42}{Proved}{Every full Grothendieck--Teichm\"uller involution has its prescribed conjugator at every prime}
Let \(\alpha=(-1,f)\in\widehat{GT}\) be a
nonidentity involution, and let \(u\in F'\) be
the unique full solution of \(u=f\alpha(u)\)
from A.37. Then
\[
 D=1+24c_2(u)\in
          (\widehat{\mathbb Z}^{\times})^2.
 \tag{A.42.1}
\]
In particular there is a full unit
\(\mu\in\widehat{\mathbb Z}^{\times}\) with
\(\mu^2=D\). For every prime \(p\), and either
local root \(\mu_p\), the prescribed pair
\((\mu_p,u_p)\) belongs to the pro-\(p\)
Grothendieck--Teichm\"uller group and its
normalized peripheral automorphism conjugates
\(c_p\) to \(\alpha_p\).

Any choice of a full square root in (A.42.1)
therefore gives a compatible conjugator on the
maximal pronilpotent quotient
\(\prod_p F_2^{(p)}\). No corresponding full
profinite conjugator or arbitrary relative-level
conjugator is asserted. In particular the statement
does not settle the dyadic case for an arbitrary
\(\mathcal P\)- or \(\widehat{GT}_0\)-involution.

\begin{proof}
A.37 proves the full pentagon for \(u\), so its
image \(u_p\) satisfies the pentagon in the
actual pro-\(p\) pure braid group. C.33 gives
\[
              1+24c_2(u_p)\in(\mathbb Z_p^\times)^2
                  \quad\text{for every }p.
\]
The coefficient is compatible with these quotient
maps. Since
\(\widehat{\mathbb Z}=\prod_p\mathbb Z_p\),
choosing a local square root at each prime yields
(A.42.1) and the stated full unit.

C.33 also gives the actual pro-\(p\) hexagons
for either root, and pro-\(p\) Frattini gives
an automorphism. Equivalently A.38 now applies
at every prime. The full identity
\(u=f\alpha(u)\) projects to the same
intertwining identity, so
\[
       g_p(x)=x^{\mu_p},\qquad
       g_p(y)=u_p^{-1}y^{\mu_p}u_p,
       \qquad \alpha_p g_p=g_p c_p.
\]
Their product is a continuous automorphism of
the maximal pronilpotent quotient and has the
claimed peripheral action and intertwining
property. This proves exactly the asserted
quotient conclusions.
\end{proof}

The additional input over A.38 is C.33's integral
metabelian hexagon identity, whose first possible
\(p\)-denominator excludes nonzero nonsquare
and nonunit parameters, together with C.32's
relative braid-torsion argument excluding zero.
Thus no local square-root obstruction remains
for this canonical word. The remaining full
condition is the third-cusp equation of A.39,
or equivalently the full hexagon for an appropriate
root. A positive answer to Furusho's Question~14
in its \emph{existential-root} variant for this
prescribed \(u\) would now suffice without a
separate square-root assumption. The literal
prescribed-root implication is refuted in C.34;
these two quantifiers must be distinguished.
The independently chosen local signs still give
no full identity: the exact \(S_4\) example
in A.39 already demonstrates that distinction.

\statement{A.43}{Proved}{The actual first scalar central refinement also carries the compatible parity-one involution}
For any finite family in A.40 or A.41, put
\[
       W_1=[W,\Pi],\qquad W_2=[W_1,\Pi].
\]
Then
\[
                         W_2=W_1.
 \tag{A.43.1}
\]
There is a full negative element of parity one
whose actual action on \(\Pi/W_1\) has order
two and retains the prescribed \(C\)-frames.
These refined involutions may be chosen as the
images of one full element for all finite
families simultaneously. Thus A.41 extends to
the actual quotient
\[
            \Pi\Big/\bigcap_J[W_J,\Pi].
 \tag{A.43.2}
\]
The square is still not shown to vanish in the
full group.

\begin{proof}
Write
\[
 E_1=\Pi/W_1=M_1\rtimes C,
 \qquad M_1=R/W_1,\qquad Z_1=W/W_1.
\]
A.40 shows that \(M_1\) is finite and perfect,
\(Z_1\) is central in \(E_1\), and the
\(C\)-action on \(M_1\) factors through
\(A_5\). Since \(M_1/Z_1=M\) is centerless,
\(Z(M_1)=Z_1\).

Consider \(E_2=\Pi/W_2\) and \(M_2=R/W_2\).
The map \(M_2\to M_1\) is a central extension
with kernel \(A=W_1/W_2\), central in all of
\(E_2\). The canonical finite perfect image
\(P\) of the universal central cover of
\(M_1\) satisfies \(M_2=PA\) and is normal
in \(E_2\). Since \(R=[R,\Pi]\),
\[
                  M_2=[M_2,E_2]\subseteq P.
\]
Thus \(M_2=P\) is finite and perfect.

The preimage \(Z_0=W/W_2\) of \(Z_1\) is
central in \(M_2\). Indeed for an element of
that preimage its commutator map on \(M_2\)
takes values in the central kernel \(A\)
and is a homomorphism; perfectness makes it
trivial. Hence \(Z_0\) is finite abelian, with
\[
                    0\longrightarrow A\longrightarrow Z_0
                         \longrightarrow Z_1\longrightarrow0.
 \tag{A.43.3}
\]

Conjugation on \(Z_0\) factors through \(C\),
and in fact through \(A_5\). To check the
second assertion, let \(v\in\ker(C\to A_5)\)
and lift its pure section element in \(E_1\)
to \(E_2\). Its conjugation on \(M_2\)
induces the identity on \(M_1\). An
automorphism of the perfect central extension
\(M_2\) inducing the identity on \(M_1\)
is the identity, by the same central-homomorphism
argument. This proves the asserted factorization.

The resulting \(A_5\)-action on (A.43.3)
is the identity on \(A\) and on \(Z_1\),
because these are central in \(E_2\) and
\(E_1\), respectively. The automorphisms of
this abelian extension which act identically
on both subquotients form the abelian group
\(\operatorname{Hom}(Z_1,A)\): their difference
from the identity is precisely such a homomorphism.
Perfectness of \(A_5\) therefore makes its
action trivial. Thus \(Z_0\) is central in
all of \(E_2\), so the image of
\([W,\Pi]=W_1\) in \(E_2\) is trivial.
This proves \(W_1\subseteq W_2\); the reverse
inclusion is automatic, giving (A.43.1).

Next use A.40's positive automorphism
\(\psi_1=\operatorname{id}_{M_1}\rtimes h_C\)
of \(E_1\), and let \(c_1\) be the genuine
standard action. Put \(b_1=\psi_1c_1\).
Its square is the identity on \(M_1\), since
\(\psi_1\) fixes that subgroup pointwise, and
on \(C\), since \(h_Cc_C=\beta_C\) is an
involution. Both \(b_1\) and \(c_1\) invert
the actual cusp, so \(b_1^2\) fixes \(\tau\).

An automorphism \(\theta\) of \(E_1\) which
fixes \(M_1\) pointwise and induces the identity
on \(C\) differs from the identity by a
homomorphism \(d:C\to Z_1\). Indeed, compare
\(\theta(i(v))\) and \(i(v)\) for the chosen
section. Their ratio centralizes \(M_1\),
hence lies in \(Z_1\), which is central in
all of \(E_1\). The multiplication rule then
makes the ratio an ordinary homomorphism.
But
\[
                  C^{\mathrm{ab}}=\Pi^{\mathrm{ab}}=C_6
\]
is generated by the image of \(\tau\), because
\(R=[R,\Pi]\). If \(\theta\) also fixes the
actual cusp, this homomorphism is trivial.
Apply the observation to \(\theta=b_1^2\)
to obtain \(b_1^2=1\) in the actual refined
quotient.

By (A.43.1), the central refinement required by
A.3 for \(\Pi/W_1\) is \(\Pi/W_1\) itself.
Thus its positive peripheral automorphism
\(\psi_1\) passes the exact lifting criterion
directly. Normalize the full ordinary lift to
fix the cusp and retain its \(C\)-frames, then
multiply by genuine \(c\). This gives the
asserted full element with refined involutory
image and parity one.

Finally \(\psi_1\) is uniquely characterized
by fixing \(M_1\) pointwise, inducing \(h_C\)
on \(C\), and fixing the actual cusp: the
ratio of two such maps is the same central
homomorphism just considered. On enlarging a
finite family the higher refined action
therefore projects to this unique lower one.
Alternatively the nonempty closed full-lift
sets have the finite-intersection property
under finite unions, exactly as in A.41.
Compactness produces one full element with
all the refined images, proving (A.43.2).
\end{proof}

A.43 eliminates every further scalar central
refinement of these particular finite perfect
families. Its argument uses the actual center
and the perfect \(A_5\)-action. A general
noncentral elementary-abelian layer need not
have this structure; its \(C_2\)-norm class
has not been computed here.

\statement{A.44}{Proved}{The canonical substitution passes all dihedral tests and the parity-relative pro-prime generation test}
Let \(\alpha=(-1,f)\in\mathcal P\) be a nonidentity
involution and let \(u\in F'\) be its canonical
parameter, \(u=f\alpha(u)\). Put
\[
 y'=u^{-1}yu,\qquad \psi(x)=x,\quad\psi(y)=y',
 \qquad V=\ker\bigl(F\to C_2,\ x,y\mapsto-1\bigr).
\]
The substitution \(\psi\) is onto every finite
dihedral quotient in which \(x,y\) are generating
reflections. For every prime \(p\), it induces
an automorphism of
\[
                         G_p=F/O^p(V),
 \tag{A.44.1}
\]
and that automorphism conjugates the induced
standard action to the induced \(\alpha\)-action.
This is a generation and unmarked automorphism
statement. It does not assert preservation of
the third cusp by \(\psi\).

\begin{proof}
Write \(E=F\rtimes\langle c\rangle\), where
the named \(c\) acts by \(\alpha\). The three
full section classes have representatives
\[
                         C=c,\qquad X=xc,\qquad Y=y'c.
\]
They are distinct by their mod-two abelian labels
\(0,x,y\). The first two cusp edges have endpoints
\(C,X\) and \(C,Y\). C.28's three-cusp cycle
therefore makes the third cusp's endpoints
\(X,Y\). Every section centralizer in \(F\)
is trivial.

\emph{The induced dihedral action.}\par
For an odd prime \(p\), use
\[
 \begin{gathered}
 q:F\twoheadrightarrow D_{2p},\qquad
 q(x)=a,\quad q(y)=at,\\
 a^2=t^p=1,\quad ata=t^{-1},\qquad
 q(f)=t^\ell,\quad q(u)=t^k.
 \end{gathered}
\]
Then \(q\alpha\) sends \(x,y\) to
\(a,at^{2\ell+1}\). This pair generates because
\(\alpha\) is an automorphism. Every generating
reflection pair of this form is obtained from the
named pair by an automorphism of \(D_{2p}\)
fixing \(a\). Thus \(\ker q\) is
\(\alpha\)-invariant, and the induced map sends
\(t\) to \(t^{2\ell+1}\). Its square is one,
so \(2\ell+1=1\) or \(-1\) modulo \(p\).
The substituted pair fails to generate exactly when
\[
                   q(y')=at^{2k+1}=a.
 \tag{A.44.2}
\]

If the induced action inverts \(t\), then
\(\ell=-1\), and the coboundary equation gives
\(2k=\ell=-1\). The induced action is
\(\operatorname{Int}(a)\). Under the isomorphism
\[
 D_{2p}\rtimes_{\operatorname{Int}(a)}C_2
       \longrightarrow D_{2p}\times C_2,\qquad
 gc^\epsilon\longmapsto(ga^\epsilon,\epsilon),
\]
both \(X\) and \(Y\) map to the central
involution \((1,1)\). Every conjugate of either
class has the same image. The third cusp's two
reflections would therefore have product one in
this quotient, although their product is a
conjugate of \(z\) or \(z^{-1}\), whose image
has order \(p\). This is impossible. Hence
\(\alpha\) acts identically on \(D_{2p}\)
and \(q(f)=1\).

\emph{The pure negative degree-\(p\) cover.}\par
Suppose (A.44.2) still holds and put
\(U=q^{-1}(\langle a\rangle)\). It is
\(\alpha\)-stable. Label \(U\backslash F\)
by \(j\in\mathbb F_p\), using \(Ut^j\).
The right actions of \(x,y,z\) are respectively
\[
                     j\mapsto-j,\qquad j\mapsto1-j,
                             \qquad j\mapsto j-1.
 \tag{A.44.3}
\]
Above each of \(x,y\) there is one width-one
cusp and \((p-1)/2\) width-two cusps. Above
\(z\) there is one width-\(p\) cusp. The
cover has \(p+2\) cusps and Euler characteristic
\(-p\), hence genus zero. Thus \(U\) is the
full profinite sphere group of rank \(p+1\).

The automorphism preserves every \(x\)- and
\(y\)-cusp individually, since its finite action
is identity and their finite frames are
\(q(1)=q(f)=1\). The \(z\)-cusp is unique,
so is preserved regardless of its frame. Every
meridian is inverted up to \(U\)-conjugacy.
Consequently C.28 applies to this pure negative
action on the \((p+2)\)-cusp group.

For clarity, when \(p=3\) the five cusp
representatives can be taken as
\[
        x,\quad yx^2y^{-1},\quad u^{-1}yu,\quad y^2,\quad z^3.
 \tag{A.44.4}
\]
The third is \(U\)-conjugate to the discrete
representative \(yxyx^{-1}y^{-1}\). Their
respective real reflections are \(c\),
\(ycy^{-1}\), \(c\), \(ucu^{-1}\), and
a local \(z\)-reflection adjusted by an even
power of \(z\) to have finite component \(a\).
They all lie in \(U\rtimes\langle c\rangle\).
The width-two frames are
\(y\alpha(y)^{-1}\) and \(f\), both with
trivial finite image; the two width-one frames
are one. The final reflection is \(v_zc\)
with \(q(v_z)=a\), giving an actual \(U\)-frame
for the unique third-cusp lift. No triviality
of that frame is assumed.

\emph{The impossible incidence graph.}\par
The full \(C\)-class splits into precisely
\(p\) \(U\)-classes, denoted \(C_j\).
They are represented by \(gcg^{-1}\) for
\(Ug=Ut^j\). Every such conjugate lies in
the subgroup, and two have the same
\(U\)-class exactly when their left cosets
agree, because \(C_F(c)=1\).
Each of the full \(X,Y\)-classes contributes
exactly one \(U\)-class: its finite component
is \(a\), and a conjugate belongs to the
subgroup precisely when its conjugator has
image in \(C_{D_{2p}}(a)=\langle a\rangle\).
The section vertices are therefore exactly
\(C_j\) together with \(X,Y\).

The width-one \(x\)-edge joins \(C_0\) to
\(X\); the width-one \(y\)-edge, represented
by \(y'\), joins \(C_0\) to \(Y\).
The unique \(z\)-edge has odd width \(p\),
so its endpoints remain in the opposite full
classes \(X,Y\), each of which has just one
\(U\)-class. An even power of \(z\) adjusts
a base reflection into the subgroup, and its
product with \(z^p\) gives the other endpoint.
These edges form a triangle.

For every nonzero pair \(\{j,-j\}\), the
width-two \(x\)-edge joins \(C_j,C_{-j}\).
Its exact local identity is
\(x^2c=xcx^{-1}\): the two conjugators are
\(g\) and \(gx\). For \(y\), use its
actual real reflection \(ucu^{-1}\), which
inverts \(y\). The two endpoints above a
width-two orbit have conjugators \(gu\)
and \(gyu\). Their labels are
\[
                         j+k,\qquad 1-j+k.
\]
By (A.44.2), these are negatives of each other.
The fixed \(y\)-orbit becomes label zero;
every other orbit gives exactly the same
nonzero pair as an \(x\)-edge. Thus the
incidence graph is one triangle together with
\((p-1)/2\) disjoint digons. C.28 requires
one cycle through every edge and vertex, a
contradiction.

This excludes (A.44.2). Failure of generation
in any finite dihedral reflection quotient
would mean that the odd number \(2k+1\)
is not a unit modulo its rotation order.
That failure is detected at an odd prime,
which has just been excluded. This proves
the first assertion.

\emph{The relative Frattini consequence.}\par
A Schreier basis of \(V\) is
\[
                         A=x^2,\qquad B=yx^{-1},\qquad C=xy.
\]
On its abelianization conjugation by \(x\)
sends \(A\) to \(A\), \(B\) to \(C-A\),
and \(C\) to \(A+B\). For odd \(p\),
the invariant plane is spanned by \(A,B+C\),
and the sign line by \(A+B-C\). These three
vectors form a basis, with determinant \(-2\).
Here \(B+C\) is the class of \(y^2\).

The substitution preserves \(V\) and
\(O^p(V)\), so induces an endomorphism of
\(G_p\). Its closed image \(H\) is generated
by \(x,y'\). The intersection with the free
pro-\(p\) kernel \(V^{(p)}\) contains
\(x^2,(y')^2\), whose Frattini images are
the two invariant basis vectors: \(u\in V\)
does not change the abelianized square.
Quotienting by that invariant plane gives
\(D_{2p}\), with \(x,y\) reflections and
\(xy\) having nonzero sign coordinate.
The dihedral result makes \(H\) onto this
quotient, so \(H\cap V^{(p)}\) also covers
the sign line. It therefore covers the entire
Frattini quotient. Frattini gives
\(H\cap V^{(p)}=V^{(p)}\), and hence
\(H=G_p\). Surjectivity makes the induced
endomorphism an automorphism.

At \(p=2\), (A.44.1) is the maximal pro-two
quotient of \(F\), since its kernel and
quotient are pro-two. Its abelianized
substitution is the identity, so ordinary
pro-two Frattini gives the same conclusion.
Finally the full identity \(u=f\alpha(u)\)
gives \(\alpha\psi=\psi c\) on the two
generators and therefore on every stated
quotient. This proves the conjugacy assertion
with precisely its stated unmarked scope.
\end{proof}

\statement{A.45}{Proved}{The canonical word determines the only possible full multiplier}
For \(\alpha,u\) as in A.44, use the universal
full dihedral quotient
\[
 q:F\twoheadrightarrow D_\infty^{\wedge}
      =\widehat{\mathbb Z}\rtimes C_2,\qquad
 q(x)=a,\quad q(y)=at,\quad ata=t^{-1}.
\]
Write \(q(u)=t^k\) and define
\[
                          d_\alpha=2k+1.
 \tag{A.45.1}
\]
Then \(k\in2\widehat{\mathbb Z}\),
\(d_\alpha\in\widehat{\mathbb Z}^{\times}\),
and \((d_\alpha)_2\equiv1\pmod4\).
The involution acts identically on
\(D_\infty^{\wedge}\). Its full conjugacy
criterion is the single equation
\[
 x^{d_\alpha}u^{-1}y^{d_\alpha}u
                   \sim_F(xy)^{d_\alpha}.
 \tag{A.45.2}
\]
If a full peripheral conjugator exists, its
multiplier is exactly \(d_\alpha\) or
\(-d_\alpha\). The same equivalence holds
in \(\widehat{GT}_0\) or \(\widehat{GT}\)
when \(\alpha\) belongs to the respective group.

For \(\alpha\in\widehat{GT}\), full conjugacy
would in particular imply
\[
                         d_\alpha^2=1+24c_2(u).
 \tag{A.45.3}
\]
Its dyadic component already holds by A.42.
The other components of (A.45.3), and the full
equation (A.45.2), are not proved here.

\begin{proof}
The quotient is obtained by imposing
\(x^2=y^2=1\). Its rotation subgroup is the
full procyclic group generated by \(t=xy\),
and its commutator subgroup is
\(2\widehat{\mathbb Z}\) inside that subgroup.
Since \(u,f\in F'\), write
\(q(u)=t^k\), \(q(f)=t^\ell\) with
\(k,\ell\in2\widehat{\mathbb Z}\).

The defining kernel is preserved by
\(\alpha\), which sends \(a\) to itself
and \(t\) to \(t^{1+2\ell}\).
Its square is one, so \((1+2\ell)^2=1\).
At every odd prime the root is positive,
because A.44 proves identity already modulo
that prime and the only roots of one in
\(\mathbb Z_p\) are \(1,-1\).
At two, \(1+2\ell\equiv1\pmod4\)
also selects the positive root. Thus
\(\ell=0\), and the entire induced action
is the identity. A.44 makes \(2k+1\)
nonzero modulo every odd prime; it is odd
at two. This proves all assertions about
\(d_\alpha\).

For any full unit \(\mu\), the normalized
substitution with parameter \(u\) sends the
two dihedral generators to
\(a,at^{2k+1}\), since a unit power of a
reflection is that same reflection. Their
product is \(t^{d_\alpha}\). The required
third-cusp power is \(t^\mu\). Rotations in
\(D_\infty^{\wedge}\) are conjugate exactly
when their exponents agree or are opposite:
its rotation subgroup is abelian and the
other coset acts by one global inversion.
Hence any full solution has
\(\mu=d_\alpha\) or \(-d_\alpha\).

The third-cusp equations for opposite signs
are equivalent, because the product of the
two inverses is conjugate to the inverse
of the product. A.39 now proves (A.45.2)
as the exact full conjugacy criterion,
including its two membership descents.
The necessary coefficient relation A.4
gives (A.45.3) when a solution exists.

Finally, for a full \(\widehat{GT}\)-involution,
A.42 supplies a prescribed peripheral
conjugator in the maximal pro-two quotient.
Project it to the pro-two dihedral quotient
\(\mathbb Z_2\rtimes C_2\). The same rotation
calculation gives \(\mu_2=\pm(d_\alpha)_2\),
while A.38 gives \(\mu_2^2=1+24c_2(u)_2\).
This proves the stated dyadic component.
\end{proof}

A.45 removes the variable full multiplier
from A.39's finite-quotient test: one may
test the fixed unit \(d_\alpha\) in every
characteristic finite quotient. It does not
replace that full conjugacy test by maximal
pro-prime tests. The comparison between its
dihedral value and A.42's other local roots
is a remaining concrete compatibility issue.

\statement{A.46}{Proved}{Generation over the full elementary abelian two quotient}
Let \(\alpha,u,\psi\) be as in A.44, and put
\[
 V_4=\ker\bigl(F\longrightarrow C_2^2\bigr)
        =\overline{F'F^2}.
\]
For every prime \(p\), the canonical substitution
\(\psi(x)=x\), \(\psi(y)=u^{-1}yu\) induces a
continuous automorphism of
\[
                         F/O^p(V_4).
 \tag{A.46.1}
\]
It intertwines the induced involution with the standard
negative involution. This asserts generation and unmarked
relative conjugacy; it does not assert that this automorphism
preserves the third cusp.

\begin{proof}
The substitution is the identity on \(F^{\mathrm{ab}}\),
so preserves \(V_4\), and then preserves its fully invariant
pro-\(p\) residual. Set \(Q=F/V_4=C_2^2\).
For odd \(p\), the Frattini module
\[
 M=V_4/\overline{V_4'V_4^p}
\]
is the homology of the finite \(Q\)-cover of the
bouquet of two circles. Its cellular complex in degrees
one and zero is
\[
 \mathbb F_p[Q]^2\longrightarrow\mathbb F_p[Q],
 \qquad (a,b)\longmapsto
       a(\bar x-1)+b(\bar y-1).
\]
All four characters of \(Q\) take values in
\(\mathbb F_p\). On the trivial character the displayed
differential is zero, and on each nontrivial character
it has rank one. Consequently
\[
             \dim M^{\mathrm{triv}}=2,
 \qquad \dim M^\chi=1\quad(\chi\ne1).
 \tag{A.46.2}
\]
The induced map \(\psi_M\) is \(Q\)-equivariant,
because \(\psi\) induces the identity on \(Q\).
The inclusion \(V_4\hookrightarrow F\) identifies
the trivial summand with \(H_1(F,\mathbb F_p)\):
transfer followed by inclusion multiplies by four,
and the two spaces both have dimension two.
Naturality and \(\psi_{\mathrm{ab}}=1\) therefore
make \(\psi_M\) the identity on this summand.

Each remaining one-dimensional summand is detected by
an explicit finite dihedral quotient. Write
\[
 D_{2p}=\langle a,t\mid a^2=t^p=1,
                          \ ata=t^{-1}\rangle.
\]
For the character \(\chi(\bar x)=\chi(\bar y)=-1\),
use \(x\mapsto a,y\mapsto at\). The image of
\(V_4\) is the nontrivial rotation subgroup,
with precisely this character for conjugation.
A.44 says that its composite with \(\psi\)
is still onto, so \(\psi_M\) is nonzero on this
character line.

For \((\chi(\bar x),\chi(\bar y))=(-1,+1)\),
use \(x\mapsto a,y\mapsto t\). Since
\(u\in F'\), its image is a rotation; it commutes
with \(t\). The composite with \(\psi\) therefore
retains both generators \(a,t\). For the remaining
character \((+1,-1)\), use
\(x\mapsto t,y\mapsto a\); the composite retains
\(t\) and a conjugate reflection, and again is
onto. In each case \(V_4\) maps onto the rotation
subgroup, since the square of its rotation generator
is nontrivial for odd \(p\). Equivariance shows that
this nonzero map on \(M\) factors through the
indicated character line. Thus \(\psi_M\) is
nonzero on every one-dimensional summand, and is
an automorphism of \(M\).

The induced endomorphism of the free pro-\(p\)
group \(V_4^{(p)}\) is now onto by the Frattini
criterion. Together with the identity on \(Q\),
this proves surjectivity of (A.46.1). A finitely
generated profinite group is Hopfian, so the induced
endomorphism is an automorphism.

At \(p=2\), the group (A.46.1) is exactly the
maximal pro-two quotient of \(F\): it is an
extension of a pro-two group by \(C_2^2\), and
every pro-two quotient of \(F\) restricts to
a pro-two quotient of \(V_4\). Its Frattini map
is the identity, which proves the assertion there.
Finally \(\alpha\psi=\psi c\) is already an
identity on the full group, so descends to every
quotient just considered.
\end{proof}

\statement{A.47}{Proved}{A genuine noncentral two layer in the compatible parity-one tower}
Let \(V\) be the augmentation submodule of the natural
permutation module \(\mathbb F_2^9\) for \(A_9\), and put
\(S=V\rtimes A_9\). Enlarge the tower of A.41 by all
tripod quotients of \(\widehat K\) onto \(S\) whose
three retained cusps form an antipodal pair and one
other vertex and have images of order seven. This
family is nonempty. Its actual joint quotient with
\(C=\Pi/R\), and its actual scalar central refinement,
carry compatible negative involutory images with
\(\Delta=1\). One full negative element realizes these
images simultaneously with the whole tower of A.41.

The same assertion holds for all such regular-head
tripod families with targets \(W\rtimes T\), where
\(T\) is a finite nonabelian simple group excluded
from A.31's congruence composition factors and \(W\)
is a nontrivial finite irreducible \(T\)-module.
No nonemptiness assertion is made for every possible
module. The explicit group \(2^8\rtimes A_9\) proves
that the enlarged tower contains a noncentral
chief layer of characteristic two. No full
involution or cofinality assertion follows.

\begin{proof}
\emph{The concrete noncentral quotient.}\par
The module \(V\) has dimension eight. For every
nonzero even-weight vector, a suitable three-cycle
changes it in exactly two positions. The orbit of
a resulting pair-vector spans \(V\), so \(V\)
is simple. It is nontrivial, and therefore has
neither invariants nor a trivial quotient. Thus
\(S\) is perfect and centerless.

Since nine is odd, \(\mathbb F_2^9\) splits as the
trivial module plus \(V\). Eckmann--Shapiro gives
\[
 H^1(A_9,\mathbb F_2^9)
       =H^1(A_8,\mathbb F_2)=0,
 \qquad H^1(A_9,V)=0.
\]
Here both alternating groups are perfect; the
cohomological comparison is
\citet[Proposition~3.3]{ABrown2010}.
Use the explicit triple from B.57,
\[
 X=(1234567),\qquad Y=(1254389),
 \qquad Z=(XY)^{-1},
\]
and its affine lifts
\[
 \widetilde X=(0,X),\qquad
 \widetilde Y=(e_1+e_2,Y),\qquad
 \widetilde Z=(e_1+e_9,Z).
 \tag{A.47.1}
\]
They multiply to one and all have order seven.
The exact certificate
\path{check_A47_affine_A9.py}
checks these identities, all 256 simultaneous
vector conjugators, and the generation of the
181440-element simple head. The two translation
coordinates in (A.47.1) are not a simultaneous
coboundary \(((1-X)v,(1-Y)v)\).

The group generated by the first two lifts projects
onto \(A_9\), and its intersection with \(V\) is
a submodule. If that intersection were zero,
it would be a complement, hence conjugate to the
standard complement by \(H^1(A_9,V)=0\). This is
exactly the excluded coboundary condition. Its
intersection is therefore \(V\), proving generation
of \(S\). The script additionally counts 1024
order-seven lift pairs and 256 simultaneous
coboundaries. It does not enumerate \(S\) itself;
the preceding argument proves its generation.

\emph{Invariant joint bundles.}\par
Take all kernels in the stated finite family.
They are permuted by \(\widehat{GT}_0\): cusp
supports and the antipodal relation are preserved,
and unit powers preserve order seven. Group the
kernels by their induced simple head kernel
\(\widehat K\to A_9\). For a fixed head, their
joint image \(P\) is a subgroup of
\(V^n\rtimes A_9\) surjecting onto \(A_9\).
Identifications of the head may twist the module;
for \(A_9\) those twists preserve the natural
module, and more generally they remain nontrivial
simple modules. Its kernel \(N\) is a submodule
of a finite semisimple sum of nontrivial modules.
Consequently
\[
                 [N,A_9]=N,
 \qquad N^{A_9}=0.
\]
It follows that \(P\) is perfect and centerless.
No splitting of \(P\) is needed.

Distinct head kernels give independent bundles.
A proper subdirect product would, by finite
Goursat, give a common nontrivial quotient. That
quotient is perfect and has a nonabelian simple
quotient. A homomorphism from a product to a
nonabelian simple group factors through one
factor, since the factor images commute and are
normal. Thus such a common quotient would identify
two head kernels. Induction proves the full product.

Every head has three nonzero cusp supports. Its
\(A_5\)-stabilizer is trivial by the antipodal-triple
calculation in B.57. Inducing one whole bundle
per head orbit therefore gives the actual finite
quotient
\[
              q:\Pi\twoheadrightarrow M\rtimes A_5,
\]
where \(M\) is a product of finite perfect
centerless bundles, regularly permuted on each
\(A_5\)-orbit. Its kernel is precisely the
intersection of the invariant family.
The quotient \(M/q(R)\) is a finite quotient of
the open normal subgroup \(\widehat K/R\) of
\(C\), so contains no \(A_9\) composition factor.
It is perfect, and its remaining composition
factors are cyclic of order two. It is therefore
trivial. Hence \(q(R)=M\), and the actual joint
quotient is \(M\rtimes C\).

\emph{The actual lift.}\par
The proof of A.40 uses only finite perfectness
of \(M\) and the regular permutation of its
bundles. Its universal central cover is finite
by the argument in A.41. For
\(W=R\cap\ker q\), the actual refinement
\(\Pi/[W,\Pi]\) has finite perfect kernel
\(M^\#\); the regular two- and three-cycles
allow its actual torsion generators to be
conjugated into the pure \(C\)-section.
Thus that proof gives
\[
 \Pi/[W,\Pi]
       =M^\#\rtimes C
       =(M^\#\rtimes A_5)\times_{A_5}C.
\]
The automorphism \(\operatorname{id}_{M^\#}\rtimes h_C\)
fixes the actual cusp and retains the full
\(C\)-frames. A.3 gives the full ordinary lift,
whose product with \(c\) has the required
negative involutory image. The proof of A.43
also applies: the bundle is centerless and
perfect, and the remaining action on its two
central subquotients factors through the perfect
group \(A_5\). Thus
\( [[W,\Pi],\Pi]=[W,\Pi]\), and the actual
first central refinement already has the
compatible order-two action.

For a general finite union of targets
\(W\rtimes T\), the kernels over a fixed simple
head lie in a semisimple sum of nontrivial
irreducible modules, possibly in different
characteristics. The same perfectness,
centerlessness, Goursat, and regular-orbit
arguments apply. A perfect quotient with only
cyclic composition factors is trivial, so the
composition-factor argument still gives
\(q(R)=M\). Joining A.41's simple families only
adds new heads or redundant simple quotients
of existing bundles. Finite unions therefore
satisfy all the same conclusions. The closed
solution sets retain exactly the prescribed
\(C\)-action and frames and the identity on
each finite quotient. Compactness, as in A.41,
provides one full element realizing their
entire compatible tower. Its full square is
still not determined.
\end{proof}

\statement{A.48}{Proved}{Nonregular orbit-one simple quotients also admit compatible images}
The parity-one tower of A.47 can be enlarged by
all finite invariant families of simple quotients
of \(\Pi=C_2\amalg C_3\) excluded from the
congruence composition factors of A.31, together
with their actual scalar central refinements.
These quotients have \(A_5\)-orbit length one
on their simple \(\widehat K\)-head kernels;
regular permutation of head factors is not
required for this assertion. One explicit
nonempty family consists of all epimorphism
kernels \(\Pi\to A_9\) for which \(\tau\)
has order fifteen.

One full negative element still realizes the
compatible involutory images with \(\Delta=1\).
This statement does not provide a full
involution or control a nonperfect noncentral
kernel over a mixed solvable head.

\begin{proof}
\emph{A central-extension lemma.}\par
Let \(R=[R,\Pi]\), \(C=\Pi/R\), and let
\(h_0\) be a full positive automorphism
preserving \(R\) and fixing \(\tau\) literally.
Write \(h_C\) for its induced action and retain
its full \(s/r\) conjugators. Consider a central
extension
\[
                 1\longrightarrow A\longrightarrow
                 C^\#\longrightarrow C\longrightarrow1
\]
with \(A\) finite abelian. Choose arbitrary
lifts \(b_s,b_r\) of the named elements of
\(C\). We claim that \(A\) injects into
\((C^\#)^{\mathrm{ab}}\), and that \(h_C\)
has a lift fixing \(A\) and \(b_sb_r\), with
the prescribed \(C\)-frames lifted.

Pull back to \(P=C^\#\times_C\Pi\). Its
normal subgroup
\[
                  R_0=\{(1,r):r\in R\}
\]
affords the same conjugation action as \(R\)
in \(\Pi\), so \(R_0=[R_0,P]\). Hence
\(P^{\mathrm{ab}}=(C^\#)^{\mathrm{ab}}\).
Set \(S=(b_s,s)\), \(T=(b_r,r)\),
\(a=b_s^2\), and \(b=b_r^3\). The central
extension of the free profinite product has
the presentation
\[
 P=\langle A\text{ central},S,T\mid
                  S^2=a,\ T^3=b\rangle.
 \tag{A.48.1}
\]
Indeed the amalgamated product of the two finite
central extensions over \(A\) has quotient
\(\Pi\); its natural map to \(P\) is the
identity on both \(A\) and that quotient,
so is an isomorphism.

The relations in the abelianization do not
kill a nonzero element of \(A\). To see this
without an integral-presentation assumption,
embed \(A\) in a divisible torsion abelian
group and choose a half of \(a\) and a third
of \(b\). They and \(A\) generate a finite
abelian group to which (A.48.1) maps, injectively
on \(A\). This proves the claimed injection.

Lift the two full \(h_0\)-conjugators from
\(\Pi\) to \(P\), conjugate \(S,T\) by
those lifts, and fix \(A\). Presentation
(A.48.1) gives an automorphism \(\phi\) of
\(P\): it covers the automorphism \(h_0\)
and is the identity on its central kernel.
For \(r\in R_0\), its image differs from
the canonical element of \(R_0\) over
\(h_0(r)\) by a central character
\(\chi:R\to A\). Equivariance makes
\(\chi\) invariant under \(\Pi\), so
\(R=[R,\Pi]\) kills it. Thus \(\phi\)
descends to an automorphism \(k\) of
\(C^\#\), fixing \(A\) and retaining the
specified \(C\)-frames.

Finally \(\phi(ST)(ST)^{-1}\in A\), since
\(h_0(\tau)=\tau\). The automorphism \(\phi\)
is the identity on \(P^{\mathrm{ab}}\):
it fixes \(A\) and conjugates \(S,T\).
Injectivity of \(A\) into this abelianization
forces that difference to be one. Hence
\(k(b_sb_r)=b_sb_r\) literally. No splitting
of \(C^\#\to C\) has been assumed.

\emph{The nonregular quotient.}\par
There is an actual epimorphism \(\Pi\to A_9\)
with
\[
 \begin{aligned}
 s&=(12)(34)(56)(78),\\
 r&=(138)(249)(576),\\
 \tau&=(149)(23758).
 \end{aligned}
 \tag{A.48.2}
\]
The exact certificate
\path{check_A48_orbit_one.py}
checks the orders and enumerates the generated
181440-element group. Its \(\widehat K\)
image is a nontrivial normal subgroup of
\(A_9\), hence all of \(A_9\); otherwise
\(A_9\) would be a quotient of \(A_5\).
The kernel on \(\widehat K\) is \(\Pi\)-normal,
so has orbit length one. Each of the twelve
cusps maps to a conjugate of \(\tau^5\),
of order three, distinguishing this head from
the tripod heads of A.47.

Take all \(\Pi\to A_9\) kernels with cusp
order fifteen. The finite family is
\(\widehat{GT}_0\)-invariant, and its joint
simple image is \(T=A_9^n\). The kernel
\(R\) maps onto \(T\), because its quotient
\(C\) has no \(A_9\) composition factor.
For \(W=R\cap\ker(\Pi\to T)\), the actual
joint quotient is therefore
\[
                             \Pi/W=T\times C.
\]
The same construction works for any finite
invariant family of the simple targets in
the statement.

\emph{The actual central product.}\par
Put \(E=\Pi/[W,\Pi]\) and
\(A=W/[W,\Pi]\). As in A.40, relative
perfectness makes
\(M=R/[W,\Pi]\) the finite perfect image
of the universal central cover of \(T\).
Let \(C^\#\) be the preimage of the pure
\(C\)-section in \(T\times C\).
An element of \(C^\#\) acts trivially on
\(T\), hence on its perfect central cover
\(M\): the difference would be a homomorphism
from \(M\) to the abelian central kernel.
Consequently
\[
                 E=MC^\#,
 \qquad [M,C^\#]=1,
 \qquad M\cap C^\#=A.
 \tag{A.48.3}
\]
Choose \(m_s,m_r\in M\) lifting the actual
finite components and write
\(s=m_sb_s\), \(r=m_rb_r\) with
\(b_s,b_r\in C^\#\). Then
\(\tau=m_sm_rb_sb_r\).
The lemma, applied to the full \(h_0\) already
provided by A.33, gives \(k\) fixing
\(A\) and \(b_sb_r\). Thus
\(\psi_E=\operatorname{id}_M\,k\) is a
well-defined automorphism of (A.48.3).
It fixes the actual cusp. Its prescribed
conjugators lie in \(C^\#\) and commute
with \(m_s,m_r\), so retain exactly the
required \(C\)-frames on the actual
\(s,r\). A.3 therefore gives the full
ordinary positive lift.

Genuine \(c\) preserves \(E,M,C^\#\).
The negative map \(b=\psi_Ec_E\) has square
one on \(M\) and on \(C\), and fixes \(A\)
after squaring. Its square on \(C^\#\)
is therefore a central character \(C\to A\).
It fixes \(\tau\) and \(M\), hence fixes
\(b_sb_r\). Since \(C^{\mathrm{ab}}=C_6\)
is generated by \(\tau_C\), the character
is zero. Thus \(b^2=1\) already on the
actual refinement \(E\).

There is no additional scalar central layer.
Writing \(W_1=[W,\Pi]\) and
\(W_2=[W_1,\Pi]\), the group \(R/W_2\)
is a finite perfect central cover of \(M\).
The preimage of \(A\) in it is central,
by perfectness. The preimage of \(C^\#\)
also centralizes that perfect cover, since
it acts trivially on \(M\). These subgroups
generate \(\Pi/W_2\), so \(W/W_2\) is
central and \(W_1=W_2\).

\emph{Compatibility with the earlier tower.}\par
For a finite regular family of A.47, its
\(\widehat K\)-head supports have size three;
the orbit-one heads have all twelve supports.
They are therefore independent. In particular
\(q(W_{\mathrm{reg}})=T\), and
\[
 q([W_{\mathrm{reg}},\Pi])=[T,T]=T.
\]
The stabilized actual refinement of the
regular family may consequently be used
as the new base in the lemma: its residual
kernel is still relative-perfect and
surjects onto \(T\). No composition-factor
exclusion for this enlarged base is needed.
The argument then retains both finite
requirements and their literal frames.
Finite unions and the same closed-set
compactness argument as A.41 give the
claimed single full element. Only its
specified quotient images have been
proved involutory.
\end{proof}

\statement{A.49}{Proved}{A nonperfect mixed-head layer has an exact Clifford refinement}
Let \(H=\operatorname{AGL}_1(\mathbb F_7)\), let
\(V\subset\mathbb F_2^7\) be its augmentation module,
and put \(G=V\rtimes H\). The complete family of
\(\Pi\to G\) epimorphism kernels consists of two
kernels, exchanged by \(c\). Write their joint
quotient as
\[
 \overline G=\overline V\rtimes\overline H,
 \qquad \overline V=V_+\oplus V_-,
 \qquad \overline H=C_7^2\rtimes C_6.
\]
For \(R,C\) as in A.31--A.34 and
\(W=R\cap\ker(\Pi\to\overline G)\), one has
\[
 q(R)=\overline V,
 \qquad
 \Pi/[W,\Pi]
   \simeq\bigl((P_+\times P_-)\rtimes\overline H\bigr)
                       \times_{\overline H}C,
 \tag{A.49.1}
\]
where \(P_\pm\simeq2_+^{1+6}\) are the specified
Clifford central covers of \(V_\pm\). Moreover
\[
                       [[W,\Pi],\Pi]=[W,\Pi].
 \tag{A.49.2}
\]
A full negative element with \(\Delta=1\) realizes
an involutory image on this actual refinement,
compatibly with the entire earlier tower of A.48.
The residual image \(\overline V\) is noncentral
and nonperfect. This is a particular mixed
solvable-head calculation, not a general norm
vanishing or a full order-two construction.

\begin{proof}
\emph{The finite heads and their cohomology.}\par
Label the seven points by \(\mathbb F_7\). There
is an actual epimorphism \(\Pi\to G\) with
\[
 s(t)=-t,\quad r(t)=2t+1,
 \qquad \widetilde s=(0,s),
 \quad \widetilde r=(e_2+e_5,r).
 \tag{A.49.3}
\]
The certificate \path{check_mixed_affine7.py}
verifies the lift orders two and three, enumerates
all \(2688=2^6\cdot42\) generated elements,
and verifies cusp order twelve.

The module \(V\) is irreducible. Under \(C_7\)
it has two irreducible three-dimensional
constituents, interchanged by a primitive
\(C_6\)-element. Its dot product is alternating
and nondegenerate, so it is self-dual. Also
\[
                         H^1(H,V)=0.
 \tag{A.49.4}
\]
Indeed \(\mathbb F_2^7=\mathbb F_2\oplus V\),
and Shapiro identifies its first cohomology
with \(H^1(C_6,\mathbb F_2)\). The trivial
summand maps isomorphically to this group
by restriction, since \(H^{\mathrm{ab}}=C_6\).
The comparison is the same
\citet[Proposition~3.3]{ABrown2010} used in A.47.

A cocycle of \(\Pi\) with this module is a pair
\((v_s,v_r)\) satisfying the order-two and
order-three norm equations. On the seven points, \(s\) has three transpositions
and one fixed point, while \(r\) has two three-cycles
and one fixed point. The norm kernels on \(V\)
therefore have dimensions three and four. Thus
the solution space has dimension seven; the
coboundaries have dimension six because \(V^H=0\). Both counts are checked exactly
in the certificate. Thus
\[
                         \dim H^1(\Pi,V)=1.
 \tag{A.49.5}
\]
Its nonzero class is (A.49.3). A proper subgroup
of \(G\) surjecting onto \(H\) would intersect
the simple \(V\) trivially, hence be a complement;
(A.49.4) makes it conjugate to the standard one.
This is precisely the zero cohomology class.

The subgroup \(V=O_2(G)\) is characteristic,
since \(H\) has no nontrivial normal two-subgroup.
There are two head kernels \(\Pi\to H\),
according as the multiplier of \(r\) is two
or four. Affine conjugation normalizes the
fixed point of \(s\) to zero and the nonzero
\(r\)-displacement to one. Every automorphism
of \(H\) is inner: its characteristic \(C_7\)
and the defining conjugation relation force
identity on \(C_6\), and the remaining scalar
and translation choices are inner. Equation
(A.49.5) therefore gives just one generating
\(G\)-kernel above each head. This is the full
finite family, so it is \(\widehat{GT}_0\)-invariant.
The formula \(c(r)=sr^{-1}s\), \(c(s)=s\)
exchanges its two members.

Their joint head has order \(294=7^2\cdot6\),
with inverse primitive \(C_6\)-characters on
its two \(C_7\)'s. The two modules \(V_+,V_-\)
are nonisomorphic simples for this head. The
joint image has the entire kernel
\(\overline V\): otherwise one simple summand
would have zero projection from that kernel,
making \(G\), of order 2688, a quotient of
the 294-element head. Semisimplicity of this
direct sum proves the asserted joint quotient.

\emph{The actual residual image.}\par
The congruence-kernel image in \(\overline H\)
is exactly \(C_7^2\). The quotient \(C_6\)
is congruence, whereas a finite solvable
congruence quotient has only two and three
as composition primes. For completeness,
\(\operatorname{SL}_2(\mathbb Z/p^n)\) is
perfect for \(p\geq5\): elementary unipotents
are commutators with a diagonal unit whose
square minus one is invertible, and they
generate the group. The remaining local
factors have only two- and three-primary
composition factors. The two different
\(C_6\)-characters then give both independent
seven-directions.

Normality forces the congruence-kernel image
in \(\overline G\) to be
\(\overline V\rtimes C_7^2\). Its nilpotent
residual is \(\overline V\): the quotient is
nilpotent, while every nilpotent image must
kill \([\overline V,C_7^2]=\overline V\),
since its normal two- and seven-parts commute.
The residual commutes with an epimorphism by
A.34's compact comparison. Hence
\(q(R)=\overline V\), as claimed.

The retained positive \(h_C\) acts identically
on \(\overline H\). At a principal congruence
level divisible by twelve the kernel's image
in \(\overline H\) is a seven-group, so this
quotient factors through the relative pro-seven
level, where A.33's positive action is identity.
The nonzero class in (A.49.5) does not inflate
from \(C\), since it is nonzero on \(R\).
Inflation is injective; therefore
\[
 H^1(C,V_+)=H^1(C,V_-)=0,
 \qquad H^1(C,\overline V^\vee)=0,
 \tag{A.49.6}
\]
where \(\vee\) denotes finite Pontryagin duality.

\emph{The actual central-kernel bound.}\par
Put \(E=\Pi/[W,\Pi]\), \(A=W/[W,\Pi]\),
and \(M=R/[W,\Pi]\). Then \(M\) is a
central extension of \(\overline V\) by
\(A\), and \(M=[M,E]\). Work in an arbitrary
finite quotient retaining \(\overline G\),
with images \(M_i,A_i\). If \(A'_i\) is
the image of \(A_i\) in \(M_i^{\mathrm{ab}}\),
duality gives an exact sequence
\[
 0\longrightarrow\overline V^\vee
  \longrightarrow(M_i^{\mathrm{ab}})^\vee
  \longrightarrow(A'_i)^\vee\longrightarrow0.
\]
The middle invariants vanish by relative
perfectness; \(A'_i\) is a trivial \(C\)-module.
The invariant long exact sequence and
(A.49.6) force \(A'_i=0\). Thus
\(A_i=[M_i,M_i]\).

Its commutator pairing is an alternating,
bilinear, \(C\)-invariant pairing on
\(\overline V\). Hence \(A_i\) is a quotient
of \((\bigwedge^2\overline V)_C\), of
dimension two. Each summand contributes one
direction: sharp two-transitivity gives
endomorphism dimension two for the seven-point
permutation module, and its splitting
\(1\oplus V\) gives
\(\operatorname{End}_H(V)=\mathbb F_2\).
Its nonzero invariant form is the dot product.
The mixed tensor summand has no coinvariants,
since the first \(C_7\) has no invariants on
\(V_+\) and is trivial on \(V_-\).
Consequently
\[
                          |A_i|\leq4,
 \qquad |M_i|\leq2^{14}.
 \tag{A.49.7}
\]
No finiteness of the original \(A\) was assumed.

\emph{The Clifford quotient attains the bound.}\par
Take seven symbols \(e_i\) with \(e_i^2=1\)
and \(e_ie_j=-e_je_i\) for \(i\ne j\).
The even-subset elements, with their central
signs, form a group \(P\) of order 128,
quotient \(V\), and square function
\[
                         q(v)=\operatorname{wt}(v)/2\pmod2.
\]
It has 36 zeros, so \(P\simeq2_+^{1+6}\).
Permuting the seven symbols is an honest
\(S_7\)-action, with the reordering sign included.
Thus \(P\rtimes H\) is an actual central
extension of \(G\).

The lifts in (A.49.3) lift further with their
same orders: \(s\) is pure and the indicated
\(r\)-lift has positive central sign. The
certificate \path{check_mixed_affine7_clifford.py}
checks the action law, generation of all
5376 elements, and cusp order twenty-four.
Pair this map with its precomposition by \(c\).
The resulting finite image is
\[
                     \widehat G_0=(P_+\times P_-)
                                           \rtimes\overline H.
\]
Its \(R\)-image is all \(P_+\times P_-\):
it covers \(\overline V\), and commutators
of lifts supported in one summand supply
its central sign independently. Therefore
\(W\) maps to the two central signs and
\([W,\Pi]\) is killed. The actual \(M\)
thus maps onto a group of order \(2^{14}\).
Bound (A.49.7) forces an isomorphism; any
additional kernel would survive in a finite
quotient retaining this one and violate the
bound. The two actual projections now give
(A.49.1), injectively by this identification
and surjectively by correcting a finite
\(P\)-coordinate with an element of \(R\).

\emph{Frames and the constructed involution.}\par
The map \((\operatorname{id}_{\widehat G_0},h_C)\)
is an automorphism of this actual fiber product
and fixes its actual cusp. The finite \(s\)
components are pure, so their centralizer maps
onto the corresponding \(\overline H\)-centralizer.
The actual order-three \(r\)-lift is
\(P_+\times P_-\)-conjugate to the pure lift
by the coprime complement theorem, giving
the same surjectivity for \(r\). The exact
certificate displays a conjugator in the
plus factor. Every retained literal
\(C\)-frame therefore glues to a finite
centralizing element, exactly as in A.40.

Genuine \(c\) exchanges the two finite
Clifford factors. The corresponding negative
automorphism on \(E\) is
\((c_{\widehat G_0},\beta_C)\), of order two.
At this point A.3 gives a full lift of its
positive companion on \(\Pi/W\); the next
paragraph is needed to retain its entire
specified action on \(E\).

\emph{Scalar stabilization and the full lift.}\par
Write \(W_1=[W,\Pi]\), \(W_2=[W_1,\Pi]\),
\(A_2=W_1/W_2\), \(M_2=R/W_2\), and
\(Z_0=W/W_2\). Again work first in finite
quotients retaining the finite Clifford image.
The same dual argument applied to
\[
 0\longrightarrow\operatorname{im}(A_2)
   \longrightarrow M_2^{\mathrm{ab}}
   \longrightarrow M^{\mathrm{ab}}=\overline V
   \longrightarrow0
\]
gives \(M_2^{\mathrm{ab}}=\overline V\).
The commutator of a lift of the old central
\(A_1=A\) with \(M_2\) is a \(C\)-invariant
pairing \(A_1\otimes\overline V\to A_2\).
Both central factors are trivial modules;
for a fixed first lift the central commutator
is a homomorphism in the second argument,
so factors through the displayed abelianization.
Since \((\overline V)_C=0\), this pairing
vanishes. Thus \(Z_0\) centralizes \(M_2\).
Since the old center is \(M'\) and
\(A_2\subseteq M_2'\), one has \(Z_0=M_2'\).
It is an elementary abelian quotient of
\(\bigwedge^2\overline V\), with action
factoring through \(\overline H\).

Let \(J=C_7^2\rtimes C_3\) be its odd
normal subgroup. The two subquotients
\(A_2,A_1\) of \(Z_0\) are trivial
\(J\)-modules. Maschke makes its action
on \(Z_0\) trivial. But
\[
                    \dim(\bigwedge^2\overline V)_J=2.
\]
In each summand the three pairs of opposite
seven-characters are cycled by \(C_3\),
leaving one direction; independent translations
kill the mixed directions. The certificate
checks alternating-square relation rank 14
using only the translation and \(r\),
without the even generator, and verifies
\(|J|=147\). The quotient \(A_1\) already
has dimension two, so \(A_2=0\).
Passing over finite quotients proves (A.49.2).
Now apply A.3 with \(N=W_1\): its actual
central refinement is itself. The specified
framed automorphism of \(E\) therefore lifts
to a full positive \(h\), and \(hc\) has
the asserted involutory image and \(\Delta=1\).

\emph{Retaining the earlier towers.}\par
A finite requirement from A.47--A.48 has
perfect finite heads, whereas \(\overline G\)
is solvable. They have no nontrivial common
quotient. On \(R\), the earlier finite
perfect image and \(\overline V\) likewise
have no common quotient; hence the earlier
kernel still maps onto \(\overline V\).
Taking its relative commutator preserves
that image, since
\([\overline V,\overline G]=\overline V\).
Its stabilized residual is relative-perfect
by A.43 as applied in A.47 and by A.48.
Using that refinement as the new base,
(A.49.5) again proves (A.49.6), since the
nonzero class still detects the residual.
Every argument above therefore retains
that finite requirement and its frames.
Finite unions and compactness give one
full element for the entire asserted tower.
Its full square is not determined.
\end{proof}

\statement{A.50}{Proved}{A nonzero torsion-frame class has an explicit correction and norm cancellation}
The tower of A.49 admits a further compatible
parity-one involutory image on an actual central
refinement whose residual module has nonzero
\(H^1(C_2,-)\). Its retained odd head is
\(H=\operatorname{Heis}_7\rtimes C_6\), and
its residual module is the pair of
21-dimensional \(\mathbb F_2\)-Schr\"odinger modules
with opposite central characters. The actual
central kernel is the additive group \(\mathbb F_8\).
A compensating \(C_{12}\)-extension of the head
is essential for the order-two torsion lift.
The scalar central refinement stabilizes,
so the specified involution is the image of a
full negative element with \(\Delta=1\),
compatibly with all preceding towers.

The retained s-frame determines a nontrivial
possible linear correction. It is explicitly
conjugation by an odd central-head lift, and
its norm under \(c\) is one. No general norm
vanishing, cofinality, or full involution is
asserted.

\begin{proof}
\emph{The actual odd head and its frame.}\par
Use symmetric Heisenberg coordinates over
\(\mathbb F_7\), with \(d\) modulo six:
\[
 \begin{aligned}
 &(x,y,z,d)(u,v,w,e)\\
 &\quad=\bigl(x+3^du,\ y+3^{-d}v,\
 z+w+(x3^{-d}v-y3^du)/2,\ d+e\bigr).
 \end{aligned}
\]
The actual \(\Pi\)-quotient has
\[
 s=(0,0,0,3),\qquad r=(1,1,2,2),
 \qquad\tau=(6,6,2,5).
 \tag{A.50.1}
\]
The certificate \path{check_heisenberg7_head.py}
enumerates all 2058 generated elements and
verifies
\[
                 s^2=r^3=1,
 \qquad \tau^6=(0,0,6,0).
\]
This head is characteristic: it is the quotient
by \(\gamma_3(\Pi')(\Pi')^7\), since the
open commutator subgroup \(\Pi'\) is free
profinite of rank two and its maximal
class-two exponent-seven quotient has order
\(7^3\). The displayed epimorphism has exactly
that order above \(C_6\).
Genuine \(c\) induces
\[
                 c_H(x,y,z,d)=(4y,2x,-z,-d).
 \tag{A.50.2}
\]
It fixes \(s\) and sends \(r\) to \(sr^{-1}s\).

The head belongs to the retained relative
pro-seven tower, so \(h_C\) is identity on
it. Its congruence-kernel image is the entire
\(\operatorname{Heis}_7\), by the solvable
congruence-prime argument of A.49. The finite
s-centralizer is
\(Z(\operatorname{Heis}_7)\times C_6\).
The retained s-frame \(\beta\) has trivial
\(\Pi\)-abelianization: A.31 chose zero complete
Kummer character, the frame formula
LM7\_PUSH E.65.1 gives zero abelianization
for its arithmetic frame, and A.33 retains
its image in the congruence quotient \(C_6\).
Consequently
\[
                           \beta_H=z^j
                         \quad(j\in\mathbb F_7).
 \tag{A.50.3}
\]
The construction below works for every \(j\);
no new choice of this retained frame is made.

\emph{The two modules and the torsion direction.}\par
Put \(k=\mathbb F_8\), choose a primitive
seventh root \(\zeta\), and index a basis
of \(k^7\) by \(\mathbb F_7\). Define
\[
 \rho(x,y,z,d)e_a
  =\zeta^{-z+xy/2+y3^da}\,e_{3^da+x}.
 \tag{A.50.4}
\]
As an \(\mathbb F_2\)-module this is \(V_+\)
of dimension 21. Its center acts by \(\zeta\),
so every invariant \(\mathbb F_2\)-subspace
is \(k\)-linear. A diagonal Heisenberg generator
has seven distinct eigenvalues, and the shift
permutes the seven eigenlines transitively.
They generate all matrix units. Hence the
module is irreducible and
\(\operatorname{End}_H(V_+)=k\).
Let \(V_-\) be its twist by \(c_H\).
The two nontrivial central-character Frobenius
orbits are distinct, so these modules are
nonisomorphic. They pair by the symmetric,
nondegenerate Fourier form
\[
                 B(v,w)=\sum_{a,b\in\mathbb F_7}
                                  v_a\zeta^{-2ab}w_b,
 \tag{A.50.5}
\]
which satisfies
\(B(\rho(g)v,\rho(c_Hg)w)=B(v,w)\).
Its trace to \(\mathbb F_2\) identifies
\(V_-\) with the finite dual of \(V_+\).

The order-two norm kernel has dimension four
and its image dimension three over \(k\).
The order-three norm kernel also has dimension
four. There are no common invariants, so
\[
 \dim_kH^1(\Pi,V_+)=1,
 \qquad \dim_kH^1(\langle s\rangle,V_+)=1.
 \tag{A.50.6}
\]
The class given by \(d_+(s)=e_0\),
\(d_+(r)=0\) has nonzero restriction to the
s-factor, hence is a generator of both spaces.
Thus the torsion-frame space is genuinely
nonzero here.

Since \(\tau^6\) has nontrivial central
character, \(\rho(\tau)-1\) is invertible.
There is a unique coboundary adjustment with
\(d_+(\tau)=0\). In the polynomial basis
\(k=\mathbb F_2[t]/(t^3+t+1)\), its exact
adjustment and normalized vectors are
\[
 \begin{gathered}
 b=(6,5,0,6,3,5,5),\\
 v_s=v_r=(1,0,5,5,5,5,0),
 \qquad B(v_s,v_s)=1.
 \end{gathered}
 \tag{A.50.7}
\]
Set \(d_-(g)=d_+(c(g))\). It has the same
normalized generator vectors and zero cusp
value, with the twisted module action.
The exact field arithmetic and all norm
equations are checked in
\path{check_heisenberg7_frame.py}.

Write \(\overline V=V_+\oplus V_-\).
The resulting joint affine quotient is
\(\overline V\rtimes H\). Indeed finite
\(H^1(H,V_+)\) vanishes by coprime cohomology
of its normal odd Heisenberg group with no
invariants. A nonzero \(\Pi\)-cocycle therefore
generates the full affine extension, as in
A.49; all nonzero classes differ by a scalar
of \(k^\times\) and have the same kernel.
The two nonisomorphic modules give independent
summands. They form an invariant two-kernel
family over the characteristic head. To see
that twisting the head introduces no other
21-dimensional module, identify its
Heisenberg restriction over \(k\); extensions
to \(C_6\) differ by a character into
\(k^\times\), and
\(\operatorname{Hom}(C_6,k^\times)=1\).

\emph{First cohomology of the retained base.}\par
The affine quotient is now proved full.
Its congruence-kernel image is
\(\overline V\rtimes\operatorname{Heis}_7\),
by normality and the nontrivial odd action.
Its nilpotent residual is \(\overline V\).
Thus \(q(R)=\overline V\), by A.34.
The nonzero class in (A.50.6) consequently
does not inflate from \(C\). Inflation is
injective and \(k\)-linear into a
one-dimensional space, giving
\[
 H^1(C,V_+)=H^1(C,V_-)=0,
 \qquad H^1(C,\overline V^\vee)=0.
 \tag{A.50.8}
\]
This sequence of deductions uses an actual
residual image, rather than treating the
finite odd quotient as a subgroup of \(C\).

\emph{The forced correction and the correct cover.}\par
The frame (A.50.3) acts on the two nonzero
s-cohomology lines by the inverse scalars
\(\zeta^{-j},\zeta^j\). The linear correction
is therefore precisely conjugation by
\(\beta_H\) on the finite affine quotient.
It fixes the normalized cusp, and its norm
under \(c\) is one. To retain this through
the actual central refinement, the following
cover keeps the torsion square explicitly.

Let \(P\) consist of \((v,w,t)\in k^7\times
k^7\times k\), with multiplication
\[
 (v,w,t)(v',w',t')
       =(v+v',w+w',t+t'+B(v,w')).
\]
It has order \(2^{45}\), center \(k\),
and quotient \(\overline V\). The paired
\(H\)-action fixes its center. The square
of the actual s-vector is one in \(k\),
so \(P\rtimes H\) alone would not permit
an order-two s-lift. Instead put
\[
 \widetilde H=H\times_{C_6}C_{12},
 \qquad
 \widehat G_1=(P\rtimes\widetilde H)
                    /\langle(1,\kappa)\rangle,
 \tag{A.50.9}
\]
where \(\kappa\) is the central involution
of \(\widetilde H\to H\), and \(1\) in
that relation is the indicated additive
central element of \(k\). The two s-square
contributions cancel. The r-lift has a
unique central correction making its order
three; it is the element 2 of \(k\) in
(A.50.7)'s coordinates.

In representatives with \(0\leq d<6\),
the additional central multiplication is
the \(C_{12}\) carry. Genuine \(c\) lifts as
\[
 (v,w,t,h)\longmapsto
 \bigl(w,v,t+B(v,w)+\mathbf1_{d(h)\ne0},c_H(h)\bigr).
 \tag{A.50.10}
\]
This is the hyperbolic flip together with
inversion on \(C_{12}\); it is an automorphism
of order two. The certificate verifies that
it fixes the actual \(s\), sends the actual
\(r\) to \(sr^{-1}s\), and inverts the
actual cusp. The specified odd-order lift
\(\widetilde\beta\) of \(z^j\) in
\(\widetilde H\) centralizes that cusp and
satisfies
\[
 c(\widetilde\beta)=\widetilde\beta^{-1},
 \qquad
 (\operatorname{Inn}(\widetilde\beta)c)^2=1.
 \tag{A.50.11}
\]
These identities are checked for all seven
possible frame values. The lift is specified
inside the central head extension; uniqueness
among all its conjugates in \(\widehat G_1\)
is not asserted.

\emph{The actual central refinement.}\par
Let \(W\) be the intersection of \(R\) with
the kernel of the joint affine quotient.
The finite Pontryagin-dual argument of A.49,
using (A.50.8), makes the actual central
kernel of \(\Pi/[W,\Pi]\) a quotient of
\((\bigwedge^2\overline V)_C\). Its dimension
is three over \(\mathbb F_2\). The self-terms
have no invariants, since the two central
character orbits contain no opposite pair.
The cross term is
\(\operatorname{End}_H(V_+)=k\).
The same dimension already holds for the
odd subgroup
\(J=\operatorname{Heis}_7\rtimes C_3\).
The exact certificate forms all 861
exterior-square coordinates of the
42-dimensional module and obtains relation
rank 858 using only its odd generators.

The map \(\Pi\to\widehat G_1\) realizes
all eight central elements on \(R\): its
image covers \(\overline V\), and commutators
of the two pure summand lifts realize every
value of the nondegenerate \(k\)-pairing.
It therefore attains the preceding upper
bound. Exactly as in A.49, the two actual
projections identify
\[
                    \Pi/[W,\Pi]
                         =\widehat G_1\times_H C.
 \tag{A.50.12}
\]
The cover (A.50.9) has thus been identified
with the required actual refinement.

Take the positive action
\((\operatorname{Inn}(\widetilde\beta),h_C)\)
on (A.50.12). Its s-conjugator
\(\widetilde\beta\) has precisely the
retained head image. For \(r\), the
centralizer of its actual order-three lift
surjects onto \(C_H(r)\), by the coprime
complement theorem for the normal two-group
\(P\). Multiplying \(\widetilde\beta\)
by a suitable such centralizer element
matches the retained r-frame. The cusp is
fixed literally. The negative companion
has order two by (A.50.11) and A.33.

The scalar refinement stabilizes by the
same argument as A.49. First the dual exact
sequence gives the next R-preimage
abelianization equal to \(\overline V\).
The pairing from the old central kernel
and \(\overline V\) into the new central
kernel vanishes because
\((\overline V)_C=0\). The next center is
therefore an exterior-square quotient.
Its action factors through \(H\), and
\(J\) acts trivially on its two central
subquotients, hence trivially by odd-order
semisimplicity. Its dimension is at most
three, already attained by the old center
\(k\). The new kernel is zero. Applied over
finite quotients retaining (A.50.9), this
proves \([[W,\Pi],\Pi]=[W,\Pi]\).
A.3 with this stabilized kernel now gives
a full lift retaining the entire action
(A.50.12), not just its quotient by the
central kernel.

\emph{Compatibility with A.49.}\par
Earlier perfect heads give no common
abelian quotient. The residual abelianization
of A.49 contributes only its two
six-dimensional modules, on which the
Heisenberg center is trivial. It cannot
surject onto either new irreducible
21-dimensional module, whose central action
is nontrivial. Finite Goursat and normality
therefore make the preceding stabilized
residual surject onto both new summands.
The inflation image in (A.50.6) is
\(k\)-linear; because its nonzero class
still detects that residual, it is zero.
Thus (A.50.8) holds for each enlarged base.
Its relative perfectness is precisely
A.49.2, together with A.43 and A.48 for
its predecessors. The entire proof retains
every prior finite requirement and its
literal frames. Finite unions and compactness
complete the stated compatible tower,
without determining the full square.
\end{proof}

\statement{A.51}{Proved}{An actual nonzero second scalar layer over a filiform odd head}
Let \(R\), \(C=\Pi/R\) be the initial congruence
pronilpotent base of A.31--A.34. There is a finite
quotient \(G\) of \(\Pi\), with a nonperfect
elementary-abelian residual image, such that, on putting
\[
 W=R\cap\ker(\Pi\to G),\qquad
 W_{i+1}=\overline{[W_i,\Pi]},\qquad W_0=W,
\]
one has
\[
 W_2\subsetneq W_1,\qquad W_1/W_2\simeq C_2,
 \qquad W_3=W_2.                                      \tag{A.51.1}
\]
Both scalar refinements are identified by explicit
finite central covers. Thus immediate scalar
stabilization is not a general consequence of
\(R=\overline{[R,\Pi]}\), or of the odd sensitivity
of the new irreducible chief modules.
This assertion does not supply a negative norm
obstruction. The single chosen head is not assumed
to be \(c\)-stable.

\begin{proof}
\emph{The actual head.}\par
Over \(\mathbb F_7\), put \(A=\mathbb F_7^4\), with
coordinates \((y,z,w,t)\), and define
\[
 D(y,z,w,t)=(0,y,z,w),\qquad E(x)=\exp(xD).
\]
The exponential is a polynomial of degree three,
whose denominators are units in \(\mathbb F_7\).
Define
\[
 N=A\rtimes C_7,\qquad
 (x,v)(u,w)=(x+u,v+E(x)w).
\]
A generator of \(C_6\) acts on \(x\) by multiplication
by \(3\) and on \(A\) by
\(B=\operatorname{diag}(5,1,3,2)\).
The equality \(BDB^{-1}=3D\) verifies the action.
Set \(H=N\rtimes C_6\), of order \(6\cdot7^5=100842\).
In coordinates \((x,y,z,w,t,d)\), choose
\[
\begin{gathered}
 s=(0,0,0,0,0,3),\qquad r_0=(0,0,0,0,0,2),\qquad
 n=(1,1,0,0,0,0),\\
 r=nr_0n^{-1}=(6,4,4,5,3,2).
\end{gathered}                                       \tag{A.51.2}
\]
These elements have orders two and three and generate
\(H\). Indeed, on \(N/[N,N]\), their translations
have nonzero components in both distinct \(C_6\)
characters \(3\) and \(5\). In either character,
a complement containing the pure \(s\) would have
zero conjugating vector, and hence zero r-translation,
a contradiction. The intersection of the generated
subgroup with \(N\) therefore surjects its Frattini
quotient and equals \(N\). The exact certificate
also enumerates all \(100842\) elements.

As in A.49, a solvable congruence quotient has no
composition factor of order seven, whereas
\(H/N=C_6\) is a congruence quotient. Thus the image
of \(K_{\rm cong}\) is \(N\). Since \(N\) is a
seven-group, \(H\) is a quotient of \(C\).

\emph{The irreducible modules.}\par
Let \(k=\mathbb F_8\), choose a primitive seventh
root \(\zeta\in k\), and put \(\ell=(1,0,0,1)\).
On \(U=k^7\), indexed by \(j\in\mathbb F_7\), set
\[
 \rho_U(x,v)e_j
       =\zeta^{\ell E(-(j+x))v}e_{j+x}.
                                                               \tag{A.51.3}
\]
The pure \(r_0\) acts by
\(e_j a\mapsto e_{2j}a^4\).
The identity \(\ell B^2=4\ell\) verifies this
semilinear action. It makes \(U\) an
\(\mathbb F_2[J]\)-module for the odd subgroup
\(J=N\rtimes C_3\).

The s-twist \(U^s\) uses the substitution
\[
 (x,y,z,w,t)\longmapsto(-x,-y,z,-w,t)
\]
and the same \(r_0\)-action. Define
\(V=U\oplus U^s\), with \(s\) interchanging the
two summands. This is a 42-dimensional
\(\mathbb F_2[H]\)-module. Its dual is obtained
by replacing \(\ell\) with \(-\ell\).

Here the irreducibility and endomorphism assertions
can be checked directly. The central top generator
\(T=(0,0,0,0,1,0)\) acts on \(U\) by \(\zeta\),
so an endomorphism commuting with it is \(k\)-linear.
The seven distinct A-characters
\(\ell E(-j)\) give every diagonal matrix unit
over \(k\), and the (x)-generator permutes them
transitively. Consequently \(U\) is irreducible
over \(\mathbb F_2\) as an N-module and
\(\operatorname{End}_N(U)=k\). The r-action
reduces this endomorphism field to \(\mathbb F_2\).
The s-twist is distinct: its central character
forces any possible Frobenius multiplier to be one;
the zero W-coordinate then forces the orbit
translation to be zero; the Y-coefficients (1)
and \(-1\) still differ. It follows that
\[
\begin{gathered}
 V\text{ is irreducible over }H,\qquad
 \operatorname{End}_H(V)=\mathbb F_2,\\
 \operatorname{End}_J(V)=\mathbb F_2\oplus\mathbb F_2.
\end{gathered}
\]
The central-character Frobenius orbits
\(\{1,2,4\}\) and \(\{-1,-2,-4\}\) are disjoint.
Thus \(V\not\simeq V^\vee\), and neither has
a self-pairing over \(J\).

Put \(\overline V=V\oplus V^\vee\).
The two trace pairings between the matching
summands span the J-invariant alternating forms,
and \(s\) interchanges them. Hence
\[
 \dim_{\mathbb F_2}(\textstyle\bigwedge^2\overline V)_J=2,
 \qquad
 \dim_{\mathbb F_2}(\textstyle\bigwedge^2\overline V)_H=1.
                                                               \tag{A.51.4}
\]
The first coinvariant space is the regular
\(\mathbb F_2[C_2]\)-module, not a trivial
two-dimensional module. Exact elimination on the
3486 exterior-square coordinates gives respective
ranks 3484 and 3485.

\emph{The actual affine quotient and inflation.}\par
The s-norm kernel on \(V\) has dimension 21,
the r-norm kernel has dimension 28, and the
coboundary space has dimension 42. Thus
\[
 \dim_{\mathbb F_2}H^1(\Pi,V)
  =\dim_{\mathbb F_2}H^1(\Pi,V^\vee)=7.
                                                               \tag{A.51.5}
\]
Normalize the s-value to zero. Actual nonzero
classes are given, in the first \(k^7\) block, by
\[
 d_+(r)=\zeta^3e_1,\qquad d_-(r)=e_1,
                                                               \tag{A.51.6}
\]
with zero in the other block. The finite certificate
checks the norm equations and proves that adjoining
each displayed vector increases the rank of the
normalized coboundary space from 21 to 22.

Since the odd normal subgroup \(N\) has no
invariant vectors,
\(H^1(H,V)=H^1(H,V^\vee)=0\).
Every nonzero class in (A.51.5) therefore gives
a full affine quotient: irreducibility makes a
subgroup over \(H\) either the whole affine group
or a complement, and a complement would give a
coboundary. The two distinct modules imply that
the paired map has image
\[
                         G=\overline V\rtimes H.
\]
Normality and \([\overline V,N]=\overline V\)
show that the image of \(K_{\rm cong}\) is
\(\overline V\rtimes N\). Its pronilpotent
residual is \(\overline V\): in any nilpotent
quotient the odd and two-primary images commute,
forcing the module to vanish. Thus the image of
\(R\) is precisely \(\overline V\).

The same argument works for every nonzero class
in (A.51.5), not only for (A.51.6). No such class
vanishes on \(R\). Inflation therefore gives
\[
 H^1(C,V)=H^1(C,V^\vee)=0,
 \qquad H^1(C,\overline V^\vee)=0.
                                                               \tag{A.51.7}
\]
For \(W=R\cap\ker(\Pi\to G)\), the joint image
contains the entire module kernel, so
\(\Pi/W=G\times_H C\).

\emph{Two explicit scalar covers.}\par
Let \(B_1,B_2\) be the trace pairings on the
matching summands of \(V\) and \(V^\vee\).
Define a two-group \(P\) by
\[
 (v,w,t)(v',w',t')
  =\bigl(v+v',w+w',
     (t_i+t'_i+B_i(v,w'))_{i=1,2}\bigr).
                                                               \tag{A.51.8}
\]
Its center is \(Z=\langle z_1,z_2\rangle\simeq C_2^2\).
The odd subgroup \(J\) acts trivially on \(Z\),
whereas \(s\) interchanges \(z_1,z_2\).
The displayed multiplication is preserved, giving
an honest action of \(H\). Set
\[
 z=z_1z_2,\qquad G_2=P\rtimes H,
 \qquad G_1=G_2/\langle z\rangle.
                                                               \tag{A.51.9}
\]

The two cocycles lift to an actual homomorphism
\(\Pi\to G_2\). The s-lift is pure and has
square one. The r-lift has cube in \(Z\);
adding that central cube to its central coordinate
makes its cube one, since three is odd. For the
specified vectors the correction is \((1,0)\).
The affine image is full. Its intersection with
\(P\) surjects \(\overline V\), and commutators
recover both central coordinates. Thus the map
onto \(G_2\) is surjective, the image of \(R\)
is \(P\), and the image of \(W\) is \(Z\).
Since
\[
              [Z,H]=\langle z\rangle\ne1,
              \qquad [\langle z\rangle,H]=1,
\]
this actual quotient already proves
\(W_2\subsetneq W_1\). No bound on an abstract
module-extension space is being substituted for
an actual quotient.

\emph{The exact first, second and third refinements.}\par
Work first in finite quotients retaining the
specified finite cover. At the first stage,
the Pontryagin-duality argument of A.49 and
(A.51.7) give abelianization \(\overline V\)
for the R-preimage. Its central kernel is a
quotient of
\((\bigwedge^2\overline V)_H\), of dimension
at most one. The cover \(G_1\) attains this bound,
so
\[
                     \Pi/W_1=G_1\times_H C.
                                                               \tag{A.51.10}
\]

The abelianization step at every following stage
has the same precise form. The kernel of the map
from the new R-preimage abelianization to the
previous abelianization is the image of the
\emph{new globally central kernel}. It is
therefore a trivial C-module, even when an older
central subgroup is no longer central in the
whole extension. Duality, (A.51.7), and relative
perfectness eliminate this new image. Inductively,
all these R-preimage abelianizations remain
\(\overline V\).

At the second stage, the preimage of the old
center centralizes the R-preimage. Its possible
central commutator is a C-invariant pairing from
the old trivial central module tensor
\(\overline V\), and
\((\overline V)_C=0\) makes it zero.
For completeness, the commutator lies in the
new central kernel, so the fourth lower-central
term vanishes; hence commutators within the
preimage of the old center vanish, which makes
the second argument of this pairing factor
through the stated abelianization.

The resulting center is an elementary-abelian
quotient of \(\bigwedge^2\overline V\), and its
action factors through \(H\). The odd \(J\)
acts trivially on both central subquotients, hence
on the extension by coprime semisimplicity.
Its dimension is at most two by (A.51.4),
already attained by \(G_2\). Thus
\[
              \Pi/W_2=G_2\times_H C,
              \qquad W_1/W_2\simeq C_2.
                                                               \tag{A.51.11}
\]

At the third stage the old center \(Z\) is a
nontrivial H-module, but it is N-trivial.
The possible new central pairing
\(Z\otimes\overline V\) therefore vanishes
on N-coinvariants. The next center is again
an exterior-square quotient on which the odd
\(J\) acts trivially. Its dimension is at most
two, already attained by the old center.
The new kernel is zero. Finite quotients
retaining (A.51.9) are cofinal for these bounds,
so the argument gives the asserted closed-group
equalities, including \(W_3=W_2\).

All finite identities, module ranks and the
actual torsion lifts are recorded in
\path{check_filiform7_scalar.py} and its output.
The strict second layer is a limitation of the
immediate-stabilization argument; it is not a
solution or obstruction to the full negative
norm equation.
\end{proof}

\statement{A.52}{Proved}{The paired filiform layer retains parity-one involutory images after two scalar steps}
The compatible family of A.40--A.50 can be enlarged
by the actual quotient in A.51 together with its
\(c\)-pullback, through both of their scalar
refinements. The first-to-second scalar kernel is
\(C_2^2\), and the third scalar step stabilizes.
There is a full negative element retaining every
preceding requirement and having an actual
parity-one involutory image on this enlarged family.
The full element is not proved to be an involution,
and the family is not proved cofinal.

\begin{proof}
Let \(q:\Pi\to G_2=P\rtimes H\) be the actual
map of A.51. Form \((q,qc)\), write \(Q\) for
its joint image, and write \(\overline H\) for
the joint head image in \(H\times H\).
Both are stable under the coordinate interchange
induced by \(c\). The head has a normal seven-group
\(\overline N\) and quotient \(C_6\); the second
head's \(C_6\)-coordinate is the negative of the
first. It is therefore an odd relative quotient
retained in \(C\), and \(h_C\) acts identically
on it.

\emph{Actual support separation.}\par
Use the commutator convention
\([g,h]=ghg^{-1}h^{-1}\). Define the actual
words in \(\Pi\)
\[
\begin{gathered}
 u=[s,r],\qquad v=[s,r^{-1}],\qquad
 a=u^3v^5,\qquad b=u^5v^3,\\
 t_1=[a,[a,[a,b]]],\qquad
 t_2=[b,[b,[b,a]]].
\end{gathered}                                       \tag{A.52.1}
\]
Their head images are computed using (A.51.2)
and \(c(s)=s\), \(c(r)=sr^{-1}s\).
On the two head abelianizations, \(a\) has
coordinates \(((1,0),(0,6))\), and \(b\) has
coordinates \(((0,1),(6,0))\). More decisively,
the full head images of the depth-four words are
\[
               (q_H,q_Hc)(t_1)=(T,1),\qquad
               (q_H,q_Hc)(t_2)=(1,T).
                                                               \tag{A.52.2}
\]
These are exact group equalities, recorded in
\path{check_filiform7_scalar.py}. No recognition
from a graded approximation is used. Both elements
are central in \(\overline N\).

The four modules
\[
                  V_+,\quad V_+^\vee,\quad
                  V_-,\quad V_-^\vee
                                                               \tag{A.52.3}
\]
are consequently mutually nonisomorphic, even
on the odd subgroup
\(\overline J=\overline N\rtimes C_3\).
The element \(t_1\) distinguishes the first
head block from the second, \(t_2\) distinguishes
the second from the first, and the disjoint
positive and negative seventh-root Frobenius
orbits distinguish each module from its dual.
There are no alternating cross-invariants between
the head blocks. The odd and full exterior-square
coinvariant dimensions are therefore four and
two, respectively; in each head block, \(s\)
interchanges the two odd invariant lines.

The image of \(R\) in each \(G_2\) is \(P\).
Its image in the direct sum of (A.52.3) is a
subdirect C-submodule, hence the whole direct
sum by their nonisomorphism. Commutators recover
all four central coordinates. It follows that
\[
            (q,qc)(R)=P\times P,
            \qquad Q=(P\times P)\rtimes\overline H.
                                                               \tag{A.52.4}
\]
The semidirect description is literal: the joint
image projects onto \(\overline H\) and contains
the entire kernel \(P\times P\), so it is the
full preimage of \(\overline H\) in the product.
In particular its s-generator is the pure head
lift, because both normalized s-vectors are zero.

\emph{The actual scalar levels.}\par
Let \(\overline G\) be the quotient of \(Q\)
by the two centers, and let
\(W=R\cap\ker(\Pi\to\overline G)\).
Use the notation \(W_i\) of A.51.
The all-nonzero-class argument there gives
\(H^1(C,V_i)=H^1(C,V_i^\vee)=0\) for every
head block. The exact center calculation is
therefore the direct sum of the two independent
calculations in A.51. The first center is
\(\mathbb F_2^2\), obtained by identifying the
two central generators within each block; the
second is \(\mathbb F_2^4\). More precisely,
\[
\begin{gathered}
 E_2=\Pi/W_2=Q\times_{\overline H}C,\qquad
 W_1/W_2\simeq C_2^2,\\
                         W_3=W_2.
\end{gathered}                                       \tag{A.52.5}
\]
At the third step the old center is
\(\overline N\)-trivial, so its tensor pairing
with (A.52.3) vanishes. The next center is an
exterior-square quotient on which the odd
\(\overline J\) acts trivially. Its dimension
is at most four, already attained. As in A.51,
the abelianization induction uses the image of
the new globally central kernel at each step.

\emph{The literal frames and the negative square.}\par
On the actual fiber product in (A.52.5) define
\[
                         \psi=(\operatorname{id}_Q,h_C).
                                                               \tag{A.52.6}
\]
The common head actions agree. This positive
automorphism fixes the actual cusp literally.
It retains the s-frame: since the finite
s-generator is pure, every element of
\(C_{\overline H}(s)\) has a pure lift in
\(C_Q(s)\). The old s-frame's head projection
can therefore be matched by such a finite
centralizer element.

The r-frame also lifts. The kernel \(P\times P\)
is a two-group and the actual r-generator has
order three. Any head centralizer element sends
it to another order-three lift over the same
head element. Coprime complement conjugacy
adjusts that lift by the kernel. Hence
\[
                C_Q(r)\longrightarrow
                C_{\overline H}(r)
                \quad\hbox{is surjective}.
\]
Choose the finite r-frame to match the old one.
These pairs of frames belong to the actual
fiber product, so they verify the required
peripheral identities for (A.52.6).

Let \(c_Q\) be coordinate interchange on the
joint image. It is the actual action induced
by \(c\), and has square one. The negative
companion of (A.52.6) is
\[
                      b=(c_Q,b_C),\qquad b_C=h_Cc_C.
                                                               \tag{A.52.7}
\]
Its head actions agree and both components
square to one. It is thus an actual involution
of \(E_2\), retaining the old normalized
parity-one data.

Apply A.3 only now, with \(N=W_2\).
By \(W_3=W_2\), its actual central refinement
is \(E_2\) itself. The verified literal frames
give a full positive lift \(h\) inducing
(A.52.6). The full \(a=hc\) has the specified
involutory image; no equality \(a^2=1\) in the
full group follows from this argument.

\emph{Retaining the preceding finite requirements.}\par
For any finite collection of the stabilized
earlier layers, let \(R_{\rm old}\) denote its
residual inside \(R\). Its relative perfectness
was proved in A.43 and A.48--A.50. The
abelianization of \(R/R_{\rm old}\) has only
the previous six-dimensional and 21-dimensional
simple modules; the perfect heads add none.
None of the new 42-dimensional simples is a
quotient of that abelianization. Equivalently,
the degree-four central seven-directions in
(A.52.2) act trivially on the old modules and
nontrivially on the corresponding new ones.

Every nonzero new cohomology class has full
affine R-image. If it killed \(R_{\rm old}\),
the resulting new simple module would be a
quotient of \((R/R_{\rm old})^{\rm ab}\),
a contradiction. Thus \(R_{\rm old}\) still
surjects onto every new simple and its central
cover. The inflation, exact-center and two-step
stabilization arguments remain valid over the
enlarged base, with the same literal frames.
Finite unions and compactness retain all
preceding requirements. This proves the claimed
compatible family, with the full square and
cofinality still undetermined.
\end{proof}

\statement{A.53}{Proved}{Every finite abelian residual layer has bounded actual scalar depth}
Let \(K_0\) be any closed normal subgroup of
\(\Pi=C_2\amalg C_3\), and put
\[
 R=R_{\rm nil}(K_0),\qquad C=\Pi/R.
\]
Let \(W\) be closed normal in \(\Pi\), contained
in \(R\), with \(V=R/W\) finite abelian. Define
\[
 W_0=W,\qquad W_{i+1}=\overline{[W_i,\Pi]}.
\]
Write \(\exp(V_2)=2^{e_2}\) and
\(\exp(V_3)=3^{e_3}\), with exponent index zero
when that primary summand is absent, and set
\[
                     m=\max\{1,e_2+1,2e_3+1\}.
                                                               \tag{A.53.1}
\]
Then the actual scalar refinements stabilize:
\[
                            W_{m+1}=W_m.
                                                               \tag{A.53.2}
\]
In particular, for an elementary-abelian
\(\ell\)-group \(V\), the respective bounds are
\[
\begin{array}{c|ccc}
 \ell&2&3&\ell\geq5\\ \hline
 \text{stabilization}&W_3=W_2&W_4=W_3&W_2=W_1.
\end{array}                                                   \tag{A.53.3}
\]

More generally, the same conclusions hold for
any closed normal \(R\) satisfying
\(R=\overline{[R,\Pi]}\), provided
\[
 H^1_{\rm cont}(C,V^\vee)=0,
 \qquad V^\vee=\operatorname{Hom}(V,\mathbb Q/\mathbb Z).
                                                               \tag{A.53.4}
\]
For every \(i\geq1\), the R-preimage
\(M_i=R/W_i\) has abelianization \(V\), has
nilpotency class at most two, and has finite
derived subgroup
\[
 Z_i=M_i'=W/W_i,
 \qquad Z_i\text{ a quotient of }
                 \textstyle\bigwedge^2_{\mathbb Z}V.
                                                               \tag{A.53.5}
\]
Its p-primary action factors through the maximal
pro-p quotient of \(\Pi\): \(C_2\), \(C_3\),
or the trivial group. These are assertions about
the actual relative central refinements.
They do not supply the peripheral automorphism
or solve its negative norm equation.

\begin{proof}
We first prove the conditional assertion and
then show that (A.53.4) is automatic over the
stated relative-pronilpotent base.

\emph{The finite action and its coinvariants.}\par
Let \(H\) be the finite action image of \(\Pi\)
on \(V\), and let
\(K_p=O^p(H)\) be the kernel of its maximal
p-group quotient. The quotient \(H/K_p\)
is generated only by the image of \(s\) when
\(p=2\), only by the image of \(r\) when
\(p=3\), and is trivial otherwise. Thus it is
trivial or cyclic of order \(p\).

Relative perfectness gives \((V_p)_H=0\).
The finite abelian p-group \((V_p)_{K_p}\)
therefore has zero coinvariants under the
p-group \(H/K_p\). It must vanish:
\[
                           (V_p)_{K_p}=0.
                                                               \tag{A.53.6}
\]
An elementary form of the needed finite-group
argument is as follows. If \(U=[U,P]\) for
a finite p-group \(P\) acting on a finite
abelian p-group \(U\), then \(U\) lies in
every lower-central term of the finite p-group
\(U\rtimes P\). Its nilpotence forces \(U=0\).

\emph{The abelianization induction.}\par
Put \(E_i=\Pi/W_i\) and
\(A_i=W_{i-1}/W_i\). The new \(A_i\) is
globally central in \(E_i\), and
\(M_i=\overline{[M_i,E_i]}\). We prove
\(M_i^{\rm ab}=V\) inductively.

Use finite quotients of \(E_i\) retaining the
given V-quotient of \(M_i\). They are cofinal:
since \(V\) is finite, compactness and
separation provide an open normal subgroup
whose intersection with \(M_i\) lies in the
kernel of \(M_i\to V\); intersect it with
any prescribed open normal subgroup.
The induced quotient of \(E_{i-1}\) also
retains \(V\), because an element projecting
into \(M_{i-1}\) already lies in \(M_i\).

Let \(B\) be the new finite R-preimage
abelianization in such a quotient. By induction
the previous one is \(V\). The kernel \(A'\)
of \(B\to V\) is the image of the
\emph{new globally central} \(A_i\), and is
a trivial C-module. Duality gives
\[
 0\longrightarrow V^\vee\longrightarrow B^\vee
   \longrightarrow (A')^\vee\longrightarrow0.
                                                               \tag{A.53.7}
\]
Relative perfectness gives
\((B^\vee)^C=0\). The invariants exact sequence
and (A.53.4) force \((A')^\vee=0\).
Hence \(B=V\), and cofinality proves the
global equality \(M_i^{\rm ab}=V\).
In particular \(Z_i=M_i'\).
This argument uses the cohomology of the
infinite retained group \(C\), not merely
the cohomology of its finite action image.

\emph{Finite derived centers at every stage.}\par
For \(i=1\), the kernel \(Z_1=A_1\) is
globally central. Its commutator pairing is
alternating and bilinear on the finite
abelianization \(V\), so \(Z_1\) is a
finite quotient of
\(\bigwedge^2_{\mathbb Z}V\). Its action is
trivial.

Suppose the assertions hold at stage \(i-1\).
The next preimage satisfies
\([Z_i,M_i]\subseteq A_i\), which is globally
central. Thus \(\gamma_4(M_i)=1\) and
\([Z_i,Z_i]=1\). The commutator descends to
a C-equivariant pairing
\[
                       Z_{i-1}\otimes_{\mathbb Z}V
                                  \longrightarrow A_i.
                                                               \tag{A.53.8}
\]
Mixed-primary tensor factors vanish. On the
p-primary part the old center's action factors
through \(H/K_p\), whereas the new \(A_i\)
is globally trivial. Equation (A.53.6)
therefore makes every primary pairing zero.
Consequently \(Z_i\) centralizes \(M_i\).

As \(M_i^{\rm ab}=V\), its derived subgroup
\(Z_i\) is again a finite quotient of
\(\bigwedge^2_{\mathbb Z}V\). The action on
it factors through \(H\), by functoriality
of the commutator pairing. The exact sequence
\[
                   0\longrightarrow A_i
                    \longrightarrow Z_i
                    \longrightarrow Z_{i-1}
                    \longrightarrow0
\]
has globally trivial new kernel. On a p-primary
part, the old action image is a p-group by
induction. The automorphisms fixing both
subquotients form a subgroup of
\(\operatorname{Hom}(Z_{i-1,p},A_{i,p})\),
a p-group. Thus the entire finite action image
on \(Z_{i,p}\) is a p-group and factors through
\(H/K_p\). This proves (A.53.5) and the
action assertion for every \(i\).

\emph{The actual augmentation filtration.}\par
Let \(T_p\) be the cyclic generator of the
action on \(Z_{i,p}\), and put \(I_p=T_p-1\),
with \(I_p=0\) for a trivial action.
Relative commutators inside this actual
abelian kernel give, for \(0\leq j\leq i\),
\[
                         (W_j/W_i)_p=I_p^j Z_{i,p}.
                                                               \tag{A.53.9}
\]
For \(p\geq5\), \(I_p=0\). For \(p=2\),
the relation \(T_2^2=1\) gives
\[
 I_2^2=-2I_2,
 \qquad I_2^{e_2+1}=(-2)^{e_2}I_2=0
                                                               \tag{A.53.10}
\]
on exponent dividing \(2^{e_2}\).
For \(p=3\), one similarly has
\[
\begin{gathered}
 I_3^3=-3I_3(1+I_3),\\
 I_3^{2e_3+1}=(-3)^{e_3}I_3(1+I_3)^{e_3}=0
\end{gathered}                                                 \tag{A.53.11}
\]
on exponent dividing \(3^{e_3}\).
Absent primary summands contribute no condition.
The choice (A.53.1) annihilates all \(I_p^m\).
Taking \(i=m+1\), \(j=m\) in (A.53.9)
proves \(W_m/W_{m+1}=1\). Equality persists
under every further closed commutator operation.

For example, in the elementary dyadic case,
\(Z_2\) is an actual
\(\mathbb F_2[C_2]\)-module and
\[
 Z_1=Z_2/(s-1)Z_2,
 \qquad W_1/W_2=(s-1)Z_2.
                                                               \tag{A.53.12}
\]
The possible nonzero second scalar layer is
exactly the contribution of the regular
two-dimensional summands. A.51 realizes such
a summand.

\emph{Automatic cohomology for the pronilpotent base.}\par
Now let \(R=R_{\rm nil}(K_0)\).
A.34 gives \(R=\overline{[R,K_0]}\), and
\(N=K_0/R\) is an actual closed normal
pronilpotent subgroup of \(C\).
Fix a prime \(p\) and write
\[
                         N=N_p\times N_{p'},
\]
using its characteristic Sylow factors.
Relative perfectness with \(K_0\) gives
\((V_p)_N=0\). The finite group
\((V_p)_{N_{p'}}\) has a finite p-group
action of \(N_p\) and has zero coinvariants.
The finite p-group argument above makes it
zero. Hence
\[
                         (V_p^\vee)^{N_{p'}}=0.
                                                               \tag{A.53.13}
\]

Let \(z:C\to V_p^\vee\) be any continuous
one-cocycle. Its restriction to the pro-p-prime
group \(N_{p'}\) is a coboundary by finite
averaging. To see the finite reduction
explicitly, take the core of an open subgroup
on which both the action and the cocycle
vanish. The resulting finite quotient has
order prime to \(p\), invertible on
\(V_p^\vee\).
Subtract the coboundary, so \(z\) vanishes
on \(N_{p'}\). For \(n\in N_{p'}\) and
\(g\in C\), normality gives
\[
 nz(g)=z(ng)=z\bigl(g(g^{-1}ng)\bigr)=z(g).
\]
All values therefore lie in the zero invariant
space (A.53.13). This proves
\(H^1(C,V_p^\vee)=0\), and summing the
primary parts gives (A.53.4).

The normal subgroup used here is the actual
\(N_{p'}\) inside \(C\). No cohomology of an
infinite extension has been inferred from a
finite head alone. This proves the unconditional
relative-pronilpotent assertion and all the
stated elementary-prime bounds.
\end{proof}

The new-base hypothesis (A.53.4) still requires
proof after an arbitrary nonabelian refinement.
Even where A.53 applies, the construction of an
equivariant peripheral automorphism with the
specified frames, and its full negative norm,
remain separate questions. The theorem bounds
the actual scalar refinement needed for those
questions; it does not settle them.

\statement{A.54}{Proved}{The exact new-base cohomology can obstruct a successive reuse of A.53}
Use the initial relative-pronilpotent base and a
stabilized layer from A.53, and put
\[
 P=R/W_m,\qquad C'=\Pi/W_m,
 \qquad P^{\rm ab}=V.
\]
Let \(T\) be a finite C-module such that, for
each prime \(p\), its p-primary part has zero
invariants under the actual normal subgroup
\(N_{p'}\) of \(N=K_0/R\).
Inflate \(T\) to \(C'\). Then restriction gives
\[
 H^j(C,T)=0\quad(j\geq0),\qquad
 H^1(C',T)\simeq\operatorname{Hom}_C(V,T).
                                                               \tag{A.54.1}
\]
In particular
\[
                        H^1(C',V)\simeq\operatorname{End}_C(V),
                                                               \tag{A.54.2}
\]
which is nonzero for \(V\ne0\).

The failure of automatic vanishing can occur
for an actual next residual layer. In the
single-head construction of A.51, let \(U\)
be its 42-dimensional simple module and let
\(R'=W_2\). There exists a closed normal
\(W'\subset R'\) with
\[
 R'/W'\simeq U^\vee,
 \qquad \dim_{\mathbb F_2}H^1(C',U)=1.
                                                               \tag{A.54.3}
\]
Thus the cohomology hypothesis for the next
application of A.53 is genuinely nonzero,
although \(R'=\overline{[R',\Pi]}\).
This is not a negative norm obstruction.

\begin{proof}
Continuous cohomology of a pro-p-prime group
with finite p-primary coefficients vanishes in
positive degree by finite averaging. The
assumed degree-zero invariants also vanish.
Apply the profinite Hochschild--Serre spectral
sequence to the actual closed normal
\(N_{p'}\subset C\), and sum over primes.
This gives \(H^j(C,T)=0\) in every degree.
The continuous inflation--restriction sequence
for
\[
                 1\longrightarrow P\longrightarrow C'
                     \longrightarrow C\longrightarrow1
\]
then has zero outer terms in degrees one and
two. Because \(P\) acts trivially on the
inflated coefficient module,
\[
 H^1(C',T)\simeq H^1(P,T)^C
       =\operatorname{Hom}_C(P^{\rm ab},T).
\]
The continuous spectral and five-term sequences
used here are stated in
\citet[Theorem~6.7.3 and Corollary~6.7.4, p.~194]{APoonen2017}.
No splitting of \(C'\) over \(C\) is assumed.

For \(T=V\), A.53 proves that the
\(N_{p'}\)-coinvariants of each \(V_p\)
vanish. Finite averaging identifies its
invariants with its coinvariants, so the
hypothesis holds and gives (A.54.2).

For the concrete next layer, the first
filiform residual abelianization is
\(U\oplus U^\vee\). The two modules are
distinct and each has endomorphism field
\(\mathbb F_2\), while A.51 gives
\[
 \dim H^1(\Pi,U)=\dim H^1(\Pi,U^\vee)=7.
\]
Equation (A.54.1) makes each corresponding
\(H^1(C',-)\) one-dimensional. Its inflation
line is spanned by the original affine
cocycle, since that cocycle factors through
the retained affine quotient of \(C'\).

In the dual module the original normalized
cocycle has s-value zero and r-value \(e_1\).
Choose instead
\[
                         d'(s)=0,
                         \qquad d'(r)=\zeta^2 e_1.
                                                               \tag{A.54.4}
\]
The updated exact certificate
\path{check_filiform7_scalar.py} checks its
r-norm equation and independence. In the
specified bit representation of \(\mathbb F_8\),
the two r-values are (1) and (4) at the
same coordinate. The normalized coboundary
space has rank 21; adjoining the first class
gives rank 22, and adjoining both gives
rank 23. Hence (A.54.4) lies outside the
inflation line.

Its restriction to \(R'\) is nonzero by
inflation--restriction. Since \(R'\subset R\)
acts trivially on \(U^\vee\), this restriction
is a continuous homomorphism. It is
\(\Pi\)-equivariant; irreducibility makes it
surjective. Its kernel \(W'\) is closed
normal in \(\Pi\), giving (A.54.3).
The dual coefficient required for this next
layer is \((U^\vee)^\vee=U\), whose
\(H^1(C',U)\) is nonzero as asserted.

One may instead return to the initial base
and retain both independent affine cocycles
simultaneously. Their joint residual
abelianization is still finite abelian, so
A.53 applies there. Thus (A.54.3) locates
the failure of an automatic stepwise
cohomology argument; it does not prove that
the next scalar refinement, peripheral lift
or negative norm is impossible.
\end{proof}


% Reconstructed from the visible, independently reviewed seventh-hour source
% after the execution environment was replaced. Not a byte-identical archive.

\statement{A.55}{Proved}{The actual second scalar center of an elementary two-layer is an explicit exterior-square coinvariant module}
Let \(K_0\) be closed normal in \(\Pi=C_2\amalg C_3\), put
\(R=R_{\rm nil}(K_0)\), \(C=\Pi/R\), and let \(W\subset R\)
be closed normal in \(\Pi\) with \(A=R/W\) finite elementary
abelian of characteristic two. Let \(H\) be the finite image of C
on A, and put
\[
 K=O^2(H),\qquad M=(\mathop{\bigwedge}\nolimits^2_{\mathbb F_2}A)_K,
 \qquad W_{i+1}=\overline{[W_i,\Pi]},\quad W_0=W.
 \tag{A.55.1}
\]
The natural commutator maps give
\[
 W/W_2\simeq M,\qquad W/W_1\simeq M_H. \tag{A.55.2}
\]
Thus an extension of C by an explicit finite class-two group attains
the actual second scalar center. If H acts trivially on M, then
\(W_1=W_2\).

The same conclusion holds for general \(R=\overline{[R,\Pi]}\)
if the actual quotient \(\Pi/W\) is identified with \(A\rtimes C\),
if \(H^1(C,A^\vee)=0\), and if H has a normal odd subgroup J with
\(A_J=0\). No peripheral automorphism or equivariant frame is asserted.

\begin{proof}
We prove the conditional version first. The maximal pro-two quotient
of H is a quotient of \(\Pi^{(2)}=C_2\). Thus H/K is cyclic of
order at most two, generated by s, and J acts trivially on M.
For \(T=\operatorname{Hom}_{\mathbb F_2}(A,M)\), we have
\(T^J=0\). Finite averaging and Hochschild--Serre give
\[
 H^i(H,T)=0\qquad(i\geq0). \tag{A.55.3}
\]
The quadratic refinements of the canonical alternating map
\(A\times A\to M\) form a nonempty affine space under T.
Equation (A.55.3) gives a unique H-equivariant refinement q.
Choose a bilinear cocycle with diagonal q. It defines a finite group
\[
 1\longrightarrow M\longrightarrow P\longrightarrow A
 \longrightarrow1,\qquad P'=M,\quad P/M=A. \tag{A.55.4}
\]
The equality \(P'=M\) follows because the commutator pairing is onto.
Any pair of automorphisms of A and M preserving q lifts to P:
choose lifts of a basis of A, whose squares and commutators are exactly
the defining relations. The automorphism kernel fixing both A and M
is T. Pulling back along H and using \(H^2(H,T)=0\) gives a strict
H-action on P, and hence a C-action through H.

Form \(P\rtimes C\) and lift the actual generators of the given
\(A\rtimes C\) quotient to \(S_0,R_0\). Their powers lie in M.
The element \(z_s=S_0^2\) is s-fixed, hence H-fixed because s
generates the action on M. It is globally central. Since
\(R\subset\Pi'\), the abelianization gives
\(\epsilon:C\to C_6\to C_2\). Pull back \(C_4\to C_2\)
and push its kernel generator to \(z_s\). Equivalently, add the carry
\[
 (c,d)\longmapsto\epsilon(c)\epsilon(d)z_s. \tag{A.55.5}
\]
This makes the actual s-lift involutory. It leaves r unaffected;
r acts trivially on M, so its cube is killed by a unique central
correction, since three is invertible on M. The repaired generators
define an actual continuous map from \(\Pi\). Its image of R maps
onto A and its commutators generate M, hence it is all P. The full
map is therefore onto the repaired extension over C.

The images of W and its next two commutator subgroups are
\[
 M,\quad IM,\quad I^2M=0,\qquad I=s-1. \tag{A.55.6}
\]
Thus the actual quotient factors through \(\Pi/W_2\) and attains M.
Conversely, conditional A.53 applies. It makes the residual preimage
class two with abelianization A, and its derived subgroup \(W/W_2\)
is an exterior-square quotient. Its action factors through H and has
a filtration of trivial scalar factors, so its action image is a
two-group and factors through H/K. Hence there is a natural
surjection \(M\to W/W_2\). Composing with the attained quotient
is identity on every commutator generator of M, proving it an
isomorphism. H-coinvariants give the second equality in (A.55.2).

For the initial base, let \(N=K_0/R\) and use its actual normal
pro-odd factor \(N_{2'}\). A.53 gives \(A_{N_{2'}}=0\), and
the invariants of A and its dual vanish by averaging. Continuous
Hochschild--Serre gives
\[
 H^i(C,A)=H^i(C,A^\vee)=0\qquad(i\geq0). \tag{A.55.7}
\]
The actual finite-kernel extension \(\Pi/W\to C\) therefore
splits by \(H^2(C,A)=0\). It has a continuous set section: an
open subgroup avoiding its finite kernel supplies local sections
on a finite clopen partition. Thus continuous extension classification
applies. The image of \(N_{2'}\) in H supplies J. All conditional
hypotheses hold, proving the unconditional assertion.
\end{proof}

\statement{A.56}{Proved}{The forced CM frame extends through an actual self-dual central layer with both involutory normalizers}
Retain the actual base C, positive action \(h_C\), negative
involution \(b_C=h_Cc_C\), and literal equivariant frames of A.33.
There is an actual further quotient over C whose finite residual
group P is class two and satisfies
\[
 \begin{gathered}
 P^{\rm ab}=V^8,\qquad \dim_{\mathbb F_2}V=108,
 \qquad\operatorname{End}_C(V)=\mathbb F_8,\\
 P'=\operatorname{Sym}^2_{\mathbb F_8}(\mathbb F_8^8),
 \qquad\dim_{\mathbb F_2}P'=108.
 \end{gathered}\tag{A.56.1}
\]
The module V is simple and self-dual. This quotient has a positive
peripheral automorphism over \(h_C\), fixing the actual cusp and
retaining the original literal C-frames, whose negative companion
and prescribed s- and r-normalizers are involutory. Its scalar
refinement is already stabilized. The complete action and frames
are realized by a full positive lift, simultaneously retaining
every previous finite family in A.47--A.52 and hence their tower.
The full negative companion has parity one and order-two images
in these quotients. For the original stored finite action and s-frame,
A.76 proves that its full square is nontrivial. Later reselected
action/frame tuples and the cofinality question remain unresolved.

The retained r-frame is specifically the one constructed in A.33:
its positive odd relative components are one, and its negative
normalizer there is \(sc\). No assertion about an unrestricted
different odd r-frame is made.

\begin{proof}
\emph{The actual head and complete affine module.}\par
Let
\[
 \begin{gathered}
 Q_6=S_3\times A_4,\quad
 s=((01),(01)(23)),\quad r=((012),(012)),\\
 K_6=\ker(\Pi\to Q_6),\quad
 H=\Pi/\overline{K_6'K_6^7},\quad
 N=K_6/\overline{K_6'K_6^7}\simeq\mathbb F_7^{13}.
 \end{gathered}\tag{A.56.2}
\]
The rank is \(72-36-24+1=13\), from the actual Bass--Serre
quotient graph. The congruence-kernel image is N: H is solvable,
whereas a solvable congruence quotient has no seven-factor.
Thus H is an actual odd relative quotient of C, and \(h_C\)
acts identically on it.

The graph certificate labels edges by \(g\in Q_6\), with boundary
\(B_{g\langle r\rangle}-A_{g\langle s\rangle}\).
Averaging the odd splitting under s and c makes s pure and gives
the actual r-vector
\[
 (3,4,6,0,5,4,3,6,1,2,6,5,3).\tag{A.56.3}
\]
The actual \(\tau^6\)-vector is the second cycle-basis vector.
On \(Q_6\), c is identity on \(S_3\) and conjugation by \((23)\)
on \(A_4\); on graph edges it is \(g\mapsto\kappa(g)s\).
These formulas specify the actual source map and action.

Put \(k=\mathbb F_8\), with \(\zeta^3+\zeta+1=0\).
The character covector
\[
 \lambda=(1,0,0,6,0,0,0,0,0,6,0,0,1)\tag{A.56.4}
\]
is fixed by s and c, has stabilizer exactly \(\langle s\rangle\),
and is moved by \(\beta_0=(1,(01)(23))\). Induce its seventh-root
character from \(N\rtimes\langle s\rangle\). The resulting V
has36 distinct N-weight lines. No Frobenius scalar two or four
normalizes the orbit; hence its underlying binary module is simple
with endomorphism field k. Opposite weights belong to the same
orbit, making it self-dual. These properties persist on
\(K=O^2(H)\), whose image \(C_3\times A_4\) is transitive on
the36 weight lines.

Normalize cocycles to have r-value zero. Exact linear algebra gives
\[
 \begin{gathered}
 E=H^1(\Pi,V),\quad\dim_kE=8,\\
 \dim_kH^1(\langle s\rangle,V)
 =\dim_kH^1(\langle\tau\rangle,V)=4,\\
 \operatorname{res}:E\xrightarrow{\sim}
 H^1(\langle s\rangle,V)\oplus H^1(\langle\tau\rangle,V).
 \end{gathered}\tag{A.56.5}
\]
All bases are specified by exact row reduction in
\path{check_Q6_forced_frame.py}. Let \(A=V\otimes_kE^\vee\)
and let d be the universal cocycle of all eight classes. The actual
normal pro-odd factor of C has no invariants, so \(H^1(C,V)=0\).
The eight restrictions to R are independent. Since A is a direct
sum of copies of the simple module V with endomorphism field k,
\(d(R)=A\). Therefore
\[
 \Pi\to A\rtimes C,\qquad g\mapsto(d(g),g_C)\tag{A.56.6}
\]
is an actual epimorphism, with \(d(r)=0\) and \(d(\tau)=d(s)\).

\emph{The attained center.}\par
Pair the18 opposite N-weight pairs. Their alternating k-form spans
the one-dimensional K-invariant form space. Conjugation by s acts
on this line by a scalar of order dividing two in \(k^\times\),
so trivially. Every K-invariant alternating binary form on A is
the trace of that pairing with a symmetric eight-by-eight k-matrix;
diagonal entries are allowed because the pairing is alternating.
Thus
\[
 M=(\mathop{\bigwedge}\nolimits^2_{\mathbb F_2}A)_K
   =\operatorname{Sym}^2_k(k^8),\qquad\dim_kM=36,\tag{A.56.7}
\]
and H acts trivially. For \(W=\ker(\Pi\to A\rtimes C)\), A.55 gives
\[
 W/\overline{[W,\Pi]}=M,\qquad
 \overline{[[W,\Pi],\Pi]}=\overline{[W,\Pi]}.\tag{A.56.8}
\]
The certificate constructs this group. In residual coordinates
\((a,z)\in A\times M\), orient the opposite pairs and use
\[
 f(a,b)=\sum_{(i,j)}\operatorname{sym}(a_i,b_j).\tag{A.56.9}
\]
Here \(\operatorname{sym}(v,w)\) has diagonal entries \(v_iw_i\)
and off-diagonal entries \(v_iw_j+v_jw_i\). A monomial head action
acquires the carry \(\operatorname{sym}(a_i,a_j)\) exactly when it
reverses an oriented pair. This is a strict H-action. The central
carry \(\epsilon(h)\epsilon(k)q_s\), with \(q_s\) the initial
actual s-square, is A.55's C4 repair. The actual pure r-lift has cube
one. The source therefore maps onto an extension of C by P, with
\(|P|=2^{864+108}\); the total quotient can be infinite.

\emph{The pure intermediate frame and exact norms.}\par
First give the intermediate s-frame H-image \(\beta_0\).
The boundary isomorphism uniquely determines the positive
cohomology action \(L_E\): \(\beta_0\) on the s-boundary and
identity on the cusp. Let \(C_E\) be the actual c-action. On A,
let L be the transpose multiplicity action and let \(\Gamma\)
combine the transpose of \(C_E\) with the c-permutation of weights.
Then \(B=L\Gamma\) has square one.

The actual c-action is \(\operatorname{Inn}(w_c)(\Gamma,c_C)\);
its inner term is solved from \(c(s)=s\), \(c(r)=sr^{-1}s\)
and has zero norm. A positive affine lift is
\(\operatorname{Inn}(w)(L,h_C)\), with conditions
\[
 (L-1)d(\tau)+(1-\tau)w=0,\qquad q=w+Lw_c.\tag{A.56.10}
\]
The initial \(q+Bq\) is nonzero. Its cusp correction space has
dimension32 over k and norm rank four; the obstruction is in that
image. A correction gives \(q+Bq=0\). The simultaneous affine
s-frame and involution equations have centralizer dimension160,
norm rank68, and the same augmented rank, so are solved as well.

Lift L and \(\Gamma\) by their symmetric-square actions and the
quadratic carries, retaining the C4 term. A central character
reconstructs c from its generator values; another makes the positive
lift fix the actual cusp. The negative companion then squares to
one on s and r. The remaining s-class residue is central. The
homogeneous cusp and s-frame correction spaces have dimensions28
and92 over k. Their central class map has binary rank83 and contains
the actual residue. The first resulting s-frame has a square outside
the center norm image, so shifting that single frame centrally fails.

The complete homogeneous class equation has277 binary parameters.
For correction lifts \(P_\tau,P_s\), and
\(d_s=\operatorname{Inn}(\beta^{-1})b\), its central square changes
modulo \((1+B)M\) by
\[
 \operatorname{Norm}_{d_s}(P_s^{-1})+
 \operatorname{Norm}_b(P_\tau).\tag{A.56.11}
\]
There is no mixed term. The polar maps are universal commutators.
Their common radical has dimension250 and linear image rank eight.
Removing it leaves an exact27-variable system. Core assignment27666
reconstructs the277-bit assignment
\[
 18677289404303020461252901246388292043973835622281379854.
 \tag{A.56.12}
\]
The complete tuples are stored in
\path{Q6_actual_central_solution.json}. Search only discovers the
witness. The separate \path{check_Q6_final_frames.py} verifies it
using the original multiplication: actual cusp, negative square,
s-conjugacy, and square of \(\beta^{-1}b\). For r use the canonical
affine conjugator \(qS\), whose affine square vanishes. Its nonzero
central square belongs to \((1+B)M\), and the recorded correction
kills it. The direct verifier checks
\[
 \begin{gathered}
 (\beta^{-1}b)^2=(J^{-1}b)^2=1,\\
 (\beta^{-1}b)(s)=s,\qquad(J^{-1}b)(r)=r^{-1}.
 \end{gathered}\tag{A.56.13}
\]
The r-head is exactly the retained one: A.33 uses positive r-frame
one at each odd relative prime, so \(A_C|_H=1\); since \(h_C|_H=1\),
its negative conjugator \(J_C=h_C(s)A_C\) has H-image s. Later
frames retain this literal C-value. The finite action therefore
glues with the actual C-action and both frames.

\emph{Every already retained s-frame.}\par
Write \(\beta_H=(n,\beta_0)\). Equivariance gives
\(n\in N^s\), \(n+\beta_0c(n)=0\). Put \(t=\beta_0n/2\),
so \((\beta_0-c)t=n\). Coprime centralizer lifting in the actual
relative-seven factor gives \(k_C\in C_C(s)\) with congruence image
one and H-image t. The intermediate frame
\[
 \beta_C^0=b_C(k_C)\beta_Ck_C^{-1}\tag{A.56.14}
\]
is equivariant for the same \(b_C\) and has pure head \(\beta_0\).
Apply the preceding construction to it.

Apply A.34/R2 to identity C with frames \((k_C,1,1)\) at
\((s,r,\tau)\), and normalize the full lift g to fix \(\tau\).
The normalizing conjugator has C-image one. Inside R, W is the
intersection of the kernels of all restricted classes in the
complete \(H^1(\Pi,V)\). Precomposition by g over identity C
permutes this entire space, hence preserves W and \([W,\Pi]\).
It induces an actual central-quotient automorphism \(g_{\rm new}\).
This uses a full R2 lift, not an assumed lift of an affine map.
Conjugating returns the original frame value
\[
 b_C(k_C)^{-1}\beta_C^0k_C=\beta_C.\tag{A.56.15}
\]
The involutory normalizer transports by conjugation; the r-frame
returns to its original value because g has r-frame one. The cusp
remains fixed.

\emph{The earlier layers and full realization.}\par
Let \(C_{\rm old}\) retain a finite family from A.47--A.52, with
finite residual \(P_{\rm old}\) over C and stabilized kernel.
Its abelianization has only the six-,21-, and42-dimensional simple
types; perfect factors contribute none. Therefore
\[
 \operatorname{Hom}_C(P_{\rm old}^{\rm ab},V)
 =\operatorname{Hom}_C(P_{\rm old}^{\rm ab},V^\vee)=0.
 \tag{A.56.16}
\]
The original base has \(H^1\) and \(H^2\) zero for these new
coefficients. Inflation--restriction gives the two needed \(H^1\)
vanishings over \(C_{\rm old}\). The complete new affine classes
remain independent on its residual kernel, so the actual quotient
is \(A\rtimes C_{\rm old}\). Conditional A.55, with the unchanged
head and odd normal subgroup, gives the same M and immediate scalar
stability. No stability of an intersection of kernels is presumed.

The actual joint quotient is the full fiber product over C.
Equivariant Goursat also proves this: a nontrivial common quotient
with the finite two-group P has a nonzero abelian quotient shared
by A and \(P_{\rm old}^{\rm ab}\), contradicting (A.56.16).
A central overlap cannot survive because \(M=P'\).
For a nonpure frame, first pair the new pure-intermediate action
with the unchanged old action over C; this intermediate action need
not be peripheral. Apply \((g_{\rm new},\operatorname{id}_{\rm old})\).
The final frames have the same original C-values and hence pair,
with involutory normalizers componentwise. No lift of \(k_C\) into
an old s-centralizer is required.

The joint kernel is scalar-stable. A.3 now realizes its complete
action and literal frames by a full positive lift. Its negative
companion retains parity one. The closed solution sets in the
compact full positive group have the finite-intersection property,
so retain the old tower and this layer simultaneously.
The remaining certificates are \path{check_Q6_central_cover.py},
\path{check_Q6_central_homogeneous.py}, and
\path{check_Q6_quadratic_core.cpp}. The search found its witness
after 18461 evaluations; direct exact verification is the proof
certificate. A.76 subsequently excludes a full involution retaining
this original stored finite action and s-frame. It leaves the later
reselected tuples and cofinality question unresolved.
\end{proof}

\statement{A.57}{Proved}{An actual four-dimensional family of nonsplit next layers is detected by an integral Bockstein}
Let \(C_{\rm new}\) be the actual quotient of A.56, with finite
residual group P over C, and let V be its simple module over
\(k=\mathbb F_8\). There is a four-dimensional k-subspace of
\[
 \ker\!\left(H^2(C_{\rm new},V)\longrightarrow H^2(\Pi,V)\right)
 \tag{A.57.1}
\]
whose nonzero elements define nonsplit extensions attained by
actual quotients of \(\Pi\). Moreover, a k-basis of this subspace
is attained simultaneously by an actual quotient with kernel
\(V^4\) over \(C_{\rm new}\). This is a family of actual next
chief-module directions; it is not asserted to exhaust the next
extensions, to be scalar-stable, or to admit a negative involution.

\begin{proof}
Put \(\mathcal O=\operatorname{GR}(4,3)\). The monomial action on
\(V=k^{36}\) lifts to the free \(\mathcal O\)-module
\(V_4=\mathcal O^{36}\): replace each seventh root of unity by its
Teichm\"uller lift and retain the same permutation matrices. The
relations hold literally, since the coefficient roots have order
seven. Thus there is an exact sequence of continuous C-modules
\[
 0\longrightarrow V\longrightarrow V_4\longrightarrow V
 \longrightarrow0. \tag{A.57.2}
\]
Both V and \(V_4\) are inflated to \(C_{\rm new}\); in particular
P acts trivially on them. The coefficient construction can be
checked with \(\mathcal O=(\mathbb Z/4)[T]/(T^3+2T^2+T+3)\).
The polynomial is a Hensel factor of
\((T^7-1)/(T-1)\), its reduction is irreducible, and the image of
T has order seven; see \path{check_A57_galois_ring.py}.

The actual normal pro-odd subgroup used in A.55 gives
\(H^i(C,V)=0\) for every i, including i zero. The five-term
sequence for \(1\to P\to C_{\rm new}\to C\to1\), together with
\(P^{\rm ab}=V^8\), therefore identifies
\[
 H^1(C_{\rm new},V)
   =\operatorname{Hom}_C(P^{\rm ab},V)=k^8.
 \tag{A.57.3}
\]
Inflation to \(H^1(\Pi,V)\) is an isomorphism: the quotient in
A.56 retains the complete eight-dimensional cocycle space, and
both sides have dimension eight.

Let \(\delta_{\rm new}\) be the connecting map associated to
(A.57.2). Its restriction to P is injective on (A.57.3).
Indeed, every homomorphism \(P\to V_4\) factors through
\(P^{\rm ab}\), which has exponent two; its image is contained
in \(2V_4\), and hence its reduction to V is zero. Exactness of
the coefficient cohomology sequence proves that
\[
 \delta_{\rm new}:k^8\longrightarrow H^2(C_{\rm new},V)
 \quad\hbox{is injective}. \tag{A.57.4}
\]
It is k-linear although \(V_4\) is not a k-vector space:
multiplication by the Teichm\"uller lift of each scalar a is a
\(\mathbb Z/4\)-module map of (A.57.2), inducing multiplication
by a at both ends. Naturality supplies the asserted linearity.

The cohomology of the profinite free product \(\Pi=C_2\amalg C_3\)
in degree two is the direct sum of the cohomologies of its two
factors. The odd factor contributes zero. On the 36 monomial
lines, s has four fixed lines and sixteen exchanged pairs.
After rescaling the exchanged pairs, these are four trivial
\(\mathcal O[C_2]\)-modules and sixteen regular modules.
Consequently the Bockstein on \(C_2\) is an isomorphism on the
four fixed k-lines and the positive-degree cohomology of the
regular pairs vanishes. Thus
\[
 H^2(\Pi,V)=k^4,\qquad
 \delta_{\Pi}=\delta_{\langle s\rangle}\circ
     \operatorname{res}_{\langle s\rangle},\qquad
 \operatorname{rank}_k\delta_{\Pi}=4.
 \tag{A.57.5}
\]
The exact joint s/cusp restriction isomorphism of A.56 identifies
\(\ker\operatorname{res}_{\langle s\rangle}\) with the
four-dimensional cusp summand. By naturality, its injective image
under (A.57.4) is precisely the intersection of
\(\operatorname{im}\delta_{\rm new}\) with (A.57.1).

For completeness these are actual quotient extensions, not only
abstract cohomology classes. Represent a nonzero class e in that
cusp summand by a continuous cocycle on \(C_{\rm new}\).
Since its s-restriction is a coboundary, choose this representative
with \(e(s)=0\) before defining the following pullback:
\[
 C_e=C_{\rm new}\mathbin{\times}_{V\rtimes C}(V_4\rtimes C),
 \qquad g\longmapsto(e(g),\bar g).
 \tag{A.57.6}
\]
Its extension class is \(\delta_{\rm new}(e)\ne0\). On \(\Pi\),
\(\delta_{\Pi}(e)=0\), so e has an actual coefficient lift
\(\widetilde e:\Pi\to V_4\). One may see this without an
existence shortcut. Lift the already normalized value e(s) by zero.
Lift e(r) arbitrarily and replace the lift by
\[
 \widetilde e(r)-\frac13(1+r+r^2)\widetilde e(r).
 \tag{A.57.7}
\]
The subtracted term reduces to zero, and the r-norm becomes zero.
The universal property of the profinite free product now gives
the required continuous cocycle. It defines a homomorphism
\(\Pi\to C_e\) over \(C_{\rm new}\).

The image of \(R_{\rm new}=\ker(\Pi\to C_{\rm new})\) in V is a
C-submodule. If it were zero, this homomorphism would factor
through a section of \(C_e\to C_{\rm new}\), contradicting
(A.57.4). Simplicity of V therefore makes that image all of V,
and the map onto \(C_e\) is surjective.

Finally choose a k-basis \(e_1,\ldots,e_4\) of the cusp summand
and form the simultaneous pullback
\[
 C_{\rm Bock}
   =C_{\rm new}\mathbin{\times}_{V^4\rtimes C}(V_4^4\rtimes C).
 \tag{A.57.8}
\]
Choose the four lifts (A.57.7). If their joint residual image
were a proper C-submodule of the semisimple module \(V^4\), some
nonzero C-linear functional \(V^4\to V\) would annihilate it.
Since \(\operatorname{End}_C(V)=k\), pushing out would split the
extension with class
\[
 \sum_{i=1}^4 a_i\delta_{\rm new}(e_i)
 =\delta_{\rm new}\!\left(\sum_{i=1}^4a_i e_i\right),
 \qquad (a_1,\ldots,a_4)\ne0.
 \tag{A.57.9}
\]
This contradicts injectivity. Hence the joint kernel is actually
\(V^4\), as claimed. No compatibility of a negative action or of
its literal normalizers has been used in this argument.
\end{proof}


% Proposed after checkpoint07; reconstructed from the complete proof sent
% during the outage and independently accepted by root and CE. Not integrated.
\statement{A.58}{Proved}{The complete finite abelian residual layer has an exact exterior-square scalar center}
Let \(K_0\) be closed normal in \(\Pi=C_2\amalg C_3\), set
\(R=R_{\rm nil}(K_0)\), \(C=\Pi/R\), and suppose
\(A=R/W\) is finite abelian for a closed normal subgroup
\(W\subset R\) of \(\Pi\). Let H be the finite C-action image
on A. For each prime p put
\[
 K_p=O^p(H),\qquad
 M_p=\left(\mathop{\bigwedge}\nolimits^2_{\mathbb Z} A_p\right)_{K_p},
 \qquad M=\prod_p M_p.
 \tag{A.58.1}
\]
Write \(W_0=W\), \(W_{i+1}=\overline{[W_i,\Pi]}\), let
\(p^{e_p}\) be the exponent of \(A_p\), taking \(e_p=0\) if
\(A_p=0\), and put
\[
 m=\max\{1,e_2+1,2e_3+1\}.
 \tag{A.58.2}
\]
The natural commutator map is an isomorphism
\[
 W/W_m\simeq M,
 \qquad (W_j/W_m)_p=I_p^jM_p\quad(0\leq j\leq m),
 \qquad W_{m+1}=W_m.
 \tag{A.58.3}
\]
Here \(I_p\) is the augmentation ideal of the cyclic action
\(H/K_p\), which is a quotient of \(C_2\) for p two, of
\(C_3\) for p three, and is trivial otherwise. In particular, if
H acts trivially on M, already \(W_1=W_2\) and \(W/W_1=M\).

The same conclusion holds for a general
\(R=\overline{[R,\Pi]}\) under the following actual hypotheses:
\(\Pi/W=A\rtimes C\), \(H^1(C,A^\vee)=0\), and, for each p,
H has a normal subgroup \(J_p\) of order prime to p with
\((A_p)_{J_p}=0\). The theorem identifies an actual quotient;
it does not provide a peripheral automorphism or a negative norm
solution on that quotient.

\begin{proof}
We first prove the conditional assertion. Since H is an actual
quotient of \(\Pi\), its maximal p-quotient has the stated cyclic
form. The normal p-prime subgroup \(J_p\) lies in \(K_p\), and
therefore acts trivially on \(M_p\). Put
\[
 T_p=\operatorname{Hom}_{\mathbb Z}(A_p,M_p),\qquad
 X_p=\operatorname{Ext}^1_{\mathbb Z}(A_p,M_p).
 \tag{A.58.4}
\]
The invariant group \(T_p^{J_p}\) is zero because
\((A_p)_{J_p}=0\). Also \(X_p^{J_p}=0\). Indeed, an invariant
abelian extension class has a strict \(J_p\)-action: the kernel
of the automorphism-lift extension is \(T_p\), and its second
cohomology with respect to the group of order prime to p is
zero. On the resulting finite abelian p-group extension, the
averaging idempotent is the identity on \(M_p\) and zero on
\(A_p\). Its image is exactly \(M_p\), so it retracts the
extension onto \(M_p\) and splits the underlying abelian
extension. This proves the claimed vanishing. Equivalently, the
two modules lie in distinct central-idempotent blocks over
\(\mathbb Z_p[J_p]\).

Coprime cohomology and the Hochschild--Serre sequence now give
\[
 H^i(H,T_p)=H^i(H,X_p)=0\qquad(i\geq0).
 \tag{A.58.5}
\]
For a trivial coefficient module \(M_p\), the universal
coefficient sequence is
\[
 0\longrightarrow X_p\longrightarrow H^2(A_p,M_p)
 \longrightarrow
 \operatorname{Hom}_{\mathbb Z}
       (\mathop{\bigwedge}\nolimits^2 A_p,M_p)
 \longrightarrow0.
 \tag{A.58.6}
\]
The last arrow records the commutator form. Its kernel consists
of abelian extensions, and its surjectivity follows directly by
an upper triangular cocycle on a cyclic decomposition of A.
The canonical quotient map in (A.58.1) thus lifts to a unique
H-invariant extension class, by (A.58.5). It defines a finite
group \(P_p\) with
\[
 1\longrightarrow M_p\longrightarrow P_p\longrightarrow A_p
 \longrightarrow1,
 \qquad P_p'=M_p,\qquad P_p^{\rm ab}=A_p.
 \tag{A.58.7}
\]
The automorphisms fixing both the quotient and the center form
\(T_p\). Equation (A.58.5) therefore lifts the H-action strictly
to \(P_p\). Set \(P=\prod_p P_p\).

We must attain this extension with the actual source relations.
In \(P\rtimes C\), choose lifts \(S_0,R_0\) of the prescribed
images of s,r in \(A\rtimes C\). Their powers lie in M.
For p different from two the s-action on \(M_p\) is trivial,
so divide the s-square residue by two and correct centrally.
For p different from three, correct the r-cube by division by
three. The remaining s-square \(z_2\in M_2\) is s-fixed and
hence H-fixed, since the action factors through the cyclic
two-quotient generated by s. If \(z_2\) has order \(2^j\),
pull back
\(C_{2^{j+1}}\to C_2\) along
\(C^{\rm ab}=C_6\to C_2\), and push its kernel generator to
\(-z_2\). The resulting central carry cancels the actual
s-square and does not alter the r-cube. The remaining r-cube
\(z_3\) is similarly H-fixed; the pullback of
\(C_{3^{j+1}}\to C_3\), with kernel sent to \(-z_3\), repairs
it without altering the s-square. The opposite-primary powers
were already repaired. These central pushouts retain P as the
kernel and give an extension of C by P with actual elements
S,R satisfying \(S^2=R^3=1\).

The source \(\Pi\) consequently maps to this extension. The
image of R surjects onto A, since it does so in the prescribed
quotient. Its commutators attain all of M by (A.58.7), so its
image is P and the homomorphism is surjective. The images of
\(W,W_1,\ldots\) in P are precisely
\[
 M,\quad IM,\quad I^2M,\ldots.
 \tag{A.58.8}
\]
On the two-primary part, \(I^2=-2I\), whence
\(I^{e_2+1}=0\). On the three-primary part,
\(I^3=-3I(1+I)\), giving
\(I^{2e_3+1}=(-3)^{e_3}I(1+I)^{e_3}=0\).
The other action ideals vanish. Thus this actual target kills
\(W_m\).

A.53 gives the reverse universal commutator quotient
\(M\twoheadrightarrow W/W_m\) and the asserted stability.
The actual target just constructed gives
\(W/W_m\twoheadrightarrow M\). Their composite is the
identity on every exterior-square commutator generator, and
therefore on M. The first map is consequently an isomorphism.
The explicit images (A.58.8) identify the entire filtration and
prove (A.58.3), rather than merely an upper bound for its order.

Finally, for the stated pronilpotent base, the actual normal
pro-p-prime factor \(N_{p'}\) of \(K_0/R\) has no invariants
or coinvariants on \(A_p\). This is the pro-p Nakayama argument
of A.53, applied to \(\overline{[R,K_0]}=R\).
Coprime averaging and continuous Hochschild--Serre therefore
give all-degree vanishing on C for \(A_p\) and its dual.
In particular \(H^1(C,A^\vee)=H^2(C,A)=0\). The actual
finite-kernel extension \(\Pi/W\to C\) has a continuous set
section: choose an open subgroup avoiding its finite kernel,
use the resulting local inverse, and choose finitely many
coset representatives. Its cohomology class is thus defined
and the vanishing of \(H^2(C,A)\) splits it. The image of
\(N_{p'}\) in H supplies the required normal p-prime subgroup
\(J_p\). All conditional hypotheses have now been checked.
\end{proof}


% Proposed after checkpoint07; independently reviewed root/CE calculation.
% Not yet incorporated into the main formal source.
\statement{A.59}{Proved}{The Bockstein quotient has an exact eighteen-dimensional first scalar refinement}
For the simultaneous actual quotient \(C_{\rm Bock}\) of A.57,
let \(W_{\rm Bock}=\ker(\Pi\to C_{\rm Bock})\). Its first
scalar refinement has kernel \(k^6\), where \(k=\mathbb F_8\),
and the following refinement is stable:
\[
 W_{\rm Bock}/\overline{[W_{\rm Bock},\Pi]}\simeq k^6,
 \qquad
 \overline{[[W_{\rm Bock},\Pi],\Pi]}
   =\overline{[W_{\rm Bock},\Pi]}.
 \tag{A.59.1}
\]
This identifies an actual additional central kernel of dimension
eighteen over \(\mathbb F_2\). No lifted negative norm or literal
normalizer assertion is included.

\begin{proof}
Retain \(\mathcal O=\operatorname{GR}(4,3)\), V and \(V_4\)
from A.57. The finite residual group of \(C_{\rm Bock}\to C\)
is the actual fiber product
\[
 P_{\rm Bock}=P\mathbin{\times}_{V^4}V_4^4.
 \tag{A.59.2}
\]
Its commutator subgroup is the old \(M=P'\), since the second
factor is abelian and the first projection is surjective.
Its abelianization is the coefficient pullback
\[
 \widetilde A=A\mathbin{\times}_{V^4}V_4^4
       \simeq V_4^4\oplus V^4.
 \tag{A.59.3}
\]
The last isomorphism is equivariant: the quotient map
\(A=V^8\to V^4\) is a split projection in the semisimple
isotypic module, with endomorphism field k. Write its integral
multiplicity module as \(T=\mathcal O^4\oplus k^4\).

Put \(K=O^2(H)\). We calculate the actual center of the abelian
layer \(\widetilde A\) using A.58. Over \(\mathcal O\), the
normal odd subgroup N has 36 distinct weights. Only opposite
weight pairs survive N-coinvariants of the exterior square.
The action of \(K/N\) is regular on these 36 weights; the
stabilizer of an unordered opposite pair is generated by
\(\beta_0\), which exchanges the two weights. In the exterior
square its action is minus-transpose. Thus one weight orbit
contributes exactly
\[
 \mathcal S(T)=
 (T\otimes_{\mathcal O}T)\big/
   \langle u\otimes v+v\otimes u\rangle.
 \tag{A.59.4}
\]
The diagonal terms have order two and must not be deleted.
For \(T=\mathcal O^4\oplus k^4\), the six free/free
off-diagonal terms contribute \(\mathcal O^6\), the remaining
22 off-diagonal terms contribute \(k^{22}\), and all eight
diagonal terms contribute \(k^8\). Hence
\[
 \mathcal S(T)=\mathcal O^6\oplus k^{30}.
 \tag{A.59.5}
\]

Here is the descent count, to avoid multiplying the answer by
an extraneous factor of three. After base change from
\(\mathbb Z/4\) to \(\mathcal O\), the original module has
three Frobenius weight orbits \(S,2S,4S\). They are distinct,
each is closed under negation, and no weight in different
orbits is opposite: such a relation would be a forbidden
Frobenius scalar normalizer of the chosen character. The
three base-changed copies of (A.59.5) therefore give
\(\mathcal O^{18}\oplus k^{90}\). Faithful unramified descent
means that the original finite abelian group is
\[
 \widetilde M=
 (\mathop{\bigwedge}\nolimits^2_{\mathbb Z}\widetilde A)_K
       \simeq C_4^{18}\oplus C_2^{90}.
 \tag{A.59.6}
\]
This is not three copies of the underlying group in
(A.59.6). The element s fixes the chosen weight and its
negative with scalar one; it commutes with the opposite-pair
transporter \(\beta_0\). Since s generates H/K, the entire
H-action on \(\widetilde M\) is trivial.

Reduction to the old elementary layer sends the six free
off-diagonal coordinates \(\mathcal O^6\) onto \(k^6\) and
is the identity on the other 30 coordinates. Consequently
\[
 \ker(\widetilde M\to M)=2\mathcal O^6=k^6.
 \tag{A.59.7}
\]
Let \(W_{\rm aff}\) be the kernel of the actual affine quotient
with residual abelianization \(\widetilde A\), and put
\(U_0=\overline{[W_{\rm aff},\Pi]}\). By A.58 and the trivial
action just proved,
\[
 W_{\rm aff}/U_0=\widetilde M,
 \qquad \overline{[U_0,\Pi]}=U_0.
 \tag{A.59.8}
\]
The actual image of \(W_{\rm aff}\) in (A.59.2) is the old
globally central M. Therefore \(U_0\subset W_{\rm Bock}\),
and (A.59.7) identifies
\(W_{\rm Bock}/U_0=k^6\). This quotient is central, giving
\(\overline{[W_{\rm Bock},\Pi]}\subset U_0\). The reverse
inclusion follows from the stability in (A.59.8). This proves
both equalities (A.59.1).
\end{proof}


% Separate post-checkpoint proposal, independently accepted root/CE/A.
% Integral argument reconstructed after the execution environment replacement.
\statement{A.60}{Proved}{The actual Bockstein quotient retains an involutory negative action}
Let h be a full positive peripheral lift realizing the A.56
quotient action and fixing \(\tau=sr\) literally. Put \(b=hc\).
Then the actual simultaneous Bockstein quotient
\(C_{\rm Bock}\) of A.57 is stable under h and c, and the
induced b-action on it has square one. Its action also has square
one on the actual scalar-stable quotient of A.59. This assertion concerns
the actual quotient kernel of A.57, independently of the chosen
integral cocycle representatives. It does not assert the literal
normalizer frames on that additional center or a full involution.

\begin{proof}
Write \(\mathcal O=\operatorname{GR}(4,3)\),
\(L=\mathcal O^{36}\), and
\(\rho:\Pi\to H\to\operatorname{Aut}_{\mathcal O}(L)\)
for the actual Teichm\"uller monomial representation. The
standard negative map satisfies
\[
 c(s)=s,\qquad c(r)=sr^{-1}s,\qquad c(\tau)=\tau^{-1}.
 \tag{A.60.1}
\]
Its monomial coefficient intertwiner C has
\(C^2=1\) and \(C\rho(g)C^{-1}=\rho(cg)\).
The positive map h induces the identity on H.

\emph{The integral cusp modules.}\par
The s-action has four fixed lines and sixteen exchanged pairs.
All entries are Teichm\"uller seventh roots. A fixed-line scalar
has order dividing two and seven, so is one; rescaling each
exchanged pair gives an ordinary permutation pair. Hence
\(L^s\) is free of rank20 and reduces onto \(V^s\).
For a monomial \(\tau\)-cycle, a nontrivial scalar product
gives an invertible \(1-\zeta\) and contributes no invariants
or coinvariants. A cycle of scalar product one can be rescaled
to an ordinary permutation cycle and contributes one free line
to each. There are four such cycles, giving
\[
 L^\tau\simeq\mathcal O^4,
 \qquad L/(\tau-1)L\simeq\mathcal O^4.
 \tag{A.60.2}
\]

The actual C normalizes \(\tau\) by inversion. Its mod-two
action on the four cusp lines is the identity in A.56. Thus it
fixes each admissible cycle. In a rescaled cycle basis it acts
by \(f_j\mapsto\zeta f_{a-j}\) for a Teichm\"uller seventh
root \(\zeta\). The same scalar acts on the invariant sum and
on the coinvariant generator. Reduction is injective on these
roots, so \(\zeta=1\). Consequently C is literally the
identity on both modules in (A.60.2). This uses the integral
monomial form; a general mod-two identity would not suffice.
The supplementary exact certificate is
\path{work_CE/code/check_Q6_integral_cusp.py}.

\emph{The free integral cusp cohomology.}\par
Set
\[
 \mathcal K_L=
 \ker\!\left(H^1(\Pi,L)\longrightarrow H^1(\langle s\rangle,L)\right).
 \tag{A.60.3}
\]
A class in this kernel has a representative d with \(d(s)=0\).
Its r-value lies in the free rank24 direct summand
\(\ker(1+r+r^2)\), since three is a unit in \(\mathcal O\).
The remaining coboundaries are the image of
\(L^s\to L\), \(a\mapsto(1-r)a\). This map is injective
modulo two: its kernel would be \(V^s\cap V^r=V^H=0\).
It therefore has a unit rank20 minor and a direct-summand
image. It follows that \(\mathcal K_L\) is free of rank four.
The cusp restriction map reduces to the isomorphism on the
four-dimensional kernel in A.56. Its source and target are
free of rank four over \(\mathcal O\); Nakayama gives
\[
 \mathcal K_L\xrightarrow{\ \sim\ }L/(\tau-1)L.
 \tag{A.60.4}
\]

Since h fixes \(\tau\) and is the identity on H, it acts as
the identity on \(\mathcal K_L\). It preserves the kernel
in (A.60.3) because it preserves the conjugacy class of s.
The coefficient-twisted c-action on cusp cohomology is
\(-C^{-1}\), from the inversion in (A.60.1). By (A.60.2)
this is minus the identity, including its integral sign.
Consequently, for each representative d of a class in
\(\mathcal K_L\), unique vectors \(w,v\in L\) satisfy
\[
 \begin{aligned}
 d(hg)&=d(g)+(1-\rho(g))w,\\
 d(cg)&=-C d(g)+(1-\rho(cg))v.
 \end{aligned}
 \tag{A.60.5}
\]
Uniqueness follows from \(L^\Pi=0\), which follows already
from the nontrivial odd weight action. Evaluating the first
equation at \(\tau\) gives \(w\in L^\tau\), so \(Cw=w\).
Applying the second equation twice gives
\((1-\rho(g))(v-Cv)=0\) for all g, hence \(Cv=v\).

Put \(q=v+w\). Since h is the identity on H, composition
in (A.60.5) gives
\[
 d(bg)=-C d(g)+(1-\rho(bg))q,
 \qquad
 d(b^2g)=d(g)+(1-\rho(g))(q-Cq)=d(g).
 \tag{A.60.6}
\]
These are equalities of actual cocycle values, not only of
cohomology classes. On \(C_{\rm new}\) the b-square is already
one. Applying (A.60.6) to four integral lifts of a basis of
the cusp kernel proves that its square is one on their joint
actual pullback. The first equations likewise show that h
and c preserve its source kernel.

\emph{The pullback is the prescribed A.57 quotient.}\par
Normalize the chosen mod-two representatives \(e_i\) by
\(e_i(s)=0\) before defining the pullback. Integral kernel
classes can be chosen to reduce to them. A lifted coboundary
restores literal equality of representatives if needed; its
restriction to \(R_{\rm new}\) is zero. The surjection
\(L^s\to V^s\) permits this while retaining the normalization.
If d and \(d'\) are two integral lifts of the same literal e,
then \(d-d'=2t\) for a V-valued cocycle t. Every class of
\(H^1(\Pi,V)\) inflates from \(C_{\rm new}\), by A.57.
An inflated cocycle and a coboundary both vanish on
\(R_{\rm new}\), so \(t|_{R_{\rm new}}=0\).
Thus all such lifts have exactly the same residual map and
the same joint source kernel. The involutory action just
proved is therefore on the actual A.57 quotient.

Finally put \(N=\ker(\Pi\to C_{\rm Bock})\). The kernel
\(N/[N,\Pi]\) of its actual first scalar refinement is central.
Since b preserves N and has square one modulo N, the map
\[
 g\longmapsto b^2(g)g^{-1}\pmod{[N,\Pi]}
 \tag{A.60.7}
\]
is a continuous homomorphism into that central kernel. It kills
\(\tau\), which generates \(\Pi^{\rm ab}=C_6\), and therefore
vanishes. A.59 identifies this refinement with the actual
scalar-stable quotient. This is the central-character argument
of A.43; it does not lift prescribed involutory normalizers.
\end{proof}


% Post-freeze proposal. Retains the earlier tower but refines the newly
% introduced A56 central frame choices; it does not retain the stored tuple.
\statement{A.61}{Proved}{The Bockstein scalar refinement admits simultaneous involutory frames after refining the new central frame choices}
The parity-one quotient action of A.56 extends to the actual
scalar-stable quotient above \(C_{\rm Bock}\) in A.59, with
simultaneously involutory s- and r-normalizers. It retains the
literal C-frames of A.33 and the earlier A.47--A.52 finite
families. A full positive peripheral lift realizes this quotient
action and these frames. Its full negative square is not asserted.

The new central frame choices introduced in A.56 may change.
In particular this assertion does not preserve the particular
stored \path{Q6_actual_central_solution.json} frame tuple, nor
a separately fixed newer cross-section transporter. The quotient
action itself is retained.

\begin{proof}
Let \(D=\Pi/U_0\) be the scalar-stable actual quotient of A.59,
and write
\[
 N=\ker(D\to C_{\rm Bock})=2\mathcal O^6\simeq k^6,
 \qquad
 \widetilde M=\mathcal O^6\oplus k^{30},
 \tag{A.61.1}
\]
with the integral descent understood as in A.59. The latter
group is central in D. Choose the explicit integral class-two
model of \path{check_A59_integral_center.py}. Its finite action
pulls back along \(C\to H\), paired with the prescribed
\(h_C\), whose H-action is the identity. The exact A.58
identification makes its source kernel \(U_0\). Thus we select
an action on this actual quotient before applying A.3; we do
not claim that every arbitrary full A.56 lift has the same
integral action. On its Bockstein projection the negative
action b has square one by the integral calculation of A.60.
It also has square one on D. Indeed, the map
\(g\mapsto b^2(g)g^{-1}\) on D is a central homomorphism
\(\Pi\to N\). It kills \(\tau\), since b inverts \(\tau\)
literally. The class of \(\tau\) generates
\(\Pi^{\rm ab}=C_6\), so this homomorphism is zero. This is
the scalar-center argument already used in A.43.

The b-action on the six free cusp/cusp coordinates
\(\mathcal O^6\) of \(\widetilde M\) is the identity.
The positive action fixes integral cusp cohomology; the
negative action changes each of its two factors by minus one.
The actual monomial coefficient involution fixes the selected
weight and its negative with scalar one. Thus it adds no
coefficient sign to their exterior-square coordinate. In
particular the central norm on these coordinates is
\[
 z\longmapsto z+b(z)=2z,
 \qquad 2\mathcal O^6=N.
 \tag{A.61.2}
\]
The action on N itself is the identity.

\emph{Separate literal frames on the Bockstein quotient.}\par
We record the integral frame argument, including the r-frame.
In one coefficient factor the action from A.60 is
\[
 b(p,g)=\bigl(-Cp+(1-\rho(bg))q,bg\bigr),\qquad Cq=q.
 \tag{A.61.3}
\]
For the s-frame \(\beta\), put \(D_s=\rho(\beta)C\).
A lift \((p,\beta)\) has the required conjugacy and norm if
and only if \(t=p-q\) is fixed by both s and \(D_s\).
These are commuting monomial involutions. Their stabilizer
characters take values in seventh roots of unity, so are
trivial. Rescaling each orbit makes a permutation module;
its integral invariants reduce onto the mod-two invariants.
Every prescribed old literal mod-two frame therefore lifts.

For r, first make a coboundary change
\(d'=d-(1-\rho)a\) with \(d'(r)=0\), possible because
three is a unit. Then \(q'=q-a-Ca\), still C-fixed. If J is
the prescribed r-normalizer, put
\(R'=\rho(b(r))\) and \(D_r=\rho(J)C\).
The lift equations are exactly
\[
 R'(p-q')=p-q',\qquad D_r(p-q')=p-q'.
 \tag{A.61.4}
\]
Here \(D_r^2=1\), \((R')^3=1\), and
\(D_rR'D_r=(R')^{-1}\). This monomial group is a quotient
of \(S_3\). Its stabilizers also have no nontrivial character
into the seventh roots, so the same integral invariant
argument applies. Undoing the coboundary change restores the
actual source r and the literal old J. This constructs both
separate involutory frames on \(C_{\rm Bock}\), without a
cross-section conjugacy assertion.

\emph{Conjugacy equations on the actual center.}\par
In the selected pure intermediate model, the original-law
certificate \path{check_A59_integral_center.py} verifies the
s-conjugacy equation before any central norm adjustment.
It includes the actual source relations, central C4 carry,
odd r-correction and transported old central frame. The
centralizer step also has an intrinsic explanation. The
difference between two affine s-conjugators with the same
old reduction lies in the doubled free cusp kernel and has
the form \((2u,0)\), with u s-fixed modulo two. The source s
has zero free cusp coordinate. Mixed commutators reduce to
zero, and quadratic terms in the doubled coordinates are
divisible by four. Thus every such difference lifts to an
actual s-centralizer on D. This explains why the chosen
joint-invariant correction preserves the actual conjugacy
equation rather than only its affine image.

Every lift of the Bockstein r-frame already has the correct
conjugacy equation in D. Any discrepancy between its conjugate
of r and the prescribed inverse is central in N; both elements
have order three. Their central discrepancy has cube one and
lies in a two-group, hence is zero. This argument does not
assume surjectivity of a two-torsion centralizer map.

Choose such conjugators \(\widetilde\beta,\widetilde J\).
The squares of their negative normalizers lie in N, since
their images on \(C_{\rm Bock}\) are involutions. Denote them
by \(n_s,n_r\). Choose independent central elements
\(z_s,z_r\in\mathcal O^6\) with
\(2z_s=n_s\), \(2z_r=n_r\). Replace
\[
 \widetilde\beta\longmapsto\widetilde\beta z_s,
 \qquad \widetilde J\longmapsto\widetilde J z_r.
 \tag{A.61.5}
\]
Centrality preserves both conjugacy equations, while
(A.61.2) makes both normalizer squares one. The choices are
independent and leave the action and the literal cusp fixed.
Their reductions change only the new cusp/cusp central
coordinates of the A.56 frames. Those changes themselves have
zero norm on the earlier quotient. This is the precise reason
that the particular stored frame need not extend although
the existential framed construction does.

\emph{Retained data and actual realization.}\par
The new residual abelianization is
\(V_4^4\oplus V^4\); all its simple constituents have
dimension108. The earlier residual abelianizations have only
the6-,21-, and42-dimensional types, and the earlier perfect
factors contribute no abelian quotient. The same equivariant
Goursat argument as A.56 therefore gives the actual full fiber
product over C. The absence of a common simple quotient also
gives the required \(H^1\)-vanishing over the retained base.
Applying the conditional A.58 theorem there identifies the
same actual center and proves its scalar stability; no
stability of an arbitrary intersection is presumed.

For a translated old s-frame, use the full identity-C R2
gauge of A.56. It preserves the full integral cusp
cohomology space and the full collection of binary s-classes,
without requiring that each binary class be fixed pointwise.
It therefore preserves the
actual affine kernel and its scalar commutator kernel.
It therefore induces an automorphism of D. Apply it to
the new component of the actual fiber product and leave
the old component fixed. This transports the pure construction
back to the original literal C-frame, including all its odd
translations. The central corrections (A.61.5) are trivial
on C and on the old component. Thus every earlier prescribed
frame and quotient action is retained.

Finally \([U_0,\Pi]=U_0\). A.3 realizes the entire action
and its literal s/r/cusp frames by a full positive lift.
Finite unions of the earlier families and compactness give
the same simultaneous retention as A.56. The retained
parity is one, because its C-image is unchanged. None of
these arguments gives a full order-two element or cofinality.
The simultaneous original-law witness and read-only verifier
are \path{A61_integral_frame_solution.json} and
\path{check_A61_final_frames.py}; their producer is
\path{check_A61_integral_frames.py}.
\end{proof}


% Complete all-exponent theorem, independently reviewed in working form.
% Do not integrate until root approves the final TeX.
\statement{A.62}{Proved}{A compatible unbounded integral cusp tower admits parity-one involutory images and simultaneous frames}
The construction of A.61 extends to every
\(\mathcal O_n=\operatorname{GR}(2^n,3)\), simultaneously.
There are actual nested scalar-stable kernels
\(U_{n+1}\subset U_n\) in \(\Pi\), with quotient
\(D_n=\Pi/U_n\) over C, whose residual abelianization and
center have types
\[
 A_n=L_n^4\oplus V^4,\qquad L_n=\mathcal O_n^{36},
 \qquad
 M_n\simeq C_{2^n}^{18}\oplus C_2^{90}.
 \tag{A.62.1}
\]
A full positive lift has parity-one negative image of order
two on every \(D_n\), with compatible involutory s/r-normalizers,
while retaining all the earlier C and A.47--A.52 data.
The newer A.56 central frames may vary; surjectivity onto
arbitrary earlier choices is not asserted.

The kernel \(U_\infty=\bigcap_nU_n\) is itself relatively
perfect. Thus the compatible data also form an actual
closed-normal quotient to which A.3 applies. The tower is
not claimed cofinal, and no full order-two element follows.

\begin{proof}
\emph{Compatible actual quotients.}\par
Let \(\mathcal O_\infty\) be the unramified degree-three
extension of \(\mathbb Z_2\). The monomial representation
on \(L_\infty=\mathcal O_\infty^{36}\) uses the same
permutations and the compatible Teichm\"uller seventh roots.
Its finite reductions are the actual \(L_n\). The integral
calculation of A.60 works over \(\mathcal O_\infty\): after
normalizing the s-value to zero, the r-norm kernel is free
of rank24, the s-fixed coboundaries have a unit rank20 minor,
and the cusp restriction is an isomorphism on the resulting
free rank-four quotient. Choose a basis of that quotient and
actual representative cocycles once over \(\mathcal O_\infty\).
This is also a direct construction with the profinite source:
their s/r-values satisfy the source relations in the
semidirect product, so define continuous cocycles on \(\Pi\).

Reduce these four cusp cocycles modulo \(2^n\), and retain
separately the four binary s-classes with zero cusp restriction.
They define an actual map
\[
 \Pi\longrightarrow A_n\rtimes C.
 \tag{A.62.2}
\]
Its residual image reduces onto the complete eight-dimensional
V-cocycle layer. Finite two-group Nakayama therefore makes
that image all of \(A_n\), so (A.62.2) is surjective.
Write \(W_n\) for its kernel. The coefficient transitions
reduce the four free coordinates and leave the four binary
coordinates unchanged, giving \(W_{n+1}\subset W_n\).
Keeping the binary coordinates separately is essential:
the torsion s-classes need not lift under reduction of the
single coefficient module \(L_{n+1}\) alone.

The actual normal pro-odd subgroup gives the initial-base
cohomology vanishings for \(A_n\) and its dual. A.58
identifies the scalar center and proves stability. In
weight notation its integral multiplicity module is
\(T_n=\mathcal O_n^4\oplus k^4\); opposite-weight
coinvariants give
\[
 (T_n\otimes T_n)/\langle u\otimes v+v\otimes u\rangle
       =\mathcal O_n^6\oplus k^{30}.
 \tag{A.62.3}
\]
The six free terms are off-diagonal; all eight diagonal
terms remain of order two. The same Frobenius descent as
A.59 gives the underlying group in (A.62.1). H acts
trivially, so with \(U_n=[W_n,\Pi]\) one has
\([U_n,\Pi]=U_n\). These are actual source kernels.
Their nestedness supplies actual quotient transitions
\(D_{n+1}\to D_n\). Equivalently, the explicit class-two
laws reduce compatibly: use the same s-carry and the
compatible correction by minus one third of the r-cube.

\emph{Compatible quotient actions.}\par
Choose one full positive lift \(h_0\) supplied by A.56,
with its older retained data, and set \(b_0=h_0c\).
Its binary action is the prescribed A.56 action. On every
integral cusp cohomology module, the positive action is
the identity and the c-action is minus the identity.
The proof of A.60 works over each \(\mathcal O_n\): the
actual coefficient C is the identity on cusp invariants
and coinvariants, the unique vectors satisfy \(Cv=v\),
\(Cw=w\), and the exact cocycle formula gives
\(d(b_0^2g)=d(g)\). Thus \(h_0,c\) preserve every
\(W_n\), and \(b_0^2\) is one on the affine quotient.
The central-character argument killing \(\tau\) gives
\(b_0^2=1\) on \(D_n\) as well. These actions are
compatible because they all come from the same full
\(h_0\). On the six free center coordinates, \(b_0\)
acts as the identity.

\emph{The integral s-centralizer gate.}\par
No unordered opposite-weight pair is exchanged by s.
Indeed the chosen weight \(\lambda\) has stabilizer
\(\langle s\rangle\), and \(\beta_0\lambda=-\lambda\).
If \(s(q\lambda)=-q\lambda\), then
\[
 q^{-1}sq\in\beta_0\langle s\rangle
             =\{\beta_0,\beta_0s\}.
 \tag{A.62.4}
\]
In \(Q_6=S_3\times A_4\), \(\beta_0\) has trivial
\(S_3\)-component, whereas s does not; \(\beta_0s\)
has trivial \(A_4\)-component, whereas the component of s
is a nontrivial double transposition. Neither can be
conjugate to s. Odd translations affect scalar phases,
not the weight permutation. The exact finite certificate
\path{check_A62_pair_orientation.py} checks all36 weights
and constructs the orientation used below.

Orient the18 unordered opposite pairs s-invariantly.
On a two-element orbit of pairs, transport an arbitrary
orientation by s; on a fixed pair, its two weights are
fixed separately by (A.62.4). The monomial lift of s in
this coordinate convention has no orientation carry.
If \(a\in2L_n^4\) is s-fixed, its mixed commutators with
the binary source coordinate vanish, since a reduces to
zero. The source s has zero free cusp coordinate.
Consequently \((a,0)\) centralizes the actual source s
in the class-two group, for every n. This proves the
required all-exponent centralizer lift. It does not rely
on a term divisible by four vanishing at higher exponent.

\emph{Nonempty simultaneous frame sets at each level.}\par
Start with the actual full conjugators of \(h_0\).
The integral invariant argument of A.61 supplies affine
s/r-frames with involutory normalizers over
\(\mathcal O_n\), reducing to their A.56 affine values.
The monomial groups are quotients of \(C_2\times C_2\)
and \(S_3\). Their stabilizers have no nontrivial
characters into the seventh roots, so orbitwise
rescaling gives permutation modules and surjective
reduction on joint invariants for every n.

For s, write the desired affine conjugator on the
\emph{right} of the full one: \(\beta_{\rm aff}
=\beta_{0,\rm aff}k\). The two conjugators have the
same image on the old affine quotient, so k lies in
\(2L_n^4\); their conjugacy equations say that k is
fixed by the original s. The preceding centralizer
lemma lifts k on the right and preserves the actual
s-conjugacy equation. For r, any lift of an affine
conjugator has the correct conjugacy equation: a
central discrepancy between two elements of order
three lies in a two-group and therefore is zero.

Choose the central coordinates to reduce to the old
A.56 frames. Both normalizer squares now lie in
\(\ker(M_n\to M_1)=2\mathcal O_n^6\).
Their binary central coordinates vanish. The central
norm is multiplication by two on these six free
coordinates. Independent central halves kill both
squares, changing neither conjugacy equation nor
the action. These halves may change newer central
frame coordinates at smaller levels; all their
normalizer equations still hold there by reduction.
They change none of the original C or earlier-family
frame projections.

For each finite union of older families the new
residual abelianization has only the108-dimensional
simple type, whereas the older abelianizations have
only the6/21/42-dimensional types. The same Goursat
and cohomology arguments as A.61 give the actual
fiber product and the hypotheses of conditional
A.58 over the retained base. Its scalar kernel is
stable. A.3 realizes the chosen action and frames
by a full positive lift.

\emph{Compatibility and the actual inverse limit.}\par
Take closed solution sets in the compact product of
the full positive group and the two full frame
groups. Require the older actions and literal
frames, the fixed \(h_0\)-action on \(D_n\), literal
cusp fixation, and order-two s/r-normalizers on
\(D_n\). Any finite list of requirements is handled
by the maximum n and the finite union of older
families just considered. Hence the sets have the
finite-intersection property. A simultaneous point
gives the compatible data asserted in the theorem.
No surjectivity between particular frame-solution
sets has been assumed.

Finally \(U_\infty\) is relatively perfect. In any
finite quotient q of \(\Pi\), the descending groups
\(q(U_n)\) stabilize. Compactness gives
\(q(U_\infty)=\bigcap_n q(U_n)\), and the stable
image equals its commutator with \(q(\Pi)\), since
each \(U_n\) is relatively perfect. This holds in
every finite quotient, proving
\[
 [U_\infty,\Pi]=U_\infty.
 \tag{A.62.5}
\]
Thus the actual inverse-limit data also satisfy
the closed-normal A.3 gate. Their residual is a
specific class-two pro-two group with the displayed
finite-rank coefficient types over C. No argument
makes this tower cofinal or its full lift an
involution.
\end{proof}


% Actual next-chief construction beyond A62. Root/CE reviewed working proof.
\statement{A.63}{Proved}{Six independent cross-cup directions give an actual nonsplit next quotient beyond the integral tower}
Let \(D_\infty\) be the actual quotient of A.62, and let
\(V^{(2)}\) denote the monomial H-module obtained from V
by squaring every scalar and retaining each permutation.
There is an actual source quotient
\[
 1\longrightarrow T\longrightarrow E\longrightarrow D_\infty
 \longrightarrow1,\qquad T=(V^{(2)})^6,
 \tag{A.63.1}
\]
with six k-linearly independent nonsplit coordinate extension
classes. They are the cross cup products of the four cusp
classes. The joint source kernel is preserved by the full
positive lift h supplied at the end of A.62 and by c.
The same extension already descends from the A.56 quotient
\(D_1\) over the initial profinite C. These quotients need
not be finite; their explicit finite group models pull back
over C. Neither a negative-square solution nor new literal
normalizer frames are asserted.

\begin{proof}
\emph{An actual finite group.}\par
Put \(A=V^4\), index the six coordinates of T by
\(1\leq i<j\leq4\), and write \(*\) for coordinatewise
multiplication on \(k^{36}\). On \(N=A\times T\) define
\[
 (a,z)(b,w)=(a+b,z+w+c(a,b)),
 \qquad c(a,b)_{ij}=a_i*b_j.
 \tag{A.63.2}
\]
Bilinearity proves associativity. The square is
\((0,c(a,a))\), and the commutator coordinates are
\(a_i*b_j+a_j*b_i\). Single-row, single-weight vectors
attain every coordinate, so
\[
 N'=T=\Phi(N),\qquad N/\Phi(N)=A.
 \tag{A.63.3}
\]
The group has order \(2^{1080}\). H acts on A by its
original monomial matrices and on T by their squared
scalars. This preserves the cocycle in (A.63.2) literally.
Thus it gives a strict action on an actual group, rather
than only an action on a graded bracket.

Coordinate Frobenius identifies \(V^{(2)}\) with V as
an \(\mathbb F_2[H]\)-module, but it is not the identity
of the literal k-representations. We keep the twist in
all coefficient maps. Its invariants and coinvariants
vanish under the same normal odd subgroup, so
\(T^H=T_H=0\).

\emph{Source generators and exact residual images.}\par
Reduce the four normalized integral cusp cocycles of
A.62 modulo two. Their s-value is zero and their r-value
\(e(r)\in A\) has r-norm zero. In characteristic two,
\(a=r e(r)\) satisfies \((1-r)a=e(r)\). Choose
\(v=(a,0)\in N\). In \(N\rtimes C\), take
\[
 S_*=(1,s_C),\qquad R_*=v(1,r_C)v^{-1}.
 \tag{A.63.4}
\]
Their orders divide two and three, respectively, and
the affine source cocycle is literally the chosen
four-tuple. Hence they define a continuous source map
\(q:\Pi\to N\rtimes C\).

Its original residual image projects onto A. Otherwise
a nonzero H-linear functional on the missing semisimple
quotient of \(V^4\) would give a nonzero combination of
the four cusp classes vanishing on the original R.
That class would inflate from C, contradicting
\(H^1(C,V)=0\). The finite two-group Frattini theorem
then gives \(q(R)=N\), proving actual surjectivity.

Let \(A_\infty=L_\infty^4\oplus V^4\) be the full
abelian residual retained in A.62, including the
separate binary s-coordinates. Form
\[
 B=A_\infty\mathbin{\times}_A N.
 \tag{A.63.5}
\]
Its derived subgroup is T, its abelianization is
\(A_\infty\), and its Frattini quotient is \(V^8\).
The two source cocycles have identical A-values,
so pair into (A.63.5). The actual R-image is onto
the complete eight-class Frattini quotient. Since
B is a finitely generated pro-two group, its closed
R-image is all B. If \(W_\infty=\bigcap_nW_n\), this
proves \(q(W_\infty)=T\) exactly.

Now
\[
 [W_\infty,\Pi]\subset U_\infty\subset W_\infty.
 \tag{A.63.6}
\]
Taking images gives
\(T=[T,H]\subset q(U_\infty)\subset T\), hence
\(q(U_\infty)=T\). Therefore the actual joint quotient
is the full fiber product
\[
 E=D_\infty\mathbin{\times}_{A\rtimes C}(N\rtimes C).
 \tag{A.63.7}
\]
For explicit surjectivity, lift a D-element by \(\Pi\)
and correct its remaining T-coordinate by an element
of \(U_\infty\). No interchange of an infinite
intersection with a commutator subgroup is used.
Since \(U_\infty\subset U_1\), the same argument gives
\(q(U_1)=T\) and the analogous actual extension over
\(D_1\).

\emph{Nonsplitting and independence.}\par
Write \(P_\infty\) for the residual group of
\(D_\infty\to C\). Its closed derived subgroup is
\(M_\infty\), globally central in \(D_\infty\).
Indeed, the closed derived subgroup surjects onto
every finite center \(M_n\), so compactness identifies
it with their inverse limit.

A continuous section of (A.63.1), followed by
projection to \(N\rtimes C\), would restrict to a
homomorphism \(\eta:P_\infty\to N\) surjective modulo
\(\Phi(N)\). It would therefore be onto N, giving
\(\eta(P_\infty')=T\). But elements of
\(M_\infty=P_\infty'\) are central in \(D_\infty\).
Their images under the section commute with its
entire image, whose H-projection is onto. Their
N-components lie in T and are H-fixed, hence are
zero. This is a contradiction. No splitting of
\(D_\infty\) over C was assumed.

For every nonzero k-linear functional
\(T\to V^{(2)}\), the commutator of its pushed-out
group is still the entire simple target. The same
argument proves that this pushout is nonsplit.
Consequently all six coordinate extension classes
are k-linearly independent. The identical argument
works over \(D_1\). This counts six actual directions,
not all possible next extensions.

\emph{Cohomology and preservation of the source kernel.}\par
Initial-base acyclicity and continuous
Hochschild--Serre give
\[
 H^1(D_\infty,V^{(2)})
  =\operatorname{Hom}_C(A_\infty,V^{(2)})\simeq k^8.
 \tag{A.63.8}
\]
Inflation to \(H^1(\Pi,V^{(2)})\) is the actual
complete eight-class isomorphism. The five-term
sequence therefore makes transgression injective
on \(\operatorname{Hom}_{D_\infty}(U_\infty,V^{(2)})\),
with image the kernel of the map from
\(H^2(D_\infty,V^{(2)})\) to \(H^2(\Pi,V^{(2)})\).

The extension cocycle in (A.63.2), after pullback,
has coordinates
\[
 e_i(g)*\bigl(g e_j(h)\bigr),\qquad i<j.
 \tag{A.63.9}
\]
These are exactly the cross cup products for the
H-equivariant coordinate multiplication
\(V\otimes V\to V^{(2)}\). The actual full framed
lift h supplied at the end of A.62 fixes the four
cusp cohomology classes. The coefficient-twisted
c-action is minus the identity integrally, hence
the identity modulo two; its intertwiner on the
target is the squared monomial intertwiner.
Naturality of cup products preserves all six
classes. Naturality and injectivity of transgression
then preserve their actual residual maps, with
these coefficient identifications, and therefore
their joint source kernel. This proves the claimed
h/c stability, including the analogous finite-base
statement over \(D_1\).

The next square is not settled by this argument.
Its discrepancy is a continuous one-cocycle on
\(D_\infty\) with values in T, since the square is
one on the base and on the coefficient module.
Literal cusp inversion kills its cusp restriction.
For the full s-frame \(\beta\) supplied by A.62,
\(b(\beta)\beta\in T\), and the discrepancy at s is
\((1-s)(b(\beta)\beta)\). Its s-restriction is therefore
a coboundary. The eight-class restriction isomorphism
then makes the discrepancy inner by a vector in
\(T^\tau\simeq k^{24}\), unique because \(T^H=0\).
A zero cohomology class does not make this literal
inner vector zero. No norm or frame conclusion is
added here.
\end{proof}

% Exact witness after varying the old affine datum. Parent/CE review requested.
\statement{A.64}{Proved}{A changed affine frame passes all six cross-cup norm equations}
There is a choice of the A.56 affine action and frames for which the
actual extension E of A.63 admits a positive peripheral automorphism \(h_E\)
over the A.62 action, with negative companion \(b_E=h_Ec\) satisfying
\[
 b_E^2=1,\qquad h_E(\tau)=\tau.
 \tag{A.64.1}
\]
There are frames \(\beta_E,J_E\) above the prescribed initial C-frames
such that
\[
 \begin{aligned}
 h_E(s)&=\beta_Es\beta_E^{-1},&
 \beta_E^{-1}b_E(\beta_E^{-1})&=1,\\
 b_E(r)&=J_Er^{-1}J_E^{-1},&
 J_E^{-1}b_E(J_E^{-1})&=1.
 \end{aligned}
 \tag{A.64.2}
\]
The choice retains the initial C-data and the older families of
A.47--A.52. It changes the discretionary A.56 affine datum; it does
not retain the original stored A.56 tuple. The construction is on
the actual quotient E, including its full integral tower. This
statement alone does not invoke the scalar-refinement lifting criterion.

\begin{proof}
\emph{The choice at the previous level.}\par
Use the normalized eight cocycles of A.57, with four cusp rows first
and four binary s-rows last. Write q for the affine inner term of the
negative action and p for the affine s-frame. We choose the four free
cusp components of
\[
 t=p-q
 \tag{A.64.3}
\]
to be zero. These are conditions on the actual affine frame, not
conditions merely on its cohomology class.

The original A.56 affine correction space has 360 binary coordinates.
The old central conjugacy equations have rank 83. Adding the four
literal vector equations in (A.64.3) gives rank 227 and leaves an
affine space of dimension 133. The exact simultaneous linear system
and its solution are recorded in
\path{records/A63_zero_cusp_frame_feasibility.json} and checked by
\path{code/check_A63_zero_cusp_frames.py}; its kernel is recovered by
exact elimination in the norm certificate below. Substitution into the
original group law checks the literal cusp, the old s-conjugacy
equation, the negative action square, and (A.64.3).

On this 133-dimensional affine space, the remaining old central
normalizer equation has zero polar form. Its linear image has
dimension ten, and its target belongs to that image. Thus it has
an exact solution. The calculation uses the original class-two
group law; it is not an inference from a graded approximation.
The resulting tuple, with both old frame norms corrected, is
\path{work_A/A63_zero_t_A56_solution.json}. Its SHA-256 is
\begin{center}\small\ttfamily
4135fabf1500006d5ec6bbe5dc0571ea91018abf911e32136cf965729d18d4e9.
\end{center}
The producer is \path{work_A/check_A63_zero_t_norm.py}; the
original A.56 witness is unchanged.

\emph{The actual six-cup action.}\par
Let \(\rho_2\) be the squared-scalar H-representation on
\(V^{(2)}\). Every coefficient cocycle used below has the explicit
form
\[
 \eta(g)=\sum_{j=1}^8\lambda_j e_j(g)^{[2]}
              +(1-\rho_2(g))v,
 \qquad (\lambda_1,\ldots,\lambda_8,v)\in k^8\oplus k^{36},
 \tag{A.64.4}
\]
where the brackets denote coordinatewise squaring. These are actual
cocycles on the quotient over C. They are not formal values assigned
only to source generators. The source-value matrix has rank 44.

Use the group N and its literal multiplication from (A.63.2).
The finite model for the extension over \(D_1\) consists of
\((g,z)\), with g in the old model and \(z\in T\), and law
\[
 (g,z)(h,w)=
 \bigl(gh,z+\rho_2(g)w+c(e(g),g e(h))\bigr).
 \tag{A.64.5}
\]
It pulls back over the full profinite C. First lift canonical c by
its monomial intertwiner and the affine inner term, and then add
six cocycles of the form (A.64.4). The 44 independent source
equations make this lift agree literally with
\(c(s)=s\), \(c(r)=sr^{-1}s\). Next lift the chosen positive
action by its affine inner term and add another six cocycles of
the same form. Literal cusp preservation contributes 36 equations
per coordinate; the s-conjugacy class contributes four. Their
combined rank is 40. All their constants are taken from (A.64.5).

The stored coefficients in
\path{work_A/A63_zero_t_cross_cup_square.json}
solve these equations. The resulting negative action has zero
literal square defect on both source generators. For the s-frame,
the coefficient centralizer has dimension 20 over k and its norm
image has dimension eight. All six obstruction classes in the
four-dimensional quotient vanish. For the r-frame, the corresponding
dimensions are twelve and five, and all six norm equations also
have solutions. The stored frames satisfy (A.64.2) under the
original multiplication (A.64.5).

The read-only verifier
\path{work_A/check_A64_cross_cup_witness.py}
does not solve for or modify the stored new witnesses. Its imported
old-model initialization reconstructs the coefficient matrices.
It checks the
source relations, the literal canonical c-action, cusp preservation,
both source square equations, both frame projections, both
conjugacy equations, and both normalizer norms. The six-cup witness
has SHA-256
\begin{center}\small\ttfamily
f5d10fa9f5491651fef0e431edab7a952bf18fa1890073cd9d87c9a8c690de50.
\end{center}
Its stdout is \path{work_A/check_A64_cross_cup_witness.out}.
Each map in this verification is an automorphism: it is a compatible
old automorphism and monomial or inner N-automorphism, followed
by a cocycle shear trivial on T. Thus checking the displayed source
identities proves identities on the actual source quotient.

\emph{Integral and older retained data.}\par
The changed free positive inner vector is still fixed by the cusp
and by the negative coefficient involution. It therefore lifts in
the integral cusp invariant module used in A.60--A.62. The free
s-frame correction (A.64.3) lifts literally by zero. The all-level
orientation argument and central-half corrections of A.62 apply
to these choices. Such corrections change only the old globally
central \(M_n\)-coordinates; they have no T-coordinate in the
fiber product (A.63.7). Consequently the stored six-cup corrections
and frames pull back compatibly through the entire integral tower.

The earlier residual abelianizations have only the simple types
6, 21, and 42, whereas this residual abelianization has only the
108-dimensional type and its integral lifts. Any common nontrivial
finite two-group quotient would have a common simple abelian
quotient. Equivariant Goursat therefore gives the same actual fiber
product with every previously retained finite family. The compact
finite-common-refinement argument of A.62 retains all those families
and the original literal C-frames. No assertion is made that the
old stored affine tuple or a later chosen transporter remains fixed.

There is also a useful full-lift consequence without a new scalar
argument. The actual full lift supplied by A.62 induces an E-action
with the same old affine action, the literal cusp, and the same
old s-frame projection. Comparing it with the displayed action gives
a T-cocycle whose cusp and s cohomology restrictions vanish. The
rank-eight restriction isomorphism makes this difference inner by
\(a\in T^\tau\). The coefficient negative action is the identity
on \(T^\tau\); hence replacing \(b_E\) by \(\operatorname{Inn}(a)b_E\)
does not change its square. Its square is therefore also one on E.
This does not by itself identify involutory lifts of the two new
frames in the full group. The scalar-refinement question is separate.
\end{proof}

% Actual scalar closure of the cross-cup layer; independent review requested.
\statement{A.65}{Proved}{Jacobi kills the potential scalar third layer and realizes the cross-cup frames fully}
For each finite integral level n of A.62, form the actual cross-cup
extension \(E_n\) of A.63, and write
\[
 W'_n=\ker(\Pi\longrightarrow E_n),\qquad
 W'_\infty=\bigcap_n W'_n.
 \tag{A.65.1}
\]
Then
\[
 [W'_n,\Pi]=W'_n\quad\hbox{for every }n,
 \qquad [W'_\infty,\Pi]=W'_\infty.
 \tag{A.65.2}
\]
Consequently the closed-kernel lifting criterion A.3 realizes the
actual action and both frames of A.64 by a full peripheral positive
automorphism of \(\Pi\), retaining the initial C-data and all the
older families. Its negative companion has square one on the entire
cross-cup extension of the integral tower. Its square in the full
group, and cofinality of the constructed tower, remain open.

\begin{proof}
\emph{The original base and the actual finite scalar problem.}\par
Keep the original quotient \(C=\Pi/R\). Write
\[
 A_n=L_n^4\oplus V^4,\qquad
 P_n'=M_n=(\bigwedge\nolimits^2_{\mathbb Z}A_n)_{O^2(H)},
 \qquad M_n\simeq\mathcal O_n^6\oplus k^{30}.
 \tag{A.65.3}
\]
Here the notation on the right specifies the abelian group and its
trivial H-action, as in A.58--A.62. The actual residual group of
\(E_n\to C\) is
\[
 Q_n=P_n\mathbin{\times}_{V^4}N,
 \qquad Q_n^{\rm ab}=A_n,
 \qquad Q_n'=M_n\oplus T.
 \tag{A.65.4}
\]
The derived subgroup assertion follows either from the explicit
commutator maps or from odd weight separation: the first summand
has trivial odd weight, and the second has the three nontrivial
Frobenius families. Both summands are attained by actual commutators.
The group \(Q_n\) is a finite two-group of class two.

Consider \(\Pi/[W'_n,\Pi]\), with its possibly infinite central
kernel over \(E_n\). It is enough to prove that every finite
quotient retaining the complete finite residual \(Q_n\) has zero
scalar kernel. Such quotients are cofinal: choose an open normal
subgroup whose intersection with the residual lies in the kernel
of its map to \(Q_n\), and intersect also with the kernel of the
finite action head. Compactness supplies this open neighborhood.

In one such quotient let \(\widetilde Q\) be the actual image of R
and let Z be the image of \(W'_n\). Then
\[
 1\longrightarrow Z\longrightarrow\widetilde Q
 \longrightarrow Q_n\longrightarrow1
 \tag{A.65.5}
\]
has globally central kernel. The abelianization argument of A.53
applies over the original C, not over the new quotient \(D_1\).
Indeed \(H^1(C,A_n^\vee)=0\) by the actual normal pro-odd factor,
while \(R=[R,\Pi]\). If \(Z_0\) is the image of Z in
\(\widetilde Q^{\rm ab}\), dualizing the exact sequence with
quotient \(A_n\) shows that every invariant character on \(Z_0\)
lifts to an invariant character on \(\widetilde Q^{\rm ab}\).
The latter has no such character by relative perfectness. Since
\(Z_0\) is a trivial C-module, this forces \(Z_0=0\). Thus
\[
 \widetilde Q^{\rm ab}=A_n,
 \qquad Z\subset\widetilde Q'.
 \tag{A.65.6}
\]
No vanishing of \(H^1(D_1,T^\vee)\) is used or asserted.

\emph{The actual third commutators.}\par
Because (A.65.5) is central and \(Q_n\) has class two,
\(\gamma_4(\widetilde Q)=1\). The third commutator gives an
H-equivariant bilinear pairing
\[
 (M_n\oplus T)\otimes_{\mathbb Z} A_n
       \longrightarrow\gamma_3(\widetilde Q)\subset Z.
 \tag{A.65.7}
\]
It is well defined: changing a lift from the derived subgroup by
Z does not change its commutator, and the second argument factors
through the abelianization because \(\gamma_4=1\).
Let J denote the original normal odd subgroup of H. Its coinvariants
on \(A_n\) vanish. Since J acts trivially on \(M_n\) and on Z,
the \(M_n\)-part of (A.65.7) is zero.

The T-part factors through \(A_n/2A_n=V^8\), since T has exponent
two. Its image therefore has exponent two. We may extend scalars
from \(\mathbb F_2\) to k in this pairing and in the third graded
commutator. Under this extension, the underlying binary module V
splits into its three Frobenius forms. Their odd weight sets are
\[
 \mathcal S,\quad2\mathcal S,\quad4\mathcal S;
 \tag{A.65.8}
\]
they are pairwise disjoint and each is closed under negation, as
verified in the actual Q6 representation in A.59. The cup bracket
has no terms between different Frobenius forms. In one form, two
vectors of the same odd weight \(\lambda\) can bracket onto every
one of the six T-coordinate lines of weight \(2\lambda\).
This follows directly from the single-row multiplication
\(a_i*b_j\) in (A.63.2).

Since the target of (A.65.7) is J-trivial, the only possible
nonzero pairing with such a T-line is against an A-vector of
weight \(-2\lambda\). This vector belongs to the next Frobenius
form, distinct from the form containing \(\lambda\). Choose a,b
of weight \(\lambda\) whose bracket attains the selected T-line,
and c of weight \(-2\lambda\). The T-components of both
\([b,c]\) and \([c,a]\) vanish because their inputs belong to
different forms. Any old \(M_n\)-components have zero further
commutator by the preceding paragraph. The integral Jacobi identity
in the third lower-central quotient therefore gives
\[
 [[a,b],c]
   =-[[b,c],a]-[[c,a],b]=0.
 \tag{A.65.9}
\]
This argument is valid on \(A_n/2A_n\): changing a representative
by twice an element changes the relevant derived value only by
an old \(M_n\)-value, whose further commutator is already zero.
All possible invariant weights have been covered. Faithful scalar
extension now proves that the T-part also vanishes over
\(\mathbb F_2\). Hence \(\gamma_3(\widetilde Q)=1\).

\emph{The exact old scalar center leaves no room for Z.}\par
We now know that \(\widetilde Q\) has class two and abelianization
\(A_n\). Its derived subgroup is therefore an actual finite
quotient of
\(\bigwedge^2_{\mathbb Z}A_n\), with action factoring through H.
In particular it is a finite two-group, so Z is also a two-group.
Odd averaging yields
\[
 E_0=(\widetilde Q')^J\longrightarrow M_n,
 \qquad\ker(E_0\longrightarrow M_n)=Z.
 \tag{A.65.10}
\]
This map is onto because \(T^J=0\). Both its subquotients are
trivial H-modules. The automorphisms fixing both subquotients embed
in \(\operatorname{Hom}(M_n,Z)\), so the H-action on \(E_0\)
has two-group image. It therefore factors through \(H/O^2(H)\).
The odd averaging projection is H-equivariant; composing it with
the exterior-square surjection shows that \(E_0\) is a quotient of
\[
 (\bigwedge\nolimits^2_{\mathbb Z}A_n)_{O^2(H)}=M_n.
 \tag{A.65.11}
\]
It also surjects onto this same finite group by (A.65.10).
Equality of finite orders forces Z=0. This proves zero scalar
kernel in every member of the cofinal finite system, and hence
\(W'_n=[W'_n,\Pi]\).

\emph{The limit, retained families, and the full lift.}\par
The source maps and transition maps are the actual compatible maps
of A.62--A.63, so \(W'_{n+1}\subset W'_n\). For any finite quotient
q of \(\Pi\), the descending sequence \(q(W'_n)\) stabilizes, and
compactness gives
\(q(W'_\infty)=\bigcap_nq(W'_n)\). Its stable image equals its
commutator with \(q(\Pi)\) by the already proved finite-level
equalities. Separation by finite quotients proves
\([W'_\infty,\Pi]=W'_\infty\).

For a finite collection of older retained families, the same proof
uses their relatively perfect source kernel as base. Put
\(P_{\rm old}=\ker(C_{\rm old}\to C)\). The abelianization of
\(P_{\rm old}\) has only the types 6, 21, and 42. None can
map to a nonzero quotient of \(A_n^\vee\), whose simple factors
all have type 108. Thus
\(\operatorname{Hom}_C(P_{\rm old}^{\rm ab},A_n^\vee)=0\).
The five-term sequence, with original-base
\(H^1(C,A_n^\vee)=H^2(C,A_n^\vee)=0\), gives the same required
\(H^1(C_{\rm old},A_n^\vee)=0\). The actual new residual is the
full product by the Goursat argument of A.64; consequently the new
residual over \(C_{\rm old}\) is exactly \(Q_n\). Every other step of
the proof uses the unchanged finite H-action. Compactness retains
all older families simultaneously.

Thus the actual extension carrying the action and frames of A.64
already equals its own scalar refinement in A.3. Apply that
closed-normal-kernel criterion, or its equivalent finite-common-
refinement compactness formulation, to realize the specified
peripheral automorphism and its projected frames by a full positive
automorphism. The quotient negative square and the two quotient
normalizer norms are exactly those proved in A.64. This realizes
the prescribed quotient data; it does not turn the full lift into
an involution. In fact A.79 shows that every full positive lift
retaining this stored A.64 finite action has a negative image of
order four in a further actual finite quotient, independently of
its frames. Later reselected actions and cofinality remain separate
questions.
\end{proof}


% Proposed after checkpoint09; actual quotient frames, not a full scalar lift.
\statement{A.66}{Proved}{A class-three twenty-direction quotient admits compatible involutory peripheral frames}
After another permitted choice of the discretionary affine data, the
integral cross-cup tower of A.65 has an actual further source quotient
\[
 1\longrightarrow (V^{(3)})^{20}\longrightarrow E^{(3)}
 \longrightarrow E\longrightarrow1.
 \tag{A.66.1}
\]
Here \(V^{(3)}\) has the cubed monomial scalars. The twenty coordinate
classes are k-linearly independent and nonsplit. The actual quotient
admits a positive peripheral automorphism, an involutory negative
companion, and involutory s- and r-normalizer frames retaining the
construction-selected initial C-data and older families. The new
positive action is the identity on its entire free-cusp group factor.
The scalar refinement above (A.66.1), and realization of these exact
new frames by a full automorphism, are not asserted here.

\begin{proof}
\emph{An actual class-three group and source map.}\par
Put \(A=V^4\), \(T=(V^{(2)})^6\),
\(U=(V^{(3)})^4\), and \(Z=(V^{(3)})^{16}\).
The indices of T are \(i<j\), those of U are \(i<j<k\), and
those of Z are ordered pairs \((r,i)\), with all indices from one
to four. On \((a,t,u,z)\) define multiplication by
\[
 \begin{aligned}
 a''&=a+b,\\
 t''_{ij}&=t_{ij}+t'_{ij}+a_i*b_j,\\
 u''_{ijk}&=u_{ijk}+u'_{ijk}+a_i*t'_{jk}+t_{ij}*b_k,\\
 z''_{ri}&=z_{ri}+z'_{ri}+a_r^{[2]}*b_i.
 \end{aligned}
 \tag{A.66.2}
\]
The product \(*\) and the power in brackets are coordinatewise.
The two triple-cup coboundaries in the u formula cancel. The other
corrections are biadditive over \(\mathbb F_2\), so (A.66.2) is
associative. H acts strictly with scalar degrees one, two, three,
and three on the four displayed blocks. Denote this finite two-group
by K. Its third lower-central subgroup is U. Its derived and
Frattini subgroups are \(T\oplus U\oplus Z\), and its Frattini
quotient is A. To check attainment, single-row brackets give T,
the triples \([a_i,[a_j,a_k]]\) give each U-coordinate, and the
ordinary brackets give
\[
 a_r^{[2]}*b_i+b_r^{[2]}*a_i.
 \tag{A.66.3}
\]
The functions \(a^2b\) and \(ab^2\) are independent as
\(\mathbb F_2\)-bilinear functions on \(k^2\), and
\(a^2b+ab^2\ne0\). Thus all sixteen ordered Z-coordinates
are attained independently. Their cross-Frobenius support
distinguishes them from the old same-family T brackets.

Use the four normalized cusp cocycles, whose s-value is zero.
In \(K\rtimes C\), take S pure and R conjugate to pure r by
an element with abelian coordinate \(r e(r)\) and all higher
coordinates zero. This gives \(S^2=R^3=1\), and it projects
literally onto the source map of A.63. Initial-base
\(H^1(C,V)=0\) makes the actual R-image onto A. The finite
two-group Frattini theorem makes it all K.

Pairing with the full old integral residual over A gives the
actual fiber product: its Frattini quotient is the complete
\(V^8\), onto which the actual R-image already surjects.
This proves (A.66.1), with no presumed finite-quotient realization.
The Frobenius weight identity \(3\mathcal S=-4\mathcal S\)
identifies the underlying binary simple type with type 108;
we nevertheless keep the literal cubic coefficient action.

For independence, push out by a nonzero k-linear functional on
the twenty coordinates. If its U-part is nonzero, the resulting
residual still has nonzero third commutators. A section from the
old class-two residual would be onto modulo Frattini and hence
onto, an impossibility. If its U-part is zero, use the ordinary
brackets (A.66.3). An H-map from old T to \(V^{(3)}\) sends a
vector supported at one physical weight to the opposite weight,
with its entries squared; this follows from the three disjoint
Frobenius weight families. An H-map from the old trivial center
is zero. On two inputs at a single physical weight, the proposed
ordinary bracket stays at that weight, so it cannot factor through
the old derived group. Independence of \(a^2b\) and \(ab^2\)
then forces every ordinary coefficient to vanish. A splitting induces an
H-equivariant map on the old residual derived group: the new
kernel centralizes the residual, so changing the chosen lifts
of C cannot alter that derived action. This proves the
twenty independent nonsplit directions.

\emph{Exact action lifts.}\par
For a six-by-four coefficient matrix L, the old shear
\(t_{ij}\mapsto t_{ij}+\sum_rL_{ij,r}a_r^{[2]}\) lifts to the
following H-equivariant automorphism of K: fix a and z and set
\[
 u_{ijk}\longmapsto u_{ijk}
   +\sum_r\left\{L_{ij,r}z_{rk}
        +L_{jk,r}(z_{ri}+a_i*a_r^{[2]})\right\}.
 \tag{A.66.4}
\]
Indeed, the coboundary of \(z_{rk}\) is \(a_r^{[2]}*b_k\),
whereas that of \(z_{ri}+a_i*a_r^{[2]}\) is
\(a_i*b_r^{[2]}\). These are precisely the change in the
u product after the t shear. The resulting maps \(F_L\)
satisfy \(F_LF_{L'}=F_{L+L'}\).

The lower canonical-c correction has only four cusp cohomology
coefficients, together with an inner term. Formula (A.66.4),
the actual monomial c-action, and these inner terms therefore
give an honest initial lift. Correct it by the complete source
one-cocycle space in coefficient \(V^{(3)}\). The intertwiner
from V is
\[
 j_3(v)_\lambda=(v_{-\lambda})^4.
 \tag{A.66.5}
\]
It gives the same 44 independent source-value columns as before.
Matching \(c(s)=s\) and \(c(r)=sr^{-1}s\) reconstructs the
literal canonical source action on the actual quotient. Its
square is one because these are the source generators.

\emph{A compatible old affine choice.}\par
Impose \(W_{\rm cusp}=0\), retaining the old central class
equations but allowing the free s-frame t to vary. The exact
system has rank 102 in 360 binary coordinates. On its
258-dimensional homogeneous space, the old central norm has
polar radical 231, radical image zero, and a 27-bit quadratic
core. The exact solution numbered 2128959 was found and checked
in the original group law. Its initial local frame was then
adjusted within that radical. The two opposite-weight local gates
have 29 possibilities each; 29 of their 841 paired possibilities
are compatible with this same norm core. The radical has local
image rank 24. The saved choice has local indices \([6,0]\).

These are finite existence certificates, not dimension heuristics.
The original-law verified tuple is
\path{work_A/A66_W0_local_compatible_partial.json}.
The relevant producers and exact coefficient records are
\path{work_A/check_A66_zero_W_free_t_norm.py},
\path{work_A/A66_zero_W_free_t_quadratic_core.json},
\path{work_A/check_A66_W0_radical_local_gate.py}, and
\path{work_A/check_A66_W0_local_adjustment.py}.
The lower-six and ordinary-sixteen direct group checks are
\path{work_A/check_A66_W0_local_lower.out} and
\path{work_A/check_A66_W0_local_ordinary.out}.
They verify their negative squares and complete s-frame norms.
Their positive correction coefficients are identically zero.

Consequently the positive map on K is literally the identity,
paired with the retained old action over C. This is an actual
automorphism because the old action fixes the finite H-head and
the four cusp cocycles. Its negative companion on the K-factor
is the canonical c just constructed. Its square on the whole
fiber product is therefore one. This does not identify a full
source lift of this selected new automorphism.

\emph{The remaining four frame norms.}\par
Combine the saved lower T-frame and ordinary Z-frame, initially
putting the four U-frame coordinates equal to zero. It already
has the correct s-conjugacy, and its norm lies in U. That norm
is fixed by both S and the coefficient involution D.

On the new kernel, canonical c has upper-triangular multiplicity
map \(I+J_{L_c}\), followed by its cubic monomial action.
The map \(J_{L_c}:k^{16}\to k^4\) has rank four. The exact
certificate gives a triangular four-by-four minor with diagonal
\((7,2,2,2)\) in the fixed F8 encoding. A homogeneous ordinary
correction Z fixed by S and D has norm
\[
 (J_{L_c}Z,0)\in U\oplus Z.
 \tag{A.66.6}
\]
A constant right inverse of \(J_{L_c}\) commutes with the common
monomial actions. It therefore lifts every joint-invariant U
residue to a joint-invariant ordinary correction. This kills all
four remaining norms while preserving the previous lower and
ordinary equations. It changes the ordinary frame only within
its homogeneous norm-zero family.

The complete direct witness is
\path{work_A/A66_W0_twenty_sframe_solution.json}; its producer
and stdout are
\path{work_A/check_A66_W0_twenty_sframe.py} and
\path{work_A/check_A66_W0_twenty_sframe.out}.
The independent coefficient check is
\path{work_CE/code/check_A66_UT4_frame_surjection.py}.
These verify source relations, canonical c, literal cusp,
negative square, s-conjugacy, and the exact s-normalizer norm.

\emph{The r-frame and the retained initial choice.}\par
We make explicit an additional closed condition available in
the A.33 construction. Its dyadic s-frame comes from the original
full Galois involution. Its dyadic r-reflection is still purely
conjugate to that involution's r-reflection: A.16 retains the
chosen reflection, and the final coprime correction conjugates it.
For the common full Galois involution, the two reflections
\(s d_s\) and \(d_r\) are purely conjugate. For example, if
\(b=\nu c\nu^{-1}\), with
\(\nu(s)=u_s s u_s^{-1}\) and
\(\nu(r)=u_r r^{\epsilon}u_r^{-1}\), their conjugator is
\(u_r^{-1}u_s\). Composing gives the dyadic transporter used here.

At each finite congruence level retain its finite transporter
image. At every odd relative prime, two involutions with that
same image are complements to a pro-odd kernel, hence are
conjugate by that kernel. The transporter therefore glues across
the relative fiber product. Add this transporter as a compact
variable in A.33. For every finite collection use a common level;
the finite-intersection property is unchanged. The initial
C-frames can thus be selected with a pure transporter from
\(s_Cd_{s,C}\) to \(d_{r,C}\). This is a selection from the
actual construction, not a claim about arbitrary norm-zero frames.

First use the new actual quotient over C before adjoining the
older semisimple families; its kernel is pro-two. The solved
s-reflection commutes with s, so \(s d_s\) is an involution.
Lift the selected C-transporter. Its conjugate is an involution
over \(d_{r,C}\). Apply A.16 in the preimage of the actual
finite subgroup \(\langle r_C,d_{r,C}\rangle\simeq S_3\),
retaining this involution. Coprime conjugacy of its order-three
complement onto the actual \(\langle r\rangle\) supplies an
involutory r-normalizer over the exact original C-frame.
Pair this r-frame with the retained older r-frames over their
identical C-frame, using the already proved actual fiber product
and separation of residual simple types. No pro-two assertion
is made about the whole joint quotient with those older families.
This is an existence proof; no coordinates of that r-frame are
asserted in the displayed computational witness.

Finally, construct the integral actions, rather than infer them
from a particular full lift. Write \(P_n^{\rm free}\) for the
four-free-cusp class-two group. The actual source comparison and
the exact centers of A.58--A.62 give
\[
 P_n\simeq P_n^{\rm free}
       \mathbin{\times}_{P_1^{\rm free}}P_1.
 \tag{A.66.7}
\]
Only the six free off-diagonal center coordinates acquire higher
two-power exponents; all mixed and binary center terms are
already modulo two. Here \(P_n\) denotes the residual group.
The same comparison of source maps extends (A.66.7) to the
corresponding actual quotients over C:
\[
 D_n\simeq D_n^{\rm free}
       \mathbin{\times}_{D_1^{\rm free}}D_1.
\]
The source C4 carry is retained in the full \(D_1\)-factor.
In \(D_n^{\rm free}\), S is pure and the chosen
opposite-pair orientation is S-invariant. These observations
verify the actual source maps; alternatively the surjective
residual map on the complete abelianization and exact center
gives the same identification.

The chosen old action is literally identity on
\(P_1^{\rm free}\). Its free abelian action is identity and
its free inner vector is zero; the remaining possible central
character vanishes on s by the strict S-frame conjugacy and
on r by order three. Pair identity on \(P_n^{\rm free}\)
with the chosen action on \(P_1\). This gives compatible
actual actions with zero free inner vector at every n.
The invariant-lifting and orientation argument of A.62 lifts
the nonzero free s-frame through those levels; it does not
deduce an integral norm from a mod-two norm. Old scalar-half
corrections have no K-coordinate. Finite common refinements
and compactness retain all levels and older families; the
latter still have different residual abelian simple types.
All statements concern the actual quotient and its frames. The
next scalar lifting criterion remains to be checked.
\end{proof}


% Actual nonzero scalar lower bound; no selected-cover automorphism assertion.
\statement{A.67}{Proved}{The next actual scalar refinement has an attained eight-dimensional binary quotient}
Let W be the source kernel of the quotient in A.66. Its actual scalar
kernel has a continuous surjection
\[
 W/[W,\Pi]\longrightarrow (\mathbb F_2)^8.
 \tag{A.67.1}
\]
The eight coordinates are attained by explicit third commutators.
Thus immediate scalar stability fails at this next layer. This is a
lower bound on the actual scalar kernel, not its complete computation.
No invariance of this selected eight-dimensional quotient under the
new automorphisms, or obstruction to their norm lifting, is asserted.

\begin{proof}
\emph{A local scalar cover.}\par
On \(k\times k\), with \(k=\mathbb F_8\), define
\[
 (a,z)(b,w)=(a+b,z+w+a^2b).
 \tag{A.67.2}
\]
This is a group G of order 64 and class two. The change of coordinates
\((a,z)\mapsto(a^2,z)\) identifies it with the type-A Suzuki
two-group defined using \(\theta(t)=t^4\) in
\citet[Section 2, p.~3]{AEaton2024}. Only that explicit definition
is used here; the following cover is proved directly.

The k-valued function
\[
 f((a,z),(b,w))=z b^4+a^2w^4
 \tag{A.67.3}
\]
is a normalized two-cocycle. Indeed, in its coboundary the only
remaining terms are \(a^2bc^4\) and
\(a^2(b^2c)^4\), which are equal in k. Consequently
\[
 (a,z,u)(b,w,v)=
 (a+b,z+w+a^2b,u+v+z b^4+a^2w^4)
 \tag{A.67.4}
\]
defines a central extension of G by the additive group k.
The scalar action \((a,z)\mapsto(\xi a,\xi^3z)\),
\(\xi\in k^\times\), fixes (A.67.3), since its two
scalar degrees are seven and fourteen.

Put
\[
 q(a,z)=\operatorname{Tr}_{k/\mathbb F_2}(a^4z)+a^7,
 \tag{A.67.5}
\]
where \(a^7\) is regarded as zero or one. Direct expansion gives
\(\operatorname{Tr}(f)=\delta q\). Hence
\(\operatorname{Tr}(u)+q(a,z)\) is a homomorphism from
(A.67.4) to \(\mathbb F_2\).

The base commutators have z-coordinate \(a^2b+ab^2\).
Their span is nonzero and invariant under multiplication by
every \(\xi^3\), hence is all k. Third commutators in the central
coordinate have the form
\[
 z a^4+a^2z^4.
 \tag{A.67.6}
\]
For fixed \(a\ne0\), substitute \(z=a^3t\). The result is
\(t+t^4\), whose image is the trace-zero plane of k.
The preceding homomorphism shows that no additional central
derived direction occurs. Since \(\operatorname{Tr}(1)=1\),
the trace-zero plane maps isomorphically onto \(k/\mathbb F_2\).
Quotient (A.67.4) by the central line \(\mathbb F_2\cdot1\).
The resulting group \(\widetilde G\) has order 256, central
kernel \(k/\mathbb F_2\simeq C_2^2\), and
\[
 |\widetilde G'|=32,\qquad
 \gamma_3(\widetilde G)=C_2^2,\qquad
 \widetilde G/\Phi(\widetilde G)=k.
 \tag{A.67.7}
\]
In particular both central directions are stem directions,
attained by third commutators. The scalar-seven action lifts
strictly and fixes this central kernel.

The complete local identities are independently checked in
\path{code/check_suzuki_local_scalar_extension.py} and
\path{work_CE/code/verify_cubic_scalar_cocycles.py}.
The first certificate checks all \(64^3\) cocycle identities,
the trace coboundary, all scalar-seven actions, and the lower
central orders \(256,32,4,1\). Equivalently, for
\(d=q+q^2\) the scalar cocycles
\[
 \operatorname{Tr}\bigl(d a^4w+q a^6b+q^2a^4b^3\bigr)
 \tag{A.67.8}
\]
give the two independent directions as \(q\) ranges modulo
\(\mathbb F_2\). The case \(q=1\) is the coboundary of
\(\operatorname{Tr}(a^7)\).

\emph{The head action and the actual source.}\par
Take one copy of \(\widetilde G\) at each of the 36 physical
weights. H acts by the actual monomial scalars and permutations;
on each central \(C_2^2\) it acts trivially before permuting
the weights. Push out the 36 central factors by their sum map
onto one \(C_2^2\). This is H-equivariant. The resulting finite
two-group covers the 36-weight product of G, has Frattini
quotient V, and still has its two central directions attained
by third commutators supported at one weight.

For each of the four cusp rows, projection onto the diagonal
ordinary coordinate \(z_{ii}\) in A.66 is exactly (A.67.2).
Use the same source map: S is pure and R is conjugate to pure r
by a lift with abelian coordinate \(r e_i(r)\). It projects
literally onto the diagonal source already used in A.66.
Pair all four covers with the actual A.66 quotient. This gives
an abstract fiber product with central kernel
\((k/\mathbb F_2)^4\).

This fiber product is an actual source quotient. To see the
essential point, take triple commutators supported in one cusp
row and one weight. They attain that row's two new central
coordinates by (A.67.6). Their old integral and ordinary
components vanish because those groups have class two. Their
old UT4 component vanishes because its triples require three
distinct cusp rows. Their components in the other three new
row covers also vanish. Thus the whole new central kernel lies
in the third commutator subgroup of the fiber product.

It follows that this fiber product has the same abelianization
and the same complete \(V^8\) Frattini quotient as the prior
residual. The actual R-image surjects onto that quotient by the
original eight source cocycles. Finite two-group Frattini, or
its finitely generated pro-two version for the integral limit,
therefore gives the full fiber product as the actual R-image.
In particular the source kernel W maps onto all eight central
coordinates. Since those coordinates are globally central,
\([W,\Pi]\) maps to zero, proving (A.67.1).

For an older retained finite family, first work over its actual
base \(C_{\rm old}\). The same distinct residual abelian simple
types as in A.66 give the full new residual product and the
required eight-class Frattini surjectivity. Apply the preceding
argument to that source kernel; its chosen elements have trivial
older components. This avoids assuming that arbitrary triple
commutators vanish in the older factors. Compactness handles
all older families simultaneously.
The resulting central quotient is nonzero, but its selected
kernel may move under an automorphism of the previous level.
The universal actual scalar refinement is invariant; descent
to this particular eight-dimensional quotient requires an
additional argument. We make no such descent or norm claim here.
\end{proof}


% Actual scalar refinement and stability; action/frame equations are separate.
\statement{A.68}{Proved}{The full cubic scalar kernel has dimension 128 and its next refinement is stable}
At a finite integral level n of the actual quotient in A.66, let
\(W_n\) be its source kernel and put
\(W_{n,1}=[W_n,\Pi]\), with closed commutators understood.
Then
\[
 W_n/W_{n,1}\simeq C_2^{128},\qquad
 [W_{n,1},\Pi]=W_{n,1}.
 \tag{A.68.1}
\]
The scalar identifications are compatible with the integral transition
maps. They also hold after adjoining any earlier finite retained family.
If \(W_\infty=\bigcap_nW_n\), then
\[
 W_\infty/[W_\infty,\Pi]\simeq C_2^{128},\qquad
 [[W_\infty,\Pi],\Pi]=[W_\infty,\Pi].
 \tag{A.68.2}
\]
These statements identify the actual scalar refinement. The action,
negative square, and normalizer equations on that refinement require
separate verification.

\begin{proof}
\emph{Reduction to actual finite central extensions.}\par
Write \(Q_n\) for the finite residual over the initial C at the
chosen level. Its abelianization is
\(A_n=L_n^4\oplus V^4\), and its derived group has blocks
\[
 M_n,\qquad T=(V^{(2)})^6,\qquad
 U=(V^{(3)})^4,\qquad Z=(V^{(3)})^{16}.
 \tag{A.68.3}
\]
Here \(M_n\) is the old globally central group of A.58--A.62.
Only T has nonzero commutators with the abelianization; their
images are U. The group has class three. Let N be the actual
normal odd weight subgroup in the finite head H. We have
\((A_n)_N=0\).

Consider any finite image of the first scalar refinement that retains
\(Q_n\) as the quotient of its R-image. Such images are cofinal for
detecting the scalar kernel: refine any finite quotient by a finite
quotient of \(\Pi/W_n\) faithful on its finite normal subgroup
\(Q_n\). Denote the resulting R-image by \(\widetilde Q\),
and its new globally central kernel over \(Q_n\) by B.
The original-base vanishing \(H^1(C,A_n^\vee)=0\), together
with \(R=[R,\Pi]\), gives
\[
 \widetilde Q^{\rm ab}=A_n.
 \tag{A.68.4}
\]
Indeed, the dual of a new abelianized central kernel is globally
fixed, whereas the dual of \(\widetilde Q^{\rm ab}\) has no
fixed vectors. The invariant exact sequence injects that new dual
into \(H^1(C,A_n^\vee)\). This is the actual Pontryagin-dual
argument of A.53. An odd-primary part of B splits as a central
factor over the finite two-group \(Q_n\), and would survive
abelianization. Hence \(\widetilde Q\) is a finite two-group.

\emph{The fourth lower-central subgroup vanishes.}\par
The entire inverse image of \(M_n\) centralizes
\(\widetilde Q\). Its commutators lie in B and define an
equivariant pairing
\(M_n\otimes A_n\longrightarrow B\). The first module
is globally trivial, and \((A_n)_N=0\), so the pairing is zero.
The same observation applies to any old globally central scalar
block in a subsequent refinement.

Since \(Q_n\) has class three,
\(\gamma_4(\widetilde Q)\subset B\) and
\(\gamma_5(\widetilde Q)=1\). Its possible generators are
commutators of lifts of U with \(A_n\). These pairings factor
through binary modules. Extend scalars to \(k=\mathbb F_8\)
and decompose by the actual N-weights. A U-weight \(3\lambda\)
can have a scalar partner only at weight
\(-3\lambda=4\lambda\).

Represent this U-weight by
\([[a_i(\lambda),a_j(\lambda)],a_k(\lambda)]\).
Put \(t=[a_i,a_j]\), of weight \(2\lambda\), and let d
have weight \(4\lambda\). The bracket \([t,d]\) projects
to zero: the old triple bracket requires the same Frobenius
family. Its lift is scalar of nontrivial weight \(6\lambda\),
so it is zero. The class-four Hall--Witt identity reduces
\([
 [[a_i,a_j],a_k],d
]\) to \([t,[a_k,d]]\), up to the irrelevant sign convention.
The inner bracket projects to the ordinary Z-block, which is
central in \(Q_n\). Its lifts have central commutators with
\(\widetilde Q\), and therefore centralize the derived subgroup.
Thus \([t,[a_k,d]]=1\). This kills \(\gamma_4\).

Furthermore \(\gamma_3(\widetilde Q)\) is elementary abelian.
If g is derived, its square projects into \(M_n\), since the
other three blocks in (A.68.3) have exponent two. The full inverse
image of \(M_n\) is central, and hence
\[
 [g,x]^2=[g^2,x]=1.
 \tag{A.68.5}
\]
This proves the exponent assertion for the generators of
\(\gamma_3\).

\emph{No additional scalar survives in degree two.}\par
Pass to \(\widetilde Q/\gamma_3(\widetilde Q)\).
Its derived group D is a quotient of
\(\bigwedge^2_{\mathbb Z}A_n\), and surjects onto
\(M_n\oplus T\oplus Z\). Coprime averaging over N makes
\(D^N\) surject onto \(M_n\); its kernel is exactly the image
of B in this class-two quotient. Both subquotients are globally
trivial. Automorphisms fixing them form a subgroup of their Hom
group, so the head action on \(D^N\) has two-group image.
The action factors through the actual H-action on \(A_n\),
and thus through \(H/O^2(H)\). Consequently \(D^N\) is a
quotient of the already exact module
\[
 \left(\bigwedge^2_{\mathbb Z}A_n\right)_{O^2(H)}=M_n.
 \tag{A.68.6}
\]
It also surjects onto \(M_n\). Equality of finite orders forces
its new scalar kernel to vanish. Therefore
\(B\subset\gamma_3(\widetilde Q)\), so B is elementary
abelian and is generated by actual third commutators.

\emph{An upper bound from actual triple pairings.}\par
A first bracket in \(M_n\) gives zero. A first bracket in T
has weight \(2\lambda\) from two equal-family inputs at
\(\lambda\); its only scalar third partner has weight
\(-2\lambda\). This is the opposite physical weight, so both
alternative brackets in Jacobi project to zero. Their lifts are
central. The T contribution is therefore zero.

An ordinary Z bracket uses free-row inputs at weights
\(\lambda,2\lambda\) at the same physical weight. Its scalar
third partner has weight \(4\lambda\). If that partner is a
binary s-row, both alternative brackets vanish and Jacobi again
kills the result. The only remaining triples use three free rows
in the three distinct Frobenius families
\(\lambda,2\lambda,4\lambda\).

For one physical weight there are three cyclic bracket tensors
on \(k^4\otimes k^4\otimes k^4\), with one independent Jacobi
relation for each ordered row triple. Their dimension is at most
\[
 3\cdot4^3-4^3=128
 \quad\hbox{over }k.
 \tag{A.68.7}
\]
The subgroup \(O^2(H)/N\), of order 36, acts regularly on the
physical weights and identifies these tensor spaces. The scalar
product has weight seven, hence is one. The remaining order-two
stabilizer acts trivially: it fixes the chosen weight line with
scalar one and does not mix the cusp multiplicities. Descent from
k preserves dimension. Thus (A.68.7) is a bound of 128 over
\(\mathbb F_2\), not 384. This argument uses actual commutator
pairings and Hall--Witt in a finite class-three group, not a
presumed realization of a graded Lie algebra.

\emph{A cover attaining all 128 directions.}\par
On the ordinary class-two quotient write
\[
 (a,z)(b,w)=
 \left(a+b,z+w+(a_i^2b_j)_{i,j}\right),
 \qquad a,b\in k^4,\quad z,w\in k^{16}.
\]
For each ordered triple define
\[
 F_{ijk}((a,z),(b,w))=z_{ij}b_k^4+a_i^2w_{jk}^4.
 \tag{A.68.8}
\]
These are k-valued central cocycles: the two coboundary terms
cancel because \((b_j^2c_k)^4=b_jc_k^4\). All 64 together
have raw central group \(k^{64}\), of binary dimension 192.
Define
\[
 (\sigma X)_{ijk}=X_{jki}^4.
 \tag{A.68.9}
\]
Then \(\sigma^3=1\), and a raw triple commutator has value
\[
 (1+\sigma)P,\qquad
 P_{ijk}=(a_i^2b_j+b_i^2a_j)c_k^4.
 \tag{A.68.10}
\]
Ordinary commutators span all \(k^{16}\); multiplying by arbitrary
third inputs spans all \(k^{64}\). Hence the raw third subgroup
is exactly \(\operatorname{im}(1+\sigma)\).

The fixed space of \(\sigma\) has binary dimension 64. Its four
diagonal index triples each contribute one fixed binary coordinate;
the remaining 60 triples form 20 cyclic orbits, each contributing
one arbitrary k coordinate. Since the order of \(\sigma\) is odd,
its fixed space complements \(\operatorname{im}(1+\sigma)\).
Quotient by that fixed space, equivalently project by the idempotent
\(\sigma+\sigma^2\). This produces a stem central cover of
binary dimension 128. Independent coordinatewise quotients
\(k/\mathbb F_2\) would be incorrect for the mixed triples:
they lose 40 directions. The raw exact commutator certificate
checks dimension 128 and the stated fixed-space complement:
\path{work_CE/code/check_generalized_scalar_triples.py} and
\path{work_CE/generalized_scalar_triple_certificate.json}.

Take the 36 physical weight copies and push out their central
factors by summation. The degrees seven and fourteen in (A.68.8)
make this H-equivariant. Pair the result with the actual \(Q_n\).
The new scalar directions are attained by third commutators.
Coprime averaging on the elementary third subgroup isolates their
N-fixed component from the old U, which has no N-fixed vectors.
Thus the complete new central kernel lies in the Frattini subgroup.
The source still surjects onto the complete original eight-cocycle
Frattini quotient. Finite Frattini therefore proves that the
fiber product is an actual source quotient. Its attained scalar
kernel has dimension 128. The upper bound proves the first
assertion of (A.68.1).

\emph{The next scalar refinement is stable.}\par
In a further scalar refinement, the entire inverse image of the
old scalar \(C_2^{128}\) centralizes the residual, by the same
pairing with \((A_n)_N=0\). In particular the ordinary Z-lifts
still lie in the second upper center: their commutators project
into that old scalar block, whose full preimage is now central.
They still centralize the derived subgroup. The fourth-commutator
argument and the degree-two exhaustion are unchanged.

The N-fixed part of the new elementary third subgroup is an
extension of the globally trivial old \(C_2^{128}\) by the
new globally trivial scalar kernel. Its head action has two-group
image and factors through the actual action on the triple pairings.
Thus it factors through \(H/O^2(H)\). The regular
\(O^2(H)/N\)-orbit identification and the trivial remaining
stabilizer give the same 128-dimensional upper bound. This group
already surjects onto the old 128-dimensional scalar group.
The new kernel must be zero. Cofinality of the finite central
quotients gives the second assertion of (A.68.1).

\emph{Integral limits and older retained families.}\par
The explicit scalar cover uses only the modulo-two cusp data.
Its central generators are therefore the same at every integral
level. The natural maps between the scalar kernels are surjective
and identify these 128 generators, hence are isomorphisms.
Set \(U_\infty=\bigcap_nW_{n,1}\). The finite-image
stabilization and compactness argument of A.62 proves
\(U_\infty=[U_\infty,\Pi]\). Compactness of the fibers over
each of the common 128 scalar values gives
\(W_\infty/U_\infty=C_2^{128}\). Consequently
\[
 U_\infty=[U_\infty,\Pi]
 \subset[W_\infty,\Pi]\subset U_\infty,
\]
which proves (A.68.2).

For an earlier finite retained family use its actual base
\(C_{\rm old}\). The established type separation and continuous
five-term sequence give \(H^1(C_{\rm old},A_n^\vee)=0\),
and the actual new residual remains \(Q_n\). All subsequent
averaging uses the unchanged finite H-action; it assumes no
new odd subgroup inside \(C_{\rm old}\). The preceding proof
therefore applies at every finite common refinement, and
compactness retains those families simultaneously. Nothing in
this scalar calculation supplies a solution of the separate
normalizer equations on the new cover.
\end{proof}


% Actual full cubic scalar frames at every integral level and in the limit.
\statement{A.69}{Proved}{The full cubic scalar cover admits simultaneous peripheral frames and a full positive realization}
At every integral level of A.66, and compatibly in their inverse limit,
the actual 128-dimensional scalar refinement of A.68 admits the
prescribed positive peripheral action,
an involutory negative companion, and involutory s- and r-normalizer
frames. These retain the construction-selected initial C-data and
all earlier finite families. The selected quotient action and frames
are realized by a full positive automorphism of \(\Pi\).
The square of its full negative companion is not determined.
The constructed tower is not asserted to be cofinal among source
quotients, and no arbitrary prescribed-frame lifting property is claimed.

\begin{proof}
\emph{The actual old center must be retained.}\par
Use the four normalized cusp rows from A.66. The corresponding
free class-two residual has opposite-weight center
\(\operatorname{Sym}^2(k^4)\), of k-dimension ten. If E is
the actual four-by-eight change of source cocycle coordinates,
the old center projection is its induced symmetric map. In explicit
coordinates it sends the old center basis indexed by \(i\leq j\)
to \(\operatorname{sym}(Ee_i,Ee_j)\), where
\[
 \operatorname{sym}(a,b)_{ij}=
 \begin{cases}a_i b_i&i=j,\\
 a_i b_j+a_j b_i&i<j.
 \end{cases}
\]
Apply the already retained inner source gauge, then subtract the
central s-value times the parity character of the head. This
makes S pure. The old C4 carry has zero projection here. The
projection is a homomorphism by the displayed natural symmetric
map, the inner gauge, and this central character; its literal
R-value is the prescribed conjugate of pure r. The exact source
certificate records all these maps and values. This retains the
old center rather than assuming the scalar cover descends after
that center has been forgotten.

Write \((\ell,\bar\ell)\) for the chosen orientation of an
opposite-weight pair. The center cocycle is the sum of
\(\operatorname{sym}(a_\ell,b_{\bar\ell})\).
Changing the orientation changes the off-diagonal central
coordinate \(m_{ij}\) by
\(a_{i,\ell}a_{j,\bar\ell}+a_{j,\ell}a_{i,\bar\ell}\).
Define
\[
 Q_{ij}(a,m)=\sum_{(\ell,\bar\ell)}
                  a_{i,\ell}a_{j,\bar\ell}
              +\begin{cases}m_{ij}&i\ne j,\\0&i=j.\end{cases}
 \tag{A.69.1}
\]
These functions are independent of the chosen orientation and
are H-invariant. Their coboundaries are the symmetric pairings
between opposite weights that occur in the scalar shear below.

\emph{An honest canonical-c lift on the full scalar cover.}\par
Let \(\pi=\sigma+\sigma^2\) be the 192-to-128 projection
of A.68. The ordinary canonical-c shear is
\[
 z_{ri}\longmapsto z_{ri}+
           \sum_sL_{ri,s}\,j_3(a_s),
 \qquad j_3(a_s)_\ell=(a_{s,\bar\ell})^4.
\]
It lifts to the full scalar cover by adding to its central coordinate
\[
 \Psi_{rij}(a,m)=\pi\left(
   \sum_s\{L_{ri,s}Q_{js}(a,m)^4
              +L_{ij,s}^4Q_{sr}(a,m)^2\}\right)_{rij}.
 \tag{A.69.2}
\]
To check this identity, substitute the shear in
\(F_{rij}=z_{ri}b_j^4+a_r^2w_{ij}^4\).
Its change is precisely the coboundary of (A.69.2), by the
opposite-weight formula for \(\delta Q_{ij}\).
The functions (A.69.1) also show that this lift commutes with
H. Its square is one because it fixes a and m.

Compose this shear with the actual monomial c-action and the
retained inner corrections. The resulting map is a homomorphism
of the explicit scalar-cover factor. Its discrepancy on S is a
central element; subtract that element times the parity character.
This makes the literal s-image correct. The r-image is then the
correct lift as well: the two candidate r-images have order three
and differ by a central two-element. The source computation
checks the exact identities
\[
 c(S)=S,\qquad c(R)=SR^{-1}S,
 \qquad c(SR)=(SR)^{-1}.
 \tag{A.69.3}
\]
The actual source is onto, so these identities also prove that
the resulting canonical c has square one. The positive action
is identity on the entire free-cusp and scalar factor, paired
with the chosen old action over C. Its negative companion is
therefore this canonical c on that factor, and is involutory
on the actual full fiber product. This constructs the quotient
action before invoking any full lifting theorem.

\emph{The simultaneous scalar and lower frame equations.}\par
Start with the verified \([6,0]\) affine and twenty-factor
s-frame from A.66. Its scalar norm is a specified vector in
the 128-dimensional center. A correction in ordinary Z can
change this norm while introducing an error in the four U
frame coordinates. These two equations must be solved together.

Let \(D_2,D_3\) denote the squared- and cubed-scalar coefficient
involutions for the s-normalizer. Retain all of the following
binary directions: 576 ordinary Z directions fixed by both S
and \(D_3\); 216 T directions fixed by S and \(D_2\);
240 U directions fixed by S; and 87 old-center directions with
zero old norm. This is a group of commuting corrections.
Its norm equations are binary linear. If
\(a=\rho(\beta)(p+C_W)\) is the effective free affine frame
coordinate, their lower U changes are
\[
 \begin{aligned}
 Z&\longmapsto J_{L_c}Z,\\
 T&\longmapsto
       (a_iT_{jk}+T_{ij}a_k)_{i<j<k},\\
 U&\longmapsto(1+D_3)U.
 \end{aligned}
 \tag{A.69.4}
\]
The new scalar change of a pure ordinary correction is
\[
 \pi\left(
  \sum_\ell
    \{Z_{ri,\ell}a_{j,\ell}^4
       +a_{r,\ell}^2Z_{ij,\ell}^4\}
       \right)_{rij}.
 \tag{A.69.5}
\]
Pure T and U corrections have zero scalar change. The old-center
change is (A.69.2) evaluated on that center. These formulas follow
from the original group multiplication, and the exact certificate
checks their source-law interpretation.

The complete system has 1119 binary variables and rank 234.
The target is in its image. The selected solution in fact uses
only ordinary Z and T directions, supported at the two joint-fixed
physical weights 8 and 11; no old-center or U correction is needed.
There is also a short explanation for the lower repair. The
ordinary scalar solution is supported at those two weights. There
\(a=(2,2,2,6)\) or \((2,2,2,7)\) in the fixed field encoding.
The pointwise map from six T coordinates to four U coordinates
in (A.69.4) has rank four when all four entries of a are nonzero.
For example, the columns \(T_{13},T_{02},T_{01},T_{23}\),
using zero-based row indices, have a nonzero triangular minor.
Thus T cancels the ordinary correction's U error without changing
its scalar norm or any lower norm.

The complete system, basis labels, coefficient columns, solution,
and corrected twenty-factor frame are recorded in
\path{work_A/A68_simultaneous_scalar_frame_solution.json}.
Its producer and stdout are
\path{work_A/check_A68_simultaneous_scalar_frame.py} and
\path{work_A/check_A68_simultaneous_scalar_frame.out}.
The direct twenty-factor group law checks the final s-conjugacy
and every lower norm; it also checks all 72 k-basis T-column
formulas. Independently, the original scalar-cover multiplication
checks the reconstructed corrected frame and the product of the
old frame with the saved correction. Both centralize the actual S
and have literal negative norm one. That independent certificate
is \path{records/full_cubic_scalar_frame_verified.json}, with
stdout \path{records/check_full_cubic_scalar_frame_verified.out}.
The source, the old-center projection, and canonical c are checked
again in that computation. No search timeout or dimension estimate
is used as a substitute for the saved solution.

\emph{The actual integral frame compatibility gate.}\par
A norm-zero central correction modulo two need not lift to a
norm-zero correction over the integers. We therefore check this
compatibility before retaining the whole tower. Let
\(\mathcal O_\infty\) be the unramified cubic extension of
\(\mathbb Z_2\). The free factor has center
\(\mathcal O_n^6\oplus k^4\): six alternating off-diagonal
coordinates and four binary diagonal coordinates. The actual
canonical c acts on the free cocycles with its integral minus
sign, together with the retained inner vector \(C_W\).

The positive monomial involutions S and
\(D=\rho(\beta)C\) commute. Their stabilizer scalars are
trivial, so orbitwise Teichmuller transport gives a joint-fixed
\(t_\infty\) lifting the selected binary t. Use the actual
integral cocycle representatives of A.62, extending the modulo-four
representatives below, and put \(p_\infty=C_W+t_\infty\).
The invariant orientation of the opposite-weight pairs makes
this an exact s-conjugator. This is a construction of integral
invariants; it does not infer them merely from binary fixedness.

In the actual free \(\operatorname{GR}(4,3)\) group, lift the
old center by Teichmuller entries and evaluate the normalizer
square. Its six free coordinates are even and its binary diagonal
coordinates vanish. In the order
\((00,01,02,03,11,12,13,22,23,33)\), the necessary central
half modulo two is
\[
 (0,2,0,2,0,3,5,0,2,0).
 \tag{A.69.6}
\]
The original integral group law verifies that multiplying by this
central half gives literal s-conjugacy and norm one.

The norm-zero subspace of the complete old 36-coordinate center
maps onto the normalized free ten-coordinate center with k-rank
10. Thus (A.69.6) has an actual old norm-zero lift, with its
complete 36 coordinates saved. Multiplying the old full binary
frame by that lift preserves its old conjugacy and norm.
Now use the actual source fiber product from A.66,
\[
 D_n\simeq D_n^{\rm free}
             \mathbin{\times}_{D_1^{\rm free}}D_1.
 \tag{A.69.7}
\]
It pairs the corrected integral free frame with that corrected
full binary frame. The source C4 carry remains in the full
\(D_1\)-factor. Consequently no untested binary or mixed
36-center norm equation is concealed in this passage.

The required old-center half changes the new scalar norm by the
actual shear (A.69.2). This change, as well as the original defect,
lies in the ordinary Z-only image of binary rank 88. A second saved
Z correction at weights 8 and 11 cancels their sum, and a lower T
correction cancels its four U errors by (A.69.4). The complete
original twenty-factor and scalar-cover laws verify the resulting
common frame. Its free-center projection is compared literally
with the binary reduction of the corrected modulo-four frame.

The integral calculation and coefficients are
\path{work_CE/code/check_full_scalar_integral_frame_gate.py}
and \path{work_CE/full_scalar_integral_frame_gate.json}.
The complete corrected witness is
\path{work_A/A69_integral_scalar_frame_solution.json},
checked in the twenty-factor law by
\path{work_A/check_A69_integral_scalar_frame.py}.
The independent original-law comparison is also recorded in
\path{work_CE/code/verify_integral_cubic_joint_frame.py} and
\path{work_CE/integral_cubic_joint_frame_verified.json}.
We independently reran the complete original scalar-cover
law on this integral-compatible witness, including its source and
common old-center projection; the results are
\path{records/full_cubic_integral_scalar_frame_verified.json} and
\path{records/check_full_cubic_integral_scalar_frame_verified.out}.
The old-center projection rank and literal half lift have an
independent certificate
\path{work_A/check_A69_old_center_half_lift.py}.

This modulo-four calculation settles every higher integral level.
Keep the same \(t_\infty\) and initial integral free-center lift.
Its central square lies in \(2\mathcal O_\infty^6\), and the
negative action is identity on that center. Divide by two and
apply the corresponding central half over
\(\mathcal O_\infty\). Its binary value is exactly the one
already determined by (A.69.6). Pair this involutory free frame
with the fixed adjusted full binary and cubic frames. Their
common binary projections agree, so they give compatible frames
at every n. Equivalently, after the first modulo-four correction,
all later central halves reduce to zero modulo two and cannot
alter the retained twenty-factor or scalar equations. The binary
diagonal center has no further higher-power coordinates.

\emph{The r-frame and actual full realization.}\par
The construction-selected initial C-frames have the transporter
proved in A.66. In the new pro-two quotient over C, the solved
s-reflection gives an involution over the prescribed r-reflection
after lifting that transporter. Apply A.16 while retaining that
involution, then use coprime conjugacy of the order-three complement
onto the actual r. This gives an involutory r-normalizer over the
literal initial C-frame. Pair it with the older retained r-frames
through the established actual fiber product. This argument is
applied before adjoining older non-two residual factors; it makes
no assertion that their whole joint kernel is pro-two.

The actual scalar-refined source kernel is relatively perfect by
A.68, including over any older finite retained-family base.
The closed-kernel form of A.3 therefore realizes the selected
peripheral action and its projected frames by a full positive
automorphism. Finite common refinements and compactness retain
the earlier families and all integral levels simultaneously.
The new r-frames are compact variables in this argument; it
does not assume that every previously selected r-frame lifts. The involutory negative
square and both norms asserted here are properties of this
quotient image. They do not establish that the full negative
companion is an involution, or that the constructed tower is cofinal among all source
quotients.
\end{proof}


% Proposed after checkpoint10; authored in work_CE.
\statement{A.70}{Proved}{Finite actual central data do not suffice for unrestricted equivariant peripheral lifting}
Let
\[
 \mathcal E=F\rtimes\langle c\rangle
   \cong C_2\amalg C_2\amalg C_2,
 \qquad N=\ker\bigl(F\longrightarrow(\mathbb Z/12)^2\bigr).
\]
There is an involutory automorphism of the actual source-central
quotient \(\mathcal E/[N,\mathcal E]\) which fixes the external
reflection, preserves all three named reflection classes, and
preserves the three named pair-product cusp classes with one common
unit multiplier. It induces an automorphism of \(\mathcal E/N\),
but has no full reflection-class-preserving peripheral lift to
\(\mathcal E\), even if the full lift is allowed infinite order.
The multiplier is mixed: its residue is \(7\pmod {24}\).
This does not refute ordinary R2, or a multiplier-one equivariant
rule for one fixed action. It makes the equivariance failure of the
mixed-sign construction in A.18 explicit on a finite actual central
refinement, using the fixed-group argument of A.6 and the S4 test
of A.39.

\begin{proof}
Put \(w=[x,y]=xyx^{-1}y^{-1}\). Since \(F'\subseteq N\),
killing \([N,F]\) makes the orientation subgroup class two.
It also imposes \(w^{12}=1\). Conversely, in the class-two group
with this relation, \(x^{12},y^{12},w\) are central and generate
the image of N, so these are exactly the required relations.
The action of c inverts \(x^{12},y^{12}\) and fixes w. Hence
killing \([N,c]\) adds \(x^{24}=y^{24}=1\) and no other
relation. The exact quotient has normal forms
\[
 x^a y^b w^d c^e,
 \qquad a,b\in\mathbb Z/24,\quad d\in\mathbb Z/12,\quad
 e\in\mathbb Z/2,
 \tag{A.70.1}
\]
with multiplication
\[
 (a,b,d,e)(u,v,k,f)=
 (a+(-1)^e u,\ b+(-1)^e v,\ d+k-b(-1)^e u,\ e+f).
 \tag{A.70.2}
\]
This explicit group realizes the converse quotient, so there are
no unproved additional relations. In particular
\[
 |\mathcal E/[N,\mathcal E]|=2\cdot24^2\cdot12=13824,
 \qquad N/[N,\mathcal E]\cong C_2^2\times C_{12}.
 \tag{A.70.3}
\]
The central kernel is generated by \(x^{12},y^{12},w\).

Define
\[
 \theta(x)=x^7,\qquad\theta(y)=y^7,\qquad\theta(c)=c.
 \tag{A.70.4}
\]
In (A.70.1) this is \((a,b,d,e)\mapsto(7a,7b,d,e)\),
because \(49\equiv1\pmod {12}\). It is an automorphism of
order two, since \(49\equiv1\pmod {24}\), and fixes the
entire central kernel in (A.70.3). Thus it really lies over an
automorphism of the original finite quotient
\((\mathbb Z/12)^2\rtimes\langle c\rangle\).
The three named reflection classes may be represented by
\(c,xc,y^{-1}c\); their images under \(\theta\) are their conjugates
by \(1,x^3,y^{-3}\), respectively. If \(z=(xy)^{-1}\),
the class-two formula gives
\[
 (xy)^7=x^7y^7w^{-21},\qquad
 \theta(xy)=x^3(xy)^7x^{-3},\qquad
 \theta(z)=x^3z^7x^{-3}.
 \tag{A.70.5}
\]
The residue 7 is the image of the genuine full unit
\(\lambda\in\widehat{\mathbb Z}^{\times}\) with
\(\lambda_2=-1\) and \(\lambda_p=1\) for every odd p.
Consequently these are common-unit peripheral data, not a
substitution by a nonunit integer.

Suppose a full lift existed. Preserving the named reflection
classes preserves the orientation subgroup F, because its parity
character is one on each of those classes. The image of c is
conjugate to c, and its conjugator can be chosen in F: multiply
an odd conjugator by c if necessary. Conjugating the full lift
by this inner element makes it fix c, without changing its
peripheral multiplier \(\Lambda\). The orientation
abelianization of (A.70.1) forces
\(\Lambda\equiv7\pmod {24}\); this remains true after the
inner normalization.

Write its x-image as \(a x^\Lambda a^{-1}\), with \(a\in F\).
Commuting with c implies
\(a^{-1}c(a)\in C_F(x)=\overline{\langle x\rangle}\),
say \(a^{-1}c(a)=x^t\). If \(a^{\rm ab}=(r,s)\),
abelianization yields
\[
 (t,0)=(-2r,-2s),\qquad s=0,\quad t=-2r.
\]
Thus \(a x^{-r}\) is literally c-fixed. The established
full identity \(C_F(c)=1\) gives \(a=x^r\). The x-image
is therefore exactly \(x^\Lambda\); the same argument
makes the y-image exactly \(y^\Lambda\). This uses the full
fixed-group theorem and does not assert trivial fixed groups
in finite quotients.

Evaluate at \(X=(2\,3\,4)\), \(Y=(1\,2\,3\,4)\) in
S4, with rightmost permutations acting first. Since
\(\Lambda\equiv7\pmod {12}\), the element
\(X^\Lambda Y^\Lambda\) has order two, whereas
\((XY)^\Lambda\) has order four. They cannot be conjugate,
contradicting the third-cusp condition. No order assumption on
the proposed full lift was used.

The self-contained certificate
\path{work_CE/research_equivariant_R2/check_finite_central_equivariant_counterexample.py}
checks all 13824 normal forms and generator edges, the 48-element
central kernel, the reflection and cusp conjugators, and the S4
obstruction. Its stdout and JSON are stored alongside it. The
presentation argument above identifies the actual source-central
quotient; the finite count alone is not used to infer that identity.

The transported-head formulation is essential. There is an
\emph{original} epimorphism from the actual three-reflection
source to \(\mathcal E/N\), and \(\theta\) transports its named
quotient tuple. The desired map on the entire actual central
refinement already exists. The theorem disproves a rule that
would infer a full equivariant peripheral lift from these data.
It leaves the same-action multiplier-one restriction undecided.
\end{proof}

\paragraph{A multiplier-one variant between different full actions.}
This scope distinction can also be made on actual full involutions.
Choose an arithmetic \(\widehat{GT}\) element
\(\psi=(\lambda,f)\), \(f\in F'\), with the mixed unit
\(\lambda\) used above; the usual arithmetic embedding and cyclotomic
surjectivity supply such an element. In every class-two quotient
\(f\) is central. Therefore \(\psi\) acts on the \emph{entire} orientation
subgroup in (A.70.1) exactly as \(\theta\), not merely on its abelian
quotient or cusp classes. Put \(\alpha=\psi c\psi^{-1}\).
This is a genuine full negative involution. Power substitution
commutes with c in a class-two group, so \(\alpha\) and \(c\) have identical
actions on \(F/\gamma_3(F)\). Since
\(\gamma_3(F)\subseteq[N,F]\), they have identical actions
on \(N/[N,F]\), and hence
\[
 [N,F\rtimes\langle\alpha\rangle]
       =[N,F\rtimes\langle c\rangle].
 \tag{A.70.6}
\]
Their actual source-central quotients are literally identified.
The identity on this quotient is an equivariant multiplier-one
map: it fixes the orientation cusp elements and external
reflection, and the other named reflection classes match by
(A.70.5) and the displayed reflection conjugators.
But a full multiplier-one conjugator g with
\(gcg^{-1}=\alpha\) would make \(\psi^{-1}g\) centralize c
in the full outer peripheral group. By A.6 this centralizer is
\(\{1,c\}\), whose multipliers are \(1,-1\), whereas
\(\lambda^{-1}\notin\{1,-1\}\). Such g cannot exist.
This concerns \emph{two different source actions}; it is not a
counterexample to a multiplier-one lifting rule for the single
fixed action c, and \(\alpha\) is already conjugate to \(c\) by \(\psi\).

\statement{A.71}{Proved}{Three-trace rigidity fixes every residually irreducible rank-two matrix quotient}
Let A be a finite commutative ring with identity and let
\(\rho:F\to\operatorname{SL}_2(A)\) be a continuous
representation. Assume its reduction at every maximal ideal of
A is absolutely irreducible. For every full peripheral
automorphism \(\Phi\) with common multiplier \(+1\) or \(-1\),
there is \(G\in\operatorname{GL}_2(A)\) such that
\[
                   \rho\Phi=\operatorname{Ad}(G)\rho.
 \tag{A.71.1}
\]
In particular \(\ker\rho\) is stable as a conclusion, without
any prior stability or semilinear-action hypothesis. Every
epimorphism \(F\to\operatorname{PSL}_2(k)\), for a finite
field \(k\) of order at least four, has a stable kernel under
such \(\Phi\), and its induced action is projective linear.
The same conjugacy conclusion holds over a profinite commutative
ring when the stated residual hypothesis holds at every finite
quotient. This is an application of classical three-trace
rigidity, with the local-ring and inverse-limit details supplied
below; no novelty claim is made for the trace theorem.
\cite{CEMcCulloughWanderley2013}

\begin{proof}
Set
\[
 X=\rho(x),\quad Y=\rho(y),\quad
 X'=\rho\Phi(x),\quad Y'=\rho\Phi(y).
\]
Conjugation preserves trace, and inversion in \(\operatorname{SL}_2\)
also preserves trace over every commutative ring. The three
peripheral conditions therefore give
\[
 \operatorname{tr}X'=\operatorname{tr}X=a,\qquad
 \operatorname{tr}Y'=\operatorname{tr}Y=b,\qquad
 \operatorname{tr}(X'Y')=\operatorname{tr}(XY)=d.
 \tag{A.71.2}
\]
Because \(\Phi\) is onto, the primed pair generates the same image as
the unprimed pair, and has the same residual irreducibility.
Cayley--Hamilton gives the identities
\[
 X^2=aX-I,\quad Y^2=bY-I,\quad
 YX=-XY+bX+aY+(d-ab)I.
 \tag{A.71.3}
\]
These hold in characteristic two as well. They imply that every
word in X and Y lies in the A-span of \(I,X,Y,XY\), and that
all multiplication in this spanning set depends only on
\(a,b,d\). Modulo any maximal ideal, absolute irreducibility
makes the generated algebra the full matrix algebra. Here is an
elementary two-dimensional verification after extension to an
algebraic closure. If a nonscalar generator has distinct
eigenvalues, its two spectral idempotents belong to the algebra;
irreducibility forces both off-diagonal entries of the other
generator to be nonzero, giving all four matrix units. If instead
it has a repeated eigenvalue, subtracting that eigenvalue and
changing basis gives \(E_{12}\). The lower-left entry of the
other generator is nonzero by irreducibility, and multiplying
by \(E_{12}\) then gives \(E_{11}\), hence the remaining
matrix units. Both generators cannot be scalar. Dimension
descends to the residue field, so the four matrices above are
a basis there.
Their coordinate determinant is consequently outside every
maximal ideal of A, hence is a unit. It follows that
\(I,X,Y,XY\) is an A-basis of \(M_2(A)\). The same holds
for the primed matrices. Equations (A.71.2)--(A.71.3) now define
an A-algebra automorphism
\[
 M_2(A)\longrightarrow M_2(A),\qquad X\mapsto X',\quad Y\mapsto Y'.
\]
This passage is a comparison of complete multiplication tables,
not a claim that traces distinguish arbitrary reducible pairs.

For completeness this algebra automorphism is inner over A.
A finite commutative ring is a product of finite local rings,
and an A-algebra map fixes the central idempotents, so work in
one local factor. The images \(P_1,P_2\) of the standard
matrix units \(E_{11},E_{22}\) are complementary idempotents.
Their image modules in \(A^2\) are rank-one projective
modules, hence free over the local ring. Choose a generator
\(v_1\) of the first and put \(v_2=\varphi(E_{21})v_1\).
The matrix-unit identities show that \(v_2\) generates the
second and \(\varphi(E_{12})v_2=v_1\). The matrix with columns
\(v_1,v_2\) is invertible and implements phi on every matrix
unit. Combining these matrices over the local factors gives
G in (A.71.1). Equality of the representations follows from
equality on the free generators.

For a projectively onto representation over k, choose arbitrary
determinant-one lifts of the two free generator images. Their
projective image is \(\operatorname{PSL}_2(k)\), so their
matrix action has no common invariant line over an algebraic
closure. The finite-field result applies to this lifted
representation, and projectivizing proves both stability of the
original kernel and projective linearity. It is unnecessary to
assert that these particular lifts generate all of
\(\operatorname{SL}_2(k)\).

Finally let A be an inverse limit of its finite quotient rings
\(A_i\), and apply the finite result to \(\rho_i\). Let
\[
 \mathcal C_i=\{G_i\in\operatorname{GL}_2(A_i):
        \rho_i\Phi=\operatorname{Ad}(G_i)\rho_i\}.
\]
Each is a nonempty finite set. Reduction maps carry
\(\mathcal C_j\) to \(\mathcal C_i\); surjectivity of
these maps is not assumed. Any finite list of compatibility
conditions is satisfied by choosing a conjugator at one common
finite refinement and reducing it. Compactness of the product
therefore gives a compatible family. Its inverse family of
matrices is compatible as well, so the limit lies in
\(\operatorname{GL}_2(A)\) and gives the required conjugacy.
This argument supplies a single conjugator in the inverse limit,
not merely unrelated conjugators at each finite level.

Over finite fields the classical nonsingular trace-triple
conjugacy result is recorded as Macbeath's theorem in
McCullough--Wanderley, Theorem 8.2, pp.~490--491.
\cite{CEMcCulloughWanderley2013}
The elementary algebra proof above gives the precise
\(\operatorname{GL}_2\) conclusion and its finite-ring
extension; it does not incorrectly assert a unique
\(\operatorname{SL}_2(k)\)-conjugacy class in odd characteristic.
Nothing here applies to the canonical unmarked substitution
before its third-cusp condition is proved, or to arbitrary
higher-dimensional finite-group quotients.
\end{proof}


\statement{A.72}{Proved}{The actual mixed ordinary-cup extension and its negative frames}
Retain the construction-selected initial C-frames and the integral,
free-cusp cubic data of A.69. There is an actual quotient of the source
adjoining 32 copies of the cubic coefficient module \(V^{(3)}\), indexed
by the ordered pairs with one free cusp index and one binary s-index.
Its negative quotient action is involutory. Its s-frame is involutory,
and an r-frame can be selected over the literal initial C-frame.
These assertions retain the old action and s-frame, including the
integral tower and the older distinct-type families. They do not yet
assert realization of these new data by a full source automorphism:
the actual scalar-refinement action is a separate condition.

\begin{proof}
\emph{The coefficient actions and actual quotient.}\par
Use the normalized eight actual source cocycles from A.57 and A.66.
The indices \(0,1,2,3\) are free cusp classes; the indices
\(4,5,6,7\) are binary s-classes. Put
\[
 F_0=\{0,1,2,3\},\qquad B_0=\{4,5,6,7\},\qquad
 I=\{(i,j):\text{exactly one of }i,j\text{ belongs to }F_0\}.
 \tag{A.72.1}
\]
In this actual normalized basis, the multiplicity permutations of the
positive action, canonical c, and their negative product are respectively
\[
 P=(4\ 5)(6\ 7),\qquad C=(6\ 7),\qquad PC=(4\ 5).
 \tag{A.72.2}
\]
The physical monomial c-action remains present on each copy of V.
The matrices, extracted from the actual old source action, are recorded
in \path{work_A/check_A70_mixed_coefficient.py} and
\path{work_A/A70_mixed_coefficient_actions.json}; the provisional
filenames precede the present statement numbering. On the 32 mixed
multiplicities the ranks of \(P-1,C-1,PC-1\) are 16, 8, and 8.

For each physical weight, define the class-two group
\[
 (a,z)(b,w)=
 \bigl(a+b,\ (z_{ij}+w_{ij}+a_i^2b_j)_{(i,j)\in I}\bigr),
 \qquad a,b\in k^8,\quad k=\mathbf F_8.
 \tag{A.72.3}
\]
The head acts with degree one on a and degree three on z. Take its
actual fiber product with the preceding quotient over the eight-cocycle
module. The normalized source s has zero free coordinates, so its new
ordinary coordinates can be zero and its square is literally one.
Construct the source r by conjugating pure r with the actual affine
lift having abelian coordinate \(r e(r)\); its order is literally three.
This defines an actual homomorphism from the source.

We verify that its image is the entire stated fiber product. A free-row
input and a binary-row input at one physical weight have zero old
opposite-weight commutator and trivial free-only cubic commutator.
Their two new ordered coordinates are
\[
 (a^2b,b^2a),\qquad a,b\in k.
 \tag{A.72.4}
\]
These values span \(k^2\) over \(\mathbf F_2\). Indeed, a trace
functional annihilating all of them would give
\(u a^2+(v a)^4=0\) for every a; the independence of the two distinct
Frobenius powers implies \(u=v=0\). Varying the row pair and weight
therefore puts all 32 new modules in the derived subgroup, and hence
in the Frattini subgroup, of the abstract joint residual. Its Frattini
quotient is unchanged. The established eight-cocycle source
surjectivity and finite Frattini generation give the actual full
fiber product. This argument also applies at every finite integral
level; compactness supplies the same actual comparison for the tower.

\emph{Constructing the actual peripheral actions.}\par
The complete coefficient-cocycle space has dimension 44 over k:
eight actual cohomology classes and 36 coboundaries. In the cubic
coefficient model the natural comparison map is
\[
 (j_3v)(\lambda)=v(-\lambda)^4.
 \tag{A.72.5}
\]
Both actions use their actual normalized affine inner terms, including
the binary coordinates. Canonical c is corrected by the full
44-dimensional source s/r evaluation system. The positive action is
corrected by the 40-dimensional system consisting of literal cusp
fixation and the s-conjugacy class with the prescribed old frame.
The evaluation ranks and the resulting complete coefficients are
saved, and the original group law checks the source relations,
canonical c, and literal peripheral conditions.

Split I into the 16 slots involving binary indices 6 or 7 and the 16
slots involving indices 4 or 5. The first half has trivial negative
multiplicity action. Its calculated negative square is the identity,
although its positive affine inner term is nonzero. The full
s-centralizer norm equation is solvable, with physical coefficient
rank eight on the 20-dimensional s-centralizer. This calculation is
\path{work_A/check_mixed_binary67_cubic.py}, with complete values in
\path{work_A/mixed_binary67_cubic_result.json}.

In the second half the first peripheral lift has a nonzero inner
negative square. Write this square as \(\operatorname{Inn}(v)\),
where \(v\) is cusp-fixed. The negative coefficient involution swaps
the paired slots and applies physical \(C_3\). Thus
\[
 v_{j'}=C_3v_j.
\]
Choose \(a_j=v_j\), \(a_{j'}=0\), one member in each pair. Then
\[
 a+B(a)=v,\qquad \tau a=a.
 \tag{A.72.6}
\]
Replacing the positive action by its composition with
\(\operatorname{Inn}(a)\) preserves the literal cusp and old quotient,
and makes the negative square one. The same paired regular-module
argument removes the s-normalizer norm in the physical s-centralizer.
Its complete image dimension is \(8\cdot20=160\) over k. The exact
initial nonzero square, correction, and final original-law checks are
\path{work_A/check_mixed_binary45_cubic.py} and
\path{work_A/mixed_binary45_cubic_result.json}.
Thus the computation records a norm repair, rather than silently
assuming that the first lift is involutory.

The two constructions agree on the actual old quotient and have
independent new kernels by (A.72.4). They therefore define the claimed
32-module action and frame. Every added coordinate and correction is
binary, so it pulls back unchanged to the old integral tower. The old
central half changes from A.69 have no new effect: the mixed law and
its corrections depend on the abelianized coordinates, while those
half changes lie in the old center. The old integral action and
literal s-frame can consequently be paired with these new coordinates.

\emph{The r-frame and retained-data scope.}\par
Apply the construction-selected initial transporter and A.16 in the
new pro-two quotient over C, retaining the solved s-reflection. It
gives an involutory r-normalizer over the literal initial C-frame.
This is done before adjoining the older non-two residual factors;
their already retained r-frames then pair through the actual fiber
product. Compact common refinements retain the integral tower and
the older distinct-type families. The new r-frames are variables in
this compactness argument. The code explicitly marks the stored
arbitrary r-coordinate representative as untested; no failure is
inferred from it and no arbitrary r-frame lifting is assumed.

The full source realization criterion still requires the action and
peripheral data on the actual scalar refinement. Neither the group
law nor the solved quotient norms dispense with that condition.
\end{proof}


\statement{A.73}{Proved}{The exact 576-dimensional scalar refinement of the mixed extension}
Let \(W_n\) be the actual source kernel after the 32 mixed ordinary
modules of A.72 have been adjoined at integral level n. Then
\[
 W_n/\overline{[W_n,\Pi]}\simeq C_2^{576}.
 \tag{A.73.1}
\]
Its next scalar refinement is stable. Together with the 128 globally
central free-cusp directions retained in A.69, the complete scalar
third subgroup has dimension 704 over \(\mathbf F_2\).
The source transitions are compatible for all n, and the intersection
of the scalar-refined kernels is relatively perfect. These conclusions
also retain the earlier distinct-type finite families. They assert an
actual source-kernel calculation, not the existence of the required
negative action or frames on the new 576-dimensional center.

\begin{proof}
\emph{The actual source and the first upper bound.}\par
Work over the initial \(C=\Pi/R\), with \(R=\overline{[R,\Pi]}\).
The residual abelianization remains
\[
 A_n=\mathcal L_n^4\oplus V^4,
 \qquad H^1(C,A_n^\vee)=0.
 \tag{A.73.2}
\]
We do not reuse this cohomological assertion over a later base where
A.54 shows that it can fail. The residual derived group has the old
\(M_n\), the free T and U modules of A.66, all 48 ordinary slots with
at least one free index, and the 128 globally central scalar
directions of A.68.

Pass to a cofinal finite actual scalar refinement. The duality
argument of A.53 preserves abelianization \(A_n\). The refined
residual is a two-group: any odd central factor would survive in
its abelianization. The full preimages of \(M_n\) and of the old
scalar 128 centralize the residual, since their new commutators
are scalar pairings with \((A_n)_N=0\), where N is the actual
normal odd subgroup used for the weight decomposition.

The fourth lower-central subgroup vanishes by the same actual
Hall--Witt argument as A.68. A possibly noncentral old third input
has the form
\(U_{3\lambda}=[[a_i(\lambda),a_j(\lambda)],a_k(\lambda)]\).
Its scalar partner has weight \(4\lambda\). The intermediate
bracket \([T_{2\lambda},A_{4\lambda}]\) projects to zero and
has nontrivial weight \(6\lambda\), so it is zero in the scalar
kernel. Jacobi reduces the remaining term to \([T,Z]\). Ordinary
Z lifts lie in the second upper center, since their commutators
project to the old scalar 128 whose preimage centralizes. They
therefore commute with the derived subgroup, giving \([T,Z]=1\).
Thus \(\gamma_4=1\).

The third subgroup is elementary abelian. The square of a derived
element projects into \(M_n\) or the old scalar 128, whose preimages
centralize, so
\( [g,x]^2=[g^2,x]=1\).
The class-two quotient has derived group a quotient of
\(\bigwedge^2_{\mathbf Z}A_n\). Its odd-invariant part surjects
onto \(M_n\), with a new globally scalar kernel. The action on
this extension has two-group image and factors through
\(H/O^2(H)\). The exact old exterior-square coinvariant from A.58
gives the reverse order bound. Hence every new scalar lies in
\(\gamma_3\).

\emph{The mixed Jacobi count.}\par
Possible scalar pairings involving T still vanish by the
opposite-physical-weight Jacobi argument of A.68. The surviving
ordinary triples use weights \(\lambda,2\lambda,4\lambda\) at one
physical weight. An ordinary bracket exists precisely when at least
one of its row indices lies in \(F_0=\{0,1,2,3\}\); write
\(B_0=\{4,5,6,7\}\). For three ordered row indices, Jacobi gives
\[
\begin{array}{c|c|c|c}
 \text{free indices}&\text{ordered triples}&
 \text{available branches minus relations}&\text{bound}\\ \hline
 3&64&3-1&128\\
 2&192&3-1&384\\
 1&192&2-1&192\\
 0&64&0&0
\end{array}
 \tag{A.73.3}
\]
In the one-free case, the absent binary--binary branch cannot acquire
a scalar replacement: its weight is nontrivial. The bound is therefore
\(128+384+192=704\) over k after scalar extension.

The N-invariant third subgroup is an extension of the globally trivial
old 128 by the new scalar kernel. Its head action has two-group image,
so it factors through \(H/O^2(H)\). The regular
\(O^2(H)/N\)-orbit identifies the 36 physical weights. The remaining
order-two stabilizer acts as one on the weight line and does not mix
row indices. Frobenius descent gives the same bound of 704 over
\(\mathbf F_2\). Since 128 directions have already been retained,
the new scalar dimension is at most 576.

\emph{An explicit cover attaining the bound.}\par
Let z contain all 48 ordered slots with at least one free index.
Use the ordinary law
\[
 (a,z)(b,w)=(a+b,z+w+(a_i^2b_j)_{ij}).
\]
Whenever both adjacent slots exist, the formula
\[
 F_{ijk}((a,z),(b,w))=z_{ij}b_k^4+a_i^2w_{jk}^4
 \tag{A.73.4}
\]
is a k-valued central cocycle. Its cancellation is
\((b_j^2c_k)^4=b_jc_k^4\). There are 320 valid triples.
The 256 triples with at least two free indices are closed under
cyclic rotation. On their raw \(k^{256}\) center use the
cyclic-Frobenius idempotent \(\sigma+\sigma^2\) from A.68; its
binary rank is 512. The other 64 triples have binary first and last
indices and a free middle index. Keep their entire \(k^{64}\)
center, of binary rank 192. This nonclosed index set has no separate
cyclic projection or coordinatewise trace quotient. The resulting
center has dimension \(512+192=704\).

The third-commutator formula on binary basis inputs is
\[
 \begin{split}
 F_{ijk}([[a,b],c])={}&
 (a_i^2b_j+b_i^2a_j)c_k^4\\
 &+c_i^2(a_j^2b_k+b_j^2a_k)^4.
 \end{split}
 \tag{A.73.5}
\]
The exact matrix on all \(24^3\) binary-basis triples has rank 704.
Its all-free projection has rank 128 and its remaining projection
rank 576. The certificate also checks the idempotent and all 512
field cancellation identities. The producer, genuine output, and
complete matrices are
\path{work_A/check_mixed_scalar_704.py},
\path{work_A/check_mixed_scalar_704.out}, and
\path{work_A/mixed_scalar_704_certificate.json}.

Take 36 physical copies and push their scalar centers out by summation.
The weight-seven and weight-fourteen identities give a strict head
action. Pair this cover with the old mixed quotient through the
identical all-free 128 projection; do not adjoin a second independent
copy of that already retained center. The new kernel is consequently
576-dimensional.

The normalized source s has zero free coordinates and zero ordinary
coordinates. Its square is literally one in this cover, because every
ordinary slot contains a free index and (A.73.4) vanishes on the two
zero-z source inputs. Construct r by conjugating pure r with the
same actual affine lift, with abelian coordinate \(r e(r)\) and
zero higher coordinates. It has order three and projects to the
already retained source r, including its all-free scalar component.
The original C4 carry is retained in the old factor; the full source
s has not been replaced by a fictitious pure head element.

The matrix (A.73.5) attains all 704 scalar directions. Coprime averaging
on the elementary third subgroup separates these fixed directions
from the nontrivial old U weight modules. In particular the new 576
kernel lies in the derived, and hence Frattini, subgroup of the abstract
residual fiber product. Its Frattini quotient is unchanged. The actual
eight-cocycle source image therefore surjects onto this entire fiber
product by finite Frattini generation. This supplies the actual
source lower bound and proves (A.73.1).

\emph{Stability, integral transitions, and older families.}\par
In the next scalar refinement, the full old 704 scalar block
centralizes its preimage by \((A_n)_N=0\). Repeat the fourth-central,
exponent, and degree-two arguments. The N-invariant third subgroup is
an extension of the globally trivial old 704 by the new scalar; its
action has two-group image and factors through \(H/O^2(H)\).
The same physical-weight identification and Jacobi tensor bound give
at most 704 directions, while the map onto the old 704 is surjective.
The new scalar kernel is zero. Thus
\[
 \overline{[W_n,\Pi]}
   =\overline{[\,\overline{[W_n,\Pi]},\Pi\,]}.
 \tag{A.73.6}
\]

All 704 cocycles use modulo-two data. Their actual source transitions
therefore agree across n and are identity on the new scalar kernel.
The compact finite-image argument of A.62 and A.68 gives the exact
constant 576 kernel at the infinite integral level and relative
perfectness of the intersection after one scalar step.
For earlier finite retained families use their actual common base
\(C_{\rm old}\). Its residual abelianization over C has only the
previous 6-, 21-, and 42-dimensional simple types; hence the five-term
argument already used in A.65 gives
\(H^1(C_{\rm old},A_n^\vee)=0\). The actual residual product follows
by the established type separation. The same finite-head averaging
then proves the asserted joint result, without postulating a normal
odd subgroup inside \(C_{\rm old}\).

The negative multiplicity now permutes binary indices 4 and 5 on the
704 center. Canonical c and the positive action must be lifted through
the actual source power relations and the complete old center.
The scalar-kernel identification itself does not solve those action,
cusp, or normalizer equations.
\end{proof}


\statement{A.74}{Proved}{Actual central data do not guarantee prescribed equivariant conjugator frames}
There is an actual finite quotient of the three-reflection source for
which the identity action, with common peripheral multiplier one, has
compatible prescribed conjugator frames on the entire actual
orientation-central refinement, including involutory quotient norms,
but has no full equivariant lift retaining those frame images.
This assertion concerns \emph{equivariant conjugator lifting} with
prescribed frames. It does not obstruct a negative cocycle
\(h\,c(h)=1\), and it does not obstruct existence when the frames
may be changed.

\begin{proof}
Put \(F=\widehat F(x,y)\), let \(c(x)=x^{-1}\),
\(c(y)=y^{-1}\), and put
\(E=F\rtimes\langle c\rangle\).
The three named reflections are \(c,xc,y^{-1}c\), and their
pair products give the three peripheral classes represented by
\(x,y,z=(xy)^{-1}\). In \(Q=\operatorname{SL}_2(\mathbf F_5)\),
take
\[
 X=\begin{pmatrix}0&4\\1&4\end{pmatrix},\qquad
 Y=\begin{pmatrix}0&1\\4&4\end{pmatrix},\qquad
 J=\begin{pmatrix}0&1\\1&0\end{pmatrix}.
 \tag{A.74.1}
\]
Both \(X,Y\) have order three, and
\(J^2=I\), \(JXJ^{-1}=X^{-1}\),
\(JYJ^{-1}=Y^{-1}\). They generate \(Q\): indeed
\(X^{-1}Y=-\left(\begin{smallmatrix}1&2\\0&1\end{smallmatrix}\right)\)
and
\(XY^{-1}=-\left(\begin{smallmatrix}1&0\\2&1\end{smallmatrix}\right)\).
Their fifth powers give \(-I\), and the resulting upper and lower
elementary matrices generate \(\operatorname{SL}_2(\mathbf F_5)\).
The actual epimorphism \(q:F\to Q\) sending \(x,y\) to \(X,Y\)
therefore has a \(c\)-stable kernel \(N\). Its induced action is
literally \(c_Q=\operatorname{Ad}(J)\). In particular,
\(E\to Q\rtimes\langle c_Q\rangle\) is an actual original
three-reflection epimorphism.

We first identify the \emph{actual} orientation-central refinement:
\[
 F/\overline{[N,F]}
   \xrightarrow{\ (q,\operatorname{ab})\ }
 Q\times\widehat{\mathbf Z}^{2}
 \quad\hbox{is an isomorphism.}
 \tag{A.74.2}
\]
The classical inputs are that \(Q\) is perfect and that
\(Q\to A_5\) is the universal central cover. Perfectness follows,
for example, by conjugating upper and lower unipotents by
\(\operatorname{diag}(2,3)\): the resulting commutators contain
all unipotents. The primary ATLAS entry records
\(H_2(A_5,\mathbf Z)=C_2\) and identifies
\(2.A_5=\operatorname{SL}_2(5)\)
\cite{CEATLASA5}. Hence this perfect double cover is the universal
one. The usual centrally closed property is also stated in
Steinberg's lectures, Section~7(vii)--(viii)
\cite{CESteinberg1967Central}; the following argument verifies the
profinite-kernel application directly.

If \(D\to Q\) is any central extension with kernel \(M\), then
\(D=M D'\) and the abstract derived subgroup \(D'\) is perfect.
The composition \(D'\to Q\to A_5\) is central: an element over
the central kernel of \(Q\to A_5\) defines a central-valued
commutator homomorphism on the perfect group \(D'\), which must
vanish. Universality supplies a homomorphism \(Q\to D'\) over
\(A_5\). Its composite with \(D'\to Q\) is the identity by
the same uniqueness property. Thus it is a section; it is continuous
because \(Q\) is finite. Its image is the entire derived subgroup,
and \(D\simeq Q\times M\). The section is unique since any two
differ by a homomorphism from the perfect group \(Q\) to \(M\).
Apply this to \(D=F/\overline{[N,F]}\). Its abelianization is
exactly \(F^{\rm ab}=\widehat{\mathbf Z}^{2}\), proving
(A.74.2). Under this canonical identification,
\[
 x=(X,e_1),\qquad y=(Y,e_2),\qquad
 c=(\operatorname{Ad}(J),-1).
 \tag{A.74.3}
\]
Killing also \([N,c]\) reduces the central
\(\widehat{\mathbf Z}^{2}\) factor modulo two. Consequently
\[
 E/\overline{[N,E]}
 \simeq (Q\times C_2^2)\rtimes\langle c\rangle,
 \qquad |E/\overline{[N,E]}|=960.
 \tag{A.74.4}
\]

On the entire group in (A.74.2), and on its extension by \(c\),
take the identity automorphism. Put
\(\beta=(-I,0)\). This element is central, \(c\)-fixed, has
order two and has zero abelianization. Prescribe \(\beta\) as
the \(x\)-cusp conjugator and prescribe one as the \(y,z\)
conjugators. The reflection \(xc\) may likewise be assigned
conjugator \(\beta\), while \(c,y^{-1}c\) have conjugator one.
Every prescribed conjugacy and pair-product equation holds literally,
and
\[
 \beta\,c(\beta)=1.
 \tag{A.74.5}
\]
These assertions hold both on the entire actual orientation-central
refinement and on the actual finite quotient (A.74.4).

Suppose a full peripheral lift retained these frame images, even
only after projection to \(Q\). Its external-\(c\) conjugator
\(a_c\) belongs to \(N\), because the prescribed frame there
is one. Normalize by conjugation with \(a_c^{-1}\). This fixes
\(c\) literally and preserves every prescribed \(Q\)-frame.
It also leaves the identity action on (A.74.2) unchanged, since
the image of \(a_c\) there is central. The normalized restriction
\(\Phi\) to \(F\) commutes with \(c\).

Write \(\Phi(x)=a x a^{-1}\). Commutation with \(c\) and
\(C_F(x)=\overline{\langle x\rangle}\) imply
\(a^{-1}c(a)=x^t\). If
\(a^{\rm ab}=r e_1+s e_2\), then
\(t=-2r\) and \(s=0\). Thus \(a x^{-r}\) is \(c\)-fixed.
The full fixed-subgroup equality \(C_F(c)=1\), already used in
A.6 and A.70, gives \(a=x^r\). Hence \(\Phi(x)=x\), and the
same argument gives \(\Phi(y)=y\). Every \(x\)-conjugator of
this identity lies in \(\overline{\langle x\rangle}\), whose
\(Q\)-image has order three. It cannot have prescribed image
\(-I\), of order two. This proves the failure of prescribed-frame
equivariant lifting.

The exact finite matrix and order960 checks are in
\begin{quote}\small
\path{work_CE/research_equivariant_R2/check_binary_icosahedral_bad_frame.py}
\end{quote}
with its stdout and certificate. They verify the original epimorphism,
perfectness, the central element and the literal quotient frame norms;
the proof of (A.74.2) supplies the source-central comparison.

The identity with all frames equal to one remains a full lift.
Moreover, requiring a full map to intertwine the fixed action \(c\)
with itself is stronger than requiring a negative cocycle
\(h\,c(h)=1\). The latter permits \(hc\) to be an involution
without requiring \(h\) to commute with \(c\). Thus neither
unframed existence, ordinary non-equivariant R2, nor the main
negative-involution conjugacy problem is refuted here.
\end{proof}


\statement{A.75}{Proved}{Framed realization through the complete mixed scalar tower}
The actual 704-dimensional scalar cover of A.73 admits a selected
positive peripheral action whose negative companion is involutory
and has involutory s- and r-frames. These data can be realized by a
full positive automorphism while retaining all integral levels and
the older finite families. The old free-only action, the literal
initial C-frames, and the required integral coordinates are retained;
discretionary higher old frames may change. The involutory negative
square asserted here concerns this quotient tower, not the full
negative automorphism.

\begin{proof}
\emph{An honest action on the actual scalar cover.}\par
Use all eight normalized actual source cocycles. Retain the complete
36-coordinate opposite-weight center and its nonzero C4 source carry.
The change \(\operatorname{Sym}^2(E)\), the actual affine gauge,
and a source parity-center character identify the normalized model
with the original source quotient. In particular this is not the
four-cusp projection with its source carry discarded. The ordinary
coordinates are the 48 allowed pairs of A.73; the scalar coordinates
are its projected 704-dimensional center. The source s has zero free
coordinates and zero ordinary coordinates, and r is constructed by
actual conjugation of pure r. The source relations and complete old
projection are checked in the original group law.

The multiplicities \(P=(4\ 5)(6\ 7)\), \(C=(6\ 7)\), and
\(B=(4\ 5)\) preserve the allowed ordinary and scalar index sets,
and commute with the cyclic-Frobenius projection. Both actual affine
inner terms are retained. Every ordinary cohomology shear L lifts
using the full old quadratic cochains
\[
 Q_{ij}(a,m)=\sum_{(\ell,-\ell)}a_{i,\ell}a_{j,-\ell}
             +\begin{cases}m_{ij},&i\ne j,\\0,&i=j,\end{cases}
\]
where the sum uses the fixed orientation of opposite weights. Its
scalar correction on a valid triple is
\[
 \Psi_{ijk}=
 \operatorname{stem}\sum_l
       \bigl(L_{ij,l}Q_{kl}^{\,4}+L_{jk,l}^{\,4}Q_{li}^{\,2}\bigr).
 \tag{A.75.1}
\]
This is the eight-row form of the actual cochain identity in A.69.
It includes the old C4 compatibility. Canonical c is corrected by
its central parity character to send s to s and r to
\(sr^{-1}s\) literally. The positive candidate is corrected by
its unique central parity character to fix the literal cusp. Its
projection is the chosen mixed32 action of A.72, including the
regular-half inner norm repair. The negative square is one on the
two actual source generators. Equivalently, its last discrepancy
is a central character killed by \(\tau\), which generates
\(\Pi^{\rm ab}=C_6\).

These finite-head maps are paired with the prescribed \(h_C\)
and \(c_C\) over \(C\to H\). The former has identity H-action,
and the latter induces the specified monomial involution. The actual
source-image theorem of A.73 identifies this pullback with the desired
source quotient. We select this quotient action before applying A.3;
we do not suppose that every arbitrary full lift induces it.

\emph{The additional action freedom and exact simultaneous repair.}\par
The stored mixed32 frame has nonzero new scalar s-class and norm.
The first 2343-bit homogeneous frame system has rank 656 and is
inconsistent. Its exact annihilator is retained in
\path{work_A/mixed_scalar_frame_system.json}. This scoped failure
is overcome by a further legitimate action family.

Let A be a zero-abelian element of the mixed ordinary subgroup whose
coefficient is fixed by both \(\tau\) and the negative coefficient
involution. This family has dimension 96 over k, or 288 over
\(\mathbf F_2\). Write
\[
 k_A=[A,\tau],\qquad
 h_A=J_{k_A}\operatorname{Inn}(A)h,
 \qquad \beta_A=A\beta,
 \tag{A.75.2}
\]
where \(J_k(g)=k^{\epsilon(g)}g\) and
\(\epsilon(s)=\epsilon(\tau)=1\),
\(\epsilon(r)=\epsilon(\beta)=0\).
The commutator \(k_A\) is central and has order at most two.
Since \(b(A)A^{-1}\) is central and \(b(\tau)=\tau^{-1}\),
\[
 b(k_A)=[b(A),b(\tau)]=[A,\tau^{-1}]=k_A.
\]
Thus (A.75.2) fixes the literal cusp and preserves the negative
square: the central character has zero norm, while
\(A b(A)\) is central. The transported frame's s-class changes
by \(k_A\), and its scalar normalizer norm changes by
\(A b(A)\). The parity character contributes nothing on the
frame because \(\epsilon(\beta)=0\). These are binary-linear
maps on the stated abelian family. Its projection to the old
free-only quotient is the identity.

We solve one simultaneous binary system containing the following
complete stated correction families:
\[
\begin{array}{c|r}
 \text{ordinary joint-fixed frame directions}&1632\\
 \text{lower free T joint-fixed directions}&216\\
 \text{free U s-fixed directions}&240\\
 \text{old36-center norm-zero directions}&87\\
 \text{independent new scalar norm directions}&168\\
 \text{mixed cusp-inner action directions}&288
\end{array}
 \tag{A.75.3}
\]
Its outputs are the old U norm, the new scalar s-class, the new
scalar norm, and the six integral free offdiagonal center changes,
the last constrained to zero. There are 2631 binary variables;
the rank is 890 and an exact solution is saved. The ordinary
s-class map is the actual commutator
\[
 \operatorname{stem}\sum_{\lambda}
     \bigl(z_{ij,\lambda}a_{s,k,\lambda}^{,4}
            +a_{s,i,\lambda}^{,2}z_{jk,\lambda}^{,4}\bigr).
 \tag{A.75.4}
\]
All scalar norm columns are computed by the original group law.
The lower U columns are the previously checked pointwise formulas,
and the added action columns are checked directly against the
original law. No exhaustion of all old affine actions is claimed
or needed for this constructive solution.

Direct reconstruction verifies the positive literal cusp, negative
square on s and r, exact s-conjugacy, and literal negative s-frame
norm one. A separate original twenty-factor implementation checks
the complete old conjugacy and norm and compares their common
old36 and ordinary projections literally. The complete system and
witness are
\path{work_A/check_mixed_scalar_action_frame_system.py} and
\path{work_A/mixed_scalar_action_frame_system.json}.
The source/action model is
\path{work_A/check_mixed_scalar_action.py}.
The direct witness verifier
\path{work_A/verify_mixed_scalar_witness.py}
preserves the input witness and records its hash and checks in
\path{work_A/mixed_scalar_witness_verified.json}; its old projection
is saved in \path{work_A/mixed_scalar_free_projection_witness.json}.
The imported historical initialization still executes its earlier
checks; this verifier does not solve or modify the new witness.

\emph{Why the six integral constraints suffice.}\par
Use the already established actual source fiber product
\[
 D_n=D_n^{\rm free}\mathbin{\times}_{D_1^{\rm free}}D_1
 \tag{A.75.5}
\]
with the full source C4 carry retained in \(D_1\), and pair it
with the binary higher-cup and scalar factors. The mixed inner
and parity action changes in (A.75.2) have zero projection to the
free integral factor and to the old all-free 128 scalar action.
The old free-only action is therefore retained exactly.

Every frame correction in (A.75.3) has zero abelianized coordinate.
Its normalized old-center change is exactly the displayed
\(\operatorname{Sym}^2(E)\) image of its old36 coordinate;
there is no affine cross term. The imposed six offdiagonal zeros
preserve the forced integral halves of A.69 literally. The remaining
four free diagonal coordinates lie in the same binary \(C_2^4\)
summand at every integral level. They are globally central, fixed
by the negative free action, and have zero norm. Lift them
constantly. All T, U, ordinary, and new scalar corrections are
binary. Their verified frame and this integral free frame have
identical common reduction, so (A.75.5) gives literal conjugacy
and norm at every n. No additional modulo-four half or full36
source carry is being omitted. The transitions are compatible,
and compactness gives the infinite tower.

This argument deliberately allows changes to discretionary higher
old frames. It retains the old free-only action, the literal
initial C-frames, older-family data, and the required integral
coordinates. It asserts no lifting of an arbitrary preselected
higher frame.

\emph{The r-frame and full positive realization.}\par
Use the construction-selected initial C-transporter. After fixing
the solved s-reflection, apply A.16 in the new pro-two quotient
over C and conjugate the order-three complement to actual r.
This selects an involutory r-frame over the literal initial C-frame.
Carry this out before adjoining the older non-two factors, then
pair their retained frames through the actual distinct-type fiber
product. The new r-frames are compact variables; surjectivity onto
arbitrary previously chosen frames is not assumed.

By A.73 the actual scalar-refined source kernel is relatively
perfect, also over every older finite retained-family base.
The closed-kernel form of A.3 realizes the selected peripheral
action and its projected frames by a full positive automorphism.
Common finite refinements and compactness retain the complete
integral tower and the earlier families. The proved negative
square and norms belong to this quotient image. The full negative
square and cofinality of the constructed tower remain open.
\end{proof}


\statement{A.76}{Proved}{An actual finite central refinement does not suffice for global-negative framed order preservation}
There are a finite actual quotient \(Q=\Pi/N\), its actual
source-central refinement
\[
 Q^\dagger=\Pi/\overline{[N,\Pi]},
 \qquad \Pi=C_2\amalg C_3,
 \tag{A.76.1}
\]
and positive peripheral data on \(Q^\dagger\), fixing the literal
cusp, with the following properties. Their negative companion has
order two and has prescribed s- and r-normalizers of order two.
An ordinary full positive framed lift exists. Nevertheless no full
positive peripheral h fixing \(\tau\) literally, inducing the
specified finite action, and retaining the specified finite s-frame
can satisfy
\[
                         (hc)^2=1.
 \tag{A.76.2}
\]
Thus the failure of framed order-preserving lifting is present for
the globally negative multiplier \(-1\), after the actual finite
central-refinement test and both quotient normalizer norms have
passed. It does not exclude reselection of the finite frame or
refute the unframed existence problem.

\begin{proof}
\emph{The original actual data and finite central refinement.}\par
Use the original A.56 witness, before the subsequent affine changes
in A.64 and A.66. Write \(D=\Pi/U\) for its actual quotient over C.
By A.56,
\[
 U=\overline{[U,\Pi]},
\]
and D has a positive peripheral action h, a negative involution
\(b=hc\), and literal s- and r-frames with both negative norms one.
The ordinary full positive framed lift exists by A.3 and A.56.

Its finite old model \(D_0\) has normal forms \((a,m,q)\),
with \(a\in V^8\), \(m\in\operatorname{Sym}^2_k(k^8)\),
and \(q\in H\), using exactly the opposite-weight law and
C4 carry in A.56. The finite head and source are (A.56.2)--(A.56.4);
in the original cocycle coordinates the source elements are
\[
 s\longmapsto(d(s),0,H_s),\qquad
 r\longmapsto(0,0,H_r).
\]
The actual polynomial positive and canonical actions are those of
\path{work_A/check_Q6_central_cover.py}, with the original saved
corrections. In particular \(D_0\) is an actual, action-stable
projection of D, of order
\[
                         |D_0|=2^{972}\,72\,7^{13}.
 \tag{A.76.3}
\]
Take \(Q=D_0\) and \(N=\ker(\Pi\to D_0)\). Since \(U\subseteq N\),
\[
 U=\overline{[U,\Pi]}\subseteq\overline{[N,\Pi]}.
 \tag{A.76.4}
\]
Therefore the actual refinement (A.76.1) is a quotient of D.
It inherits the negative involution and both norm-one frames,
not merely corresponding data on a conveniently chosen cover.
The original full framed lift also realizes these descended data.

The refinement is finite. Indeed its central subgroup
\(Z=N/\overline{[N,\Pi]}\) has finite index \(m=|Q|\)
and is topologically finitely generated by the finite-index
Schreier argument. The continuous transfer to Z restricts on Z
to \(z\mapsto z^m\). Its image has exponent dividing six,
because the source generators have orders two and three.
Consequently Z has exponent dividing \(6m\). A finitely generated
abelian profinite group of bounded exponent is finite, proving
the assertion. This also gives a direct finite construction;
no unspecified limiting quotient is used in (A.76.1).

The old witness is
\path{work_A/Q6_actual_central_solution.json}.
Its SHA-256 begins \texttt{b042bf92e479019c}; the complete hash,
together with those of all obstruction inputs, is recorded in
\path{work_A/A76_original_witness_hashes.json}.
In particular this theorem concerns the original stored datum,
not one of the later repaired tuples.

\emph{Every candidate preserves the actual six-cup quotient.}\par
Put \(T_0=V^{(2)}\), interpreted as the underlying binary
H-module with its displayed squared monomial action. The finite
old residual \(P_0\) satisfies \(P_0^{\rm ab}=V^8\).
The normal odd subgroup of H has no invariants on \(T_0\), so
coprime Hochschild--Serre gives
\[
 H^1(H,T_0)=H^2(H,T_0)=0,
 \qquad H^1(D_0,T_0)
       \simeq\operatorname{Hom}_H(P_0^{\rm ab},T_0)
       \simeq k^8.
 \tag{A.76.5}
\]
Here the Hom is over the underlying additive binary modules;
the Frobenius comparison supplies the displayed k-dimension.
The source cocycles are the complete eight classes, so inflation
in (A.76.5) is an isomorphism onto \(H^1(\Pi,T_0)\).
The five-term sequence therefore makes the actual transgression
\[
 \operatorname{Hom}_{D_0}(N,T_0)
              \longrightarrow H^2(D_0,T_0)
 \tag{A.76.6}
\]
injective.

The six literal cusp cups \(e_i\smile e_j\), \(0\leq i<j\leq3\),
are the independent actual extension classes of A.63. Their finite
source quotient is
\[
 D_0\mathbin{\times}_{V^4\rtimes H}(N_6\rtimes H),
 \tag{A.76.7}
\]
where \(N_6\) has the coordinatewise cross-cup law and kernel
\(T_0^6\). Its actual source attainment is the same finite
Frattini argument as A.63. The retained old positive action fixes
the four cusp cohomology classes, hence all six cup classes.
Equivariance and the injectivity in (A.76.6) force preservation
of their six actual kernel maps and force identity on their
coefficient quotient. Canonical c similarly has its prescribed
squared monomial action. Thus every proposed full lift of the
retained data descends to (A.76.7); kernel stability and coefficient
action are conclusions, rather than extra restrictions on candidates.

\emph{Exhausting all source cochains and frame changes.}\par
Two descended positive actions with the same old \(D_0\)-action
and coefficient action differ by a \(T_0^6\)-valued cocycle on
\(D_0\). By (A.76.5), its source inflation ranges over exactly
the complete 44-dimensional cocycle space in each coefficient
copy: eight cohomology classes and 36 coboundaries. The literal
cusp condition kills its cusp restriction. Keeping the same old
s-frame makes its s-restriction a coboundary. The actual joint
restriction isomorphism (A.56.5), with the Frobenius comparison,
therefore forces the difference to be
\[
 \operatorname{Inn}(a),\qquad a\in(T_0^6)^\tau.
 \tag{A.76.8}
\]
The negative coefficient action is identity on this cusp-fixed
space. If \(b'=\operatorname{Inn}(a)b\) and the frame is
transported to \(\beta'=a\beta\), direct multiplication gives
\[
 \beta'^{-1}b'(\beta'^{-1})
                 =\beta^{-1}b(\beta^{-1}).
 \tag{A.76.9}
\]
Every remaining frame over the same old frame differs by an element
of \((T_0^6)^s\), and changes the square by its complete
\(1+D_s\) norm. Old globally central frame changes have no new
\(T_0\)-component in (A.76.7). Thus (A.76.8)--(A.76.9) exhaust
all admissible source cochains and all new frame variations for
this prescribed old datum.

\emph{The exact nonzero invariant.}\par
In each coefficient copy, \((T_0)^s\) has k-dimension 20;
the frame norm has rank eight and invariant target dimension 12.
At zero-based physical weight 23, the frame coefficient involution
fixes the weight literally with scalar one. Evaluation there
annihilates the entire norm image. In cup order
\((01,02,03,12,13,23)\), the original-law s-square has values
\[
                           (0,0,5,0,5,5)\in k^6.
 \tag{A.76.10}
\]
The value five in the 03 coordinate is nonzero. It therefore
cannot be removed by any of the variations just exhausted.

The independent direct certificate
\path{work_CE/code/verify_A63_s_norm_direct.py}
checks all 360 binary s-centralizer norm columns, all 216
coefficient-action columns, and all 72 literal-cusp inner
correction columns. The complete results are
\path{work_CE/A63_s_norm_independent_certificate.json}
and \path{work_CE/review_A63_fixed_affine_s_norm.txt}.
Their hashes are bound in the witness record above. This is an
exact norm-image annihilator, not a search timeout or a statement
about one unvaried representative.

\emph{Why a full negative involution would contradict it.}\par
Suppose a full positive peripheral h fixes \(\tau\) literally,
retains the specified finite action and s-frame, and has
\(b=hc\) involutory. Let \(\beta\) be its full s-conjugator
with that prescribed projection. Since
\(b(s)=\beta s\beta^{-1}\), involutivity gives
\[
              \beta^{-1}b(\beta^{-1})\in C_\Pi(s)=\langle s\rangle.
 \tag{A.76.11}
\]
The free-factor centralizer equality is the full profinite
free-product theorem.\cite{A_HR1985}
The retained S3 quotient has trivial b-action, and \(\beta\)
has trivial image there. Hence the element in (A.76.11) maps to
one in S3, while s has nontrivial image. It must therefore be
one in \(\Pi\). Its image in the actual quotient (A.76.7) would
be a norm-one s-frame, contradicting (A.76.10).

The multiplier of b in this argument is the global unit \(-1\);
the positive lift has multiplier one. The literal cusp and the
specified finite frame are part of the lifting problem. No claim
is made for unnormalized lifts with only an unrelated quotient
cusp framing, or for existential reselection of the old finite
affine datum. The later constructions A.64 and A.75 change that
datum and remain fully compatible with this counterexample.
\end{proof}


\statement{A.77}{Proved}{One finite relation module detects the prescribed-frame obstruction that every scalar iteration misses}
Let \(\Pi=C_2\amalg C_3=\langle s,r\rangle\), let
\(q:\Pi\twoheadrightarrow Q\) be a finite quotient injecting both
factors, and put \(L=\ker q\). Define
\[
 \Phi_p(L)=\overline{[L,L]L^p},\qquad
 G_p=\Pi/\Phi_p(L)\quad(p=2,3),\qquad
 G_6=\Pi/\overline{[L,L]L^6}.
 \tag{A.77.1}
\]
Then
\[
 \operatorname{im}\bigl(C_{G_2}(s)\longrightarrow Q\bigr)
       =\langle s_Q\rangle,
 \qquad
 \operatorname{im}\bigl(C_{G_3}(r)\longrightarrow Q\bigr)
       =\langle r_Q\rangle.
 \tag{A.77.2}
\]
Both equalities hold with \(G_6\) in place of the corresponding
\(G_p\). If \(Q\) retains the standard \(S_3\) quotient, an
automorphism of \(G_6\) preserving the kernel of
\(G_6\to Q\) therefore determines uniquely the
\(Q\)-images of its normalized conjugators of \(s\) and
\(r^{-1}\), whenever these conjugators exist: the normalizations
are images \(1\) and \(s\), respectively, in \(S_3\).
If the automorphism is involutory and acts trivially on \(S_3\),
both lower normalizer norms equal one.

For the original datum of A.76, let \(E\) be its actual finite
six-cup quotient (A.76.7). Every full normalized positive lift
retaining the original finite action and s-frame has negative
square nontrivial already on
\[
 \mathcal G=\Pi/\Phi_2\bigl(\ker(\Pi\to E)\bigr).
 \tag{A.77.3}
\]
The orders of these specified finite groups are
\[
 |E|=2^{1620}\cdot72\cdot7^{13},\qquad
 |\mathcal G|=|E|\,2^{1+|E|/6}.
 \tag{A.77.4}
\]
Nevertheless, all iterated actual scalar refinements of the old
finite datum admit the original compatible involution and both
norm-one frames.

\begin{proof}
\emph{The first relative module controls the centralizer image.}\par
For \(p=2,3\), write \(a=s,r\), respectively. Let \(m\in\Pi\)
lift an element of \(C_{G_p}(a)\), and set
\[
                       J=q^{-1}(\langle a_Q\rangle).
\]
The two elements \(a\) and \(mam^{-1}\) lie in \(J\), although
\(m\) need not lie in \(J\): their images in \(Q\) are equal.
They coincide modulo \(\Phi_p(L)\), which is contained in
\(\Phi_p(J)\), and consequently have the same image in the
mod-\(p\) abelianization of \(J\).

The open profinite Kurosh theorem writes \(J\) as a finite free
profinite product of copies of \(C_p\) and a finite-rank full
free profinite group. There are no factors from the other source
cyclic group, since \(q\) injects it and \(J/L\) has order \(p\).
The free-product torsion theorem conjugates each element of order
\(p\) into one of the cyclic factors. Different factors and
different nonidentity powers of a fixed factor generator have
different images in the mod-\(p\) abelianization. Hence
\(a\) and \(mam^{-1}\) are conjugate in \(J\). These are the
open-subgroup and torsion statements in the full profinite
category; see \citet{A_HR1985}.

Choose \(j\in J\) with \(j^{-1}mam^{-1}j=a\). The full
free-factor centralizer is \(C_\Pi(a)=\langle a\rangle\), so
\(j^{-1}m\in\langle a\rangle\). Thus
\(q(m)\in\langle a_Q\rangle\). The reverse inclusion is
immediate. Since \(G_6\) maps to both \(G_p\), (A.77.2) also
holds for \(G_6\).

The kernel \(L\) is free profinite of rank \(d=1+|Q|/6\), by
open Kurosh and the Euler characteristic of \(C_2\amalg C_3\).
Consequently
\[
            |G_p|=|Q|p^d,\qquad |G_6|=|Q|6^d.
 \tag{A.77.5}
\]
This specifies finite quotients directly; no enumeration of their
elements is required.

\emph{Extracting the lower frames.}\par
Suppose \(b\in\operatorname{Aut}(G_6)\) and
\[
          b(s)=\beta s\beta^{-1},\qquad
          b(r)=Jr^{-1}J^{-1}.
 \tag{A.77.6}
\]
Two conjugators for either equation differ on the right by the
corresponding centralizer. Equation (A.77.2) bounds their
\(Q\)-ambiguity by the relevant cyclic source factor. If their
\(S_3\)-images agree, that ambiguity is trivial. This proves the
asserted uniqueness. Whenever an induced action on \(Q\) is used,
stability of the kernel of \(G_6\to Q\) is retained explicitly.

If \(b^2=1\), then \(b(\beta)\beta\) centralizes \(s\) and
\(b(J)J\) centralizes \(r\). Their \(S_3\)-images are one under
the stated normalizations and trivial \(S_3\)-action. Their
\(Q\)-images therefore lie in the intersection of the pure
\(S_3\)-kernel with the corresponding cyclic factor, and equal
one. This is a necessary lifting condition; existence of such an
action on \(G_6\) is not asserted.

\emph{A single actual noncentral detector for A.76.}\par
A.76 proves that every full candidate retaining its data preserves
\(E\) and has the prescribed action on its six coefficient copies.
Canonical \(c\) preserves the same kernel. Its characteristic
subgroup \(\Phi_2(\ker(\Pi\to E))\) is therefore preserved by
both, so the negative companion \(b=hc\) acts on the same finite
\(\mathcal G\) for every candidate.

Let \(\beta\) be the candidate's full s-conjugator retaining the
original finite frame, and put
\[
                         n=\beta^{-1}b(\beta^{-1}).
 \tag{A.77.7}
\]
Before \(b^2=1\) has been established, this notation denotes a
group element, not a square in a semidirect product by \(C_2\).
Its image in \(E\) is a normalizer square because the induced
negative action on \(E\) is involutory. That image lies in the
six-cup kernel and is nonidentity by the exhaustive invariant
(A.76.10). It therefore lies outside \(\langle s_E\rangle\).
Equation (A.77.2), now with \(Q=E\), shows that \(n\) cannot
centralize \(s\) in \(\mathcal G\). But direct substitution in
\(b(s)=\beta s\beta^{-1}\) gives
\[
                           b^2(s)=n^{-1}sn.
 \tag{A.77.8}
\]
Thus \(b_{\mathcal G}^2(s_{\mathcal G})\ne s_{\mathcal G}\).
The six attained coefficient copies add \(6\cdot108=648\)
binary dimensions to A.76's old finite quotient. This proves
(A.77.4) and the claimed finite detection.

\emph{Every scalar iteration still passes.}\par
Put \(N_0=\ker(\Pi\to D_0)\), for A.76's old quotient, and
\(N_{j+1}=\overline{[N_j,\Pi]}\). The original A.56 kernel
satisfies \(U=[U,\Pi]\) and \(U\subseteq N_0\). Inductively,
\[
                            U\subseteq N_j\quad(j\ge0).
 \tag{A.77.9}
\]
Hence every \(\Pi/N_j\) is an actual quotient of the same
\(D=\Pi/U\). Its original negative involution and both
norm-one frames descend compatibly to all these quotients.
The subgroups are stable because \(N_0\) is stable and the
operation \(N\mapsto[N,\Pi]\) is equivariant.

Each quotient is finite. Inductively, if \(m=[\Pi:N_j]\), the
subgroup \(Z=N_j/[N_j,\Pi]\) is central of index \(m\) in
\(\Pi/[N_j,\Pi]\). It is finitely generated. Transfer restricted
to \(Z\) is the \(m\)-th power map, and it factors through the
abelianization of the ambient group, a quotient of
\(\Pi^{\rm ab}=C_6\). Thus \(Z\) has exponent dividing \(6m\)
and is finite. The existing full positive lift realizes all these
scalar data simultaneously. Its prescribed-frame negative square
is nevertheless detected by (A.77.3). Arbitrarily many scalar
iterations therefore do not replace the noncentral lifting gate.
\end{proof}

The successful reselections below and in A.64--A.75 remain
compatible with this result. The finite detector retains A.76's
literal cusp and prescribed finite frame. Neither an unframed
nonexistence theorem nor a full peripheral conjugacy is inferred.


\statement{A.78}{Proved}{The original finite action passes the six-cup norm test after frame reselection}
Keep the original finite positive action of A.56 and its induced
actual six-cup action used in A.76. There is a different finite
s-frame, with the same finite \(H\)-projection as the initial frame, which has normalizer
norm one both in the old quotient and in the actual six-cup
quotient. Thus the six-cup obstruction in A.76 does not exclude
every frame for that fixed action. No full involution or retention
through every later quotient is asserted.

\begin{proof}
We use the original coordinate normalization and original actions
without changing their parameters \(\mathrm{WE}\) and
\(\mathrm{PSHIFT}\). Let \(v_0,\ldots,v_{91}\) be the recorded
\(k\)-basis of the joint affine s-centralizer and negative-norm
space, where \(k=\mathbb F_8\). In the old class-two group define
\[
             Q=\bigl(7(v_{89}+v_{91}),0,1\bigr),
             \qquad \beta'=\beta Q.
 \tag{A.78.1}
\]
The field element seven is the element encoded as seven in the
fixed binary basis of \(\mathbb F_8\); it is not an integer
multiplicity in characteristic two. The basis vectors and their
full old coordinates are in
\path{A76_fixed_action_frame_space.json}.

The original group law checks
\[
 QsQ^{-1}=s,\qquad
 \beta's\beta'^{-1}=h(s),\qquad
 \beta'^{-1}(hc)(\beta'^{-1})=1.
 \tag{A.78.2}
\]
No old central correction is needed for (A.78.1). In normalized
free cusp rows, the right correction has the following four
Tate-coordinate values:
\[
\begin{array}{c|cccc}
 \text{physical weight}&8&11&20&23\\ \hline
 \text{four-row value}&(0,0,0,0)&(0,0,0,0)&(0,0,0,7)&(0,0,0,7).
\end{array}
 \tag{A.78.3}
\]

For clarity, the finite calculation locating the witness retains
all affine frame variables with the old action fixed. The old
s-class equations have rank 54 in 276 binary variables. At each
of the four coordinates in (A.78.3), the six new norm equations
have exactly 29 solutions among the \(8^4=4096\) four-row values.
Their \(29^4=707281\) combinations give 841 locally allowed
coordinate quadruples compatible with the old class equations.
These are not all old frames: fixing such a quadruple leaves
a 198-dimensional binary affine fiber before the old central
quadratic norm is imposed. The first recorded affine representative
is (A.78.3) and already satisfies
the old full norm by (A.78.2). These counts explain the complete
finite search space; existence relies on the following direct
verification of the saved witness.

Lift \(\beta'\) to the actual six-cup quotient. The six
conjugacy discrepancies lie in the image of \(1-s\), so the
recorded linear corrections give an exact s-conjugator. The
remaining six squares lie in the complete s-centralizer norm
image. Its recorded corrections yield \(\widetilde\beta\) with
\[
 \widetilde\beta s\widetilde\beta^{-1}=h(s),\qquad
 \widetilde\beta^{-1}(hc)(\widetilde\beta^{-1})=1.
 \tag{A.78.4}
\]
Both equations are checked using the original six-cup source
multiplication, inverse, and automorphisms. The source relations,
literal cusp, and negative quotient square are checked at the
same time. The old action parameters are compared exactly with
the original file, not reconstructed from a weaker invariant.

The complete old and six-cup frames, corrections, and direct checks
are stored in
\path{A76_fixed_action_frame_reselection.json},
with SHA256
\[
\begin{gathered}
 \texttt{a927329a694b42a68562354e6311e16a2}\\
 \texttt{47e0f6e4e36e902812d1a599f6a1662}.
\end{gathered}
 \tag{A.78.5}
\]
The producer is
\path{check_A76_fixed_action_frame_reselection.py}.
An independent verifier, which reads the witness without searching,
is
\path{work_CE/code/verify_original_A56_frame_reselection.py};
its output and
\path{work_CE/original_A56_frame_reselection_independent.json}
verify (A.78.2)--(A.78.4), the identical old action, and the same
finite head. This proves the asserted reselection.

The enumerated residual frame family keeps the initial finite
\(H\)-head. Possible changes of that head are not claimed
exhausted. The exhibited solution already suffices to show why
the fixed-frame obstruction cannot be promoted to an all-frame
obstruction at this level. Passing this finite test does not
settle existence of a full negative involution.
\end{proof}

\statement{A.79}{Proved}{A specified finite action forces a noncentral order-four negative image, independently of every frame choice}
There are an actual finite quotient \(Q=\Pi/N_0\),
\(\Pi=C_2\amalg C_3\), and a positive peripheral automorphism
\(h_0\) of Q fixing the literal cusp, with the following properties.
Put
\[
 N_{j+1}=\overline{[N_j,\Pi]},\qquad Q_j=\Pi/N_j.
 \tag{A.79.1}
\]
Every \(Q_j\) is finite. The specified action extends compatibly
to all these actual source-central refinements, with involutory
negative companions and norm-one s- and r-normalizer frames.
One ordinary full positive peripheral automorphism realizes all
these data simultaneously.

Nevertheless, no full positive peripheral h satisfying
\[
                  h(\tau)=\tau,\qquad h_Q=h_0
 \tag{A.79.2}
\]
has \((hc)^2=1\). No s-frame or r-frame is prescribed in this
nonexistence assertion. In fact every such full h preserves a
specified actual finite source quotient \(E_{\rm ord}\) and its
negative companion has exact order four there. The orders are
\[
 |Q|=2^{972}\,72\,7^{13},\qquad
 |E_{\rm ord}|=2^{2700}\,72\,7^{13}.
 \tag{A.79.3}
\]
Thus actual central refinements, even all their compatible
iterations, do not suffice for order-preserving lifting of this
globally negative finite action. The literal-cusp normalization
and the specified finite action remain hypotheses. The later
\(W=0\) actions are not excluded, and the existential
\(\Delta=1\) involution problem remains open.

\begin{proof}
\emph{The specified actual action and the compatible central tests.}\par
Use the changed affine tuple of A.64, before the subsequent
\(W=0\) reselection in A.66. Its exact data are
\path{work_A/A63_zero_t_A56_solution.json}; its SHA256 is
\[
\begin{gathered}
 \texttt{4135fabf1500006d5ec6bbe5dc0571ea}\\
 \texttt{91018abf911e32136cf965729d18d4e9}.
\end{gathered}
 \tag{A.79.4}
\]
The actual finite old model Q has normal forms
\((a,m,q)\), with \(a\in V^8\),
\(m\in\operatorname{Sym}^2_k(k^8)\), \(q\in H\),
\(k=\mathbb F_8\), and \(V=k^{36}\). It uses the exact
opposite-weight group law and C4 source carry of A.56.
The source is the actual pair
\[
 s\longmapsto(d(s),0,H_s),\qquad
 r\longmapsto(0,0,H_r).
 \tag{A.79.5}
\]
In \path{work_A/check_Q6_central_cover.py}, the positive action
is the fixed linear action followed by conjugation by
\texttt{WE}, and the central parity correction
\texttt{PSHIFT}. Substitution of precisely the two entries from
(A.79.4) defines \(h_0\). This specifies an actual finite
automorphism, rather than only a class of conjugators or an
outer action. Canonical c is the actual source automorphism in
the same model. The residual \(P_0=\ker(Q\to H)\) satisfies
\[
              P_0^{\rm ab}=V^8,
 \qquad P_0'=M,
 \qquad M\text{ is H-trivial}.
 \tag{A.79.6}
\]

For the central-refinement assertion, use the actual quotient
\(D=\Pi/U\) of A.65 with the A.64 tuple. Its kernel satisfies
\(U=\overline{[U,\Pi]}\). Its positive action fixes the literal
cusp, its negative companion is involutory, and both quotient
normalizer norms are one. A.65 supplies a full ordinary positive
lift realizing these data. This prior construction is independent
of the ordinary-cup obstruction below.

The finite Q is an action-stable quotient of this D. Since
\(U\subseteq N_0\), induction gives
\[
 U=\overline{[U,\Pi]}\subseteq N_j\quad(j\geq0).
 \tag{A.79.7}
\]
Consequently every actual quotient in (A.79.1) is a quotient of D
and inherits its compatible action and both norm-one frames.
The same full ordinary lift realizes all of them.

Finiteness holds at every step. If \(N_j\) is open, the central
subgroup
\(Z_j=N_j/\overline{[N_j,\Pi]}\) of \(Q_{j+1}\)
has index \(m_j=|Q_j|\) and is topologically finitely generated.
The continuous transfer \(Q_{j+1}\to Z_j\) restricts on this
central subgroup to \(z\mapsto z^{m_j}\). Its image has exponent
dividing six, since \(Q_{j+1}^{\rm ab}\) is a quotient of
\(\Pi^{\rm ab}=C_6\). Thus \(Z_j\) has exponent dividing \(6m_j\),
so it is finite. This proves the induction and, in particular,
the finiteness of the actual first refinement
\(Q_1=\Pi/\overline{[N_0,\Pi]}\).

\emph{The actual ordinary-cup detector.}\par
Let \(a:Q\to V^4\rtimes H\) be the four normalized free
cusp cocycles, writing \(a(g)\) for their additive coordinate.
They satisfy \(a(s)=0\) literally. Put
\(T_0=V^{(3)}\) with its cubed monomial H-action and
\(T=T_0^{16}\). On \(V^4\times T\) use the class-two law
\[
 (a,z)(b,w)=(a+b,z+w+\mathcal C(a,b)),
 \qquad \mathcal C(a,b)_{ij}=a_i^{[2]}*b_j,
 \tag{A.79.8}
\]
for all sixteen ordered pairs \(0\leq i,j\leq3\), with
coordinatewise products and squares. H acts strictly. Pairing
with Q gives the group whose normal forms and product are
\[
 \begin{aligned}
 E_{\rm ord}&=\{(g,z):g\in Q,\ z\in T\},\\
 (g,z)(h,w)&=
 \bigl(gh,\ z+\rho_3(g)w+
       \mathcal C(a(g),\rho(g)a(h))\bigr).
 \end{aligned}
 \tag{A.79.9}
\]
The actual source elements are
\[
 S=(s,0),\qquad
 R=\bigl(r,\mathcal C(a(r),\rho(r)^2a(r))\bigr).
\]
The latter is obtained by conjugating pure r by
\((\rho(r)a(r),0)\) in the ordinary factor, so \(R^3=1\).
This is the source map already used in A.66 and implemented in
\path{work_A/check_A66_ordinary_cubic.py}, with the stored
action (A.79.4).

This is the full actual source image. At a single physical
weight, brackets between distinct free rows give
\((u^2v,v^2u)\), whose binary span is \(k^2\). A trace
annihilator would make a polynomial with distinct powers two
and four vanish on all of k, forcing both coefficients to zero.
The same-row brackets \(u^2v+v^2u\) similarly span k.
The old center pairs opposite physical weights, so these
single-weight brackets have zero old-center component.
They attain all sixteen new copies independently. The already
attained old affine image and the finite Frattini argument of
A.66 therefore give precisely (A.79.9). Its kernel is all T,
which proves the second order formula in (A.79.3).

\emph{Every candidate preserves this detector and has the forced
coefficient action.}\par
The actual normal odd subgroup of H has no invariants on
\(T_0\). Coprime Hochschild--Serre and (A.79.6) give
\[
 \begin{gathered}
 H^1(H,T_0)=H^2(H,T_0)=0,\\
 H^1(Q,T_0)\simeq
       \operatorname{Hom}_H(V^8,T_0)\simeq k^8.
 \end{gathered}
 \tag{A.79.10}
\]
The Hom is over the underlying additive binary modules.
The literal intertwiner is
\(j_3(v)_\lambda=v_{\operatorname{partner}(\lambda)}^4\);
we retain this Frobenius comparison instead of identifying the
two k-representations without it. The eight actual source
classes and 36 coboundaries have rank 44. The source cocycle
space has dimension \(20+24=44\), from the order-two and
order-three source relations, and its coboundary space has
dimension 36. Thus inflation in
(A.79.10) is an isomorphism onto \(H^1(\Pi,T_0)\).
The continuous five-term sequence therefore makes
\[
 \operatorname{Hom}_Q(N_0,T_0)
                 \longrightarrow H^2(Q,T_0)
 \tag{A.79.11}
\]
injective; see \citet{APoonen2017}.

The sixteen extension classes in (A.79.9) are the literal
Frobenius cup products \(e_i^{[2]}\smile e_j\).
They are independent by the actual source attainment above.
The action \(h_0\) fixes the four cusp cohomology classes;
canonical c preserves them with its actual coefficient
intertwiner. Equivariance and (A.79.11) force every full h in
(A.79.2) to preserve all sixteen actual kernel maps and to act
identically on T. Canonical c likewise preserves their joint
kernel and acts by the cubed monomial operator \(C_3\).
Thus descent and the coefficient action are forced for every
candidate; they are not additional restrictions imposed on h.

\emph{Centralizer surjectivity removes every frame choice.}\par
For every \(g\in C_Q(s)\), the literal lift \((g,0)\)
centralizes S in (A.79.9). Indeed both cup terms in the products
with S vanish because \(a(s)=0\), and the old products agree.
Consequently
\[
                  C_{E_{\rm ord}}(S)\longrightarrow C_Q(s)
                   \quad\text{is surjective}.
 \tag{A.79.12}
\]
This includes every centralizing head and every residual affine
or central coordinate.

Let two automorphisms of \(E_{\rm ord}\) cover \(h_0\), act
identically on T, and preserve the conjugacy class of S. Choose
their S-conjugators \(B_1,B_2\). The image of
\(B_1^{-1}B_2\) centralizes s in Q. By (A.79.12), cancel a
lift of this centralizer from \(B_2\); the remaining difference
from \(B_1\) lies in T. Hence the two S-images are T-conjugate.
Their difference cocycle has zero s-restriction class, without
prescribing either conjugator or even its finite head.

\emph{Exhausting all actions over the specified finite action.}\par
Let \(h_*\) be the actual finite action in the saved ordinary
certificate. It fixes \(\tau\) literally and preserves the
S-conjugacy class. Every descended candidate differs from it by
a T-valued 1-cocycle on Q. The preceding paragraph kills its
s-restriction class; the literal cusp kills its tau restriction.
After pullback to the actual source cyclic subgroups, the
complete joint restriction isomorphism is
\[
 H^1(Q,T_0)\longrightarrow
 H^1(\langle s\rangle,T_0)
       \oplus H^1(\langle\tau\rangle,T_0),
 \qquad 8=4+4.
 \tag{A.79.13}
\]
This is the full source restriction of A.56.5, transported by
\(j_3\), also verified by the actual 44/40 coefficient matrices.
The subgroups in (A.79.13) belong to \(\Pi\); no finite cyclic
cohomology group is substituted for the cusp subgroup.

It follows that every candidate has the form
\[
                  h_E=\operatorname{Inn}(v)h_*,
                  \qquad v\in T^\tau.
 \tag{A.79.14}
\]
The vector is unique because \(T^H=0\). The actual negative
coefficient \(C_3\) fixes \(T^\tau\) pointwise. Therefore,
putting \(b_*=h_*c\), we have
\[
 (\operatorname{Inn}(v)b_*)^2
       =\operatorname{Inn}(v+b_*(v))b_*^2=b_*^2.
 \tag{A.79.15}
\]
This exhausts every peripheral action lift of \(h_0\), including
all choices of s- and r-frames.

\emph{The exact nonzero square and verification.}\par
The original-law certificate
\path{work_A/A66_ordinary_cubic_result.json} gives
\[
 b_*^2=\operatorname{Inn}(t),\qquad
 t\in T^\tau,\qquad
                  t_{00,\,\mathrm{weight}\,0}=1\in k.
 \tag{A.79.16}
\]
The equality is checked on the actual S and R. Since \(T^H=0\),
this nonzero t induces a nonidentity inner automorphism.
As T has exponent two, the square in (A.79.16) has order two,
so \(b_*\), and by (A.79.15) every candidate negative image,
has exact order four. A full negative involution would induce
an involution on this preserved finite quotient, a contradiction.

The independent earlier source-law projection is
\path{work_CE/A66_ordinary_projection_independent.json}.
The separate fresh verifier
\path{work_A/code/verify_fixed_action_ordinary_square.py}
preserves the stored witness and checks the actual source and
canonical-c identities, literal cusp, ranks 44/40/36, and all 192
binary cusp-inner directions. It verifies the coordinate one
in (A.79.16) and that the square moves an actual source generator.
Its genuine stdout is the corresponding \path{.out} file;
the complete results and input hashes are in
\path{work_A/fixed_action_ordinary_square_verified.json}.
Its imported legacy group builder performs existing exploratory
algebra but writes no artifacts. No enumeration of the enormous
group in (A.79.3) or search-timeout inference is used.

The finite action (A.79.4) is essential. A.66 and A.75 use
changed actions with zero free cusp inner term and are compatible
with this obstruction. No assertion is made about every finite
action with \(\Delta=1\), an unnormalized full cusp image,
or the universal conjugacy question.
\end{proof}

\statement{A.80}{Proved}{An odd relation-module lift and a finite negative action with no full involutory realization, without prescribing a cusp normalization}
Let \(\Pi=C_2\amalg C_3=\langle s,r\rangle\),
\(\tau=sr\), and let \(q:\Pi\twoheadrightarrow Q\) be a finite
quotient injecting the two factors. Write \(s_0=q(s)\),
\(r_0=q(r)\), and \(t=q(\tau)\). Suppose that
\(b_0\in\operatorname{Aut}(Q)\) and \(\beta_0,J_0\in Q\) satisfy
\[
 \begin{gathered}
 b_0^2=1,\qquad b_0(t)=t^{-1},\\
 b_0(s_0)=\beta_0s_0\beta_0^{-1},\qquad
 b_0(r_0)=J_0r_0^{-1}J_0^{-1},\\
 b_0(\beta_0)\beta_0=b_0(J_0)J_0=1.
 \end{gathered}
 \tag{A.80.1}
\]
For every prime \(p\nmid2|Q|\), put
\[
 L=\ker q,\qquad
 \Phi_p(L)=\overline{[L,L]L^p},\qquad
 G_p=\Pi/\Phi_p(L).
 \tag{A.80.2}
\]
There are an involution \(b_p\in\operatorname{Aut}(G_p)\)
above \(b_0\) and lifts \(\beta_p,J_p\) of the prescribed
\(\beta_0,J_0\) such that
\[
 \begin{gathered}
 b_p(\tau)=\tau^{-1},\qquad
 b_p(s)=\beta_ps\beta_p^{-1},\qquad
 b_p(r)=J_pr^{-1}J_p^{-1},\\
 b_p(\beta_p)\beta_p=b_p(J_p)J_p=1.
 \end{gathered}
 \tag{A.80.3}
\]
These are actual first relation-module quotients, of order
\[
                      |G_p|=|Q|p^{1+|Q|/6}.
 \tag{A.80.4}
\]

Apply this assertion to the specified \(Q,h_0\) of A.79,
\(b_0=h_0c_Q\), and \(p=5\). The resulting finite negative
involution \(b_5\) is not induced by any full negative peripheral
involution of \(\Pi\), even if the full cusp is only required to
be inverted up to conjugacy. Every permitted semisimple extension
chosen in the proof gives such a finite action. Equivalently, the finite positive
peripheral action \(h_5=b_5c_{G_5}\) has no full positive
peripheral lift whose negative companion has order two. Its
order-two realization fails without prescribing any full cusp
normalization or any full normalizer frame.

The existence of an ordinary full positive lift of this new
\(h_5\), or of compatible actions on its actual central
refinements, is not asserted. Thus this conclusion does not
extend A.79's central-refinement assertion to a newly chosen
coefficient action. It also does not exclude the surviving
\(W=0\) actions or settle the main existential question.

\begin{proof}
\emph{An index-two semisimple extension lemma.}\par
We first prove the precise positivity fact needed below. Let
\(Q\lhd E\) have index two, let k be a field whose
characteristic does not divide \(|E|\), and let X be an actual
finite-dimensional kQ-module. If its class in \(K_0(kQ)\)
belongs to the restriction image of \(K_0(kE)\), then X is
the restriction of an actual kE-module. Here \(K_0\) is the
free abelian group on the simple modules. Semisimplicity over
the stated field follows from the general-field form of
Maschke's theorem; see \citet[Theorem~7.4 and
Proposition~7.9]{AMilne2025}.

For completeness, the induction adjunction used in the proof
is the explicit bijection
\[
 \begin{aligned}
 \operatorname{Hom}_{kE}(kE\otimes_{kQ}S,T)
       &\longrightarrow\operatorname{Hom}_{kQ}(S,T),\\
 F&\longmapsto\bigl(u\mapsto F(1\otimes u)\bigr).
 \end{aligned}
 \tag{A.80.5}
\]
Its inverse sends f to \(g\otimes u\mapsto g f(u)\).
The tensor relation makes the inverse well defined, so this
is valid over k without any complex-character assumption.

Fix a simple kQ-module S and let \(S^b\) denote its twist
by an element of \(E\setminus Q\). If \(S^b\not\simeq S\),
then
\[
 \operatorname{Res}_Q^E\operatorname{Ind}_Q^E S
           \simeq S\oplus S^b.
\]
Every E-endomorphism preserves the S-summand, and its action
there determines its action on the other summand. Thus its
endomorphism ring is \(\operatorname{End}_{kQ}(S)\), a
division ring. Semisimplicity makes \(\operatorname{Ind}S\)
simple. By (A.80.5), every E-simple whose restriction contains
S is this one. On this two-element orbit the restriction
semigroup therefore contains exactly the nonnegative multiples
of \([S]+[S^b]\), which are precisely the positive classes in
its restriction lattice.

If \(S^b\simeq S\), decompose
\[
 \operatorname{Ind}_Q^E S\simeq\bigoplus_i T_i^{n_i},
 \qquad \operatorname{Res}_Q^E T_i\simeq S^{e_i}.
\]
Restriction gives \(\sum_i n_i e_i=2\). Every E-simple
supported on S occurs among these \(T_i\), by (A.80.5).
Consequently all available restriction multiplicities are one,
or the sole available multiplicity is two. Every nonnegative
integer in their restriction lattice is again realized by an
actual nonnegative sum of E-simples. The Q-isotypic support of
an E-simple is one orbit under \(E/Q\), since an orbit sum is
E-stable. Applying the preceding two cases independently to
these disjoint orbits proves the lemma. This is an index-two
argument, not a general positivity assertion for restriction.

\emph{The actual relation module has a virtual extension.}\par
Set \(k=\mathbb F_p\), \(M=L/\Phi_p(L)\), and
\(E=Q\rtimes_{b_0}\langle\epsilon\rangle\), where
\(\epsilon^2=1\). The two elements
\[
                 d_s=\beta_0^{-1}\epsilon,
                 \qquad d_r=J_0^{-1}\epsilon
 \tag{A.80.6}
\]
are involutions. The subgroups
\(A_s=\langle s_0,d_s\rangle\simeq C_2\times C_2\)
and \(A_r=\langle r_0,d_r\rangle\simeq S_3\) have
intersections \(\langle s_0\rangle\) and
\(\langle r_0\rangle\) with Q, respectively. This is
where both actual norm-one frame hypotheses are used.

The Bass--Serre cover of the source has edge set Q and vertex
sets \(Q/\langle s_0\rangle\) and
\(Q/\langle r_0\rangle\). Its first homology is exactly M,
so its augmented cellular sequence gives
\[
 \begin{split}
 0\longrightarrow M\longrightarrow k[Q]\longrightarrow
 k[Q/\langle s_0\rangle]\oplus k[Q/\langle r_0\rangle]
 \longrightarrow k\longrightarrow0.
 \end{split}
 \tag{A.80.7}
\]
One can equally form this finite graph for the discrete modular
group and complete its finite-index free subgroup; the induced
topology on an open subgroup is full. Thus (A.80.7) computes
the actual source relation module. Its dimension is
\(1+|Q|/6\), proving (A.80.4), as in A.77.

Each term on the right of the corresponding virtual identity
\[
 [M]=[k]+[k[Q]]-[k[Q/\langle s_0\rangle]]
                   -[k[Q/\langle r_0\rangle]]
 \tag{A.80.8}
\]
is the restriction of an E-module. Use the trivial module,
\(k[E/\langle\epsilon\rangle]\), \(k[E/A_s]\), and
\(k[E/A_r]\), respectively. It follows that \([M]\) lies
in the restriction image. We are not asserting that the
particular differential in (A.80.7) is E-equivariant.

\emph{Extend the actual cusp submodule with negative sign.}\par
Put \(m=\operatorname{ord}(t)\). The torsion-free discrete
cover underlying L is an oriented surface with
\([Q:\langle t\rangle]\) cusps, all of width m. The
positive boundary classes have only their total-sum relation.
Hence their Q-submodule P in M is
\[
 P\simeq
 k[Q/\langle t\rangle]\big/
 k\!\left(\sum_{g\langle t\rangle}e_{g\langle t\rangle}\right),
 \qquad e_{\langle t\rangle}\longmapsto[\tau^m].
 \tag{A.80.9}
\]
This is the same elementary boundary-homology calculation used
in C.29. When there is one cusp the quotient is zero.
Because \(b_0(t)=t^{-1}\), the rule
\[
       \epsilon e_{g\langle t\rangle}
              =-e_{b_0(g)\langle t\rangle}
 \tag{A.80.10}
\]
gives an E-module extending this Q-module and preserves its
sole relation. It sends \([\tau^m]\) to its negative.

Maschke supplies a Q-splitting \(M=P\oplus X\).
Since \([X]=[M]-[P]\) is in the restriction image, the
index-two lemma extends X to an actual E-module. Combine
that extension with (A.80.10). On the actual M this gives
an invertible k-linear map B satisfying
\[
 B(qv)=b_0(q)B(v),\qquad B^2=1,
 \qquad B[\tau^m]=-[\tau^m].
 \tag{A.80.11}
\]
The action on P has been retained literally under the chosen
Q-splitting; it is not merely an equality of virtual classes.

\emph{The cusp equation and the odd norm on the actual group.}\par
The finite extension \(1\to M\to G_p\to Q\to1\)
splits because \(p\nmid|Q|\). This follows, for example,
by coprime averaging in \(H^2(Q,M)\). Choose a splitting
and write \(G_p=M\rtimes Q\), with the actual source
cusp \(\tau=(a,t)\). Set
\[
        \bar a=\frac1m\sum_{j=0}^{m-1}t^j a.
        \qquad \tau^m=m\bar a\in M.
 \tag{A.80.12}
\]
Equation (A.80.11) says \(B\bar a=-\bar a\).
Therefore the right side of
\[
                   (1-t^{-1})k=-t^{-1}a-Ba
 \tag{A.80.13}
\]
has zero t-average. Coprime averaging identifies that kernel
with \((1-t^{-1})M\), so choose a solution k. The map
\[
 b_*(v,q)=\bigl(Bv+(1-b_0(q))k,\ b_0(q)\bigr)
 \tag{A.80.14}
\]
is an automorphism inducing \(b_0\), and (A.80.13) is
exactly \(b_*(\tau)=\tau^{-1}\). Its square is
\(\operatorname{Inn}(v_0)\), where \(v_0=k+Bk\).
The equality \(b_*^2(\tau)=\tau\) gives
\(v_0\in M^t\), and \(Bv_0=v_0\). Thus
\[
                        b_p=\operatorname{Inn}(-v_0/2)b_*
 \tag{A.80.15}
\]
still inverts the literal cusp and has square one. This
correction is allowed because p is odd and lies in the actual
cusp centralizer. No odd-order assertion about the entire
automorphism kernel has been used.

\emph{Retain both prescribed lower frame values.}\par
Work in \(\widetilde E=G_p\rtimes_{b_p}C_2\), which
maps to E. For \(A=\langle s\rangle\) or
\(\langle r\rangle\), the normalizer maps onto the
normalizer of \(A_Q\). Indeed a lift of an element of
the latter conjugates A to another complement in the
preimage of \(A_Q\). Coprime complement conjugacy moves
it back to A by an element of M. In this elementary-abelian
case this is just \(H^1(A_Q,M)=0\), proved by averaging.
The kernel of the normalizer map is \(C_M(A)\): a
normalizing element of M acts trivially on A because
\(A\to A_Q\) is injective.

Choose a normalizer lift z of the prescribed \(d_s\) or
\(d_r\) in (A.80.6). Its square is a vector w in
\(C_M(A)\), fixed by conjugation by z. Replacing z by
\((-w/2)z\) makes its square one without changing its
image or its normalizing property. Its action on A is the
required one, since that action is already prescribed on
the faithfully mapped \(A_Q\). Writing the two adjusted
elements as \(\beta_p^{-1}\epsilon\) and
\(J_p^{-1}\epsilon\) proves all of (A.80.3), with
the literal prescribed lower values. These are two separate
normalizer choices; no additional common transporter is
asserted. The resulting \(b_p\) is nontrivial: if Q is
cyclic it is C6 and cusp inversion is already nontrivial
there; otherwise (A.80.9) makes the actual cusp have
nontrivial p-part, which is inverted.

\emph{The odd relation quotient detects every cusp-conjugator head.}\par
For any noncyclic factor-faithful Q and any odd p, not
necessarily prime to \(|Q|\), the same boundary calculation
gives
\[
          \operatorname{im}\bigl(C_{G_p}(\tau)\to Q\bigr)
                        =\langle t\rangle.
 \tag{A.80.16}
\]
There are \(r=|Q|/m\geq2\) positive boundary classes.
They are distinct over \(\mathbb F_p\): if \(r\geq3\),
the difference of two coordinate vectors is not a multiple
of the total-sum relation; if \(r=2\), the two nonzero
classes are opposite and p is odd. A centralizer of
\(\tau\) fixes \([\tau^m]\), so its deck action fixes
the identity cusp coset and its Q-image lies in
\(\langle t\rangle\). Powers of the actual cusp give
the reverse inclusion.

In particular,
\[
                   C_\Pi(\tau)=\overline{\langle\tau\rangle}
                              \simeq\widehat{\mathbb Z}.
 \tag{A.80.17}
\]
For the centralizer equality, refine every finite quotient
by the standard S3 quotient and apply (A.80.16);
compactness of the closed procyclic subgroup gives the
claim. Its full \(\widehat{\mathbb Z}\)-type can be
checked without a torsion assumption: modulo every integer
n the matrices
\[
 \begin{pmatrix}0&-1\\1&0\end{pmatrix},\qquad
 \begin{pmatrix}0&-1\\1&1\end{pmatrix}
\]
define a source quotient in
\(\operatorname{SL}_2(\mathbb Z/n)/\{\pm I\}\), and
their product is the class of
\(\left(\begin{smallmatrix}1&1\\0&1\end{smallmatrix}\right)\),
of exact order n.

\emph{Involution-preserving normalization and the A.79 consequence.}\par
Suppose that a full negative involution b induces the
constructed \(b_5\), and write
\[
                        b(\tau)=u\tau^{-1}u^{-1}.
\]
Since \(b^2=1\), the element \(z=b(u)u\) centralizes
\(\tau\). On the one hand \(b(z)=u z u^{-1}\);
on the other hand (A.80.17) and cusp inversion give
\(b(z)=u z^{-1}u^{-1}\). Hence \(z^2=1\), and
(A.80.17) forces \(z=1\).

The image of u in \(G_5\) centralizes the actual cusp,
because \(b_5\) inverts it literally. By (A.80.16),
\(u_Q\in\langle t\rangle\). Choose
\(v\in\widehat{\mathbb Z}\) with \(u_Q=t^v\), and
replace u by \(u\tau^{-v}\). This changes neither its
cusp-conjugacy equation nor its norm-one property, and now
\(u_Q=1\). The automorphism
\[
                           b'=\operatorname{Inn}(u^{-1})b
 \tag{A.80.18}
\]
remains an involution, inverts \(\tau\) literally, and
still induces \(b_0\) on Q. Then \(h=b'c\) is a full
positive peripheral automorphism fixing \(\tau\) literally
and inducing \(h_0\). This contradicts A.79.

All hypotheses for the construction hold for that Q: its
order is \(2^{972}\,72\,7^{13}\), its retained quotient
frames satisfy (A.80.1), and \(5\nmid|Q|\). Canonical c
preserves \(L\) and its characteristic subgroup
\(\Phi_5(L)\), so \(c_{G_5}\) and \(h_5=b_5c_{G_5}\)
are actual automorphisms of the same finite source quotient.
This proves the stated finite nonrealization assertion.
The semisimple choice in (A.80.11) does not supply an
ordinary full realization or a lift to a further source-central
quotient, and none has been used.
\end{proof}



\statement{A.81}{Proved}{Actual central compatibility and ordinary full realization do not restore an involution, even without cusp or frame normalization}
Let \(q:\Pi\twoheadrightarrow Q\), \(h_0\), and
\(b_0=h_0c_Q\) be the specified actual finite data of A.79,
and put \(L=\ker q\). For infinitely many primes
\(p\nmid2|Q|\), there are positive peripheral actions
\(h_p\) on the actual finite quotient
\[
 N_p=\overline{[L,L]L^p},\qquad G_p=\Pi/N_p
 \tag{A.81.1}
\]
with all the following properties. Their negative companions
\(b_p=h_pc_{G_p}\) are involutions fixing the literal
negative cusp and admitting both norm-one normalizer frames
above the prescribed Q-frames. The action extends to an
involution on the \emph{actual first source-central refinement}
\[
                         D_p=\Pi/\overline{[N_p,\Pi]},
 \tag{A.81.2}
\]
again with literal cusp inversion and both norm-one frames.
The corresponding positive action satisfies the ordinary
lifting criterion: some full positive peripheral h induces
\(h_p\), and h may fix \(\tau\) literally.

Nevertheless, no full negative peripheral involution induces
\(b_p\), even if its full cusp is only inverted up to
conjugacy and no full normalizer frame is prescribed. Thus
none of the ordinary full lifts of \(h_p\) can have
involutory negative companion. This assertion concerns a
constructed family at infinitely many eligible fresh primes;
no particular such prime or enormous coefficient matrix is
claimed to have been computed. Compatibility at every
iterated central refinement is not asserted. The surviving
\(W=0\) choices and the main existential problem remain open.

\begin{proof}
The proof identifies the new central obstruction exactly and
chooses the coefficient action to remove it. Throughout,
commutators are \([x,y]=xyx^{-1}y^{-1}\).

\emph{The compact-surface module admits an anti-symplectic
involution over algebraic numbers.}\par
Work first over \(k=\overline{\mathbb Q}\). Let
\(\mathscr M\) be the first homology of the actual
torsion-free modular surface cover, \(\mathscr B\) its
boundary submodule, and
\(\mathscr H=\mathscr M/\mathscr B\) its compact-surface
homology with the Q-invariant nondegenerate intersection
form. By the actual source sequences in A.80,
\[
 [\mathscr H]=2[1]+[kQ]
 -[k(Q/\langle s_0\rangle)]-[k(Q/\langle r_0\rangle)]
 -[k(Q/\langle t\rangle)],
 \tag{A.81.3}
\]
where \(t=q(\tau)\). The trivial multiplicity is zero.
For a nontrivial irreducible Q-module S its multiplicity is
\[
 m_S=\dim S-\dim S^{s_0}-\dim S^{r_0}-\dim S^t.
 \tag{A.81.4}
\]
Duality and twisting by \(b_0\) commute on the simple
modules. The fixed-space dimensions in (A.81.4) are unchanged
by either operation. We construct a \(b_0\)-semilinear
involution \(B_{\mathscr H}\) which negates the actual
intersection form, treating their joint orbits.

On four distinct constituents \(S,S^*,S^{b_0},(S^*)^{b_0}\),
choose a twisted isomorphism between the two dual pairs,
use the negative adjoint on the dual partner, and use the
inverse on the return maps. This gives an anti-symplectic
involution. If \(S\not\simeq S^*\) and \(b_0\) fixes
both classes, choose an intertwiner A on S and rescale it
to have \(A^2=1\). Use A on its multiplicity block and
\(-(A^{-1})^*\) on its dual block.

Suppose next that \(S\not\simeq S^*\) and
\(S^{b_0}\simeq S^*\). There is a nondegenerate twisted
bilinear form F, unique up to scalar, satisfying
\[
                    F(qx,b_0(q)y)=F(x,y).
 \tag{A.81.5}
\]
Its transpose is either F or \(-F\). In the symmetric
case, tensor it with a symmetric multiplicity form. In the
alternating case we need \(m_S\) even, as follows.
The actual norm-one lower frames make
\[
 F_s(x,y)=F(x,\beta_0y),\qquad F_r(x,y)=F(x,J_0y)
 \tag{A.81.6}
\]
alternating and nondegenerate: transposition uses
\(b_0(\beta_0)=\beta_0^{-1}\) and
\(b_0(J_0)=J_0^{-1}\). Moreover, \(s_0\) is an
isometry for \(F_s\), whereas \(r_0\) is self-adjoint
for \(F_r\) and t is self-adjoint for F. For example,
\[
 F_r(r_0x,r_0^{-1}y)=F_r(x,y),\qquad
 F(tx,t^{-1}y)=F(x,y).
\]
The finite-order operators are semisimple. Their fixed
subspaces are nondegenerate for these alternating forms:
in the self-adjoint cases distinct eigenspaces are
orthogonal, and in the isometric case the eigenvalue-one
space pairs only with itself. Thus all four dimensions on
the right of (A.81.4) are even, and \(m_S\) is even.
Tensor F with an alternating multiplicity form. In both
cases the resulting twisted form on \(X=S^{m_S}\) is
symmetric. If \(A:X\to X^*\) is its matrix, then
\[
                     (x,\phi)\longmapsto(A^{-1}\phi,Ax)
 \tag{A.81.7}
\]
on \(X\oplus X^*\) is the required anti-symplectic
semilinear involution.

Suppose S is self-dual and its class is fixed by \(b_0\).
Let its invariant form F have symmetry
\(\varepsilon\in\{1,-1\}\). Rescale an intertwiner A
so that \(A^2=1\). Its multiplier on F is then
\(\nu\in\{1,-1\}\). The actual symplectic isotypic
block has form \(F\otimes G\), where G has symmetry
\(-\varepsilon\). Seek \(A\otimes J\). If
\(\nu=-1\), take \(J=1\). If \(\nu=1\), take an
involutory anti-isometry of G. For alternating G this
exists on a symplectic basis. For symmetric G it exists
over k when its dimension is even, by a hyperbolic basis.
In this latter case F is alternating; every fixed space
of a finite-order Q-element is nondegenerate alternating,
so (A.81.4) again proves the required even dimension.
Finally, when \(b_0\) exchanges two distinct self-dual
classes, choose a twisted anti-isometry between their
isotypic blocks and its inverse. The invariant forms of
each given symmetry and dimension are isometric over k;
a square root of \(-1\) changes an isometry to an
anti-isometry if necessary. These cases exhaust the joint
duality and twisting orbits.

All constructions use algebraic numbers. Extend
\(\mathscr B\) by the signed cusp permutation of A.80
and choose a rational Q-splitting of
\(\mathscr M\to\mathscr H\). We obtain B on
\(\mathscr M\), defined over a number field, with
\[
 B^2=1,\qquad B(qv)=b_0(q)Bv,
 \qquad B|_{\mathscr B}\text{ is the negative cusp permutation},
 \tag{A.81.8}
\]
and with anti-symplectic induced action on \(\mathscr H\).
Semisimplicity here is over the stated field, as in
\citet[Theorem~7.4]{AMilne2025}; the preceding form and
parity argument is additional to A.80's module extension.

\emph{Choose eligible fresh primes.}\par
Take a finite Galois number field K containing the entries
of B and the finitely many basis changes just used. Avoid
their denominators and determinants, and avoid \(2|Q|\).
For clarity, existence of infinitely many completely split
primes needs no density estimate. Choose an integral
primitive element of K with monic minimal polynomial f
and \(f(0)\ne0\). For any finite excluded set of primes
\(\mathcal S\), enlarge it to contain the prime divisors
of \(\operatorname{disc}(f)\) and the preceding bad
integers. Put \(m_{\mathcal S}=\prod_{\ell\in\mathcal S}\ell\)
and \(A=f(0)m_{\mathcal S}\). The nonconstant polynomial
\[
                 g(X)=f(AX)/f(0)\in\mathbb Z[X]
                 \quad\text{satisfies}\quad
                 g(X)\equiv1\pmod{m_{\mathcal S}}.
\]
A prime divisor p of an integer \(|g(n)|>1\) is outside
\(\mathcal S\), and f has a root modulo p. Since
\(p\nmid\operatorname{disc}(f)\), this gives an unramified
degree-one prime of K over p. Galoisness makes all residue
degrees one, so p splits completely. Applying this to
arbitrary finite excluded sets proves the assertion;
compare the standard stronger statement in
\citet[Theorem~8.31 and Corollary~8.32]{AMilneANT2020}.
Reduction at such a prime gives the actual module
\[
 M=H_1(L,\mathbb F_p),\qquad
 P=\langle Q\cdot[\tau^m]\rangle\subset M,
 \qquad H=M/P,
 \tag{A.81.9}
\]
where \(m=\operatorname{ord}(t)\), and an operator B
satisfying (A.81.8) modulo p, with anti-symplectic
action on H. The integral boundary submodule and compact
homology are torsion-free, so the avoided basis
determinants ensure these are the actual mod-p source
modules and forms. This is an existence construction at
infinitely many fresh primes, not a numerical selection
of p or a stored matrix.

\emph{Identify the actual new central p-kernel.}\par
Put \(N=N_p\), \(Z=N/\overline{[N,\Pi]}\), and
let \(E_p\) be the actual quotient of \(D_p\) obtained
by killing the prime-to-p factors of Z. A.79's transfer
argument shows that \(D_p\) and Z are finite. Write
\(A=\ker(E_p\to G_p)\), and let \(\widetilde P\)
be the preimage of M. Its extension by A is central.
The source module has \(M_Q=0\), since its trivial
multiplicity in (A.80.8) is zero.

The kernel of \(\widetilde P^{\rm ab}\to M\) is
Q-trivial. Coinvariants are exact on finite abelian
p-groups when \(p\nmid|Q|\). A nonzero kernel would
therefore give a nonzero Q-invariant character of
\(\widetilde P\), hence of L, to \(C_p\), contrary
to \(M_Q=0\). It follows that
\[
              \widetilde P^{\rm ab}=M,
              \qquad A=\widetilde P',
              \qquad A\text{ is a quotient of }
                    W=(\bigwedge^2_{\mathbb F_p}M)_Q.
 \tag{A.81.10}
\]
The central commutator group thus has exponent p. The
pth-power map \(M\to A\) is additive for odd p,
since the class-two cross term has coefficient
\(p(p-1)/2\). It is Q-equivariant with trivial target,
so it vanishes. Hence \(\widetilde P\) has exponent p.

The upper bound is attained by an actual source quotient.
Define \(\mathcal E(M)=M\times W\) with law
\[
                 (a,z)(b,w)=(a+b,z+w+\tfrac12(a\wedge b)_Q),
 \tag{A.81.11}
\]
and give Q its natural action on M and trivial action on
W. Split \(G_p=M\rtimes Q\), and write its actual
source elements as \((a_s,s_0)\) and \((a_r,r_0)\).
In \(\mathcal E(M)\rtimes Q\) use the source images
\[
 S=(a_s,0,s_0),\qquad
 R=(a_r,-\tfrac16(a_r\wedge r_0a_r)_Q,r_0).
 \tag{A.81.12}
\]
They have orders two and three. Indeed the cube of
\((a_r,z,r_0)\) has central coordinate
\(3z+(a_r\wedge r_0a_r)_Q/2\). The source image
surjects \(G_p\); its L-image surjects M, so its
commutators attain the whole W. Thus the source image is
the entire displayed group. Together with (A.81.10),
this identifies the actual marked quotient
\[
                         E_p=\mathcal E(M)\rtimes Q.
 \tag{A.81.13}
\]

\emph{Remove precisely the central cusp discrepancy.}\par
An inner change of the entire marked model identification
over Q arranges
\[
             \tau=(a_0,z,t),\qquad ta_0=a_0,
             \qquad \tau^m=(ma_0,mz,1).
 \tag{A.81.14}
\]
Lift the required affine conjugation to \(\mathcal E(M)\)
and transport both source generators together. This changes
the model identification, not the actual source kernel;
the strict target law (A.81.11) remains the same.
Here \(a_0=[\tau^m]/m\) spans the identity boundary
line and \(P=\langle Q a_0\rangle\). The chosen B
sends \(a_0\) to \(-a_0\).

Let \(B_W\) be its induced exterior-square action on W.
The homogeneous action
\((a,w,q)\mapsto(Ba,B_Ww,b_0(q))\) is involutory,
and its central cusp discrepancy is
\(z_0=(1+B_W)z\). To locate it, kill P and
\((P\wedge M)_Q\). In this quotient the cusp is
\((0,\bar z,t)\). Every positive boundary generator
of L, conjugate to \(\tau^m\), has the same central
coordinate \(m\bar z\). The actual compact surface
relation therefore gives
\[
                 \Omega+|Q|\bar z=0
                 \quad\text{in }(\bigwedge^2H)_Q,
 \tag{A.81.15}
\]
where \(\Omega=\sum_i A_i\wedge B_i\) is the
symplectic bivector. This fixes the scalar and the
orientation convention; it does not discard boundary
cross terms before passing to the stated quotient.
The anti-symplectic action on H sends \(\Omega\) to
\(-\Omega\). Hence
\[
                         z_0\in I=(P\wedge M)_Q.
 \tag{A.81.16}
\]

The map
\[
                 M^t\longrightarrow I,
                 \qquad v\longmapsto(a_0\wedge v)_Q
 \tag{A.81.17}
\]
is onto. Move a Q-translate of \(a_0\) back to
\(a_0\) in coinvariants, and then average the other
argument over \(\langle t\rangle\). The map is
equivariant for \(-B\) on its source and \(B_W\)
on its target. Since \(B_Wz_0=z_0\), averaging
over \(C_2\) supplies \(v\in M^t\) with
\[
                         Bv=-v,
                 \qquad (a_0\wedge v)_Q=-z_0.
 \tag{A.81.18}
\]
Compose the homogeneous action with conjugation by
\((v,0,1)\). Its square is one, since
\(v+Bv=0\) and \(v\wedge Bv=0\). Direct use of
(A.81.11) changes the cusp discrepancy by
\(-(v\wedge a_0)_Q=(a_0\wedge v)_Q\).
Equation (A.81.18) therefore makes the actual cusp
literally inverse. Denote this involution on \(E_p\)
by \(\widetilde b_p\), and its image on \(G_p\)
by \(b_p\). The correction may change the finite
\(G_p\)-action; that action was a permitted choice
within the constructed family.

The lower s- and r-frames lift with their prescribed
Q-values. First lift them to \(G_p\) by A.80's
normalizer argument. Any further lift to \(E_p\)
gives the exact conjugacy equation: the two source
elements being compared differ by a central p-element
and have the same order two or three, so that element
is trivial. The normalizer square lies in W and is
fixed by \(\widetilde b_p\); its half corrects the
frame without changing its \(G_p\)-image. Thus both
frames are norm one in the actual \(E_p\).

\emph{Match the old central primes and apply the full
ordinary criterion.}\par
Let \(Q_1=\Pi/\overline{[L,\Pi]}\), the actual
old first central refinement. It is finite, and its
kernel over Q has exponent dividing \(6|Q|\), by
transfer. Thus \(|Q_1|\) is prime to p. For every
prime \(\ell\ne p\), restriction across
\(N\lhd L\), with \(L/N=M\) a p-group, gives
\[
 \operatorname{Hom}_\Pi(L,\mathbb Q_\ell/\mathbb Z_\ell)
       \simeq
 \operatorname{Hom}_\Pi(N,\mathbb Q_\ell/\mathbb Z_\ell).
 \tag{A.81.19}
\]
The positive-degree cohomology of M on these trivial
\(\ell\)-primary coefficients vanishes by averaging;
the inflation--restriction isomorphism is equivariant
for the outer \(\Pi/L\)-action, so taking invariants
gives (A.81.19). These are continuous characters.
Equivalently, the prime-to-p central parts of the two
actual source refinements agree. Together with
(A.81.13), this proves the actual fiber-product identity
\[
                         D_p=E_p\times_Q Q_1.
 \tag{A.81.20}
\]
The source image is full: over Q its two residual
kernels have coprime orders, and a subdirect product
of those kernels is their full product. If
\(d=1+|Q|/6\) and \(w=\dim_{\mathbb F_p}W\), this
also gives the explicit finite order
\[
                              |D_p|=|Q_1|p^{d+w}.
 \tag{A.81.21}
\]

A.79 independently supplies the old involutory
negative action on \(Q_1\), with literal cusp and
both norm-one frames above \(\beta_0,J_0\). Pair
it in (A.81.20) with \(\widetilde b_p\). Their
source actions agree on Q, and their two frames pair
over the same literal Q-values. This gives the
claimed negative involution and both frames on the
actual \(D_p\).

Canonical c preserves L, its characteristic subgroup
N, and \(\overline{[N,\Pi]}\). Composing the
constructed action on \(D_p\) with c gives a
positive peripheral automorphism fixing \(\tau\)
literally. The actual central criterion A.3, in its
closed-normal modular form audited after A.34, now
supplies a full ordinary positive automorphism inducing
\(h_p=b_pc_{G_p}\) on \(G_p\). An inner correction
in N fixes the full cusp literally if required.
This does not assert that the full lift induces the
entire chosen action on \(D_p\); the retained action
in that criterion is the one on \(G_p\).

Finally, (A.80.16)--(A.80.18) apply at every eligible
odd p. A full negative involution inducing \(b_p\)
would have cyclic cusp-conjugator image in Q and
could be normalized, preserving both its involutivity
and its Q-action. A.79 excludes the result. The
nonrealization conclusion is therefore independent
of all full cusp and frame normalizations, whereas
the new ordinary full positive realization has just
been proved separately. Only the first central
refinement (A.81.2) has been used.
\end{proof}


\statement{A.82}{Proved}{The actual reflection augmentation module is projective at every prime and free at the dyadic prime}
Let \(E=F\rtimes C_2\), where \(F=\widehat F_2\), be
one of the actual pants extensions already considered: the pure
negative extension of C.28, or the regular extension of C.43.
Let \(c_1,\ldots,c_n\) represent its involution classes, put
\(C_i=\langle c_i\rangle\), and let \(\varepsilon\)
be its orientation character. Here \(n=3\) in the first case
and \(n=1\) in the second. For every prime p put
\[
 \Lambda_G=\mathbb Z_p[[G]],\qquad
 P_p=\bigoplus_{i=1}^n\mathbb Z_p[[E/C_i]],\qquad
 R_p=\ker(P_p\longrightarrow\mathbb Z_p),
 \tag{A.82.1}
\]
where the last map sums the augmentations. These are actual
compact modules over the full group E. For every finite discrete
E-module A, the restriction maps satisfy
\[
 \begin{aligned}
 H^1(E,A)&\longrightarrow\bigoplus_{i=1}^n H^1(C_i,A)
                    &&\text{is onto},\\
 H^2(E,A)&\longrightarrow\bigoplus_{i=1}^n H^2(C_i,A)
                    &&\text{is injective}.
 \end{aligned}
 \tag{A.82.2}
\]
Consequently \(R_p\) is finitely generated projective for
every p, and
\[
 R_p\oplus(R_p\otimes\varepsilon)\simeq\Lambda_E^{\,n+1},
 \qquad
 R_2\simeq\mathbb Z_2[[E]]^{(n+1)/2}.
 \tag{A.82.3}
\]
In particular the two dyadic ranks are two and one, respectively.
This does not produce a finite graph of groups or a generating
marked reflection system.

\begin{proof}
C.28 supplies the three classes in the pure negative case.
For the regular case we use the complete C.43 hypotheses:
\(F=\widehat F(x,y)\) is the full free subgroup of index
two, the odd involution interchanges the classes of x and y
with exponent \(-1\), and it preserves the class of
\(z=(xy)^{-1}\) with exponent \(-1\). C.43 proves
that there is exactly one F-conjugacy class of sections
and that its centralizer in F is trivial. Since F is
torsion-free, these are exactly the involution classes
in E, and in both cases \(C_E(c_i)=C_i\).
The index-two free subgroup and these self-centralizers
make E real projective, by the exact criterion in
\citet[Theorem~2.1]{A_PalQuickFormal2025}.

\emph{Prescribed involution classes are permitted.}\par
We use the allowed-class formulation of real projectivity in
\citet[Definition~1.1(4), Remark~1.2(b), and Proposition~2.3]{AHaran1990}.
We recall why the finite strengthening introduces no recognition
hypothesis. Given a finite epimorphism \(B\to Q\) and
a conjugacy-closed set \(\mathcal C\) of involutions in B,
form a finite-rank free profinite product D with one \(C_2\)
factor mapping to a representative of each class in
\(\mathcal C\), and enough free generators to map onto B.
The full free-product torsion theorem puts every involution
of D over \(\mathcal C\); see
\citet[Theorem~A']{A_HR1985}.

Let S be the preimage of \(B\setminus(\mathcal C\cup\{1\})\).
The compact set \(\{s^2:s\in S\}\) avoids one. Choose
an open normal U in \(\ker(D\to B)\) disjoint from that
set. Then \(D/U\to B\) lifts every element of
\(\mathcal C\) to an involution and has no involution
above a forbidden element of B. If the original finite
problem for E has nonidentity involution heads lying under
\(\mathcal C\), real projectivity solves the composite
problem through \(D/U\). Its projection to B sends all
involutions into \(\mathcal C\). This gives an ordinary
solution with the prescribed allowed classes; it does not
assert a proper solution.

\emph{Refine the finite action head before prescribing local
cocycles.}\par
Fix a finite E-module A and a finite action quotient
\(q_0:E\to Q_0\). There is a finer finite quotient
\(q:E\to Q\) in which all \(q(c_i)\) are nonidentity
and pairwise nonconjugate, and such that
\[
 \operatorname{im}\bigl(C_Q(q(c_i))\longrightarrow Q_0\bigr)
                         \subseteq q_0(C_i)
                 \quad(1\le i\le n).
 \tag{A.82.4}
\]
Nonidentity follows by retaining the orientation quotient.
Distinct compact conjugacy classes remain distinct in a
sufficiently fine finite quotient. For the centralizers,
as N ranges over open normal subgroups contained in
\(\ker q_0\), the closed sets
\[
                  X_{i,N}=\{g\in E:[g,c_i]\in N\}
\]
decrease to \(C_i\). If none had its \(q_0\)-image
inside \(q_0(C_i)\), compactness would give an element
of \(C_E(c_i)\) with a forbidden \(q_0\)-image.
This is impossible. A common refinement gives (A.82.4).
Only the projected centralizers are controlled here;
centralizers in arbitrary finite quotients need not be cyclic.

Choose arbitrary local classes represented by
\(a_i\in A\) with \((1+c_i)a_i=0\). In
\(B=A\rtimes Q\), each
\[
                            v_i=(a_i,q(c_i))
\]
is an involution. Use the union of their B-conjugacy classes
as the allowed set. The preceding finite problem has a
solution \(\psi:E\to B\), and we can write
\(\psi(g)=(d(g),q(g))\) for a continuous cocycle d.
The pairwise nonconjugate heads force \(\psi(c_i)\)
to be conjugate to \(v_i\), rather than another
\(v_j\). A conjugator \((a,k)\) has
\(k\in C_Q(q(c_i))\), and its formula is
\[
                         d(c_i)=ka_i+(1-c_i)a.
 \tag{A.82.5}
\]
By (A.82.4), k acts on A as one or as \(c_i\).
Either action leaves \([a_i]\in H^1(C_i,A)\)
unchanged. Thus d has all prescribed local classes,
proving the first assertion of (A.82.2). The centralizer
refinement is essential: allowed conjugacy classes alone
would not give the displayed local cohomology classes.

\emph{Degree two and compact projectivity.}\par
A continuous cocycle representing a class in \(H^2(E,A)\)
factors through a finite quotient through which the action
also factors. Refine that quotient to retain all involution
heads. If its restrictions to the \(C_i\) vanish, the
corresponding finite extension has involution lifts of all
these heads. A solution of the real embedding problem
splits its pullback to E. This proves the degree-two
injectivity in (A.82.2).

Continuous Ext and Shapiro applied to (A.82.1) give
\[
 \begin{gathered}
 H^1(E,A)\longrightarrow\bigoplus_iH^1(C_i,A)
       \longrightarrow\operatorname{Ext}^1_E(R_p,A)\\
       \longrightarrow H^2(E,A)
       \longrightarrow\bigoplus_iH^2(C_i,A).
 \end{gathered}
 \tag{A.82.6}
\]
For every finite p-primary A, the middle Ext group is
therefore zero. This is the compact projectivity criterion,
not merely a calculation with trivial coefficients. Indeed,
every map from \(R_p\) into a finite module quotient
lifts through a finite epimorphism by the vanishing of its
pullback extension class. The restriction computation below
makes \(R_p\) finitely generated. The finite lifting sets
are consequently finite, and compactness over the finite
quotients of any compact module epimorphism supplies a
continuous lift. Hence \(R_p\) is projective.

\emph{Split the induction counit, including at two.}\par
Each \(E/C_i\) restricts to the regular F-set, whence
\[
 \operatorname{Res}_F R_p
       \simeq\Lambda_F^{\,n-1}\oplus I_F
       \simeq\Lambda_F^{\,n+1}.
 \tag{A.82.7}
\]
The last equality uses the rank-two free augmentation ideal.
Induction and the tensor identity give the exact sequence
\[
 0\longrightarrow R_p\otimes\varepsilon
  \longrightarrow\operatorname{Ind}_F^E\operatorname{Res}_F R_p
  \longrightarrow R_p\longrightarrow0,
 \qquad
 \operatorname{Ind}_F^E\operatorname{Res}_F R_p
       \simeq\Lambda_E^{\,n+1}.
 \tag{A.82.8}
\]
The last map is the induction counit, and its kernel comes
from the augmentation ideal of \(\mathbb Z_p[E/F]\).
The sequence splits by the projectivity just proved.
At \(p=2\) this is not a splitting of
\(\mathbb Z_2[C_2]\) into sign eigenspaces.

Modulo two the twist disappears. Thus, with
\(\overline R=R_2/2R_2\),
\[
             \overline R\oplus\overline R
                     \simeq\mathbb F_2[[E]]^{\,n+1}.
 \tag{A.82.9}
\]
At every finite quotient \(E/U\), finite-length
Krull--Schmidt cancellation makes the reduction of
\(\overline R\) free of rank \((n+1)/2\).
For completeness, this cancellation uses the finite-length
Fitting decomposition: an endomorphism of an indecomposable
module is invertible or nilpotent, so its endomorphism ring
is local; the usual summand-exchange argument then gives
unique indecomposable multiplicities. Doubling those
multiplicities in (A.82.9) permits cancellation.

The finite sets of bases have the inverse-limit finite
intersection property: a basis at a common finer quotient
induces compatible bases at any finite collection of levels.
No arbitrary prechosen bases are being retained. Compactness
gives \(\overline R\simeq\mathbb F_2[[E]]^{(n+1)/2}\).
Lift such a basis to \(R_2\). The cokernel of the resulting
map from \(\Lambda_E^{(n+1)/2}\) is a compact pro-two
abelian group C with \(C=2C\), hence is zero by its finite
quotients. The surjection splits by projectivity; its kernel
has zero reduction modulo two and vanishes by the same
argument. This proves (A.82.3).
\end{proof}

\emph{The actual incidence map still has to be invertible.}\par
In the three-class case let \(e_i\) be the identity-coset
vector of \(\mathbb Z_2[[E/C_i]]\). The particular map
\[
 d:\mathbb Z_2[[E]]^2\longrightarrow R_2,
       \qquad d(a,b)=a(e_2-e_1)+b(e_3-e_1)
 \tag{A.82.10}
\]
has cokernel
\[
 \operatorname{coker}d\simeq
 \ker\bigl(\mathbb Z_2[[E/H]]\longrightarrow\mathbb Z_2\bigr),
       \qquad H=\overline{\langle C_1,C_2,C_3\rangle}.
 \tag{A.82.11}
\]
Indeed identifying the three identity vertices in the
permutation modules imposes exactly the relations of H;
this is also the component module of the two-edge coset
graph. The reduced component module vanishes exactly when
\(H=E\), since a proper closed H is contained in a proper
open subgroup and its nontrivial finite coset module detects
the difference. For the one-class case fix an actual \(x\in F\),
put \(H_x=\overline{\langle C_1,x\rangle}\), and
write the entire one-loop map as
\[
 \begin{gathered}
 d_x:\mathbb Z_2[[E]]\longrightarrow R_2,
       \qquad d_x(a)=a(xe_1-e_1),\\
 \operatorname{coker}d_x\simeq
 \ker\bigl(\mathbb Z_2[[E/H_x]]\longrightarrow\mathbb Z_2\bigr).
 \end{gathered}
 \tag{A.82.12}
\]
Indeed the relations \(ge_1=gxe_1\), together with the
existing \(gce_1=ge_1\) for \(c\in C_1\), identify
exactly the cosets \(gH_x\) of E. This proves the
cokernel description directly and makes surjectivity
equivalent to \(H_x=E\).

In either case a surjective displayed incidence map is
an isomorphism: after choosing the target basis from
(A.82.3), it is a surjective endomorphism of a finite-rank
free completed group-ring module. Its reductions over
\((\mathbb Z/2^a)[E/U]\) are surjections of finite
sets of equal order, hence injective; these quotients
separate the original module. Thus target freeness
leaves the selected incidence map's surjectivity---the
required generation assertion---unproved.
The full completed group algebra is not local; a unit test
after augmentation or a pro-two specialization does not
supply that assertion. No marked peripheral conjugator is
claimed.


\statement{A.83}{Proved}{Actual odd finite-head twisted duality fixes all self-dual multiplicities}
Let E be the pure negative pants extension in A.82, and
let K be an open normal subgroup of E contained in F.
Put \(Q=E/K\), let \(k=\mathbb F_p\) for an odd
prime p, and define the actual finite module
\[
             \overline R=k\otimes_{\mathbb Z_p}(R_p)_K.
 \tag{A.83.1}
\]
In the Grothendieck group of finite \(k[Q]\)-modules,
\[
                    [\overline R^*]=[\overline R\otimes\varepsilon].
 \tag{A.83.2}
\]
If \(p\nmid|Q|\), this is an actual module isomorphism.
Under this same prime-to-order hypothesis, A.82 and the
isomorphism force every self-dual simple S to
occur in \(\overline R\) with multiplicity
\[
                         2\dim_{\operatorname{End}_Q(S)}S.
 \tag{A.83.3}
\]
In particular, if every simple \(k[Q]\)-module is
self-dual and \(p\nmid|Q|\), then
\(\overline R\simeq k[Q]^2\).
The nonself-dual paired blocks and the modular isomorphism
problem are not removed by this statement.

\begin{proof}
Write
\[
 A=\bigoplus_{i=1}^3 k[Q/q(C_i)],\qquad M=H_1(K,k).
\]
Restriction of every summand of \(P_p\) to K is free,
since K meets its torsion stabilizers trivially. The
homology sequence of (A.82.1), first reduced modulo p,
gives the actual exact sequence
\[
                 0\longrightarrow M\longrightarrow\overline R
                    \longrightarrow A\longrightarrow k
                    \longrightarrow0.
 \tag{A.83.4}
\]
Thus \([\overline R]=[M]+[A]-1\). The module
\(\overline R\) is projective by A.82, although
semisimplicity will only be used under the stated
prime-to-order hypothesis.

\emph{Retain the actual boundary and its sign.}\par
Let \(D_j\) be the actual dihedral normalizer of the
jth base cusp in E, and write \(\overline D_j=q(D_j)\).
Its rotation subgroup is the image of that cusp; its order
may be any positive integer, including one. The actual
cusp permutation module is
\[
                         C=\bigoplus_{j=1}^3 k[Q/\overline D_j].
 \tag{A.83.5}
\]
These cosets count precisely the positive oriented
boundary circles of the K-cover. A negative element
reverses every boundary orientation. The boundary
submodule of M is therefore the quotient of
\(C\otimes\varepsilon\) by its one signed total-sum
relation. If H is the compact-surface homology quotient,
\[
                         [M]=[H]+\varepsilon([C]-1).
 \tag{A.83.6}
\]
This is a quotient by the diagonal total-sum vector,
not a claimed splitting of an augmentation map; it
remains valid when p divides the number of cusps.

The intersection form on H is nondegenerate and changes
by \(\varepsilon\), giving
\[
                                H^*\simeq H\otimes\varepsilon.
 \tag{A.83.7}
\]
The sign is an actual surface-pair calculation, as in C.6
and A.81, rather than a Nielsen realization assumption.
The boundary of the relative degree-two orientation class
is the total sum of the oriented boundary generators.
Literal exponent \(-1\) on them makes the induced
relative degree negative; naturality of the compact
intersection pairing gives (A.83.7). Positive deck
transformations have degree one. This uses the source
cyclic boundary groups and the actual compact quotient.

\emph{Use the finite dihedral polygon before summing
endpoints.}\par
For \(D=\overline D_j\), let \(C_+,C_-\) be the
images of the two local reflection subgroups. A
barycentrically subdivided polygon has two vertex orbits,
one free edge orbit, and orientation character
\(\varepsilon\). Its cellular sequence is
\[
 0\longrightarrow k(\varepsilon)\longrightarrow k[D]
 \longrightarrow k[D/C_+]\oplus k[D/C_-]
 \longrightarrow k\longrightarrow0.
 \tag{A.83.8}
\]
When the rotation image is trivial, the same sequence is
the two-edge polygon, with both vertex stabilizers equal
to D. Thus no separate nondegeneracy assumption on the
finite cusp image is required. In the representation
Grothendieck group, (A.83.8) implies
\[
 (1-\varepsilon)\bigl([k[D/C_+]]+[k[D/C_-]]\bigr)
                              =2(1-\varepsilon)[k].
 \tag{A.83.9}
\]
Here the regular module is unchanged by the sign twist,
and \((1-\varepsilon)^2=2(1-\varepsilon)\).

Induce (A.83.9) to Q and sum over the three actual
cusps. C.28's incidence triangle says that each global
involution class occurs exactly twice among the six
local endpoints. Consequently
\[
                     2(1-\varepsilon)[A]
                              =2(1-\varepsilon)[C].
 \tag{A.83.10}
\]
The group \(G_0(k[Q])\) is free abelian on its simple
modules, so the integer two can be cancelled. We obtain
\((1-\varepsilon)([A]-[C])=0\).

\emph{The precise duality and multiplicity consequence.}\par
Both A and C are permutation modules and hence self-dual.
Using (A.83.4), (A.83.6), and (A.83.7), we find
\[
 \begin{split}
 [\overline R^*]-[\overline R\otimes\varepsilon]
              &=(1-\varepsilon)([A]-[C])\\
              &=0.
 \end{split}
 \tag{A.83.11}
\]
This proves (A.83.2) for every odd p. If \(p\nmid|Q|\),
Maschke's theorem makes equality of these simple
multiplicities an isomorphism; see
\citet[Theorem~7.4 and Proposition~7.9]{AMilne2025}.
A.82 also reduces to
\[
             \overline R\oplus(\overline R\otimes\varepsilon)
                                      \simeq k[Q]^4.
 \tag{A.83.12}
\]
For self-dual S, (A.83.11) equates its two orientation
multiplicities. Their sum in (A.83.12) is four times
the regular multiplicity. This gives (A.83.3), where
\(\operatorname{End}_Q(S)\) is the finite division
field. Over a nonsplitting field the regular multiplicity
is \(\dim_{\operatorname{End}_Q(S)}S\), not
\(\dim_k S\).

For a nonself-dual pair, the duality can instead exchange
opposite orientation imbalances. No such imbalance has
been constructed in an actual \(\overline R\), but
these identities do not exclude it. When p divides
\(|Q|\), the proof asserts the composition-factor
identity (A.83.2) and does not invoke an additional
Cartan-map theorem or a compact inverse-limit duality.
In particular it supplies no unproved full integral
self-duality, finite tree, or marked generation assertion.
A.84 below supplies the additional Cartan-map and compact
projective-dual arguments, upgrading the modular comparison
without removing the nonselfdual ambiguity.
\end{proof}


\statement{A.84}{Proved}{The actual compact augmentation module has orientation-twisted projective duality}
Let E be the pure negative pants extension of A.82. For every prime p
put \(\Lambda=\mathbb Z_p[[E]]\) and let
\(\iota:\Lambda\longrightarrow\Lambda\) be the continuous
anti-involution induced by \(g\mapsto g^{-1}\). For a finitely generated
projective compact left \(\Lambda\)-module P define its projective dual by
\[
 P^\#=\operatorname{Hom}_{\Lambda,\mathrm{cont}}(P,\Lambda),
 \qquad (a\cdot f)(v)=f(v)\iota(a).
 \tag{A.84.1}
\]
This is a compact left module; it is not the linear or Pontryagin dual
of the underlying infinite compact abelian group. The actual module of
A.82 satisfies
\[
                         R_p^\#\simeq R_p\otimes\varepsilon.
 \tag{A.84.2}
\]
For every odd p and every finite head \(Q=E/K\), where
\(K\trianglelefteq_o E\) and \(K\subseteq F\), the finite module
\(\overline R=\mathbb F_p\otimes_{\mathbb Z_p}(R_p)_K\)
therefore satisfies the actual isomorphism
\[
                         \overline R^*\simeq
                                    \overline R\otimes\varepsilon,
 \tag{A.84.3}
\]
including when \(p\mid |Q|\).

More precisely, let S be a simple \(\mathbb F_p[Q]\)-module and
\(P(S)\) its projective cover. If \(S\simeq S^*\), its projective
cover occurs as a direct summand of \(\overline R\) exactly
\[
                          2\dim_{\operatorname{End}_Q(S)}S
 \tag{A.84.4}
\]
times. When every simple is self-dual, this gives
\(\overline R\simeq\mathbb F_p[Q]^2\), without the
prime-to-order restriction. In the modular case (A.84.4) concerns
projective-cover summands, not composition-factor multiplicities.
Neither (A.84.2) nor (A.84.4) removes the nonselfdual multiplicity
ambiguity or proves the required incidence map surjective.

\begin{proof}
\emph{The finite projective comparison.}\par
We use the injectivity of the Cartan map for finite group algebras,
in the following precise form: two finite projective modules over a
splitting field of positive characteristic are isomorphic if their
composition-factor multiplicities agree. The proof of
\citet[Proposition~10.2.1, Theorem~10.2.2, and Corollary~10.2.4]{AWebb2016}
was read at this step. Its character pairing identifies the pairing of
a projective-cover character with a simple character as a Kronecker
delta, so projective characters are independent. Additivity on
composition series then gives the asserted comparison. Only
injectivity is used; an integral surjectivity assertion for the
Cartan map is not required.

The comparison also holds over the actual field \(k=\mathbb F_p\).
Choose a finite splitting extension \(k'/k\). Scalar extension
preserves projectivity and composition-factor equality, so the two
extended projectives are isomorphic. On restriction of scalars this
is an isomorphism of d copies of the original modules, where
\(d=[k':k]\). Finite-length Krull--Schmidt cancellation makes their
indecomposable multiplicities equal and hence gives an isomorphism
over k. This is an explicit descent argument, rather than an
assumption that the head is split over \(\mathbb F_p\).

By A.82, \(\overline R\) is projective. Its ordinary dual is
projective as well: the dual of the regular finite group algebra is
isomorphic to its regular module, by the group basis and its dual
basis, and a dual direct summand is again a direct summand.
The orientation twist preserves projectivity. A.83 gives equality
of the composition-factor classes of these two projectives, even
at modular odd heads. The preceding comparison proves (A.84.3).

\emph{Use a dual that commutes with the actual finite base changes.}\par
For \(a\geq1\) and \(K\trianglelefteq_o E\), \(K\subseteq F\), put
\[
 \Lambda_{a,K}=(\mathbb Z/p^a)[E/K],\qquad
 P_{a,K}=\Lambda_{a,K}\widehat\otimes_\Lambda P.
 \tag{A.84.5}
\]
Every such quotient ideal is stable under \(\iota\) and under the
orientation twist. A finite idempotent-matrix description of P shows
that \(P^\#\) is finitely generated projective and gives canonical
base-change identifications
\[
 (P^\#)_{a,K}\simeq(P_{a,K})^\#.
 \tag{A.84.6}
\]
Here the symbol \(\#\) on the right uses definition (A.84.1)
over the finite ring \(\Lambda_{a,K}\).
For a finite free module this is coordinatewise reduction of
the Hom matrix, and a projective module is the image of an idempotent
matrix; reducing that same idempotent proves the assertion for its
direct summand. No flatness of the finite quotient is assumed.

When \(a=1\), the finite group algebra admits the coefficient map
\(t:k[Q]\to k\) taking the coefficient of the identity. The map
\[
 f\longmapsto t\circ f,
 \qquad
 \ell\longmapsto\left(v\longmapsto
                  \sum_{g\in Q}\ell(g^{-1}v)g\right)
 \tag{A.84.7}
\]
gives mutually inverse maps between the right-hand side of
(A.84.6) and the ordinary contragredient dual of \(P_{1,K}\).
For example the coefficient of g in a \(k[Q]\)-linear \(f(v)\)
is \(t(f(g^{-1}v))\). Also
\(t(f(v)g^{-1})=t(f(g^{-1}v))\), so the stated left action is
exactly the contragredient action. Thus the finite duality used in
(A.84.3) is the reduction of the projective dual (A.84.1).

\emph{Lift each finite isomorphism, then make the choices compatible.}\par
Fix an odd p. Write \(M=R_p^\#\) and
\(N=R_p\otimes\varepsilon\). We have just proved
\(M_{1,K}\simeq N_{1,K}\) for every K. Both \(M_{a,K}\) and
\(N_{a,K}\) are finite projective \(\Lambda_{a,K}\)-modules.
Projectivity of the source lifts the composite
\(M_{a,K}\to M_{1,K}\simeq N_{1,K}\) through
\(N_{a,K}\twoheadrightarrow N_{1,K}\). The resulting map
\(M_{a,K}\to N_{a,K}\) is an isomorphism: its cokernel has zero
reduction modulo p and is killed by \(p^a\), hence is zero;
the surjection splits by target projectivity, and the kernel has
the same zero-reduction property and is zero too.

Let \(\mathcal I_{a,K}\) be the finite nonempty set of these
module isomorphisms. The canonical base changes (A.84.6) give
reduction maps between these finite sets. Their compatibility
conditions have the finite intersection property: for any finite
collection of indices and transition conditions, take a common
refinement \((b,L)\), with b at least all the a and L contained
in all the K, choose an isomorphism there, and reduce it to the
specified levels. Compactness of the product of the finite sets
therefore supplies a compatible system. Finite projective modules
over \(\Lambda\) are the inverse limits of their reductions
(A.84.5), as follows from the same idempotent description.
The compatible system and its inverses give (A.84.2).

The quotients with \(K\subseteq F\) are cofinal among the finite
group heads of E, so this is a full completed-group-ring assertion.
We have not taken an inverse limit of ordinary vector-space duals,
whose natural transition maps would point in the opposite direction.
At \(p=2\), (A.84.2) follows directly from A.82: both the projective
dual and the orientation twist of its rank-two free module are
rank-two free.

\emph{Read projective summands, with the correct modular meaning.}\par
For \(A=k[Q]\), the regular module is
\[
 A\simeq\bigoplus_S P(S)^{d_S},\qquad
                         d_S=\dim_{\operatorname{End}_Q(S)}S.
 \tag{A.84.8}
\]
This follows by passing to the semisimple head \(A/\operatorname{rad}A\)
and then taking projective covers. Duality carries \(P(S)\) to
\(P(S^*)\). To check the latter without a semisimplicity assumption,
write \(P(S)=Ae\) for a primitive idempotent e. The group-algebra
coefficient pairing identifies its contragredient dual with
\(A\iota(e)\). The head of this last indecomposable projective is
\(S^*\): the action of \(\iota(e)\) on \(S^*\) is the dual of the
action of e on S, and its action on another simple T vanishes unless
\(T^*\simeq S\). Twisting by \(\varepsilon\) similarly carries
\(P(S)\) to \(P(S\otimes\varepsilon)\).

Let \(n_S\) be the multiplicity of \(P(S)\) in \(\overline R\).
Equation (A.84.3) and the reduction of A.82 give
\[
 n_{S^*}=n_{S\otimes\varepsilon},\qquad
                  n_S+n_{S\otimes\varepsilon}=4d_S.
 \tag{A.84.9}
\]
For self-dual S these imply \(n_S=2d_S\), proving (A.84.4).
If all simples are self-dual, (A.84.8) gives the asserted freeness
at that finite head. For a nonselfdual pair, (A.84.9) still permits
opposite imbalances. We have constructed no such imbalance for the
actual module. The conclusion remains a module theorem: it does not
choose generators that are differences of peripheral vertices and
does not solve the marked generation or full conjugacy questions.
\end{proof}

